\documentclass[11pt]{article}
\usepackage[a4paper,margin=27mm]{geometry}
\usepackage{amsmath,amssymb,amsthm,mathtools}
\usepackage[T1]{fontenc}
\usepackage{lmodern,microtype}
\pdfextension mapfile {+lm.map}
\pdfextension mapfile {+cm.map}
\pdfextension mapfile {+symbols.map}
\usepackage{booktabs,longtable,array,enumitem}
\usepackage{xurl}
\usepackage[hidelinks]{hyperref}
\setlength{\emergencystretch}{2em}
\newtheorem{theorem}{Theorem}[section]
\newtheorem{proposition}[theorem]{Proposition}
\newtheorem{lemma}[theorem]{Lemma}
\newtheorem{corollary}[theorem]{Corollary}
\newtheorem{remark}[theorem]{Remark}
\newtheorem*{remark*}{Remark}
\newcommand{\CP}{\mathbb{C}P}
\newcommand{\Z}{\mathbb{Z}}
\newcommand{\Q}{\mathbb{Q}}
\newcommand{\F}{\mathbb{F}}
\newcommand{\Tor}{\operatorname{Tor}}
\newcommand{\Ext}{\operatorname{Ext}}
\title{Finite categorical models of $\CP^2$\\and two-dimensional resolutions of $S^2$}
\author{AI-produced by GPT-6 Astra (OpenAI, ChatGPT Pro mode)}
\date{25 November 2026\\Working research record --- INCOMPLETE}
\begin{document}
\maketitle
\begin{abstract}
We independently audit the supplied finite-category calculations and pursue
the minimum-size and bounded-resolution questions. We prove the unrestricted interval $18\le n\le22$ with fresh complete scope joins. We verify the explicit
22-arrow category and 23-element monoid constructions, including integral
homology and the primitive cup square. The two supplied witnesses in each
size are isomorphic. We construct and independently check a complete
ordinary-isomorphism census of monoids through order seven. The least
monoid order is exactly 23; there is one minimizer up to isomorphism and
anti-isomorphism, or two ordinary isomorphism classes. For acyclic
categories, we prove that the least object number is five, by explicit
47-arrow, 42-arrow, 39-arrow and 38-arrow models with independent integral and fundamental-group
certificates. The least arrow count at five objects lies between 32 and 38;
we also exclude infinite families of five-object profiles. The converse
assertions about finite monoids with bounded homology are independently
proved, with the chronology of proof exposure stated accurately. The
minimum arrow number, minimum acyclic arrow count at five objects, and
cardinal-gap questions are still open
in this working version. No unreceived certificate stream is treated as
verified. We make no priority claim.
\end{abstract}
\tableofcontents
\section{Part A: finite categorical models}
\label{CP2:partA}
We count identities among arrows. A finite category has finitely many
arrows; its classifying space is the geometric realization of its nerve.
An acyclic category has only identity endomorphisms and no directed cycle
of distinct objects. A category with acyclic support may have nontrivial
corner monoids. These are different hypotheses.

The supplied FINMON return claims $17\leq n\leq22$ and monoid minimum
$23$. Its 14 GB of output and certificate streams were not attached.
Our reconstructed upper bounds and complete minimum-monoid classification
are proved below. Every category through seventeen arrows is independently
excluded. The broader audit of inherited reductions for eighteen through
twenty-one arrows remains open. Every partial computation is scoped to
its actual completed profiles.

\subsection{Comparison maps and elementary obstructions}

Throughout this subsection, $B\mathcal C$ denotes the realization of the
ordinary nerve. Products in a monoid are written in composition order:
$ab$ means that $b$ is applied first. Our chain formulas instead list
arrows chronologically. These conventions are converted explicitly in
the certificate programs.

\begin{lemma}[Proved: deleting an object retract]
\label{CP2:A-retract-deletion}
Suppose $x\ne y$ are objects of a finite category $C$, and there are
arrows $r:x\to y$ and $s:y\to x$ with $sr=1_x$. Let $D$ be the full
subcategory obtained by deleting $x$. The inclusion $j:D\to C$ induces
a homotopy equivalence. It removes at least three arrows, including
identities. If $C$ is connected with $k$ objects and $n$ arrows, the
excess $n-(2k-1)$ drops by at least one.
\end{lemma}
\begin{proof}
For $c\in D$, the comma category $(j\downarrow c)$ has terminal object
$(c,1_c)$. In $(j\downarrow x)$ choose $(y,s)$. An arbitrary object is
$(d,f:d\to x)$, and the arrow $rf:d\to y$ gives a morphism
$(d,f)\to(y,s)$ because $srf=f$. For a comma morphism
$h:(d,f)\to(d',f')$ we have $f'h=f$, so $(rf')h=rf$. These arrows
are therefore the components of a natural transformation from the
identity functor to the constant functor at $(y,s)$. They contract the
classifying space of this comma category. Quillen's Theorem~A,
\S1 of~\cite{CP2-Quillen1973}, applies to the specified inclusion.
The arrows $1_x,r,s$ are distinct and all disappear. Connectedness
follows either from the homotopy equivalence or directly from the
retract paths. Subtracting $2(k-1)-1$ from the new arrow count proves
the excess assertion.
\end{proof}

\begin{remark*}
It would be incorrect to define a functor deleting $x$ by conjugating
arrows with $r,s$ without checking identities: that proposed map sends
$1_x$ to $rs$, which need not be $1_y$. The comma-category proof avoids
this error. The statement concerns homotopy equivalence of classifying
spaces, not equivalence of categories.
\end{remark*}

\begin{lemma}[Proved: two blocks and their cup products]
\label{CP2:A-two-block}
Let $C$ be a finite category whose objects are partitioned into full
subcategories $D,E$, with no arrows from $E$ to $D$. There is a
double-bar space $T$ and a natural homotopy-pushout description
\[
 BC\simeq BD\mathop{\cup}\limits_T^h BE.
\]
If two positive-degree cohomology classes over a field both restrict
to zero on $BD$ and $BE$, their product is zero. In particular, if
$H^2(BD;\mathbb Q)=H^2(BE;\mathbb Q)=0$, all rational degree-two
products on $BC$ vanish.
\end{lemma}
\begin{proof}
Use the bisimplicial set whose $(p,q)$-simplices consist of a string
of $p$ arrows in $D$, an arrow from its last object to the first
object of a string of $q$ arrows in $E$, and that latter string.
Composition supplies the two actions and all faces; identities supply
degeneracies. Its iterated realization is $T$. Projection to the two
strings gives maps $T\to BD$ and $T\to BE$.

A simplex of $NC$ either belongs to one block or crosses exactly once
from $D$ to $E$. For a mixed simplex, sum the barycentric coordinates
on the $D$ vertices. Away from the two endpoints, division by this
sum and its complementary sum gives the two simplex coordinates and
one interval coordinate. At an endpoint the irrelevant block is
forgotten. These descriptions commute with faces and degeneracies.
Realizing them identifies $BC$ with the double-mapping-cylinder model,
up to the usual subdivision between diagonal and iterated realization
of the bisimplicial set. This proves the natural comparison asserted.

In that model put $A=BD\sqcup BE$. The inclusion is a cofibration, and
collapsing all of $A$ to one point gives $\Sigma(T_+)$. A class whose
restriction to $A$ is zero lifts to relative cohomology by the exact
sequence of the pair. All positive-degree products in a suspension
are zero: cover the suspension by its two contractible cones, lift
the two factors relative to the respective cones, and use the
relative cup product into the cohomology relative to their union,
which is zero. Applying the quotient identification and then the map
from relative to absolute cohomology proves the assertion.
\end{proof}

\begin{corollary}[Proved: a fixed incoming arrow]
\label{CP2:A-fixed-arrow}
Let $E$ consist of one object $v$ with corner monoid $M$. Suppose
$H^2(BD;\mathbb Q)=0$, and suppose some $f:d\to v$ satisfies
$mf=f$ for every $m\in M$. Then every rational degree-two product
on $BC$ is zero. The dual assertion holds for a fixed outgoing arrow.
\end{corollary}
\begin{proof}
The fixed arrow defines a simplicial section $BM\to T$ of the
projection $T\to BM$: use the object $d$ and arrow $f$ in the
$D$ and correspondence coordinates and retain the complete $M$
string. The equations $mf=f$ verify the face identities. Pullback
$H^2(BM;\mathbb Q)\to H^2(T;\mathbb Q)$ is therefore injective.
Mayer--Vietoris and the vanishing on $D$ imply that every degree-two
class of the pushout restricts to zero on both endpoint spaces.
Lemma~\ref{CP2:A-two-block} applies.
\end{proof}

\begin{lemma}[Proved: elementary facts about the minimal ideal]
\label{CP2:A-small-minimal-ideal}
Let $M$ be a finite monoid, $I$ a minimal nonempty two-sided ideal,
and $e\in I$ an idempotent. Put $G=eMe$. Then $G$ is a group.
If $Me=G$ or $eM=G$, the specified inclusion of the group corner in
the Karoubi envelope gives a homotopy equivalence $BG\simeq BM$.
If neither equality holds and $1\notin I$, then $|M|\ge5$.
At order five this last case is uniquely the identity-adjoined
$2\times2$ rectangular band.
\end{lemma}
\begin{proof}
Every finite ideal contains an idempotent: among the eventually
periodic powers of one of its elements choose a sufficiently large
multiple of the period. For $a\in eMe$, minimality gives $MaM=I$,
so $e=uav$ for suitable $u,v\in M$ and
$e=(eue)a(eve)$. In a finite monoid a product equal to the identity
forces its factors to be units: a one-sided inverse makes a finite
translation surjective and hence injective, giving the other
inverse. Applying this fact twice inside $eMe$ proves that it is a
group.

Let $K$ be the full two-object subcategory of $\operatorname{Kar}(M)$
on $1,e$ (one object when $e=1$). The object $e$ is a retract of $1$,
so Lemma~\ref{CP2:A-retract-deletion} identifies $BK$ with $BM$.
If $Me=G$, include the object $e$ into $K$. The comma category over
$1$ has objects the elements of $Me=G$ and arrows given by the free
transitive right $G$ action; there is exactly one arrow between
any ordered pair of objects. The comma category over $e$ has the
same description. Both are contractible, so Theorem~A gives
$BG\simeq BK$. If $eM=G$, use the dual comma categories.

Suppose neither equality holds. Choose $a\in Me\setminus G$.
Since $ea\in G$, replace $a$ by $a(ea)^{-1}$, obtaining
$ae=a$ and $ea=e$. Similarly choose $b\in eM\setminus G$ with
$eb=b$ and $be=e$. The four elements $e,a,b,ab$ are distinct:
the pairs of their left and right products with $e$ are respectively
$(e,e),(e,a),(b,e),(b,a)$. They lie in $I$, while $1$ lies outside.
Any element of $G$ other than $e$ is different from these four, since
its two products with $e$ both equal itself. Thus at order five
$G=\{e\}$ and $I=\{e,a,b,ab\}$. Now $ba\in eMe$ equals $e$,
and the relations just displayed determine the rectangular-band
table on these four elements. The remaining element is the
identity. Finally, if $1\in I$, then $I=M$ and the preceding
unit argument, with $e=1$, says that $M$ itself is a group.
\end{proof}

\begin{lemma}[Proved: rectangular bands with an identity]
\label{CP2:A-band-spheres}
For nonempty finite sets $I,J$ of sizes $r,c$, let
$M=\{1\}\sqcup\{x_{ij}:i\in I,j\in J\}$ with
$x_{ij}x_{kl}=x_{il}$. Then
\[
 BM\simeq\bigvee^{(r-1)(c-1)}S^2,
\]
where a wedge of zero spheres means a point.
\end{lemma}
\begin{proof}
Choose distinguished indices $0$ and put $R=\mathbb ZM$, $e=x_{00}$,
$P=Re$, $Q=R(1-e)$ and $D=R/\mathbb Z[I\times J]$. These are left
$R$-modules; $P,Q$ are projective because their defining elements
are orthogonal idempotents. The quotient $D$ is infinite cyclic as
an abelian group, with all band elements acting by zero.

The augmentation $P\to\mathbb Z$ has kernel with basis
$x_{i0}-x_{00}$ for $i\ne0$. Each band element annihilates these
differences, so this kernel is $D^{r-1}$ as an $R$-module.
The module $Q$ has the abelian basis
\[
 1-e,\qquad x_{ij}-x_{i0}\quad(i\in I,\ j\ne0).
\]
The map $Q\to D$ induced by the quotient has kernel the direct sum,
over $j\ne0$, of copies of $P$, with the map on the $j$th copy
$x_{i0}\mapsto x_{ij}-x_{i0}$. These maps are $R$-linear by the
rectangular-band multiplication rule. They are injective, and
their displayed bases prove exactness over the integers. Splicing
gives the complete projective resolution
\[
0\longrightarrow P^{(r-1)(c-1)}\longrightarrow Q^{r-1}
 \longrightarrow P\longrightarrow\mathbb Z\longrightarrow0.
\]
Tensor with the trivial right $R$-module $\mathbb Z$. The terms
$P,Q$ become respectively $\mathbb Z,0$, because the augmentations
of $e,1-e$ are $1,0$. Thus integral homology is $\mathbb Z$ in
degree zero, $\mathbb Z^{(r-1)(c-1)}$ in degree two, and zero
in every other positive degree. The universal group is trivial,
since every nonidentity element is idempotent. Hurewicz supplies a
map from the stated wedge inducing an isomorphism on $H_2$.
It is an integral homology equivalence between simply connected
CW spaces, hence a homotopy equivalence.
\end{proof}

\begin{corollary}[Proved: the order-five homotopy assertion in R7]
\label{CP2:A-five-types}
Every finite monoid of order at most five has a classifying space
homotopy equivalent to that of a finite group, except for the
five-element rectangular-band monoid, whose classifying space is
$S^2$. In particular, every corner of order at most four is
rationally acyclic, and every corner of order at most five has
zero third rational homology.
\end{corollary}
\begin{proof}
Apply Lemmas~\ref{CP2:A-small-minimal-ideal} and
\ref{CP2:A-band-spheres}. A finite group has zero positive rational
homology, since averaging splits its trivial rational module.
No finite enumeration is needed for this homotopy classification;
the separate census below verifies the table counts and identities.
\end{proof}

\begin{lemma}[Proved: exact projectivity test]
\label{CP2:A-projectivity-test}
For the preceding $M,e,G$, the trivial left $\mathbb QM$-module is
projective if and only if $Me=G$. The trivial right module is
projective if and only if $eM=G$.
\end{lemma}
\begin{proof}
Projectivity on the left is equivalent to a left-module splitting
of the augmentation $\mathbb QM\to\mathbb Q$: it suffices because
$\mathbb QM$ is free, and it is necessary by the defining lifting
property. Such a splitting is an invariant vector of augmentation
one. If $Me=G$, left multiplication by $m\in M$ on $G$ is group
translation by $me\in G$. The vector
$|G|^{-1}\sum_{g\in G}g$ supplies the splitting.

For the converse, first note that $eM$ is a minimal nonempty right
ideal. If $x\in eM$, write $e=uxv$ by minimality of the two-sided
ideal. Since $ex=x$, we have $e=(eue)xv$, so the inverse of
$eue\in G$ gives $e\in xM$. Thus $xM=eM$. The same argument
applies to every idempotent in $I$.
If $Me\ne G$, choose the normalized idempotent $a$ from the proof
of Lemma~\ref{CP2:A-small-minimal-ideal}. The minimal right ideals
$eM$ and $aM$ are disjoint: if they met they would be equal, and
then $a\in eM\cap Me=eMe=G$, a contradiction.
An invariant vector fixed by left multiplication by $e$ is
supported in $eM$, while one fixed by $a$ is supported in $aM$.
Hence an invariant vector must be zero and cannot have augmentation
one. This proves the left assertion. Opposition proves the right.
\end{proof}

\begin{lemma}[Proved: semisimple diagonal and support height]
\label{CP2:A-semisimple-height}
Suppose $C$ has acyclic support, and every corner algebra over
$\mathbb Q$ is semisimple. If every strict chain of objects has at
most $h$ arrows, then $H_i(BC;\mathbb Q)=0$ for $i>h$.
For a two-object category of acyclic support with at most five
arrows, $H^2(BC;\mathbb Q)=0$.
\end{lemma}
\begin{proof}
Let $A=\mathbb QC$, let $B$ be the product of the corner algebras,
and let $J$ be the span of all strict arrows. Acyclic support makes
$J$ an ideal and $A=B\oplus J$ an algebra splitting. The normalized
relative bar resolution of the constant left module $L$ has terms
$A\otimes_B J^{\otimes_B p}\otimes_B L$. Its standard insertion-of-
identity contraction proves exactness before augmentation. Because
$B$ is semisimple every $B$-module is projective, so each displayed
term is projective over $A$. A tensor in $J^{\otimes_B p}$ requires
a strict composable chain of $p$ objects changes; the tensor is
zero for $p>h$. Tensoring with the constant right module computes
the usual categorical nerve homology, giving the asserted bound.

For the last statement, write the two corner orders as $a,b$ and
the nonempty cross-hom size as $s$. If $s=1$, the double bar is
$BM\times BN$ and the collage is the join $BM*BN$. Both factors
are connected and have zero rational first homology (their
universal groups are finite), so the join has zero rational second
homology. If $s\ge2$, then $a+b\le3$, so both corner algebras
are semisimple: an order-two monoid is either $C_2$ or
$\{1,z:z^2=z\}$, both having rational algebra $\mathbb Q\times
\mathbb Q$. Apply the height-one case just proved.
\end{proof}

\begin{lemma}[Proved: singleton cross-homs]
\label{CP2:A-singleton-support}
Let $C$ be connected, have acyclic support and at least two objects,
and have one arrow in each nonempty cross-hom. If $P$ is the
support poset, the functor $q:C\to P$ induces an isomorphism on
rational second cohomology. This is the pullback of an actual
functor and therefore respects products.
\end{lemma}
\begin{proof}
Filter normalized rational chains by the number $p$ of strict
object changes. The internal complex for a chain
$x_0<\cdots<x_p$ is the tensor product of the corner bar complexes:
all intervening hom-sets are singleton trivial bimodules. Its
internal-degree-zero homology is $\mathbb Q$, and all internal
degree-one homology is zero, since every finite monoid has finite
universal group. Thus the bottom row of the resulting first-quadrant
spectral sequence is the chain complex of $BP$, and the next row
is zero.

In internal degree two, an edge contributes the direct sum of the
second homology of its two corners. Its horizontal boundary maps
each summand, with its boundary sign, onto that same corner
summand in horizontal degree zero. Every object is incident to
an edge by connectedness, so $E^2_{0,2}=0$. The only possible
survivor of total degree two is $E^2_{2,0}=H_2(BP;\mathbb Q)$.
No later differential enters it, and its only possible outgoing
differential lands in the zero internal-degree-one row.
The filtration has finitely many horizontal columns, so it
converges in each degree. The map induced by $q$ is the identity
on the bottom row, proving the homology and then cohomology
isomorphism. Naturality gives the assertion about products.
\end{proof}

\begin{lemma}[Proved: the support graph budget]
\label{CP2:A-cycle-budget}
Let $P$ be a finite connected poset, with comparability graph $G$
and cycle rank $R=|E(G)|-|V(G)|+1$.
If $R\le6$, then $H_4(BP;\mathbb Q)=0$.
If $R\le3$ and $H_1(BP;\mathbb F_2)=0$, then $G$ is chordal.
In the latter range its nontrivial clique blocks are triangles
in a clique tree, with at most three triangle additions, or one
complete four-vertex block.
\end{lemma}
\begin{proof}
The order complex is the clique complex of $G$. A four-simplex
requires a complete graph on five vertices, of cycle rank six.
If one occurs and $R\le6$, contracting it leaves a tree; all
remaining edges are bridges. There is exactly one four-simplex,
whose simplicial boundary is nonzero, and there is no
five-simplex. If none occurs, there are no four-chains at all.
Either case gives the first assertion.

For the second, triangle boundaries span the graph cycle space
over $\mathbb F_2$. Choose a basis of at most three triangle
boundaries. An induced cycle of length at least four has at most
one edge in any triangle: two would force a chord. No sum of
three basis triangles could contain all its edges, a
contradiction. Thus there is no induced cycle of length at least
four. Chordality gives the clique-tree decomposition, which can
also be read directly here by successive deletion of a simplicial
vertex. A four-clique uses all three cycles. Otherwise every
nontrivial maximal clique is a triangle, and adding a triangle
along an edge or a vertex adds one cycle. This proves the stated
small-range list.
\end{proof}

\subsection{A complete independent monoid census through order seven}

\begin{theorem}[Proved by complete finite computation]
\label{CP2:A-monoid-census}
The complete numbers of ordinary isomorphism classes of monoids of
orders one through seven are
\[
1,\ 2,\ 7,\ 35,\ 228,\ 2237,\ 31559.
\]
Up to isomorphism and anti-isomorphism the corresponding numbers are
\[1,2,6,27,156,1373,17730.\]
Our ordinary catalogs agree exactly,
after relabelling, with every supplied RESIDUE table in orders
one through five. These are counts of monoids with identity,
not counts of all semigroups or literal oriented table lists.
\end{theorem}
\begin{proof}
Our program \path{code/monoid_census.cpp} fixes the identity label
at zero and treats the other $(n-1)^2$ entries as unknown.
For each triple $(a,b,c)$ whose inner products are known it
enforces $(ab)c=a(bc)$, either rejecting an unequal pair of
known values or forcing one unknown value from the other.
Every such step is a necessary consequence of associativity.
It then takes the first remaining row-major unknown entry and
branches over every element of the monoid. Induction on the
search tree proves that every identity-zero associative table
occurs, and occurs once. The seven order-seven jobs partition
the search by the value of $1\cdot1$ and cover all seven
values disjointly.

At a complete table, intrinsic element invariants partition
possible relabellings. They include principal-ideal sizes,
translation fixed-point counts, commuting-element counts and
the power index and period. An isomorphism preserves these
invariants. We sort their values, enumerate every permutation
within an equal-invariant block, and retain the least resulting
word. Thus we remove exactly ordinary isomorphism duplicates;
opposition is not used in this enumeration.

The independent checker \path{code/check_monoid_census.cpp}
does not use this invariant partition. For every retained
table it enumerates all identity-fixing permutations, verifies
all associativity equations, checks that no two representatives
are isomorphic, computes the automorphism order, and verifies
that the recorded multiplicity is
$(n-1)!/|\operatorname{Aut}(M)|$. It also computes opposition
classes from the full permutation orbits. Finally
\path{code/verify_census.py} independently puts the supplied
RESIDUE tables into unrestricted lexicographic normal form and
compares the complete sets in each order at most five.

The labelled totals, with identity fixed, are
\[
1,\ 2,\ 11,\ 156,\ 4122,\ 208672,\ 18507440.
\]
The independent checker verifies $11,338,875$ associativity
equations in the representative tables. Every check passes;
the exact tables, automorphism-orbit checks, run counts,
source hashes and output hashes are supplied in
\path{outputs/monoid_census/}. Their ordinary and opposition
counts agree with the opened OEIS records
\cite{OEIS_A058129,OEIS_A058133} and the published distinction
between these two equivalence relations. The algorithmic
exhaustion proof above, not agreement of counts alone,
establishes completeness.
\end{proof}

\begin{remark*}[Catalog scope]
The number $1032$ used in the inherited order-eight search is
a restricted possible-positive-$H^2$ oriented cover, not a
complete ordinary order-eight monoid catalog. The published
ordinary count in order eight is $1,668,997$. Any use of that
restricted cover therefore requires its specific homological
reduction and an exhaustive proof within the reduced family.
The new order-one-through-seven catalogs above do not certify
that separate order-eight reduction.
The fresh restricted catalogs used below have $510$ ordinary classes
with nonzero rational $H_2$ at order eight and $8046$ at order nine;
their separate exhaustion proofs are
Propositions~\ref{CP2:A-positive-corners-eight} and
\ref{CP2:A-positive-corners-nine}. For comparison the full published
counts are
\[
\begin{array}{c|r|r}
\text{order}&\text{ordinary isomorphism}&\text{isomorphism and opposition}\\\hline
8&1,668,997&858,977\\
9&3,685,886,630&1,844,075,697.
\end{array}
\]
The ordinary counts are the opened OEIS record A058129; the opposition
counts are A058133 and the abstract of Distler--Kelsey at the authors'
university repository\cite{OEIS_A058129,OEIS_A058133,CP2-DistlerKelsey2009}.
Thus neither restricted catalog is advertised as a full catalog in its
order. Our reduction to nonzero rational second homology is part of
every use of these restricted lists.
\end{remark*}

\subsection{What the literal cup certificates establish}

\begin{proposition}[Proved: the certificate obstruction]
\label{CP2:A-cup-certificate}
Let $C$ be a finite connected category. Suppose normalized
nerve chains over $\mathbb F_2$ give $H_1=0$ and
$\dim H_2=1$. If a nonboundary two-cocycle $u$ admits a
three-cochain $v$ with $\delta v=u\smile u$, or if
$H_3(C;\mathbb F_2)\ne0$, then $BC$ is not homotopy equivalent
to $\CP^2$.
\end{proposition}
\begin{proof}
For chronological nonidentity arrows $[a_1|\cdots|a_j]$, the
normalized boundary deletes the first or last arrow or composes
an adjacent pair, with the usual alternating signs. A term
containing an identity after composition is degenerate and is
zero. These are exactly the face formulas of the nerve, so
their chain homology is that of $BC$.
The Alexander--Whitney formula on a four-simplex is
\[
(u\smile u)[a|b|c|d]=u[a|b]\,u[c|d].
\]
A literal equality $\delta v=u\smile u$ therefore says that
the square of the generator of $H^2(C;\mathbb F_2)$ is zero.
The corresponding generator of $H^2(\CP^2;\mathbb F_2)$ has
nonzero square. The other obstruction follows from the
vanishing of $H_3(\CP^2;\mathbb F_2)$. Both are invariant
under homotopy equivalence.
\end{proof}

In the new linear-profile search, the initial filter and the
cochain producer use separate bar calculations. The independent
checker then reconstructs each typed category from its literal
table, verifies every associativity equation, recomputes the
first two Betti numbers, checks that the supplied cocycle is
not a coboundary, and evaluates its square-primitive equation
on every normalized four-simplex. For a positive-third-homology
certificate it recomputes the necessary fourth-boundary rank.
No finite-degree calculation is used to assert vanishing of an
uncomputed homology tail.


\subsection{Increasing a corner by one arrow}
\begin{lemma}[Proved: identity padding preserves the classifying space]
\label{CP2:A-identity-padding}
Let $C$ be a nonempty finite category and $x$ one of its objects.
There is a category $C^+$ with the same objects and support, one
additional endomorphism at $x$, and $BC^+\simeq BC$. Thus the
existence of a model with $n$ arrows and $k$ objects implies
existence with every larger arrow count and the same $k$.
The assertion concerns support; it does not preserve acyclicity
of the category when that term requires trivial endomorphisms.
\end{lemma}
\begin{proof}
In $C^+$ retain every old arrow and old composite, but declare
the old $1_x$ to be an additional idempotent endomorphism and
adjoin a new identity $1_x^+$. The new arrow is a two-sided
identity for every arrow incident with $x$. Associativity
follows from that of $C$: a triple using a new identity reduces
to an old composition, and all other triples are unchanged.
There is exactly one new arrow, only the $x$ corner changes,
and no hom-set becomes empty or nonempty.

Here is a comparison using actual functors. Adjoin to $C$ a
new object $x'$ and form $D$ as follows. Put $r(x')=x$ and
$r(y)=y$ on old objects. For every ordered pair $(a,b)$ except
$(x',x')$, put $D(a,b)=C(r(a),r(b))$, with distinct typed copies
when a new object occurs. Put
$D(x',x')=C(x,x)\sqcup\{1_{x'}\}$. Old-labelled products use
the composition of $C$, and the new $1_{x'}$ is an identity.
This again defines a category. The map $R:D\to C$ sends every
old-labelled arrow to its label and $1_{x'}$ to $1_x$; it is
a functor. The inclusion $I:C\to D$ is a functor and $RI=1_C$.
The old identity label gives an arrow $u:x\to x'$ in $D$.
With component $u$ at $x'$ and identity components elsewhere,
we have a natural transformation $IR\Rightarrow1_D$.
Naturality is exactly the identity law in $C$ for old-labelled
arrows, and the new identity causes no additional equation.
Thus $BI$ and $BR$ are homotopy inverses.

Let $v:x'\to x$ be the other typed copy of $1_x$. Then
$vu=1_x$, whereas $uv$ is the old identity label in
$D(x',x')$. Hence $x$ is a retract of $x'$. The full
subcategory $D\setminus\{x\}$ has the same classifying
space as $D$ by Lemma~\ref{CP2:A-retract-deletion}; after
renaming $x'$ as $x$, it is exactly $C^+$.
This proves the claimed homotopy equivalence and all count
assertions. Iterating at the chosen corner gives every larger
arrow count. No inclusion $C\to C^+$ that sends old arrows
to themselves is asserted to preserve identities.
\end{proof}


\subsection{Proper retracts across an arbitrary support cut}

\begin{theorem}[Proved: a cut retraction preserves projective-plane type]
\label{CP2:A-cut-retract}
Let $C$ be a finite category whose objects are the disjoint union of
full subcategories $D,E$, with no arrows from $E$ to $D$.
Write $S(d,e)=C(d,e)$ for its connecting profunctor, a functor
$D^{\mathrm{op}}\times E\to\mathbf{FinSet}$.
Suppose $r:S\to S$ is an idempotent natural transformation.
Keep $D,E$ and replace each $S(d,e)$ by $rS(d,e)$, obtaining the
finite subcategory $Y$. If $BC\simeq\CP^2$, then $BY\simeq\CP^2$.
If any component of $r$ is not the identity, $Y$ has strictly fewer
arrows and the same objects and support as $C$.
\end{theorem}
\begin{proof}
Naturality means that every $r_{d,e}$ commutes with precomposition
by arrows of $D$ and postcomposition by arrows of $E$. Its images
therefore form a subprofunctor and define $Y$. There cannot be two
composable crossing arrows. Hence the maps which are identities
on $D,E$ and equal to $r$ on crossing arrows give an actual functor
$R:C\to Y$. For the inclusion $i:Y\hookrightarrow C$, idempotence
gives $Ri=\operatorname{id}_Y$.

Use the natural two-block mapping-cylinder model of
Lemma~\ref{CP2:A-two-block}. Its bridge spaces for $Y$ and $C$ are
$T_Y\subset T_C$, obtained by taking the two-sided bars of the
respective profunctors and then realizing. They are CW subcomplexes.
The two corner spaces are identical. Collapsing the model of $BY$
in that of $BC$ consequently gives
\[
 \frac{T_C\times I}
 {T_C\times\partial I\;\cup\;T_Y\times I}
   =\Sigma(T_C/T_Y).
\]
The models and their inclusions are natural in the profunctor, so
this computes the relative cohomology and its cup product for
$(BC,BY)$. Every product of positive-degree relative classes is
zero, because the displayed quotient is a suspension.

Assume $BC\simeq\CP^2$. The retraction makes each
$H_j(BY;\mathbb Z)$ a direct summand of $H_j(BC;\mathbb Z)$.
Thus $BY$ is connected, its only possible positive integral
homology groups are summands of $\mathbb Z$ in degrees two and
four, and its fundamental group is trivial as a retract of
$\pi_1(BC)=1$. If $H^2(BY;\mathbb Z)=0$, the generator
$u\in H^2(BC;\mathbb Z)$ lifts to relative cohomology; its square
would vanish, a contradiction. Therefore $H^2(BY;\mathbb Z)$ is
$\mathbb Z$ and the split injection
$R^*:H^2(BY;\mathbb Z)\to H^2(BC;\mathbb Z)$ is an isomorphism.
Choose $v$ with $R^*v=u$. Then $R^*(v^2)=u^2$ is the integral
four-dimensional generator. The split degree-four map is also
an isomorphism, and $v^2$ is primitive. The integral recognition
argument of Theorem~\ref{A:monoid-upper23} gives $BY\simeq\CP^2$.

A finite idempotent self-map is a bijection only when it is the
identity. Thus a nonidentity component has a proper image.
Each originally nonempty component remains nonempty, so support
and object set are preserved while the total arrow count drops.
\end{proof}

\begin{corollary}[Proved: endomorphisms of a minimal connecting profunctor]
\label{CP2:A-minimal-cut-endomorphisms}
In a category realizing $\CP^2$ with the least possible total
number of arrows, every natural endomorphism of every connecting
profunctor across a support cut is a componentwise bijection.
The same holds for an acyclic model with a fixed number of objects
and the least arrow count among such models.
\end{corollary}
\begin{proof}
For a natural endomorphism $f$, choose a common positive exponent
$k$ which is at least the transient length of every finite component
map and is divisible by every eventual cycle length. Then $f^k$
is a natural idempotent. If any component of $f$ is not bijective,
its eventual image is proper and that component of $f^k$ differs
from the identity. The theorem gives a strictly smaller category
of the same homotopy type. In the acyclic case, the object set,
trivial corners and support are unchanged, so it is also a smaller
acyclic model at the same object count. Both contradict minimality.
\end{proof}


\subsection{A fresh restricted corner census through order eight}
\begin{proposition}[Proved: positive rational second homology corners]
\label{CP2:A-positive-corners-eight}
Up to ordinary isomorphism, the numbers of monoids of orders
$5,6,7,8$ with nonzero rational second homology are respectively
$1,5,43,510$.
\end{proposition}
\begin{proof}
The minimal-ideal and projectivity reductions above imply that both the
row and column counts of the minimal ideal must be at least two.
If its maximal group were nontrivial, the ideal would have at least
eight elements; adjoining the identity would require at least nine.
The group case has zero positive rational homology. Thus, through order
eight, every required ideal is a rectangular band, of shape $2\times2$,
$2\times3$, or $3\times2$.

Fix a labelled band and its identity extension. Every additional element
acts by a pair of transformations of the row and column sets. The program
\path{code/positive_corners.cpp} enumerates all multisets of such pairs;
sorting them loses only a relabelling of the additional elements.
For each product outside the band, it tries every element whose actor
equals the composite actor. This includes all outside elements with that
actor, the identity when applicable, and the unique band point for a
constant actor. Products involving the band are forced. Propagation uses
only associativity equalities, and every remaining domain is exhausted.
This proves coverage, including nonfaithful action representations.

The complete family counts, before imposing positive second homology, are
\[
\begin{array}{c|rrrr}
\text{shape}&5&6&7&8\\\hline
2\times2&1&9&80&1181\\
2\times3&0&0&1&12\\
3\times2&0&0&1&12
\end{array}
\]
The independent program \path{code/check_positive_corners.cpp}
recomputes every minimal ideal, verifies all $11983584$ associativity
equations, and computes the balanced augmentation quotient exactly over
$\mathbb Q$. This quotient is $H_2$ by the band Tor formula proved below;
its ambient dimension is at most two, so integer determinants settle all
rank decisions. It canonicalizes using unrestricted identity-fixing
permutations, verifies uniqueness, and obtains the asserted positive
counts. Through order seven its positive canonical sets agree exactly
with the separately generated full ordinary census. The scope and literal
tables are in \path{outputs/positive_corners/}.

The old restricted order-eight cover contained $516$ entries, including
six explicitly marked rational-$H_2$-zero entries. Our count $510$
therefore refines that cover; it does not assert that those six entries
were omissions or mistakes. This is not a census of all order-eight
monoids.
\end{proof}


\subsection{Right band acts and the seven-state double-band gate}
\begin{lemma}[Proved: complete right-$P_5$ act normal forms]
\label{CP2:A-P5-normal-forms}
Let $P_5=\{1,e_{00},e_{01},e_{10},e_{11}\}$, with
$e_{ij}e_{kl}=e_{il}$. Every finite right $P_5$-set is determined,
up to an equivariant bijection, by $a$ singleton core orbits,
$b$ two-element core orbits with labelled columns, and a multiset
of $t$ ordered pairs of these $a+b$ orbits. Its cardinality is
$a+2b+t$. The only ambiguity is permutation of the singleton
orbits, permutation of the double orbits, and permutation of
equal-type transient points. Consequently the numbers on
$1,\ldots,7$ points are $1,3,6,17,44,130,383$.
\end{lemma}
\begin{proof}
For a right act $X$, put $Y=X\{e_{ij}\}$. If $y=xe_{ij}$,
then $ye_{kl}=xe_{il}$ is independent of $k$. Thus its orbit
consists of one point or of two points canonically labelled
by the column $l$. Two such orbits meeting at a point coincide,
since applying either column operation to that point recovers
both labelled outputs. These are precisely the core orbits.
For $x\in X\setminus Y$, each row $i$ selects the orbit
$\{xe_{i0},xe_{i1}\}$, and the column selects its labelled
point. These choices are unrestricted and satisfy every act
equation. Transient points with the same ordered orbit pair
are interchangeable; core orbit sizes are preserved. This
proves both existence and completeness of the normal form.

The producer \path{code/p5_right_acts.py} enumerates every
$a,b,t$ with $a+2b+t=s$, all multisets of ordered orbit pairs,
and their $S_a\times S_b$ orbits. An independent count uses
Burnside's lemma. For a permutation of the orbit set, let its
action on ordered pairs have cycle lengths $d_1,\ldots,d_r$.
The fixed multisets of size $t$ are counted by the coefficient
of $x^t$ in $\prod_i(1-x^{d_i})^{-1}$. The separate program
\path{code/check_p5_normal_forms_burnside.py} averages these
integer coefficients and checks every parameter case against
the literal normal forms, as well as every action equation.
It obtains the stated counts. The input tables, all parameter
counts and digests are under \path{outputs/p5_bisets/}.
\end{proof}

\begin{proposition}[Proved: the $P_5$--$P_5$ seven-state case]
\label{CP2:A-P5-P5-seven}
No category with two objects, both corner monoids isomorphic
to $P_5$, seven arrows in one direction between the objects,
and no arrows in the reverse direction has type $\CP^2$.
The total arrow count, including identities, is seventeen.
\end{proposition}
\begin{proof}
Such a category is exactly a commuting left and right $P_5$
action on seven states. Fix a representative of the right act
from Lemma~\ref{CP2:A-P5-normal-forms}. The transpose of the
band indices identifies $P_5^{\mathrm{op}}$ with $P_5$.
Taking all state relabellings of all $383$ normal forms therefore
gives every labelled left action. The new program
\path{code/p5_bisets.cpp} constructs exactly $1038598$ such
left actions. It verifies that the normal-form state orbits are
disjoint and that their cardinality is
$\sum 7!/|\operatorname{Aut}(X)|$.

For each fixed right act it computes every transformation of
the seven states commuting with its four band operations.
It tries every labelled left action whose four maps are in
that centralizer. The automorphism group of the right act
acts on these choices; taking the least member of each orbit
loses only an equivariant state relabelling. Its full group
is obtained by testing every one of the $7!$ state permutations.
The program separately verifies that the number of compatible
labelled actions equals the sum of these orbit cardinalities.
In the trivial-right-action case the already complete $383$
left normal forms are used, with the same orbit-weight check.
No monoid automorphism quotient, opposite-category quotient,
retraction cut or conditional corner reduction is used.

There are $1904633$ compatible labelled left actions over the
right-act representatives and $56597$ resulting state orbits
of bisets. Exactly $15606$ fail $H_1(-;\mathbb F_2)=0$;
another $24704$ fail $\dim H_2(-;\mathbb F_2)=1$.
Every remaining one of the $16287$ literal categories has a
nonzero two-cocycle whose square is the coboundary of a supplied
three-cochain. The independent program
\path{code/check_p5_bisets.cpp}, which imports no producer
code, reconstructs all typed tables and normalized chains.
It checks $12622425$ associativity equations, $7557168$
cocycle equations and $37525248$ square-primitive equations.
All records pass. Since the nonzero class of $H^2(\CP^2;
\mathbb F_2)$ has nonzero square, every case is excluded.
The exact scope, complete literal stream, digests and checker
receipt are in \path{outputs/p5_bisets_live/} in the checkpoint.
This proposition asserts only the double-$P_5$ case; mixed
order-five corner monoids remain a separate gate.
\end{proof}

\begin{corollary}[Proved: the double-$P_5$ cases through seventeen arrows]
\label{CP2:A-P5-P5-through17}
A two-object category with both corners $P_5$, arrows only in one
cross direction, and at most seventeen total arrows cannot have
classifying space $\mathbb{CP}^2$.
\end{corollary}
\begin{proof}
The preceding construction and independent literal check apply unchanged
with $s=1,\ldots,6$ states. The numbers of state-isomorphism orbits are,
respectively, $1,6,32,181,1124,7667$. After the two mod-two homology
filters, the retained counts are $0,4,10,60,360,2320$. Every retained
category has a supplied and independently checked square primitive.
The zero-state case is disconnected. Exact streams and independent
receipts are in \path{outputs/p5_bisets_small/}; the separate checker
is \path{code/check_p5_bisets_all_sizes.cpp}. Together with the
seven-state proposition this covers every possible cross size.
\end{proof}


\subsection{Independently certified upper bounds}

\begin{theorem}[Proved]\label{A:monoid-upper23}
There is an aperiodic monoid of order $23$ whose classifying space is
homotopy equivalent to $\CP^2$. The two supplied constructions numbered
$1412$ and $1651$ are isomorphic, and each has trivial automorphism group.
\end{theorem}
\begin{proof}
Let $I=\{0,1,2\}$ and $J=\{0,1,2,3,4\}$. We multiply transformation
pairs by
\[
 (\sigma,\tau)(\alpha,\beta)=(\sigma\alpha,\beta\tau).
\]
Take the identity, the fifteen constant pairs $x_{ij}$, and the seven
actors in the first pair of columns below. The second pair specifies the
other supplied construction, using its own labels $a,\ldots,g$.
\[
\begin{array}{c|cc|cc}
&\multicolumn{2}{c|}{1412}&\multicolumn{2}{c}{1651}\\
&\sigma&\tau&\sigma&\tau\\\hline
a&002&02200&010&04034\\
b&002&41144&020&04400\\
c&010&04034&020&21122\\
d&202&03300&020&30033\\
e&202&12211&022&01100\\
f&202&40044&022&32233\\
g&212&00334&212&00334
\end{array}
\]
The indicated sets are closed under this associative multiplication. Our
fresh constructor and separate verifier check all $529$ products and
$12\,167$ associativity triples for each table. The literal tables use
labels $0=1$, $x_{ij}=1+5i+j$, and $a,\ldots,g=16,\ldots,22$.
They are in \path{monoid/certificates/cp23_candidate_1412.json} and
\path{monoid/certificates/cp23_candidate_1651.json}.

We give the homotopy certificate for the first table. Put $R=\Z M$ and
define projective left $R$-modules
\[
 P=Rx_{00},\quad E=R(c-x_{00}),\quad F=R(1-c),\quad G=R(a-x_{00}).
\]
Each displayed generator is a ring idempotent. The exact resolution is
\[
0\longrightarrow P\xrightarrow{d_4}G\oplus G
 \xrightarrow{d_3}P\oplus F\xrightarrow{d_2}E
 \xrightarrow{d_1}P\xrightarrow{\varepsilon}\Z\longrightarrow0,
\]
where the $R$-linear maps are determined by
\begin{align*}
d_1(c-x_{00})&=x_{10}-x_{00},\\
d_2(x_{00},0)&=x_{03}-x_{00},\\
d_2(0,1-c)&=c+d-g-x_{00},\\
d_3(a-x_{00},0)&=(x_{00}-x_{20},b+e-f-x_{01}),\\
d_3(0,a-x_{00})&=(0,a-x_{00}),\\
d_4(x_{00})&=(x_{02}-x_{00},-x_{02}+x_{00}).
\end{align*}
For $P$ use the basis $x_{00},x_{10},x_{20}$. Bases for the other modules
are $z-zx_{00}$ for $E,G$ and $z-zc$ for $F$, where respectively
\begin{align*}
E:\quad&z\in\{x_{03},x_{04},x_{13},x_{14},x_{23},x_{24},c,d,f,g\},\\
F:\quad&z\in\{1,x_{01},x_{02},x_{11},x_{12},x_{21},x_{22},a,b,e\},\\
G:\quad&z\in\{x_{02},x_{12},x_{22},a\}.
\end{align*}
The verifier reconstructs these bases inside $\Z^{23}$ and checks every
translate of each generating idempotent has integer coordinates in its
claimed basis. Thus these are the whole principal projectives. It checks
all module actions, all differential equivariance equations, and all zero
composites. The abelian ranks of the terms are $(3,10,13,8,3)$; the ranks
of $\varepsilon,d_1,d_2,d_3,d_4$ are $(1,2,8,5,3)$; every nonzero Smith
invariant is $1$. Each differential image is consequently saturated,
contained in the next kernel, and of the same rank. It equals that kernel.
The last differential is injective. This proves exactness in every degree.

For a ring idempotent $e$, $\Z\otimes_R Re$ is the image of multiplication
by $\varepsilon(e)$ on $\Z$. Consequently the coinvariant terms are
$\Z,0,\Z,0,\Z$, with all differentials zero. We obtain integral homology
$\Z$ in degrees $0,2,4$, and zero in every other degree.

Extend the matrix with rows $(1,1),(0,1),(-1,0),(1,0)$ by zero row and
column at index zero, and call it $L$. The normalized cochain
\[
u(m,n)=L\bigl(\tau_m(0),\sigma_n(0)\bigr)
\]
is an integral two-cocycle, as checked on every triple. We compare the
projective resolution with the normalized left bar resolution. Its
contracting homotopy on a term $z[m_1|\cdots|m_d]$ inserts $z$ at the
front with coefficient $1$ if $z\ne1$, and gives zero if $z=1$.
Starting from $F_0(x_{00})=x_{00}[\,]$, we set
\[
 F_d(e)=e\,s(F_{d-1}d_d(e))
\]
on each projective generator and extend $R$-linearly. Idempotent support
makes the formula well defined. The separate verifier checks every chain
equation and every equivariance equation of the stored comparison.
The comparison has $(3,40,172,448,324)$ sparse terms in degrees $0$--$4$.
Integral evaluation on the coinvariant generators gives
\[
\langle u,F_2(x_{00},0)\rangle=-1,
\qquad \langle u\smile u,F_4(x_{00})\rangle=1.
\]
Thus the degree-two class and its square are primitive.

Every band element is idempotent, hence trivial in the universal group;
the relation $mx_{00}\in K$ then makes every element of $M$ trivial there.
Thus $BM$ is simply connected. We recall the recognition argument, also
covered by \cite[Proposition~4C.3]{CP2-Hatcher2002}. A Hurewicz generator
$S^2\to BM$ has a $3$-connected mapping-cylinder pair: relative
Hurewicz in degree three applies after the degree-two isomorphism and
the vanishing of relative $H_3$. Relative Hurewicz in degree four then
represents the sole relative homology generator by a map of a four-cell.
This gives an integral homology equivalence
$S^2\cup_h e^4\to BM$ between simply connected CW complexes, hence a
homotopy equivalence. The primitive square makes the Hopf invariant of
$h$ equal to $\pm1$. Since $\pi_3(S^2)$ is generated by a Hopf map,
this cofiber is $\CP^2$. The verifier separately checks stabilization of
every element's powers, proving aperiodicity.

The row permutation $(0\ 2)$ and column permutation $(1\ 2)(3\ 4)$
conjugate the first actor set to the second. Their element permutation is
\[
\begin{split}
(&0,11,13,12,15,14,6,8,7,10,9,1,3,2,5,4,\\
 &20,21,22,17,18,19,16).
\end{split}
\]
Our independent isomorphism program checks all $720$ row/column
permutations for every ordered pair of the two models. An isomorphism
must preserve the characteristic minimal band; its restriction there is
a row and column permutation, and faithfulness of the actor representation
determines its values elsewhere. This search is therefore exhaustive. It
finds exactly one isomorphism for each ordered pair, including each
automorphism pair. Neither listed monoid is anti-isomorphic to either
listed monoid, because opposition swaps the unequal minimal-band row and
column counts. Their opposites belong to the same class when opposition
is allowed.
\end{proof}

\begin{lemma}[Proved: constant minimal-band compression]
\label{A:band-action-compression}
Let $M$ have a proper rectangular-band minimal ideal $K=I\times J$,
and let $S$ be a nonempty finite right $M$-set on which every $k\in K$
acts constantly, with value $c_k$. Let $C$ have objects $a,b$, corners
$M,\{1_b\}$, cross hom-set $C(a,b)=S$, and $C(b,a)=\varnothing$.
There is a monoid
\[
 N=(M\setminus K)\amalg(I\times S),\qquad BC\simeq BN.
\]
\end{lemma}
\begin{proof}
Let $\rho:K\to I$ be the row map and let $\sigma_m:I\to I$ be the
left row action. Set $\eta(k)=(\rho(k),c_k)$. Preserve old products
outside $K$, replace old products in $K$ by their $\eta$-images, and set
\[
m(i,s)=(\sigma_m(i),s),\quad (i,s)m=(i,sm),\quad
(i,s)(j,t)=(i,t).
\]
The action equations give
$\eta(mk)=(\sigma_m\rho(k),c_k)$ and
$\eta(km)=(\rho(k),c_km)$; these verify all associativity cases
meeting the new band. The other cases are inherited from $M$. Thus
$q:M\to N$, equal to the identity outside $K$ and to $\eta$ on $K$,
is a unital homomorphism.

Fix $e=(i_0,s_0)$ and let $D$ be the full subcategory of
$\operatorname{Kar}(N)$ on $1,e$. The object $e$ retracts from $1$.
The inclusion $N\to D$ induces a homotopy equivalence by
Quillen's Theorem A \cite[\S1, Theorem~A]{CP2-Quillen1973}.
Indeed, its comma category over $e$ has objects $x\in eN$ and
morphisms $m:x\to y$ satisfying $ym=x$. The components
$x:x\to e$ define a natural transformation from the identity to
the constant functor at $e$: one has $ex=x$, and naturality is
exactly $ym=x$. The comma category over $1$ has a terminal object.
Thus both comma categories are contractible. Define $F:C\to D$ by
$a\mapsto1$, $b\mapsto e$, $m\mapsto q(m)$ and
$s\mapsto(i_0,s)$. The displayed product rules verify its functor laws.

The comma category $(F\downarrow e)$ has terminal object $(b,e)$.
In $(F\downarrow1)$ the objects over $a$ are $n\in N$, with arrows
$n\to n'$ labelled $m\in M$ when $n'q(m)=n$; the objects over $b$
are one sink for each $i\in I$. There is one arrow $(i,s)$ to sink $i$.
For a right $M$-set $X$, write $A(X)$ for the category of arrows
$x\to y$ labelled $m$ with $ym=x$. Its degree-$d$ nerve is naturally
$X\times M^d$, so it preserves colimits and injections. As a right
$M$-set, $N=M\amalg_K(I\times S)$. Each row category is a sieve,
so the comma nerve is obtained by attaching a cone to each row nerve.
Its realization is the pushout
\[
 BA(M)\amalg_{BA(K)}\coprod_{i\in I}\operatorname{Cone}(BA(S)).
\]
The map $BA(K)\to BA(M)$ is a simplicial subcomplex inclusion, hence
a cofibration, so this is a homotopy pushout. The category $A(M)$ has
terminal object $1$. Each component $A(e_iM)$ of $A(K)$ contracts by
the natural transformation $x\to e_i$ labelled $x$: naturality is
exactly $ym=x$. Therefore $BA(K)\to I$ is a homotopy equivalence.
Contracting the cones now shows that the comma category is contractible.
Quillen's theorem gives $BC\simeq BD\simeq BN$.
\end{proof}

\begin{theorem}[Proved]\label{A:category-upper22}
There is a two-object category with $22$ arrows, identities included,
and classifying space $\CP^2$. Both supplied $22$-arrow constructions
are isomorphic and have been independently checked.
\end{theorem}
\begin{proof}
In either monoid $N$ from Theorem~\ref{A:monoid-upper23}, retain the
eight elements outside the band and the ten band elements in rows $0,2$.
This is an $18$-element submonoid $P$: every outside actor preserves
$\{0,2\}$, and a product of two outside actors cannot be the constant
row map with value $1$. Indeed an outside actor taking the value $1$
has singleton inverse image of $1$, so a composite of two such
nonconstant maps cannot be constantly $1$.

Set $M=P^{\mathrm{op}}$ and let $S=I$ act on the right by the original
row transformations. Every element of its $5\times2$ minimal band
acts constantly. The resulting act-cone category has $18+1+3=22$
arrows. Compression in Lemma~\ref{A:band-action-compression} replaces
the $5\times2$ band by a $5\times3$ band and recovers exactly
$N^{\mathrm{op}}$. Its four product types agree with the specified
multiplication. Thus $BC\simeq B(N^{\mathrm{op}})\simeq\CP^2$.

The files \path{monoid/certificates/cp22_candidate_1412.json} and
\path{monoid/certificates/cp22_candidate_1651.json} give every source,
target, identity and composition, as well as the full comparison target
and functor. Undefined composites are recorded as $-1$. The independent
checker verifies $972$ action equations, $382$ composable-pair functor
equations and $6862$ composable associativity triples for each category.
The isomorphism of the source monoids preserves the retained row set;
its restriction on the corner and its row permutation on the cross
arrows give an isomorphism of the two categories, checked on every
composition.

The actual comparison functor has target the two-object category
$\operatorname{Kar}(N^{\mathrm{op}})\{1,x_{00}\}$, with $32$ arrows.
This corrects the prompt's wording ``onto the $23$-element monoid.''
The comparison is the proved span of homotopy equivalences; the upper
bound is unaffected. Neither theorem in this subsection asserts a
matching lower bound or a global classification of all minimal models.
\end{proof}


\subsection{The monoid lower bound: completed reductions and exact scope}
\label{A:monoid-lower-audit}

We distinguish a complete classification inside an explicitly specified
action hull from the classification of all finite monoids. This subsection
proves the three-row gates and the fixed-hull classification; the completed
large-band and four-by-four arguments follow in
Sections~\ref{A:large-band-target-proof} and \ref{A:monoid-global-answer}.
Together they establish the global answer and classification in
Theorem~\ref{A:monoid-minimum23-classification}.

\begin{lemma}[Proved: explicit coordinates on the minimal ideal]
\label{A:minimal-ideal-coordinates}
Let $K$ be the minimal ideal of a finite monoid $M$ and choose
an idempotent $e\in K$. Then $G=eMe$ is a group, and there
are finite sets $I,J$ identifying $K$ with $I\times G\times J$,
with multiplication
$(i,g,j)(i',g',j')=(i,gp_{ji'}g',j')$ for $p_{ji'}\in G$.
If $G=1$ this is a rectangular band. If $K=M$, then $M$ is a group.
\end{lemma}
\begin{proof}
The group assertion for $eMe$ is
Lemma~\ref{CP2:C-minimal-corner}; we give the coordinates explicitly.
Put $I=\{a\in Me:ea=e\}$ and
$J=\{b\in eM:be=e\}$. For $x\in Me$, the element $ex$
lies in $G$, and $x=a g$ uniquely with $g=ex$ and
$a=xg^{-1}\in I$. Dually $y\in eM$ is uniquely $h b$,
where $h=ye\in G$ and $b=h^{-1}y\in J$.
Minimality gives $K=MeM$, so every $z\in K$ can be written
$a g b$. Conversely, such a factorization recovers
$g=eze$, $a=ze\,g^{-1}$ and $b=g^{-1}ez$, proving uniqueness.
The product is
$a g b a'g'b'=a g(ba')g'b'$, and $ba'\in eMe=G$.
This proves the formula with $p_{ba'}=ba'$.
If $1\in MeM$, a factorization $1=men$ in a finite monoid
forces $e$ to be a unit: a one-sided inverse in a finite
monoid is two-sided, and a product is a unit only if each
factor is a unit. An idempotent unit is $1$, so $e=1$ and
$M=eMe=G$.
\end{proof}

\begin{lemma}[Proved: two augmentation shifts]
\label{A:band-Tor}
Let $M$ be a finite monoid with proper rectangular minimal ideal
$K=I\times J$, and let $k$ be $\Z$ or a field. Set
$R=kM$, $A=R/kK$, $V=\ker(kI\to k)$ as a left $R$-module,
and $W=\ker(kJ\to k)$ as a right $R$-module. There are natural
isomorphisms
\[
 H_{n+2}(BM;k)\cong\operatorname{Tor}^{A}_{n}(W,V)
 \quad(n\geq0).
\]
In particular $H_2(BM;k)$ is the quotient of $W\otimes_kV$ by
$wm\otimes v-w\otimes mv$ for $m\in M$.
\end{lemma}
\begin{proof}
Choose $e=(i_0,j_0)\in K$. Multiplication identifies the $R$-bimodule
$kK$ with $Re\otimes_k eR$: on the standard bases this is the
bijection $(i,j)\mapsto(i,j)$. Consequently $kK$ is projective on
both sides. Also $K$ annihilates $V$ and $W$, and
$kK\otimes_R V=Re\otimes_k eV=0$.
The right-projective resolution $0\to kK\to R\to A\to0$
therefore gives $\operatorname{Tor}^R_q(A,V)=0$ for $q>0$ and
$A\otimes_R V=V$. Filtering the two-sided double bar complex by
the $A$ resolution yields
\[
 E^2_{p,q}=\operatorname{Tor}^A_p
   (W,\operatorname{Tor}^R_q(A,V))
 \Longrightarrow\operatorname{Tor}^R_{p+q}(W,V).
\]
Only $q=0$ remains. Finally use the exact augmentation sequences
$0\to V\to Re\to k\to0$ and
$0\to W\to eR\to k\to0$.
Their middle terms are projective, and $eR\otimes_RV=eV=0$;
the long exact Tor sequences give the stated two shifts, including
$H_2=W\otimes_AV$. All constructions commute with homomorphisms
which preserve the identified band and its actions.
\end{proof}

\begin{lemma}[Proved: preservation of $H_2$ and surjection on $H_3$]
\label{A:band-quotient-H3}
Suppose $q:M\twoheadrightarrow N$ is a surjective monoid
homomorphism whose restriction identifies their rectangular minimal
bands bijectively. Then $q$ induces an isomorphism on $H_2(-;k)$
and a surjection on $H_3(-;k)$, for $k=\Z$ and for every field.
\end{lemma}
\begin{proof}
The identified row and column sets have the same actions: every
element of $N$ has a preimage, and multiplication on the band
determines the action. Their balanced presentations of $H_2$ are
therefore identical. For the second assertion use the natural
two-sided bar complexes computing the Tor groups in
Lemma~\ref{A:band-Tor}. The chain map is the identity in degree
zero, whose group is $W\otimes_kV$, and is surjective in degree
one because $kM/kK\to kN/kK$ is surjective. If $z$ is a target
one-cycle and $y$ is any chain lifting it, then $dy=0$ because
its image is $dz=0$ and the degree-zero map is injective. Thus
every homology class in degree one is represented by the image
of a source cycle. Naturality of the two shifts proves the claim.
This argument does not claim an isomorphism in higher degrees.
\end{proof}

\begin{lemma}[Proved: rational minimal-group collapse]
\label{A:minimal-group-collapse}
For a finite monoid whose minimal ideal has Rees coordinates
$I\times H\times J$, collapse the $H$ coordinate inside the
ideal and leave every other element distinct. The resulting
quotient $q:M\to\overline M$ is a monoid homomorphism, its
minimal ideal is $I\times J$, and $Bq$ is a rational homology
equivalence and hence a rational cohomology-ring isomorphism.
\end{lemma}
\begin{proof}
Left multiplication changes $(i,h,j)$ to
$(\sigma(i),a_i h,j)$ and right multiplication changes it to
$(i,hb_j,\tau(j))$. These descriptions follow by commuting with
all translations from the opposite side. They show directly
that forgetting $h$ is a congruence, including a nontrivial
Rees sandwich matrix.

Put $R=\mathbb QM$ and choose an idempotent $e$ in the minimal
ideal. The multiplication map identifies
$\mathbb QK$ with $Re\otimes_{\mathbb QH}eR$.
The module $Re$ is free over $\mathbb QH$ on the right, so
$\mathbb QK$ is projective as a right $R$-module. Averaging
the group coordinate in each $H$ class defines an $R$-bimodule
idempotent on $\mathbb QK$, since left and right translations
permute those coordinates. Its kernel is exactly
$J_0=\ker(R\to\mathbb Q\overline M)$, which is consequently
right projective. On $\mathbb QK\otimes_R\mathbb Q$ the
averaging operator is the identity; hence
$J_0\otimes_R\mathbb Q=0$. The resolution
$0\to J_0\to R\to\mathbb Q\overline M\to0$ now gives
vanishing higher Tor with the trivial module. The change of
rings argument used in Lemma~\ref{A:band-Tor} identifies the
homology through the specified map $q$. The induced map on
cohomology respects cups, proving the ring statement.
\end{proof}

\begin{lemma}[Proved: a rank-one balanced form has square zero]
\label{A:rank-one-square}
For a finite monoid with rectangular minimal band, every rational
degree-two class represented by a balanced matrix of rank at most
one has zero square. No finiteness assumption on $BM$ is needed.
\end{lemma}
\begin{proof}
Write the nonzero matrix as $L=b\otimes a$, with $a_0=b_0=0$.
Balance means $B_m^tL=LA_m$, where $A_m,B_m$ are the induced
augmentation operators. Equality of these two rank-one tensors
gives a common scalar character $\chi_m$ and affine laws
\[
 a_{\sigma_m(i)}=a_m+\chi_m a_i,\qquad
 b_{\tau_m(j)}=b_m+\chi_m b_j.
\]
Here $a_m=a_{\sigma_m(0)}$ and $b_m=b_{\tau_m(0)}$.
The composition rules are
$a_{mn}=a_m+\chi_m a_n$ and
$b_{mn}=b_n+\chi_n b_m$.
A nonzero slope permutes a finite nonconstant rational value set,
so it is $1$ or $-1$. A slope-one translation is zero, and all
slope-minus-one maps have a common center on each value set.
Choose those centers $A,B$, or zero if no reflection occurs, and set
\[
 a'_m=a_m+A(\chi_m-1),\qquad
 b'_m=b_m+B(\chi_m-1).
\]
Then $\chi_m a'_m=\chi_m b'_m=0$ for every $m$.
The affine two-cocycles $u(m,n)=b_m a_n$ and
$u'(m,n)=b'_m a'_n$ obey
\[
 u'-u=-\delta\bigl(Ab+Ba+AB(\chi-1)\bigr).
\]
For the normalized three-cochain
$w(m,n,p)=b'_m a'_n b'_n a'_p$, substitution in the five terms
of the coboundary gives
\[
 (\delta w)(m,n,p,q)
 =u'(m,n)u'(p,q)
   +b'_m a'_q\chi_n b'_n\chi_p a'_p
 =u'(m,n)u'(p,q).
\]
Thus $[u]^2=0$.

For completeness, balanced forms do give all rational degree-two
classes. By Lemma~\ref{A:band-Tor}, their vector space is the
dual of the balanced quotient of $W\otimes V$. The formula
$u_L(m,n)=L_{\tau_m(0),\sigma_n(0)}$, with zero row and
column at index zero, satisfies the cocycle equation by balance.
The band cycles
\[
 -[x_{ij}|x_{ij}]-[x_{00}|x_{00}]
 +[x_{i0}|x_{i0}]+[x_{ij}|x_{00}]
\]
pair with $u_L$ as $-L_{ji}$. These cycles are the images of
the elementary tensors under the two augmentation connecting
maps. The formula therefore realizes exactly the dual
identification, rather than just a subspace of cohomology.
\end{proof}

\begin{corollary}[Proved: a census-free exclusion through order nine]
\label{A:monoids-nine}
No finite monoid of order at most nine has classifying space
$\CP^2$. More generally, a minimal ideal with at most two rows
or at most two columns forces every rational degree-two square
to vanish, independently of the maximal group and of finite
homotopy type.
\end{corollary}
\begin{proof}
Apply Lemma~\ref{A:minimal-group-collapse}. With one row or
one column, one augmentation module is zero, so $H_2$ vanishes.
With two rows or two columns every balanced matrix has rank
at most one, and Lemma~\ref{A:rank-one-square} applies.
If both row and column counts are at least three, the minimal
ideal has at least nine elements. The identity lies outside a
proper minimal ideal, giving at least ten elements in the
monoid. If the minimal ideal contains the identity, the monoid
is a group; its positive rational homology vanishes. These
cases exhaust all monoids of order at most nine.
\end{proof}

\paragraph{Action hulls and the exact finite scope.}
For $r$ rows and $s$ columns, put $d=(r-1)(s-1)$. An actor is
a pair of transformations $(\sigma,\tau)$; its multiplication
is $(\sigma,\tau)(\alpha,\beta)=(\sigma\alpha,\beta\tau)$.
Let $R_{\sigma,\tau}\subset\mathbb Q^d$ be the space of
coefficient relations in $B_\tau^tL=LA_\sigma$.
For a proper sum $R$ of such spaces define
\[
 H(R)=\{(\sigma,\tau):R_{\sigma,\tau}\subset R\}.
\]
Balance is closed under multiplication, so this is a monoid.
If $R$ is generated by a monoid's actors, the hull contains all
of them and has precisely $R$ as its relation space. The action
homomorphism into $H(R)$ preserves the band and induces an
isomorphism on rational $H_2$, and a surjection on integral
$H_2$ whenever the band is unchanged. Thus a hypothetical
$\CP^2$ source has $\dim R=d-1$ and its degree-two generator
is pulled back rationally from that hull.

The new program \path{monoid/code/relation_lattice.py}
forms every actor relation space by rational row reduction and
closes proper sums modulo $S_r\times S_s$. Every new orbit is
marked in full. It suffices to join individual actor spaces to
each orbit representative, since the actor set is invariant
under the permutation group. A hyperplane has no larger proper
join, so those joins are settled by dimension, without a
numerical rank estimate. The independent checker
\path{monoid/code/relation_lattice_check.py} reconstructs the
relations as double differences of matrix entries and computes
permutation actions by expanding augmentation basis tensors.
It verifies every orbit, every proper join and every hull's
actor list and spanning equality. The resulting counts are
\[
\begin{array}{c|r|r|r|r}
 (r,s)&r^rs^s&\text{proper spaces}&S_r\times S_s\text{ orbits}
       &\dim H^2=1\text{ orbits}\\\hline
 (3,3)&729&111&13&6\\
 (3,4)&6912&2767&53&21\\
 (3,5)&84375&68719&214&71
\end{array}
\]
The $3\times3$ count is thirteen because opposition has not
been divided out; the inherited eleven-orbit count used it.

\begin{proposition}[Proved: the $3\times3$, $3\times4$, and
 $3\times5$ gates for $\CP^2$ through order twenty-three]
\label{A:three-row-classification}
No finite monoid of type $\CP^2$ has a minimal ideal with
three rows and three columns, with any maximal group. None has
a $3\times4$ rectangular minimal band. Among monoids of order
at most twenty-three with $3\times5$ rectangular minimal band,
the only $\CP^2$ model, up to isomorphism, is the monoid of
Theorem~\ref{A:monoid-upper23}. All three statements also hold
after opposition.
\end{proposition}
\begin{proof}
Only the one-dimensional balanced-form spaces in the preceding
table can occur. Their exact partition in our orbit numbering is
\[
\begin{array}{c|r|r|r|l}
 (r,s)&\operatorname{rank}L=1&\text{duplicate values}
 &\text{new hulls}&\text{new hulls with }H_4(\mathbb Q)\ne0\\\hline
 (3,3)&3&0&3&11\ (|H|=30)\\
 (3,4)&6&4&11&38\ (25),\quad40\ (23)\\
 (3,5)&8&37&26&164\ (34),\quad198\ (30),\quad200\ (30)
\end{array}
\]
The first column is covered by Lemma~\ref{A:rank-one-square}.
For duplicate values, identify column points on which every
balanced form agrees; balance preserves this equivalence
relation under all column actions. The action map descends to
the smaller set, and the displayed affine cocycles factor
through it. The augmentation quotient is surjective. Hence
the smaller action image has the same nonzero one-dimensional
rational balanced quotient and is covered inductively by the
smaller-shape argument. A repeated row is handled in the same
way, or by opposition. All source finiteness assumptions below
are imposed on the original monoid, not on this intermediate
image.

For each new hull the program
\path{monoid/code/hull_prefixes.py} constructs an integral
projective resolution prefix through degree five. Its projectives
are principal summands generated by explicitly checked idempotents
or differences of nested idempotents. The certificate stores their
bases, the entire action, differentials, and tensor differentials.
The separate checker
\path{monoid/code/projective_prefix_check.py} reconstructs the
table and every projective, checks all equivariance equations,
zero composites, rational ranks and integral Smith invariants.
Saturated images and the complementary kernel ranks prove exactness.
Thus its calculation of $H_4$ depends on an exact resolution,
not on a guessed truncation. Every new hull omitted from the
last column has verified zero rational fourth homology, which
kills a pulled-back square.

We explain the cyclic-character exclusions used for the few
remaining hulls. For an integral affine cocycle $u$ and an actor
$a$, let $G_a$ be its eventual cyclic group, $e$ its identity,
and $g=ea$. Its restricted degree-two class is
\[
 c_a=\sum_{h\in G_a}u(h,g)\pmod{|G_a|}.
\]
Indeed $u(e,h)=u(h,e)=u(e,e)$; subtracting this constant
normalizes the cocycle and changes the displayed sum by a
multiple of $|G_a|$. This is the elementary cyclic cocycle
formula for $H^2(BG_a;\Z)$.
If an actor with $c_a\ne0$ were the image of an element of a
finite monoid $M$, its eventual group would surject onto $G_a$.
The nonzero character would remain nonzero under inflation,
since $H^2(BG;\Z)=\operatorname{Hom}(G,\mathbb Q/\Z)$ for
finite groups and precomposition with a surjection is injective.
Such a restriction cannot come from a line bundle on a finite
CW model of $BM$. The reduced line-bundle class is nilpotent
by Lemma~\ref{CP2:C-Knilpotence}. Choose a cyclic prime-power
subgroup on which the restricted character is nontrivial.
Atiyah's completion theorem \cite[Theorem~7.2]{CP2-Atiyah1961}
identifies its K-theory with the completed representation ring.
Evaluation at a primitive prime-power root of unity maps the
reduced character to a nonzero element $\zeta^a-1$ of the
corresponding $p$-adic cyclotomic integer domain. The
augmentation ideal maps into its maximal ideal, so evaluation
extends to the completion, and no power of this nonzero
element vanishes. This contradicts nilpotence. A local group is interpreted as
the endomorphisms of its idempotent in the Karoubi envelope,
so the classifying-space map is well typed even when $e\ne1$.

The independent program
\path{monoid/code/character_prune_check.py} recomputes all
these residues, grows allowed submonoids by generated closures,
and compares every maximal target with its independently
checked integral projective prefix. The results are
\[
\begin{array}{c|r|r|l}
 \text{hull}&\text{forbidden actors}&\text{maximal allowed targets}
 &\text{verified conclusion}\\\hline
 (3,3),11&8&9\text{ of order }16&H_4(\mathbb Q)=0\\
 (3,4),38&2&1\text{ of order }23&H_4(\mathbb Q)=0\\
 (3,5),198&4&1\text{ of order }26&H_4(\mathbb Q)=0\\
 (3,5),200&4&1\text{ of order }26&H_4(\mathbb Q)=0
\end{array}
\]
For the first line the producer tested all $4096$ subsets,
finding $154$ closed sets; the independent closure-growth
checker recovers exactly these sets and the nine maxima.
In the last three lines the entire allowed set is already
closed, which proves directly that it is the unique maximum.
The degree-two class factors through one of these targets,
so its square vanishes on the original source.

The exceptional $3\times4$ hull 40 has order twenty-three.
The newly checked resolution prefix has abelian ranks
$(3,12,23,22,15,12)$ in degrees zero through five and
coinvariant ranks $(1,0,1,0,1,0)$.
The exact integral bar comparison in
\path{monoid/certificates/hull_3x4_40_cup.json} gives
\[
 H_2=\Z,\quad H_3=0,\quad H_4=\Z,\qquad x^2=2y,
\]
where $x$ and $y$ are primitive. The checker verifies the
comparison on projective generators and all their translates,
the cocycle on every triple, and evaluations on integral cycle
bases. These conclusions need only exactness through degree
four; we make no assertion here about an unverified periodic
tail. A band source with $H_2=\Z$ maps surjectively on integral
$H_2$ to this hull, so its primitive two-class is pulled back
up to sign. Its square is then even, which excludes $\CP^2$.
The same argument applies after identifying duplicate values,
because the corresponding augmentation map is still surjective.

It remains to treat the $3\times5$ hull 164. Its matrix, with
zero row and column suppressed, is
\[
 L=\begin{pmatrix}1&1\\0&1\\-1&0\\1&0\end{pmatrix}.
\]
It has thirty-four actors; four forbidden characters leave a
closed thirty-element allowed hull. The complete lattice of
closed subsets containing the identity and all fifteen band
elements has $2186$ members. Exactly $1711$ have order at most
twenty-three; among these $1160$ have integral $H_2=\Z$.
The two newly authored programs
\path{monoid/code/cp23_hull.py} and
\path{monoid/code/cp23_hull_check.py} obtain the same scope
by subset enumeration and by generated closures, respectively.
The latter reconstructs the integer balanced presentation,
its Smith form and the three two-sided-bar boundary matrices
over $\mathbb F_5$ for every relevant image. The resulting
counts are
\[
\begin{array}{c|r|r|r|r}
 |N|&\text{all images}&H_2(N;\Z)=\Z
  &\dim H_4(N;\mathbb F_5)=0
  &(b_2,b_3,b_4)=(1,0,1)\\\hline
16&1&0&0&0\\17&12&0&0&0\\18&60&0&0&0\\
19&166&12&12&0\\20&309&140&140&0\\
21&416&314&314&0\\22&421&380&376&0\\
23&326&314&310&2
\end{array}
\]
The four exceptions of order twenty-two have profile $(1,1,2)$;
the other two exceptions of order twenty-three have profile
$(1,2,3)$. They are excluded for every band-preserving source
by the surjection on $H_3(-;\mathbb F_5)$ in
Lemma~\ref{A:band-quotient-H3}. All images with zero fourth
homology over $\mathbb F_5$ have zero rational fourth homology,
so their pulled-back degree-two squares vanish.

The remaining two images are indices 1412 and 1651 in the
certificate. Both have the complete integral $\CP^2$ certificate
of Theorem~\ref{A:monoid-upper23}, and the displayed isomorphism
there identifies them. A source of order at most twenty-three
surjecting onto either twenty-three-element image must be
isomorphic to it. This proves the asserted classification.
For the $3\times3$ assertion with a nontrivial maximal group,
use Lemma~\ref{A:minimal-group-collapse}; the rational square
obstructions persist and all forbidden-character restrictions
are applied to the original finite-type source. This does not
assume that the collapsed quotient has finite CW type.
\end{proof}

\paragraph{Relation to the completed global proof.}
The preceding proposition classifies every monoid of order at most
twenty-three with a $3\times5$ rectangular minimal band, including
nonfaithful action sources. The $4\times4$ action computation is
completed independently in Proposition~\ref{A:four-by-four-through23}.
The remaining large bands are excluded by the target-specific
algebra argument of Section~\ref{A:large-band-target-proof}; it
avoids the stronger inherited arbitrary-module contracted-algebra
claim and its binary-pencil classification. The global minimum and
all monoids attaining it are therefore established in
Theorem~\ref{A:monoid-minimum23-classification}.


\subsection{The remaining large minimal bands}
\label{A:large-band-target-proof}

The argument in this subsection is specific to the projective-plane
target. It avoids the stronger inherited assertion about arbitrary
modules over four-dimensional contracted monoid algebras. Throughout,
algebras and modules in this subsection are finite dimensional over
$\mathbb Q$, and all modules are unital.

\begin{lemma}[Proved: local-algebra facts]
\label{A:local-target-facts}
Let $A$ be local, with radical $J$ and residue division algebra $D$.
For nonzero right and left modules $W,V$, respectively, there is a
surjection
\[
 W\otimes_A V\twoheadrightarrow
 (W/WJ)\otimes_D(V/JV).
\]
Consequently, if $\dim_{\mathbb Q}(W\otimes_AV)=1$, then
$D=\mathbb Q$ and both $W,V$ are cyclic. Moreover, for a
nonprojective finite right module $W$,
$\operatorname{Tor}_i^A(W,D)\ne0$ for every $i\geq1$;
the analogous assertion holds with the sides interchanged.
\end{lemma}
\begin{proof}
Quotienting both tensor factors gives the displayed surjection.
If the tops have dimensions $a,b$ over $D$, its target has
rational dimension $ab\dim_{\mathbb Q}D$. Nonzero finite
modules have nonzero tops by nilpotence of $J$. A one-dimensional
tensor product therefore forces $a=b=1$ and $D=\mathbb Q$.
A lift of a top generator generates the whole module, again
by nilpotence of $J$.

Finite projectives over a local finite algebra are free: choose
a basis of their tops, lift it to a surjection from a free
module, and split the surjection; the complementary projective
has zero top and hence is zero. Construct a minimal free
resolution by taking top bases at each stage. Its differentials
have entries in $J$. Such a resolution cannot end with a
nonzero free module in positive degree. Indeed a nonzero
element in the last nonzero radical power is annihilated by
$J$ on both sides and supplies a nonzero vector in the kernel
of that last differential. Tensoring the nonterminating minimal
resolution with $D$ makes every differential zero, proving
nonvanishing in every positive degree.
\end{proof}

\begin{lemma}[Proved: nonlocal algebras through dimension four]
\label{A:nonlocal-four}
If $A$ is nonlocal and $\dim_{\mathbb Q}A\leq4$, then
$\operatorname{Tor}_2^A(W,V)\ne0$ implies nonzero Tor in
arbitrarily high degrees.
\end{lemma}
\begin{proof}
Write $S=A/J$. Semisimple algebras have zero positive Tor.
We record the elementary small-dimension analysis, including
the nonsplit possibilities. A nonzero radical bimodule
incident with a field factor of degree $d$ has dimension at
least $d$, since it is a unital vector space over that field.
A matrix factor of dimension four similarly cannot act on
a one-dimensional radical. Thus a nonsplit nonlocal algebra
of dimension at most four with nonzero radical can only have
a quadratic field factor, a rational factor and a
one-dimensional radical supported on the rational factor.
It is the product of that field and the dual numbers.
This has the asserted Tor behavior.

In the remaining nonsemisimple cases $S=\mathbb Q^r$,
$r\geq2$. Lift the primitive orthogonal idempotents. Such
lifts can be obtained successively modulo the radical powers;
for a square-zero error $\delta=p^2-p$, the correction
$p+(1-2p)\delta$ is idempotent. Here $\dim J\leq2$.
If $J^2\ne0$, then $r=2$ and $J/J^2,J^2$ are
one dimensional. A generator of $J/J^2$ must be a loop:
a single off-diagonal component cannot multiply with itself.
Lifting that loop gives $x^2\ne0$, $x^3=0$, and the algebra
is $\mathbb Q[x]/(x^3)\times\mathbb Q$.

For $\mathbb Q[x]/(x^m)$, $m\geq2$, Jordan decomposition
expresses every finite module as a sum of $A/(x^a)$.
For $a<m$ its free resolution alternates multiplication by
$x^a$ and $x^{m-a}$. Tensoring with $A/(x^b)$, $b<m$,
has positive-degree homology dimension
\[
 \min(a,b)+\min(m-a,b)-b
 =\min\{a,b,m-a,m-b\}>0.
\]
Free summands contribute no positive Tor. This proves the
claim for the truncated-polynomial products above.

It remains that $J^2=0$. The idempotent decomposition makes
$A$ the radical-square-zero algebra of a directed quiver,
with $r$ vertices and $d=\dim J$ arrows, where $r+d\leq4$.
First syzygies $\Omega W,\Omega V$ are semisimple because
they lie in radicals of projective covers. Dimension shifting
on both sides gives
\[
 \operatorname{Tor}_n^A(W,V)
   \cong\operatorname{Tor}_{n-2}^A(\Omega W,\Omega V)
 \quad(n\geq2).
\]
For $n=2$ this follows from the long exact sequence and the
zero map $\Omega W\otimes_A\Omega V\to
\Omega W\otimes_A P_V$: the second syzygy factor maps into
$JP_V$, while $\Omega WJ=0$.
Nonzero degree-two Tor therefore means that a vertex simple
occurs in both syzygies. That vertex has an incoming and an
outgoing arrow. With at most two arrows and $r+d\leq4$,
it lies on a loop or a two-cycle; a line through three
distinct vertices would require at least five dimensions.
The minimal resolution of a vertex simple in a
radical-square-zero algebra has basis the directed paths
starting at that vertex. Repeating this cycle gives positive
Tor of the two matching simples in unbounded degrees.
Direct sums and the displayed shift finish the proof.
\end{proof}

\begin{lemma}[Proved: a small local summand cannot supply the target]
\label{A:local-four-small-target}
Let $A$ be local of dimension four. There are no nonzero modules
$\dim V\leq2$, $\dim W\leq5$ satisfying
\[
 \bigl(\dim\operatorname{Tor}_0,
       \dim\operatorname{Tor}_1,
       \dim\operatorname{Tor}_2,
       \dim\operatorname{Tor}_3\bigr)=(1,0,1,0).
\]
\end{lemma}
\begin{proof}
Lemma~\ref{A:local-target-facts} makes the residue field
$\mathbb Q$ and both modules cyclic. Hence $\dim W\leq4$.
If either module has dimension one, its positive Tor with a
nonprojective opposite module never vanishes. If $\dim W=4$,
the cyclic cover $A\to W$ is an isomorphism, so positive Tor
vanishes. If $\dim W=3$, its first syzygy is the
one-dimensional simple module, and
$\operatorname{Tor}_3(W,V)=\operatorname{Tor}_2(\mathbb Q,V)
\ne0$ because $V$ is nonprojective. Only
$\dim W=\dim V=2$ remains.

Put $I=\Omega V$. The cyclic cover has kernel of dimension
two. Since $\operatorname{Tor}_1(W,V)=0$,
$W\otimes I$ embeds in $W$, with image the kernel of
$W\to W\otimes V$; its dimension is one. The quotient
onto $\operatorname{top}W\otimes\operatorname{top}I$
therefore makes $I$ cyclic. Put $L=\Omega^2V$, also of
dimension two. If $L$ were not cyclic it would be the sum
of two simples, giving
$\operatorname{Tor}_3(W,V)=\operatorname{Tor}_1(W,L)\ne0$.
Thus $L$ is cyclic. Its syzygy $I'=\Omega^3V$ has dimension
two. The equality $\operatorname{Tor}_1(W,L)=0$ embeds
$W\otimes I'$ into $WJ$, of dimension one, so $I'$ is
cyclic as well.

The first four projectives in the minimal resolution of $V$
are consequently $A$. After tensoring with $W$, each has
dimension two. Since $\operatorname{Tor}_0$ has dimension
one, the first differential has rank one. Vanishing of
$\operatorname{Tor}_1$ makes the second differential rank
one. Dimension one in $\operatorname{Tor}_2$ makes the
third differential zero. But the fourth differential has
image in $WJ$, of dimension one, so it cannot cover the
two-dimensional kernel of the third. This contradicts
$\operatorname{Tor}_3=0$.
\end{proof}

\begin{lemma}[Proved: the nonlocal five-dimensional alternatives needed here]
\label{A:nonlocal-five-target}
Let $A$ be nonlocal of dimension five. Either nonzero
$\operatorname{Tor}_2(W,V)$ persists in arbitrarily high
degrees for all finite module pairs, or $A$ is one of the
following types:
\begin{enumerate}
\item $B\times\mathbb Q$ with $B$ a four-dimensional local algebra;
\item the radical-square-zero algebra of the directed line on
three vertices;
\item the two-vertex algebra with basis $e_1,e_2,a,b,c$,
$a\in e_1Ae_2$, $b\in e_2Ae_1$, $ab=c\ne0$, $ba=0$,
and $cJ=Jc=0$.
\end{enumerate}
\end{lemma}
\begin{proof}
We again use $S=A/J$. Nonsplit semisimple factors require
their degree in every incident radical component. The only
new nonsemisimple possibility relative to
Lemma~\ref{A:nonlocal-four} is a quadratic field and a
rational factor with a two-dimensional radical. A component
joining the two factors is a single triangular arrow and
gives a hereditary algebra: its bimodule is projective on
both sides over semisimple fields, and the standard
one-arrow projective resolution has length one.
A radical component on the quadratic factor gives a local
square-zero algebra over that division residue algebra;
its first syzygies are semisimple and
Lemma~\ref{A:local-target-facts} gives Tor persistence.
A radical supported on the rational factor gives an algebra
of dimension at most three, already covered by the
square-zero or truncated-polynomial arguments. Larger field
factors and matrix factors leave at most one radical
dimension and give no additional case. Separable field
factors split across the nilpotent ideal by lifting a
primitive element and correcting its minimal polynomial;
its derivative is invertible, so the usual successive
Newton correction applies.

Now assume $S=\mathbb Q^r$, $r\geq2$, and lift the
orthogonal idempotents. If $r\geq4$, at most one radical
arrow remains and the preceding argument applies. If
$r=3$ and $J^2\ne0$, a sole loop generates the radical,
giving $\mathbb Q[x]/(x^3)\times\mathbb Q^2$.
If $r=3$ and $J^2=0$, there are two arrows. The
semisimple-syzygy proof of Lemma~\ref{A:nonlocal-four}
applies unless a vertex has an entering and a leaving
arrow but lies on no cycle. The only exception is the
directed line on three vertices, giving alternative (2).

For $r=2$ we have $\dim J=3$. If $J^2=0$, every vertex
with both entering and leaving arrows lies on a loop or
a two-cycle, so the same proof gives persistence.
If $\dim J^2=2$, the single generator of $J/J^2$ is a
loop and gives $\mathbb Q[x]/(x^4)\times\mathbb Q$,
which is covered by (1). Suppose $\dim J^2=1$; then
$J^3=0$ and there are two arrows modulo $J^2$.
Two parallel cross arrows cannot generate a nonzero
square. A loop at each vertex gives a product of
$\mathbb Q[x]/(x^3)$ and $\mathbb Q[y]/(y^2)$.
Two opposite cross arrows give (3), because their
composites lie at distinct vertices and exactly one
can span the one-dimensional square. Two loops at
one vertex give (1).

The remaining configuration, up to opposition, consists
of $x\in e_1Je_1$ and $a\in e_1Je_2$.
The potentially nonzero products $x^2$ and $xa$
lie in different idempotent components, so precisely
one spans $J^2$.
If $x^2=c$ and $xa=0$, write $P_i=e_iA$ and
$T_i$ for right simples. Then
\[
 \operatorname{rad}P_1=L\oplus P_2,\quad
 L=\langle x,c\rangle,\quad P_2=T_2,\qquad
 \Omega L=T_1\oplus P_2,\quad
 \Omega T_1=L\oplus P_2.
\]
Every first right syzygy is a sum of $L,T_1,P_2$:
the cross arrow acts trivially there, and on the
first-vertex component this is Jordan decomposition
for a square-zero operator. Its nonprojective
summands have positive Tor periodic with period two.
If $xa=c$ and $x^2=0$, the same decomposition of
$\operatorname{rad}P_1$ holds, but now $\Omega L=L$.
On $L$ the cross arrow is an isomorphism between
its two one-dimensional vertex spaces. Every submodule
of a sum of $L$ and $P_2$ has injective cross-arrow
map and splits as a sum of those two modules.
Thus first syzygies have positive Tor periodic with
period one. In both cases a nonzero degree-two Tor
group recurs indefinitely. Opposition interchanges
the two tensor factors. These cases exhaust the
two-arrow configurations and prove the alternatives.
\end{proof}

\begin{proposition}[Proved: the $3\times6$ band with five outside elements]
\label{A:three-by-six-CP2}
A finite monoid with a $3\times6$ rectangular minimal band
and exactly five outside elements cannot have classifying
space $\CP^2$. The same holds for its opposite.
\end{proposition}
\begin{proof}
Let $A=\mathbb QM/\mathbb QK$. Its dimension is five;
the row and column augmentation modules have dimensions
$\dim V=2$, $\dim W=5$. The band Tor formula would give
\[
 \dim\operatorname{Tor}_n^A(W,V)=
 \begin{cases}1,&n=0,2,\\0,&n=1\text{ or }n\geq3.\end{cases}
\]
If $A$ is local, Lemma~\ref{A:local-target-facts}
makes $W$ cyclic. Its cover $A\to W$ has equal
source and target dimensions, hence is an isomorphism,
contradicting positive Tor. In the nonlocal case
Lemma~\ref{A:nonlocal-five-target} leaves its three
exceptional alternatives. For $B\times\mathbb Q$,
positive Tor comes entirely from the $B$ summands;
their tensor product is nonzero and hence has dimension
one, since the total tensor dimension is one.
Lemma~\ref{A:local-four-small-target} applies to
those summands and excludes this case.

For the line algebra write $a\in e_1Ae_2$,
$b\in e_2Ae_3$, $ab=0$.
Splitting the image of $b$ inside the kernel of $a$
decomposes every left module into the three simples
$S_1,S_2,S_3$ and the two arrow modules $Ae_2,Ae_3$;
dually the right indecomposables are
$T_1,T_2,T_3,e_1A,e_2A$.
The two nontrivial right resolutions are
\[
 0\to e_3A\xrightarrow{b}e_2A\xrightarrow{a}e_1A
   \to T_1\to0,\qquad
 0\to e_3A\xrightarrow{b}e_2A\to T_2\to0.
\]
Thus degree-two Tor one requires one $T_1$ in $W$
and one $S_3$ in $V$. The remaining dimension of
$V$ must be $S_1$: $S_2$ would give nonzero degree-one
Tor, and a second $S_3$ would give degree-two Tor
dimension at least two. Degree-one vanishing excludes
$T_2$ from $W$. Hence
\[
 V=S_1\oplus S_3,\qquad
 W=T_1\oplus T_3^{\oplus s}\oplus(e_1A)^{\oplus p}
                    \oplus(e_2A)^{\oplus q},\quad
 s+2p+2q=4.
\]
The tensor dimension is $1+s+p$, so it equals one
only for $(s,p,q)=(0,0,2)$.
Every balanced functional for this pair is supported
on the single $T_1\times S_1$ coordinate and has
matrix rank one. Rank is invariant under changes
of module bases; therefore the actual band balanced
form also has rank one. Lemma~\ref{A:rank-one-square}
contradicts the required nonzero square.

For the two-vertex algebra, the left indecomposables
are $S_1,S_2,L=Ae_1/\mathbb Qc,Ae_1,Ae_2$, of
dimensions $1,1,2,3,2$, and the right ones are
$T_1,T_2,R=e_1A/\mathbb Qc,e_1A,e_2A$.
To see completeness, split off the length-three
chains representing a basis of $\operatorname{im}(ab)$.
On the complementary representation both composites
vanish; splitting the two arrow images inside the
opposite kernels leaves length-two chains and simples.
The relevant right resolutions are
\[
\begin{split}
 &0\to e_2A\xrightarrow{a}e_1A\to T_1\to0,\\
 &0\to e_2A\xrightarrow{a}e_1A\xrightarrow{b}e_2A
       \to T_2\to0,\\
 &0\to e_2A\xrightarrow{a}e_1A\xrightarrow{c}e_1A
       \to R\to0.
\end{split}
\]
Evaluating them on the five possible two-dimensional
left modules shows that only $V=L$ can have degree-one
Tor zero and degree-two Tor one: two copies of a simple
give an even degree-two dimension, $S_1\oplus S_2$
forces nonzero degree-one Tor, and $Ae_2$ gives zero
degree-two Tor. With $V=L$, one must have exactly
one $T_2$, and neither $T_1$ nor $R$ can occur in
$W$. Thus
$W=T_2\oplus(e_1A)^{\oplus p}\oplus(e_2A)^{\oplus q}$
with $3p+2q=4$, forcing $p=0,q=2$.
But $T_2\otimes L=0$ and each $e_2A\otimes L$
has dimension one, so the tensor dimension is two,
contradicting the target.

These small calculations are also regenerated from
literal multiplication tables and the displayed
resolutions by
\path{monoid/code/five_dimensional_pairs.py}.
Its complete module-pair output is
\path{monoid/certificates/five_dimensional_pairs.json}.
The proof above establishes their decomposition
scope and uses no list of monoid multiplicative bases.
\end{proof}

\begin{corollary}[Proved: the other large-band shapes through order twenty-three]
\label{A:large-bands-through23}
No finite monoid of order at most twenty-three with
rectangular minimal band of shape $3\times6$, $3\times7$
or $4\times5$ has type $\CP^2$. The same holds after
opposition.
\end{corollary}
\begin{proof}
Let $d=|M\setminus K|=\dim A$. For $3\times6$,
either $d=5$, treated by Proposition~\ref{A:three-by-six-CP2},
or $d\leq4$. For $3\times7$ we have $d\leq2$;
for $4\times5$, $d\leq3$. In each of the latter
cases $A$ has dimension at most four and one of
$V,W$ has dimension strictly greater than $d$.
If $A$ is local, tensor dimension one would make
both modules cyclic by Lemma~\ref{A:local-target-facts},
which is impossible for a module of dimension greater
than its algebra. If $A$ is nonlocal,
Lemma~\ref{A:nonlocal-four} contradicts the target's
vanishing Tor in all degrees above two.
\end{proof}


\subsection{The four-by-four gate and the complete monoid answer}
\label{A:monoid-global-answer}

\begin{proposition}[Proved: every $4\times4$ action image through order twenty-three is excluded]
\label{A:four-by-four-through23}
If a finite monoid has a $4\times4$ rectangular minimal band
and classifying space $\CP^2$, its image in the monoid of
row and column transformations has at least twenty-four elements.
In particular the source monoid has at least twenty-four elements.
\end{proposition}
\begin{proof}
The complete relation-space method of
Section~\ref{A:monoid-lower-audit} now has nine coordinates.
Our new constructor enumerates all $65536$ transformation
pairs. Their distinct proper individual relation spaces have
counts
\[
\begin{array}{c|rrrrrrr}
 \text{dimension}&0&3&4&5&6&7&8\\\hline
 \text{count}&1&12&1269&504&7732&1590&205.
\end{array}
\]
Their proper sums give $169323$ spaces, in $577$ orbits under
$S_4\times S_4$. The independent dual-form checker verifies
all $4841964$ explicit representative joins; the remaining
$1685637$ joins are settled because one operand is a
hyperplane and no strictly larger proper space exists.
It also checks every orbit, every actor list and the equality
between the specified relation space and the span of its
hull's actors. Its completed result is
\path{monoid/certificates/relation_lattice_4x4_verified.json}.
The earlier interrupted states are superseded by this complete
certificate; none is counted as a completed lattice.

Exactly $149$ orbits have one-dimensional balanced-form
space. Twelve have rank-one matrix and are excluded by
Lemma~\ref{A:rank-one-square}. Sixty-one have repeated row
or column values. Identifying equal values gives a map into
a smaller-shape action monoid, with the same rational
degree-two class and a surjection on its integral balanced
quotient. The relative smaller-shape proof in
Proposition~\ref{A:three-row-classification} applies:
zero products pull back; forbidden local characters are
tested on the original finite source; and an even-square
target pulls its primitive two-class back primitively.

For the remaining $76$ hulls we construct integral
projective resolution prefixes through degree five. All
$76$ pass the separate projective-prefix checker; combined
with the earlier shapes and maximal allowed targets the
batch has $129$ distinct certificates and no remaining
failure. In $66$ of these $76$ hulls the rational fourth
homology is zero. This count includes orbit 210, whose
binary or odd-primary homology alone need not reveal the
rational vanishing. The remaining ten are
\[
\begin{array}{c|r|l}
 \text{orbit}&\text{order}&\text{degree-two square gate}\\\hline
216&60&\text{bounded images below}\\
217&140&\text{bounded images below}\\
221&82&\text{bounded images below}\\
223&54&\text{zero on every integral four-cycle}\\
359&56&\text{content }2\\
448&29&\text{content }2\\
454&39&\text{content }2\\
455&29&\text{content }2\\
495&29&\text{content }2\\
499&29&\text{content }2.
\end{array}
\]
Here content means the greatest common divisor of the
evaluations on a saturated integral basis of degree-four
cycles in the coinvariant projective prefix. The new
program \path{monoid/code/attach_bar_comparison.py}
attaches a literal normalized-bar comparison to each of
these seven prefixes. The independent checker verifies
all idempotent projectives, differential equations,
integral exactness, comparison chain equations and basis
translates, then checks the cocycle on every triple.
It obtains $H_2=\Z$ and a primitive degree-two class
in each case. Saturated kernel bases are obtained with
unimodular Hermite transformations, and the displayed
contents are checked on the resulting full cycle lattices.
Thus orbit 223 has zero rational square. In every
content-two case, the image of any integral four-cycle
from a source has even evaluation. The integral
$H_2$ surjection pulls the target two-class back
primitively, so a $\CP^2$ source would contradict its
unit square. This leaves exactly the first three hulls.

Fix one of these hulls and write $H$ for the actual
action image of the source. The image contains the
identity and all sixteen constants. Suppose $|H|\leq23$;
there are at most six additional actors.
An actor with nonzero cyclic restriction of the integral
hull cocycle is forbidden by the line-bundle argument
in Proposition~\ref{A:three-row-classification}.
The complete literal input tables have respectively
$9,47,19$ forbidden actors in hulls $216,217,221$.
Their full groups of row and column permutations,
including opposition when it preserves the hull,
have orders $2,16,4$.

The new producer
\path{monoid/code/bounded4_enumerate.cpp} starts from
the identity and constants, adjoins one allowed actor,
and takes its multiplicative closure. A closure is
discarded only if it meets a forbidden actor or has
more than six additional actors. Every wanted image
is reached: adjoining its generators produces only
intermediate closures contained in that image.
Products with constants remain in the band, so the
closure computation need only track additional actors.
The independent input checker tests all $65536$ pairs
against the balanced matrix, reconstructs every product,
cyclic character and symmetry, and checks the exact
text input consumed by the two C++ programs.

The rational relation-rank test is exact. Each raw
relation row has squared Euclidean norm at most eight.
Hadamard's inequality bounds every nonzero minor of
order at most nine by $8^{9/2}<65521$. Therefore
reduction modulo the prime $65521$ preserves its
rational rank. Rank eight is precisely the required
one-dimensional rational $H_2$.
The complete counts through order twenty-three are
\[
\begin{array}{c|r|r|r|rrrr}
\text{hull}&\text{closed sets}&H_2\text{-rank-one sets}
 &\text{orbits}&20&21&22&23\\\hline
216&15606&10841&5426&68&630&1736&2992\\
217&128639&102602&6495&52&519&1789&4135\\
221&33274&24962&6297&76&659&1901&3661
\end{array}
\]
The last four columns count orbits by order; none has
order below twenty. All closed sets, including those
discarded by the rank test, are retained in the
compressed scope streams. The independent checker
\path{monoid/code/bounded4_check.cpp} checks each set's
closure, verifies that every allowed one-actor extension
and generated closure is present, and recomputes all
symmetry orbits. This proves the search scope directly
from its literal output, independently of the producer's
work queue.

For each representative, with $q=|H|-17$, the two-sided
bar complex in Tor degree $j$ has dimension $9q^j$.
The producer constructs its first three differentials
over $\mathbb F_2$, checks consecutive composites and
uses column elimination. The independent checker
reconstructs the matrices as rows, checks all
consecutive composites, uses row elimination, and
separately proves independence of every listed pivot
column witness. The results are
\[
\begin{array}{c|r|r|r}
\text{hull}&H_4(\mathbb F_2)=0
 &H_4\ne0,\ H_3\ne0&H_4\ne0,\ H_3=0\\\hline
216&5228&188&10\\
217&6187&291&17\\
221&5995&286&16.
\end{array}
\]
Every representative has binary $H_2$ dimension one.
Zero fourth binary homology implies zero rational
fourth homology and kills the pulled-back square.
Positive third binary homology also excludes a
possibly nonfaithful source, because the quotient
onto its action image preserves the band and
Lemma~\ref{A:band-quotient-H3} gives a surjection
on $H_3(-;\mathbb F_2)$.

The remaining $43$ representatives comprise $39$
of order twenty-two and $4$ of order twenty-three.
For every one,
\path{monoid/code/bounded4_residual_cups.py}
constructs an integral projective prefix and actual
bar comparison through degree four. Exactness through
degree three is sufficient: extend the projective
prefix to a resolution, then extend the comparison
into the acyclic bar resolution. The independent
checker finds content one in degree two and content
zero for the square on all degree-four cycles.
Thus every residual square vanishes rationally.
Their complete literal tables and compressed
comparisons, with all $43$ separate verification
records, are indexed by
\path{monoid/certificates/bounded4_residual_cups.json}.
No possible image of order at most twenty-three
remains, proving the proposition.
\end{proof}

\begin{theorem}[Proved: minimum monoid order and all minimal monoids]
\label{A:monoid-minimum23-classification}
The least order of a finite monoid $M$ with $BM\simeq\CP^2$
is exactly $23$. Up to isomorphism and anti-isomorphism,
there is exactly one such monoid of order $23$, namely
the aperiodic monoid of Theorem~\ref{A:monoid-upper23}.
Up to ordinary isomorphism there are exactly two classes,
represented by that monoid and its opposite. Each has
trivial automorphism group.
\end{theorem}
\begin{proof}
The upper bound and its integral homotopy-type certificate
are Theorem~\ref{A:monoid-upper23}.
Suppose $|M|\leq23$ and $BM\simeq\CP^2$.
Use the explicit coordinates of
Lemma~\ref{A:minimal-ideal-coordinates}.
A group has zero positive rational homology and cannot
occur. A minimal ideal with one row or column gives
zero rational $H_2$ after minimal-group collapse;
with two rows or columns every rational square is
zero by Lemma~\ref{A:rank-one-square}.
The $3\times3$ shape is excluded for every maximal
group by Proposition~\ref{A:three-row-classification}.

If the remaining minimal maximal group were nontrivial,
the ideal alone would have at least $3\cdot4\cdot2=24$
elements. Thus the minimal ideal is a rectangular
band. The identity lies outside it. Up to opposition,
the only shapes with at least three rows and columns
and at most twenty-two band elements are
\[
 3\times3,\quad3\times4,\quad3\times5,\quad
 3\times6,\quad3\times7,\quad4\times4,\quad4\times5.
\]
The first two are excluded by the three-row proposition;
the $4\times4$ shape is excluded by
Proposition~\ref{A:four-by-four-through23}; and the
three large-band shapes are excluded by
Corollary~\ref{A:large-bands-through23}.
Only $3\times5$ remains. Its complete classification
through order twenty-three in
Proposition~\ref{A:three-row-classification} gives
exactly the displayed twenty-three-element monoid.
This simultaneously proves the lower bound and
exhaustiveness at equality, including nonfaithful
action representations.

The two order-twenty-three witnesses in that search
are explicitly isomorphic. Exhausting all
$3!5!=720$ row and column permutations shows that
each has only the identity automorphism and that
there is exactly one isomorphism between the two
listed labelings. This exhausts monoid isomorphisms:
the minimal ideal is characteristic, an automorphism
of its rectangular band permutes rows and columns,
and the faithful outside actions determine every
remaining element. The opposite has minimal-band
shape $5\times3$, and cannot be isomorphic to the
$3\times5$ monoid, since the cardinalities of the
principal left ideals in the characteristic minimal
ideal differ. Nerve reversal gives a homeomorphism
$BM\cong B(M^{\mathrm{op}})$, so the opposite has
the same homotopy type. The literal opposite table
is included among the minimal-model certificates.
\end{proof}


\subsection{Mixed five-element corners with seven connecting arrows}

\begin{lemma}[Proved: the relative-cup obstruction]
\label{CP2:A-collage-relative-cup}
Let $X$ be the classifying space of a category with two objects,
corner monoids $P,N$, arrows $S$ from the $P$ object to the $N$
object, and no arrows in the other direction. If a class
$\alpha\in H^2(X;\mathbb Q)$ restricts to zero on both corner
spaces, then $\alpha^2=0$.
\end{lemma}
\begin{proof}
The normalized nerve realizes the homotopy pushout
$BP\leftarrow E\to BN$, where $E$ is the two-sided action bar
space of the commuting left $N$ and right $P$ actions on $S$.
Use its double mapping cylinder. It has an open cover $U,V$
which deformation retracts onto $BP,BN$, respectively. The
vanishing restrictions lift $\alpha$ to both $H^2(X,U;\mathbb Q)$
and $H^2(X,V;\mathbb Q)$. Their relative cup product lies in
$H^4(X,U\cup V;\mathbb Q)=0$ and maps to $\alpha^2$.
\end{proof}

\begin{lemma}[Proved: rational homology of order-five corners]
\label{CP2:A-order-five-corner}
Every five-element monoid other than
$P_5=(\{0,1\}\times\{0,1\})^1$ has vanishing positive-degree
rational homology. For $P_5$, rational homology is $\mathbb Q$
in degrees zero and two and zero in all other degrees.
\end{lemma}
\begin{proof}
Collapse the local group of the minimal ideal by
Lemma~\ref{A:minimal-group-collapse}. If both dimensions
of the resulting rectangular band exceed one, cardinality
five forces the $2\times2$ band and its identity, and the local
group was trivial. Otherwise one augmentation module in
Lemma~\ref{A:band-Tor} vanishes, so all homology in degrees
at least two vanishes. A finite monoid has a finite universal
group: its image in its universal group is a finite submonoid
of a group, hence a subgroup, and it already generates that
group. Thus rational first homology also vanishes.

For completeness, write $R=\mathbb QP_5$, $e=e_{00}$ and
$v=e_{10}-e_{00}$. The following is a projective resolution of
the trivial left module:
\[
0\longrightarrow (Re)^2\xrightarrow{d_2}R
\xrightarrow{d_1}Re\xrightarrow{\epsilon}\mathbb Q
\longrightarrow0,
\qquad d_1(r)=rv,
\qquad d_2(e_{i0},0)=e_{i0},\quad
 d_2(0,e_{i0})=e_{i1}.
\]
Indeed, $\ker\epsilon=\mathbb Qv$; $\ker d_1$ is exactly
$\mathbb Q\{e_{ij}\}$; and $d_2$ identifies its two summands
with the two band columns. Tensoring with the trivial right
module gives $\mathbb Q^2\xrightarrow{(1,1)}\mathbb Q
\xrightarrow{0}\mathbb Q$, which proves the assertion and
also shows $\operatorname{pd}_R\mathbb Q\le2$.
\end{proof}

\begin{lemma}[Proved: three analytic mixed-corner reductions]
\label{CP2:A-mixed-five-reductions}
Let $N$ be a finite monoid and let
$S$ be a nonempty finite left-$N$, right-$P_5$ biset. The associated
two-object category is not of type $\CP^2$ in any of the following
cases:
\begin{enumerate}
\item $BN$ is rationally acyclic and the right $P_5$ action has a global fixed point;
\item the minimal ideal of $N$ has one column;
\item $BN$ is rationally acyclic, the right $P_5$ action has no fixed point,
and its band core consists of one two-element orbit.
\end{enumerate}
\end{lemma}
\begin{proof}
A fixed point $s$ gives a natural transformation from the corner
inclusion $P_5\to\mathcal C$ to the constant functor at the other
object: its naturality equations are $sm=s$. Hence every positive
cohomology class restricts to zero on $BP_5$. Its restriction to
$BN$ is also zero by the rational acyclicity hypothesis. Apply
Lemma~\ref{CP2:A-collage-relative-cup}.

For the second case choose an idempotent $e$ in the minimal ideal
and let $G=eNe$. The normalized coordinates of
Lemma~\ref{A:minimal-ideal-coordinates} have a single column.
For every $n\in N$, right multiplication on any fixed row is
multiplication by a group element on its $G$ coordinate. Thus
$p=|G|^{-1}\sum_{g\in G}g$ satisfies $pn=p$ and
$\epsilon(p)=1$. The trivial right $\mathbb QN$ module is the
projective summand $p\mathbb QN$.
The rational chain complex of $E$ computes
\[
\mathbb Q\mathbin{\otimes^{\mathbf L}_{\mathbb QN}}
\mathbb Q[S]\mathbin{\otimes^{\mathbf L}_{\mathbb QP_5}}
\mathbb Q.
\]
Projectivity removes the first derived tensor, and the displayed
resolution in Lemma~\ref{CP2:A-order-five-corner} makes its
homology zero above degree two. A one-column minimal ideal makes $BN$
rationally acyclic by Lemmas~\ref{A:minimal-group-collapse} and
\ref{A:band-Tor}; thus both corners have zero rational
homology above degree two, so the mapping-cylinder exact sequence
gives $H_4(X;\mathbb Q)=H_3(E;\mathbb Q)=0$.

For the third case let the core points be $c_0,c_1$, labelled by
column. Every $s\in S$ satisfies $se_{ij}=c_j$. In particular
swapping the two band rows defines a category automorphism $F$
which fixes $N$ and every member of $S$. We compute its action,
including on the bridge, rather than assuming that a visible
graph gives an equivariant model.
Tensor the displayed left-$R$ resolution with $\mathbb Q[S]$.
Its degree-zero term is $\mathbb Qc_0$, its degree-one term is
$\mathbb Q[S]$, and $d_1=0$. Its degree-two term maps
isomorphically onto $\mathbb Q\{c_0,c_1\}$ in degree one.
The quotient by this acyclic subcomplex is the $\mathbb QN$
complex
\[
\mathbb Q[0]\ \oplus\ T[1],\qquad
T=\mathbb Q[S]/\mathbb Q\{c_0,c_1\},
\]
with zero differential. The first summand is the trivial
$N$ module, since every $N$ endomorphism of the right act fixes
the unique core orbit pointwise.

Write $\sigma(e_{ij})=e_{1-i,j}$. A $\sigma$-semilinear lift
of $F$ to the left resolution sends $e_{i0}$ in degree zero
to $e_{1-i,0}$, sends $r\in R$ in degree one to $-\sigma(r)$,
and sends $e_{i0}$ in each degree-two summand to
$-e_{1-i,0}$. Each displayed differential commutes with these
maps; the sign in degree one follows from $\sigma(v)=-v$.
After tensoring with $S$, the induced map on the quotient is
$+1$ on $\mathbb Q[0]$ and $-1$ on $T[1]$.
Applying a free right $\mathbb QN$ resolution preserves the
quasi-isomorphism just used. Consequently
\[
H_k(E;\mathbb Q)=H_k(N;\mathbb Q)\oplus
\operatorname{Tor}^{\mathbb QN}_{k-1}(\mathbb Q,T),
\]
where the second term is absent for $k=0$. Since the first
summand vanishes for $k>0$, $F_*$ is $-1$ on every positive
homology group of $E$, in particular on $H_3(E)$ and hence
on $H_4(X)$.

On $H_2(BP_5)$ the same semilinear chain map is $-1$.
If $X$ had type $\CP^2$, a generator
$\alpha\in H^2(X;\mathbb Q)$ either restricts to zero on
$BP_5$, when the relative-cup lemma gives $\alpha^2=0$, or
restricts nontrivially, when $F^*\alpha=-\alpha$.
In the latter case multiplicativity gives
$F^*(\alpha^2)=\alpha^2$, whereas the action on $H^4(X)$
is $-1$. Thus $\alpha^2=0$ over $\mathbb Q$ in both cases,
a contradiction.
\end{proof}

\begin{proposition}[Proved: complete mixed-corner gate]
\label{CP2:A-five-five-seven}
No category on two objects with five endomorphisms at each
object, seven arrows in one direction, and no arrows in the
other direction has classifying space of type $\CP^2$.
\end{proposition}
\begin{proof}
If neither corner is $P_5$, use the relative-cup lemma and
Lemma~\ref{CP2:A-order-five-corner}. If both are $P_5$, use
Proposition~\ref{CP2:A-P5-P5-seven}. In the remaining case,
opposition, which preserves the homotopy type of the nerve,
allows us to take $P_5$ as the source corner and $N$ as the
target; $P_5\cong P_5^{\mathrm{op}}$.

Use the complete right-act normal forms of
Lemma~\ref{CP2:A-P5-normal-forms}. Of the $383$ seven-state
forms, precisely $13$ have no global fixed state. The first
has one double core orbit and is excluded analytically above.
The remaining twelve have centralizer sizes, in the recorded
normal-form order,
\[
266,88,88,175,47,40,63,47,43,22,63,21.
\]
These centralizers consist of all endomorphisms of the fixed
right act, including noninvertible maps. The complete ordinary
order-five monoid catalog contains $228$ tables. Excluding
$P_5$ and the one-column minimal ideals leaves $42$ target
monoids: $33$ with minimal band $1\times2$, seven with
minimal band $1\times3$, one with minimal band $1\times4$,
and one with a $1\times2$ minimal ideal whose local group
has order two.

The new producer \path{monoid/code/mixed5_bisets.cpp} enumerates
all homomorphisms from each of these $42$ targets into each
of the twelve endomorphism monoids. It fixes the identity map,
assigns an undetermined image, and repeatedly enforces
$f_{ab}=f_a\circ f_b$ whenever both factors are assigned.
Every consistent complete assignment is therefore reached.
It takes state orbits under the full automorphism group of the
right act and checks their orbit-weight sum against the raw
assignment count. This gives $802964$ labelled assignments
and $383050$ state orbits across the $504$ scopes.

An independent scope checker
\path{monoid/code/mixed5_scope_check.cpp} constructs the
endomorphisms by a different procedure. An equivariant map
chooses a map of double core orbits; it must preserve the
column in each core, and each transient point can go to
exactly those states with the resulting ordered orbit pair.
This description is necessary and sufficient, so it constructs
every endomorphism without the producer's transformation
scan. The checker assigns the four target-monoid images in
fixed order, rejecting a partial assignment only when a fully
assigned multiplication equation fails. It matches every
resulting state orbit against the producer stream, checks its
stabilizer and all orbit weights, and recomputes every degree-one
and degree-two mod-two homology filter independently. Its $3274729$ search nodes yield
exactly the same $802964$ assignments and $383050$ orbits. The
independent filters exclude $36612$ cases in degree one and
$224921$ in degree two.

Of the state orbits, $121517$ pass the necessary tests
$H_1(-;\mathbb F_2)=0$ and $H_2(-;\mathbb F_2)=\mathbb F_2$.
For $121509$ of these the producer supplies a literal nonzero
two-cocycle and a three-cochain whose coboundary is its square.
The other eight have positive degree-three mod-two homology.
The separate root checker \path{code/check_p5_bisets.cpp}
imports neither action search nor cup producer and checks these
literal tables, cocycles, primitives and degree-three exclusions.
It checks $94175675$ associativity equations, $56383888$ cocycle
equations and $279956736$ square-primitive equations; all pass.
Thus no scope remains. Exact receipts and the complete compressed
action and cochain streams are retained under
\path{monoid/certificates/mixed5/}; only newly generated data
are included there.
\end{proof}


\subsection{Five complete triangular sixteen-arrow profiles}

\begin{lemma}[Proved: cancelling a trivial action side in degree two]
\label{CP2:A-trivial-biset-side-H2}
If one side of a nonempty finite biset acts trivially, collapsing
that corner to the trivial monoid defines an actual functor from
its two-object collage to the corresponding ordinary act cone.
It induces an isomorphism on rational second cohomology.
\end{lemma}
\begin{proof}
Suppose the right $P$ action is trivial and put
$Z=B(N\ltimes S)$. The bridge is $BP\times Z$. Since every
finite monoid has zero rational first homology, K\"unneth gives
\[
H^1(BP\times Z;\mathbb Q)=H^1(Z;\mathbb Q),\qquad
H^2(BP\times Z;\mathbb Q)=
H^2(BP;\mathbb Q)\otimes H^0(Z;\mathbb Q)
\oplus H^2(Z;\mathbb Q).
\]
The first-corner pullback is diagonal into the first summand
and is injective because $Z$ is nonempty. The other corner
pullback lies in the second summand. The Mayer--Vietoris
short exact sequence in degree two therefore has the same
end terms as for the ordinary left-$N$ act cone: $H^1(Z)$ on
the left and $\ker(H^2(BN)\to H^2(Z))$ on the right. The
collapse induces the identity on both, hence an isomorphism
in the middle. The left-trivial case is dual. The comparison
is an actual functor, so preserves cup products.
\end{proof}

\begin{proposition}[Proved: all five $P_5$-source profiles at total sixteen]
\label{CP2:A-P5-source-sixteen}
No two-object collage with source corner $P_5$ and profile
\[
(5,7,4),\quad(5,6,5),\quad(5,5,6),\quad(5,4,7),\quad(5,3,8)
\]
has classifying space of type $\CP^2$. Here the entries give
source corner order, target corner order and cross-hom size.
The assertion includes every target monoid of the indicated
order, with no assumption on its rational second homology.
\end{proposition}
\begin{proof}
The source is fixed and the target ranges over the complete
ordinary catalogs of orders three through seven. A target
whose minimal ideal has one column is excluded by
Lemma~\ref{CP2:A-mixed-five-reductions}. The target $P_5$
in the third profile is covered by the separately regenerated
double-$P_5$ six-state gate. All other target tables are retained
at this stage; their counts, ordered by target order three
through seven, are
\[
1,6,42,445,5852.
\]
In particular a target with zero $H_2$ but possible higher
rational homology is not omitted merely on that basis.

The right-$P_5$ acts are the complete core-orbit normal forms.
The new eight-state constructor
\path{monoid/code/p5_eight_normal_forms.py} obtains $1212$
classes. Its independent verifier extracts the core graphs
from the literal maps, checks that their canonical signatures
are distinct, and matches each $(a,b,t)$ stratum with the
Burnside coefficient formula of
Lemma~\ref{CP2:A-P5-normal-forms}. It checks $155136$ literal
action equations. Exactly $31$ forms have no fixed state.
For smaller state counts the complete independently checked
root catalogs are used.

There are four further complete-scope omissions. A trivial
right side reduces by Lemma~\ref{CP2:A-trivial-biset-side-H2}
to an ordinary left-act cone over a monoid of order at most
seven; Lemma~\ref{CP2:A-high-seven-act}, applied to its
opposite, makes its rational degree-two products zero.
If the set of right-fixed states is a singleton, every
commuting target map fixes its unique member. This common
fixed cross arrow makes both corner inclusions null-homotopic,
so Lemma~\ref{CP2:A-collage-relative-cup} applies.
Finally, a target with one minimal row is rationally acyclic,
by minimal-group collapse and the band Tor formula. For those
targets only, the fixed-right-state and the unique-double-core
reductions of Lemma~\ref{CP2:A-mixed-five-reductions} apply.
They are not used for targets having more than one minimal
row and column.

For every remaining pair, the producer
\path{monoid/code/p5_small_targets.cpp} first constructs all
endomorphisms of the right act. It assigns values of a
transformation and propagates exactly the equations
$f(xm)=f(x)m$; each still undetermined value ranges over every
state. It then enumerates all target-monoid homomorphisms into
this endomorphism monoid, fixes the identity, and propagates
only the multiplication equations. Full right-act automorphism
groups are obtained by testing every state permutation. The
producer checks the sum of all state-orbit weights against the
number of labelled action homomorphisms.

The independent checker
\path{monoid/code/p5_small_scope_check.cpp} instead uses the
core graph to build all endomorphisms: a singleton core can
map only to a singleton core, a double core can map to either
kind with columns preserved, and each transient can map to
any state of the resulting ordered orbit-pair type. Its
invertible members are precisely the full automorphism group.
It assigns the target maps in a fixed order, without the
producer's multiplication propagation, and checks every orbit,
stabilizer and degree-one/two mod-two homology filter against
the complete output stream. Its counts agree in every scope.

The complete results are
\[
\begin{array}{c|r|r|r|r|r}
\text{profile}&\text{corner/action scopes}&\text{labelled actions}
 &\text{state orbits}&u^2=0&H_3(\mathbb F_2)>0\\\hline
(5,7,4)&6672&83471&50506&19839&16\\
(5,6,5)&1151&60812&38630&18314&0\\
(5,5,6)\text{ mixed}&294&41660&26241&9907&0\\
(5,4,7)&72&54634&27792&9105&0\\
(5,3,8)&30&54498&15891&6310&0.
\end{array}
\]
All other orbits fail a necessary first- or second-homology
condition. The root literal checker
\path{code/check_p5_bisets_all_sizes.cpp} independently checks
every retained table, nonboundary cocycle, square primitive
or positive-third-homology certificate. All pass. These are
new streams generated from the complete ordinary target
catalogs; none is an inherited or incomplete historical prefix.
The exact scope receipts, full compressed action and cup
streams, and hashes are under
\path{monoid/certificates/p5_small_targets/}.
\end{proof}

\begin{remark*}[Precise remaining triangular baseline]
This proposition closes these five profiles, including the
previously separate $(5,7,4)$ control. It does not by itself
exclude all two-object collages with at most sixteen arrows.
The positive-source-corner profiles with source orders six,
seven and eight, and the analytic large-corner reductions,
remain separate obligations unless completed elsewhere in
this manuscript.
\end{remark*}


\subsection{A universal two-constant action obstruction}

\begin{theorem}[Proved: the identity and two constants]
\label{CP2:A-two-constant-image}
Let $S$ be a nonempty finite left-$N$, right-$M$ biset, with
$M,N$ finite monoids. Suppose its right transformation image is
exactly
\[
\{\operatorname{id}_S,c_0,c_1\},
\]
where $c_i$ is constant at $s_i$ and $s_0\ne s_1$.
Then its two-object collage is not of type $\CP^2$.
The same conclusion holds with the two sides interchanged.
There is no bound on the cardinalities of $M,N,S$.
\end{theorem}
\begin{proof}
Commutation with the two constants forces every left actor to
fix $s_0$ and $s_1$. Put $U=\{s_0,s_1\}$ and let $X_U\subseteq X_S$
be the two collage spaces. Their corner spaces are identical.
The subset $U$ is a subbiset. Write
\[
\chi(m)=\begin{cases}1&\text{if $m$ acts identically on $S$},\\
0&\text{if $m$ acts constantly on $S$}.
\end{cases}
\]
This is a multiplicative zero-one character. On
$L=\ker(\mathbb Q[S]\to\mathbb Q)$ the right action is scalar
$\chi$. The vector $u=[s_1]-[s_0]$ spans a trivial left
$N$ submodule. Put $V=L/\mathbb Qu$ and let
\[
\gamma\in\operatorname{Ext}^1_{\mathbb QN}(V,\mathbb Q)
\]
be the class of
\begin{equation}\label{CP2:A-two-constant-extension}
0\longrightarrow\mathbb Q\xrightarrow{\,u\,}L
\longrightarrow V\longrightarrow0
\end{equation}
as left modules. In this display the commuting right action
on all three modules is scalar $\chi$.

Set
\[
D_N=\mathbb Q\mathbin{\otimes^{\mathbf L}_{\mathbb QN}}V,
\qquad
D_M=\mathbb Q_\chi\mathbin{\otimes^{\mathbf L}_{\mathbb QM}}\mathbb Q.
\]
The relative bridge coefficient is
$\mathbb Q[S]/\mathbb Q[U]=V\otimes\mathbb Q_\chi$.
The double mapping-cylinder chains therefore identify the
reduced relative complex of $(X_S,X_U)$ with
$(D_N\otimes_{\mathbb Q}D_M)[1]$.
Since $\chi$ takes the value zero, $H_0(D_M)=0$.
In particular K\"unneth supplies the summand
\[
\operatorname{Tor}^{\mathbb QN}_1(\mathbb Q,V)
\otimes H_1(D_M)
\ \subseteq\ H_3(X_S,X_U;\mathbb Q).
\]

We identify its boundary, so no comparison of dimensions alone
is being used. The coefficient extension for $U\subset S$ is
the pushout of \eqref{CP2:A-two-constant-extension}, tensored
on the right by $\mathbb Q_\chi$, along the inclusion
$u:\mathbb Q_\chi\to\mathbb Q[U]$.
Choose any vector-space splitting of $L\to V$. In the resulting
normalized double-bar chains the off-diagonal left-action face
is exactly a cocycle for $\gamma$, the right-action faces are
scalar $\chi$, and the vertex augmentations vanish on these
augmentation-module lifts. Hence the connecting chain map is
$\gamma\otimes\operatorname{id}$ followed by the actual
inclusion of $u$ in the core complex.

The left $N$ action on $U$ is trivial. By
Lemma~\ref{CP2:A-trivial-biset-side-H2}, collapsing that corner
induces an isomorphism on rational second homology. The
augmentation sequence
\[
0\longrightarrow\mathbb Q_\chi\xrightarrow{\,u\,}
\mathbb Q[U]\longrightarrow\mathbb Q\longrightarrow0
\]
identifies the reduced ordinary core act-cone complex with
$D_M[1]$, naturally with respect to the displayed inclusion.
Consequently the relative boundary on the preceding K\"unneth
summand is, up to the harmless global boundary sign,
\begin{equation}\label{CP2:A-two-constant-boundary}
\gamma\otimes\operatorname{id}:
\operatorname{Tor}^{\mathbb QN}_1(\mathbb Q,V)
\otimes H_2(X_U;\mathbb Q)
\longrightarrow H_2(X_U;\mathbb Q).
\end{equation}
The duality being used here is the exact vector-space duality
\[
\operatorname{Ext}^1_{\mathbb QN}(V,\mathbb Q)
\cong
\operatorname{Hom}_{\mathbb Q}
\bigl(\operatorname{Tor}^{\mathbb QN}_1(\mathbb Q,V),\mathbb Q\bigr).
\]
It follows by applying $\operatorname{Hom}_{\mathbb Q}(-,\mathbb Q)$
to a free resolution of $V$ tensored with the trivial right
module; the resulting cochain complex is
$\operatorname{Hom}_{\mathbb QN}(-,\mathbb Q)$.

If $\gamma\ne0$, it is a nonzero, hence surjective, functional.
Equation~\eqref{CP2:A-two-constant-boundary} is therefore onto.
The map $H_2(X_U;\mathbb Q)\to H_2(X_S;\mathbb Q)$ vanishes,
so every rational degree-two class on $X_S$ is relative to $X_U$.
Both corners are contained in $X_U$; collapsing it in the
mapping cylinder gives the reduced suspension of the quotient
of the two edge-bar spaces. All positive-degree products in
that quotient vanish. Their pullbacks to $X_S$ therefore vanish,
which is incompatible with the square on $\CP^2$.

If $\gamma=0$, take an $N$-linear retraction
$q:L\to\mathbb Q$ with $q(u)=1$ and put
$q(s)=q([s]-[s_0])$. It satisfies
$q(s_0)=0$, $q(s_1)=1$ and $q(ns)=q(s)$.
Let $\rho:\mathbb Q\to U$ send $1$ to $s_1$ and every other
rational number to $s_0$. Then $r(s)=\rho(q(s))$ fixes $U$
and is left equivariant. It is also right equivariant, since
every right actor is identity or constant at one of the two
fixed points. Thus $r:S\to U$ is an actual biset retraction,
and $X_U$ is an actual retract of $X_S$.

Every rational degree-two product on $X_U$ is zero: its
trivial-side collapse is a degree-two cohomology isomorphism
to the ordinary two-state $M$-act cone, and that act embeds
in the affine line as $\{0,1\}$, with identity and constant
affine maps. Apply Lemma~\ref{CP2:A-high-affine} and pull back.
If $X_S$ were a projective plane, the split injection
$H^2(X_U;\mathbb Q)\to H^2(X_S;\mathbb Q)$ would have rank
zero or one. In the zero case its degree-two generator is
relative to $X_U$ and has square zero by the suspension
quotient. In the rank-one case the generator is pulled back
from $X_U$ through the retraction and again has square zero.
Both are contradictions. This proves the theorem; opposition
proves its left-handed version.
\end{proof}

\begin{remark*}[Computational scope of this gate]
Testing the literal transformation image against exactly the
three maps in the theorem is sufficient. No corner monoid is
replaced by that image, and no homotopy-type assumption on
an individual corner is made. The split case asserts the
projective-plane obstruction; it does not assert that all
products on an arbitrary larger-rank collage vanish.
\end{remark*}


\subsection{The remaining finite triangular profiles at sixteen arrows}

\begin{lemma}[Proved: complete right acts of a rectangular band]
\label{CP2:A-rectangular-band-act-normal-forms}
Let $K=I\times J$, with $(i,j)(k,l)=(i,l)$, and let $X$ be
a finite nonempty right $K^1$-set. Each orbit in its core $XK$
is a quotient of the column set $J$, with a specified partition
of the labelled set $J$. A transient state in $X\setminus XK$
is specified by an $I$-tuple of these core orbits. These data
are unrestricted. Two such descriptions give isomorphic acts
exactly when core orbits with the same column partition are
permuted and transient states with the same row tuple are
permuted.
\end{lemma}
\begin{proof}
If $x=se_{ij}$ lies in the core, then
$xe_{kl}=se_{il}$ is independent of $k$. Thus its orbit is
parametrized by $l\in J$, with a definite partition identifying
columns that give the same state. Any two such core orbits
that meet agree: applying all column operations to the meeting
point recovers both parametrized orbits. For an arbitrary state
$s$, its row-$i$ images $se_{ij}$ lie in one of these orbits,
and $j$ selects the labelled column point. Conversely the rule
just described satisfies
$(se_{ij})e_{kl}=se_{il}$ for every choice of the row tuple.
An equivariant bijection preserves the labelled column
partition of each core orbit and preserves every transient's
row tuple. This gives precisely the stated equivalence.

The new constructor \path{monoid/code/triangular_band_inputs.py}
enumerates every multiplicity of each column-partition type,
then every multiset of row tuples, modulo permutations of core
orbits of identical type. It also verifies coverage independently
by Burnside's lemma: for a permutation of the typed core orbits,
let $d_1,\ldots,d_a$ be its cycle lengths on $I$-tuples. The
number of invariant transient multisets of cardinality $t$ is
the coefficient of $z^t$ in
$\prod_j(1-z^{d_j})^{-1}$. The group average matches every
multiplicity stratum. A separate extraction from the literal
transformation maps yields pairwise distinct canonical core
signatures and checks all action equations.
The complete state-isomorphism counts needed below are
\[
\begin{array}{c|rrrr}
(|I|,|J|)\backslash |X|&4&5&6&7\\\hline
(2,2)&17&44&130&383\\
(2,3)&36&108&360&\text{not needed}\\
(3,2)&38&194&1272&\text{not needed}.
\end{array}
\]
\end{proof}

\begin{proposition}[Proved: the nine remaining finite profiles]
\label{CP2:A-triangular-sixteen-finite-profiles}
No two-object collage with one of the profiles
\[
\begin{gathered}
(8,4,4),\ (7,5,4),\ (6,6,4),\ (8,3,5),\ (7,4,5),\\
(6,5,5),\ (7,3,6),\ (6,4,6),\ (6,3,7)
\end{gathered}
\]
has classifying space of type $\CP^2$. The entries specify
the source corner order, target corner order and number of
arrows from source to target; the reverse hom-set is empty.
\end{proposition}
\begin{proof}
If both corners have zero rational second cohomology, use
Lemma~\ref{CP2:A-collage-relative-cup}. A target of order at
most four is rationally acyclic. An order-five target with
nonzero rational second cohomology is $P_5$; its cases here
are already covered, after opposition, by
Proposition~\ref{CP2:A-P5-source-sixteen}.
For the equal-order $(6,6,4)$ profile, use opposition if needed
to choose a positive-$H_2$ source. Thus all other cases require
only the complete positive-source catalogs of
Proposition~\ref{CP2:A-positive-corners-eight}: five monoids
in order six, $43$ in order seven and $510$ in order eight.
The target always ranges over its complete ordinary catalog,
including all targets with zero second homology.

A target whose minimal ideal has one column has projective
trivial right rational module. Exact coinvariants replace its
left action on $S$ by the actual finite orbit set $N\backslash S$.
This is an actual quotient functor to an ordinary right-source
act cone. In the target bar direction its edge comparison is
a rational homology equivalence, and $BN\to *$ is also one;
the double mapping-cylinder comparison therefore induces an
isomorphism of rational cohomology rings in every degree.
If the original collage were a projective plane, the resulting
ordinary act cone would have bounded rational homology.
Lemma~\ref{CP2:A-baseline-eight-act} excludes it, since the
source has order at most eight. This explicitly transfers the
boundedness hypothesis along an all-degree comparison.

After that gate, a target of order three is uniquely the
identity-adjoined two-element right-zero band. Its $(6,3,7)$
case is excluded for every state count by
Theorem~\ref{CP2:A-baseline-right-zero-three-six}.
This analytic theorem, rather than an interrupted computation,
is the premise for that ninth profile. We now give the complete
new finite computation for the other eight.

Each source's minimal ideal is a rectangular band of one of
the shapes in Lemma~\ref{CP2:A-rectangular-band-act-normal-forms}.
The source input includes its literal table and normalized band
coordinates. These are recomputed from its minimal ideal and
checked directly. Restrict a right source action to this band
and choose its complete normal form. For a source element $a$
outside the band, its unknown map $f_a$ must satisfy the two
families of equations
\[
 f_a\circ f_k=f_{ka},\qquad f_k\circ f_a=f_{ak}
 \quad(k\in K).
\]
The right-hand maps are known band maps. The first family
forces values on core states; the second lists exactly the
permitted image of each remaining state. The producer tries
every resulting map for every outside element and propagates
only the equations $f_{ab}=f_b\circ f_a$. Thus every source
action extending the band form occurs.

There are three whole-band omissions. If all band maps are
the identity, so are all source maps, because $ka\in K$ and
$f_a=f_a\circ f_k=f_{ka}$; the trivial-side comparison reduces
to an ordinary target act cone of order at most six.
If the band has exactly one fixed state, it is also the unique
source-fixed state, and commutation makes it target-fixed.
If the target has one minimal row, it is rationally acyclic,
so any source-fixed state is sufficient for the relative-cup
obstruction. These omissions are verified from literal maps
and the target's minimal coordinates. No arbitrary zero-$H_2$
target is treated as rationally acyclic.

For each remaining source action the producer computes its
complete state automorphism group and its full transformation
centralizer. It enumerates every unital homomorphism from the
actual target monoid into that centralizer, then takes orbits
under exactly that source-action automorphism group.
It omits an entire source-action branch if its literal image
is exactly identity and two distinct constants, using
Theorem~\ref{CP2:A-two-constant-image}. That right-action
record retains the actual maps and the gate. No other
source-image replacement, subgroup restriction, proper
idempotent retraction or assumed smaller-category bound is
used by this computation.

The new program is \path{monoid/code/triangular_bisets.cpp}.
Its independent checker
\path{monoid/code/triangular_scope_check.cpp} imports none of
the producer. It scans all $|S|^{|S|}$ transformations to
construct every possible outside-actor map and every centralizer,
uses all $|S|!$ permutations to obtain automorphisms, and assigns
source and target images in fixed order without multiplication
propagation. It matches every source and biset state orbit,
checks its full stabilizer, and verifies the raw/orbit-weight
identity at both stages. It also recomputes every first- and
second-homology filter independently. All eight scopes pass.

The complete results after the stated analytic gates are
\[
\begin{array}{c|r|r|r|r}
\text{profile}&\text{labelled left actions}&\text{biset state orbits}
 &u^2=0&H_3(\mathbb F_2)>0\\\hline
(8,4,4)&312830&270768&170768&1506\\
(7,5,4)&168779&141646&93396&1512\\
(6,6,4)&139637&109909&70358&1584\\
(8,3,5)&486579&337863&170948&924\\
(7,4,5)&253790&160559&82189&702\\
(6,5,5)&149282&86721&47200&0\\
(7,3,6)&307533&133930&59815&250\\
(6,4,6)&217931&71295&30578&78\\\hline
\text{total}&2036361&1312691&725252&6556.
\end{array}
\]
There are $49612$ state-isomorphism classes of right source
acts over the retained band forms, with $121907$ completed
source-action/target-monoid scopes. The raw source-action count
before quotienting by band-form automorphisms is $225352$.
Of the biset orbits, $26991$ fail the first mod-two homology
condition and $553892$ fail the second; all $731808$ remaining
ones have the displayed obstruction.

The root literal checker
\path{code/check_p5_bisets_all_sizes.cpp}, whose computation
is for arbitrary typed two-object tables despite its filename,
reconstructs all normalized chains and checks each nonboundary
cocycle, square primitive or positive-third-homology certificate.
The additional independent program
\path{monoid/code/join_biset_streams.py} rebuilds the category
from every retained action record and matches its source,
target, complete multiplication word and reason byte for byte
with the corresponding literal cup record. It also requires
exact simultaneous end of file, so a truncated or unjoined
stream cannot pass. The finalized receipts, literal data and
hashes are retained under
\path{monoid/certificates/triangular_inputs/} and
\path{monoid/certificates/triangular_runs/}.
\end{proof}

\begin{remark*}[The stopped ninth computational control]
An optional direct run for $(6,3,7)$ was stopped by its explicit
centralizer memory guard before completion. No prefix of that
stream is used in this proposition, and no completed receipt
is claimed for it. Its diagnostic and partial bytes are
retained with an explicit incomplete-status record. The
separate all-state graph theorem cited in the proof supplies
the complete mathematical exclusion of that profile.
\end{remark*}


\subsection{The complete one- and two-object baseline}

\begin{theorem}[Proved: at most two objects require seventeen arrows]
\label{CP2:A-two-object-baseline-sixteen}
A finite category with at most two objects and classifying
space of type $\CP^2$ has at least seventeen arrows,
identities counted.
\end{theorem}
\begin{proof}
The one-object case follows from the complete minimum-monoid
classification, Theorem~\ref{A:monoid-minimum23-classification}.
A two-object model must be connected. If arrows exist in
both directions, Proposition~\ref{CP2:A-baseline-strong-two}
excludes even seventeen arrows. Otherwise it is a triangular
collage. By identity padding,
Lemma~\ref{CP2:A-identity-padding}, any model with fewer than
sixteen arrows would give one with exactly sixteen arrows,
with the same two objects and the same support. It therefore
suffices to exclude triangular collages at that exact total.

Write $(m,n,s)$ for the source corner order, target corner
order and cross-hom size. Thus $m+n+s=16$, with all three
entries positive. If both corners have zero rational second
cohomology, the relative-cup lemma excludes the projective-plane
square. Otherwise choose a corner with positive second
cohomology as the source, taking the opposite category if
necessary. The order-five corner classification gives
$m\geq5$. This choice does not assume $m\geq n$.

A singleton cross-hom realizes the join of the two connected
corner spaces and has zero rational second homology.
Two cross arrows are excluded by
Lemma~\ref{CP2:A-high-two-cross}; three by
Theorem~\ref{CP2:A-three-cross-reduction}. Hence $s\geq4$.

If $n=1$, the collage is an ordinary right act cone and
$m+s=15$. For $m\leq9$ use
Theorem~\ref{CP2:A-bounded-nine-acts}: the assumed
projective-plane type supplies bounded rational homology.
The only remaining possibilities are $(m,s)=(10,5)$ and
$(11,4)$. They are excluded, respectively, by
Proposition~\ref{CP2:A-ten-six-state-rational} and
Corollary~\ref{CP2:A-four-state-small}.

If $n=2$, both possible target monoids have projective
trivial right rational module. Indeed they are $C_2$ or
$\{1,z:z^2=z\}$; averaging in the first case and the
idempotent $z$ in the second give the splitting. The actual
orbit-act quotient functor is a rational homology and
cohomology-ring equivalence, by the exact-coinvariant
argument in Proposition~\ref{CP2:A-triangular-sixteen-finite-profiles}.
Thus it preserves the rational projective-plane target and
all-degree homological boundedness. Now $m+s=14$ and
$s\geq4$, so $m\leq10$. The bounded-nine theorem handles
$m\leq9$, and $(m,s)=(10,4)$ is covered by the
rational ten-element/six-state proposition. No actual
homotopy equivalence is inferred from this rational comparison.

It remains that $n\geq3$ and $s\geq4$. The possibilities
are exactly
\[
\begin{array}{c|l}
 m& (n,s)\\\hline
5&(3,8),(4,7),(5,6),(6,5),(7,4)\\
6&(3,7),(4,6),(5,5),(6,4)\\
7&(3,6),(4,5),(5,4)\\
8&(3,5),(4,4)\\
9&(3,4).
\end{array}
\]
The first row is Proposition~\ref{CP2:A-P5-source-sixteen},
since a positive-$H_2$ order-five source is $P_5$.
The next three rows are
Proposition~\ref{CP2:A-triangular-sixteen-finite-profiles},
including its separately proved graph-theorem treatment of
$(6,3,7)$. The last row is
Proposition~\ref{CP2:A-right-zero-four-states}, which includes
the projective opposite-corner alternatives explicitly.

For an arithmetic check independent of the table searches,
\path{monoid/code/two_object16_inventory.py} enumerates all
$55$ positive integer triples with $m\geq5$ and sum sixteen.
It verifies that the preceding gates are exhaustive and that
exactly the displayed fifteen profiles remain after the
small-cross and small-target reductions. The exact inventory
is \path{monoid/certificates/two_object16_inventory.json}.
Every case has now been excluded; padding gives the statement
for all smaller totals as well.
\end{proof}


\subsection{A literal retract gate for the seventeen-arrow searches}

\begin{lemma}[Proved: a literal biset idempotent supplies a forbidden smaller retract]
\label{CP2:A-seventeen-central-actor-retract}
Let a two-object collage have exactly seventeen arrows,
and let $S$ be its nonempty connecting biset. Suppose a
transformation $r:S\to S$ is idempotent, differs
from the identity, and commutes with both monoid actions. Then the collage cannot have type $\CP^2$.
\end{lemma}
\begin{proof}
Thus $r:S\to rS$ is an actual equivariant
biset retraction. Its image is a proper subset: a bijective
idempotent is the identity. Keep both corners and replace
$S$ by $rS$ to obtain a category $Y$ with at most sixteen
arrows and an actual categorical retraction from $X$.

The relative quotient of the two double mapping cylinders
is a suspension, since both corner spaces are retained
and only the bridge is replaced by its sub-bridge.
If $X$ had type $\CP^2$ and $H^2(Y;\mathbb Z)$ vanished,
its degree-two generator would lift from that suspension
quotient and would have square zero. Therefore the split
map in degree two has rank one. All integral homology
of $Y$ is a direct summand of the homology of $X$;
naturality of the primitive square now forces the full
integral projective-plane homology and primitive square
on $Y$. Its fundamental group is a retract of the trivial
group. The integral recognition theorem consequently
gives $BY\simeq\CP^2$, contradicting
Theorem~\ref{CP2:A-two-object-baseline-sixteen}.

The new seventeen-arrow searches using this gate retain
the literal transformation code for $r$. A right actor
commuting with all right actors automatically commutes
with the left action. More generally, a proper idempotent
in the center of the full right-action endomorphism monoid
commutes with every possible left action. These two cases
allow a whole right-action branch to be omitted. After
a particular left action is assigned, any proper idempotent
commuting with both actions allows that individual orbit
to be omitted. The complete action stream retains each
such orbit and its witness. The independent scope checker
separately verifies $r^2=r$, $r\ne1$, and all required
commutation equations, including the full endomorphism
monoid when a whole branch is omitted. The sixteen-arrow searches
use no such gate. Stopped computations are not treated
as completed scopes.
\end{proof}


\subsection{The complete positive-second-homology source scope at order nine}

\begin{proposition}[Proved: the order-nine source catalog]
\label{CP2:A-positive-corners-nine}
There are exactly $8046$ ordinary isomorphism classes of
nine-element monoids with nonzero rational second homology.
Of these, $8044$ have a rectangular minimal ideal and two
have a maximal subgroup of order two in their minimal ideal.
The rectangular-band counts are
\[
\begin{array}{c|r|r}
\text{minimal shape}&\text{all ordinary classes}&H_2(\mathbb Q)\ne0\\\hline
2\times2&21534&7784\\
2\times3&191&129\\
3\times2&191&129\\
2\times4&1&1\\
4\times2&1&1.
\end{array}
\]
\end{proposition}
\begin{proof}
The minimal-ideal coordinates and rational averaging argument
of Lemmas~\ref{A:minimal-ideal-coordinates} and
\ref{A:minimal-group-collapse} show that a positive second
homology group requires at least two minimal rows and two
minimal columns. A proper ideal uses at most eight elements.
Thus a trivial maximal subgroup gives exactly the five
rectangular shapes displayed. A nontrivial maximal subgroup
forces the shape $2\times C_2\times2$ and leaves only the
identity outside the ideal.

For each rectangular shape, normalize its band as
$x_{ij}$ and write every outside element as its row and
column transformation pair. Every multiset of these pairs,
including the forced identity, occurs in the enumeration.
Products involving a band element are forced. Each product
of two other outside elements is allowed exactly those
labels having the required composite actor pair. Testing
associativity on every completion therefore exhausts all
monoids with the normalized band. Every isomorphism preserves
the minimal ideal and restricts there to independent row
and column permutations; outside labels may be permuted
arbitrarily. These are precisely the canonicalization
permutations, so the ordinary isomorphism scope is exact.

The new order-nine extension
\path{monoid/code/positive_corners_nine.cpp} completes this
scope. The largest family has $3876$ actor multisets,
$905$ closed actor multisets, and $297345$ labelled associative
tables. The other nontrivial families have $5886$ actor
multisets each and $1672$ labelled associative tables each.
The new independent program
\path{monoid/code/check_positive_nine.cpp} assigns outside
products in reverse order, without the producer's domain
propagation. It checks only fully specified associativity
equations during the search and reaches exactly every
reported canonical table. In the $2\times2$ family alone
it checks $3785831415$ such equations. It also reconstructs
the balanced augmentation tensor over $\mathbb Q$, by exact
integer row operations in dimension at most three. These
ranks agree with the producer's modular bar computation
for every table, and it checks both membership and full
coverage of each positive-homology stream. All five
complete checks pass.

For the nontrivial-group case, normalize the sandwich
matrix to
$\left(\begin{smallmatrix}1&1\\1&g^b\end{smallmatrix}\right)$,
where $g^2=1$ and $b\in\{0,1\}$. Thus the two monoids have
identity together with elements $(i,a,j)$,
$i,a,j\in\{0,1\}$, and multiplication
\[
 (i,a,j)(k,d,l)=(i,a+d+bjk\pmod2,l).
\]
These are the two normalized sandwich possibilities. They
are nonisomorphic: the zero sandwich bit leaves the local
group visible in the universal group, whereas for bit one
$g=e_{01}e_{10}$ is a product of idempotents and the
universal group is trivial. Rational averaging collapses
each to $P_5$, so both have rational second Betti number one.
This gives the stated total.
\end{proof}

\begin{lemma}[Proved: complete small acts of the two Rees-group sources]
\label{CP2:A-C2-Rees-nine-acts}
For the two monoids in the preceding proof, the numbers of
right acts up to state isomorphism on four and five states
are respectively
\[
\begin{array}{c|rr|rr}
 &\multicolumn{2}{c|}{\text{state-isomorphism classes}}
 &\multicolumn{2}{c}{\text{labelled acts}}\\
 b&4&5&4&5\\\hline
 0&37&121&626&10371\\
 1&18&47&281&3676.
\end{array}
\]
The delivered lists are complete literal actions.
\end{lemma}
\begin{proof}
Write $E,H,A,B$ for the right transformations of
$(0,0,0),(0,1,0),(1,0,0),(0,0,1)$. Let $U=\operatorname{im}E$.
The multiplication table gives the following exact description.
Choose any idempotent $E$ on the state set; any involution
$\sigma$ of $U$, extended by $H=\sigma E$; any section
$s:U\to S$ of $E$, extended by $B=sE$; and any retraction
$A:S\to U$ satisfying
\[
 A\,B=H^b,
 \qquad H^0:=E.
\]
The action of $(i,a,j)$ is
\[
 B_j\,H_a\,A_i,
 \quad (A_0,A_1)=(E,A),\quad
 (B_0,B_1)=(E,B),\quad (H_0,H_1)=(E,H),
\]
where juxtaposition denotes composition of transformations.
Indeed $A_kB_j=H^{bkj}$, and this identity verifies every
right action product. Conversely the four named elements
generate the ideal and force exactly these choices. Thus
scanning every labelled idempotent, involution, section and
retraction gives every right action once.

There is a separate structural count. The core is
$U\cup B(U)$. A fixed point of $\sigma$ gives a one-state
fixed core or a two-state column core. A free two-point
$\sigma$-orbit gives a four-state core when neither of
its two $U$ points is fixed by the section. If $b=0$, it
may also give a three-state core when exactly one point
is fixed by the section, or a two-state group core when
both are fixed. If $b=1$, the equation $AB=H$ forbids
both of these overlaps. The one-state, two-state column
and three-state cores have trivial automorphism group;
the two-state group and four-state cores have automorphism
group $C_2$. There are no further core types.

A transient state is specified freely by the ordered pair
$(Ex,Ax)\in U^2$. Thus, for each multiplicity of each core
type, transients form a multiset of such ordered pairs,
modulo the typed core permutations and the stated internal
automorphisms. Burnside's lemma computes its cardinality:
for the cycle lengths $d_i$ of a core automorphism on
$U^2$, take the coefficient of $z^t$ in
$\prod_i(1-z^{d_i})^{-1}$ and average over that finite
automorphism group. The independent core count agrees
with the labelled-generator enumeration in every stratum.
The latter also checks every literal action equation and
all stabilizer/orbit-size identities. The new constructor
and both count mechanisms are in
\path{monoid/code/rees_c2_nine_acts.py}.
\end{proof}

\begin{proposition}[Proved: the nontrivial minimal-group portions]
\label{CP2:A-C2-Rees-nine-bisets}
A collage with profile $(9,3,5)$ or $(9,4,4)$ and a
nine-element source whose minimal ideal has nontrivial
maximal subgroup cannot have type $\CP^2$.
\end{proposition}
\begin{proof}
If the opposite corner has one minimal column, exact
coinvariants give the ordinary source act cone, preserving
all rational homology and its ring. Boundedness is therefore
preserved, and Lemma~\ref{CP2:A-bounded-nine-acts} excludes it.
Every remaining opposite corner in orders three and four
has one minimal row, so is rationally acyclic. A right-fixed
state gives the usual relative-cup obstruction; a unique
right-fixed state is also left-fixed. A trivial right action
reduces to an ordinary opposite-corner act cone. These are
the only analytic omissions.

For every other right act in Lemma~\ref{CP2:A-C2-Rees-nine-acts},
\path{monoid/code/rees_c2_bisets.cpp} computes the entire
transformation centralizer and every unital opposite-corner
action, modulo the exact right-act automorphism group.
Its separate checker
\path{monoid/code/rees_c2_biset_check.cpp} scans all $S^S$
transformations for the centralizer and assigns target
images in fixed order without product propagation.
It verifies every orbit, stabilizer and first/second
mod-two homology filter. The independent literal bar
checker and the further exact stream linker also pass.
The complete counts are
\[
\begin{array}{c|c|r|r|r|r}
 b&\text{profile}&\text{labelled actions}&\text{state orbits}
 &u^2=0&H_3(\mathbb F_2)>0\\\hline
0&(9,3,5)&653&249&43&26\\
0&(9,4,4)&696&431&86&40\\
1&(9,3,5)&383&103&43&2\\
1&(9,4,4)&382&218&86&6.
\end{array}
\]
All $332$ retained tables have the displayed obstruction.
The complete source, action, scope and cup records are
under \path{monoid/certificates/positive_nine/} in the
checkpoint archive. This proposition does not assert that
the rectangular-band portions of the same two profiles
have finished.
\end{proof}


\subsection{The P5 corner with four connecting arrows}

\begin{lemma}[Proved: a full regular column core]
\label{CP2:A-P5-full-regular-core}
Let $R=P_5$ and let $S$ be an arbitrary finite nonempty
left-$N$, right-$R$ biset. Suppose that $S$ is a disjoint union
of two-point regular column orbits for the minimal $2\times2$
band of $R$, with no transient or fixed states. Then the
classifying space of the collage has $BR$ as a wedge summand.
In particular it cannot have the homotopy type of $\CP^2$.
\end{lemma}
\begin{proof}
Write the states as $(a,j)$, where $a$ indexes the column
orbits and $j\in\{0,1\}$ is the labelled column. Thus
$(a,j)e_{kl}=(a,l)$. Every map commuting with the right action
preserves the labelled column, so $n(a,j)=(na,j)$ for a
well-defined action of $N$ on the orbit set. Set
$h_{(a,j)}=e_{0j}\in R$. For every $m\in R$ and $n\in N$,
\[
 h_s m=h_{sm},\qquad h_{ns}=h_s.
\]
The first identity gives a natural transformation between the
projection of the right action category to $BR$ and a constant
functor: its components are the displayed $h_s$. Equivalently,
inserting $h_s$ supplies the usual bar homotopy from that
projection to the constant map. The second identity makes
this homotopy $N$-equivariant. It therefore descends through
the other bar direction and proves that the bridge projection
$E=B(N,S,R)\longrightarrow BR$ in the double mapping-cylinder
model is nullhomotopic.

The homotopy pushout of $BR\leftarrow E\to BN$ consequently
splits as
\[
 BR\vee\operatorname{hocofib}(E\longrightarrow BN).
\]
The rational second cohomology of $BR$ is one-dimensional,
with zero square. If the wedge had the cohomology of $\CP^2$,
its only degree-two generator would be the pullback from this
summand, and would have zero square. This is impossible.
\end{proof}

\begin{lemma}[Proved: an acyclic opposite corner on at most four states]
\label{CP2:A-P5-four-state-acyclic-opposite}
Let $N$ be any finite monoid with rationally acyclic classifying
space. A left-$N$, right-$P_5$ biset on at most four states
cannot give a collage of type $\CP^2$.
\end{lemma}
\begin{proof}
If a state is fixed by the right $P_5$ action, apply the
fixed-state case of Lemma~\ref{CP2:A-mixed-five-reductions}.
Otherwise all core orbits have two points, by
Lemma~\ref{CP2:A-rectangular-band-act-normal-forms}.
There is either a unique such core, covered by the
semilinear row-interchange case of
Lemma~\ref{CP2:A-mixed-five-reductions}, or exactly two cores.
In the latter case the four states are all core states, so
Lemma~\ref{CP2:A-P5-full-regular-core} applies.
\end{proof}

\begin{proposition}[Scope reduction for the $(5,8,4)$ profile]
\label{CP2:A-P5-eight-corner-scope}
To exclude every collage with source $P_5$, opposite corner
of order eight, and four connecting arrows, it suffices to
check the $1205$ ordinary monoids of order eight having at
least two minimal rows and two minimal columns. These are
exactly the $1181$, $12$ and $12$ tables in the complete
$2\times2$, $2\times3$ and $3\times2$ minimal-band catalogs,
respectively. All tables in those catalogs must be used,
including the ones with zero rational second homology.
\end{proposition}
\begin{proof}
A one-column minimal ideal makes the trivial right rational
$N$-module projective. Exact left coinvariants replace the
biset by its actual orbit set, and the resulting ordinary
$P_5$ act cone has the same rational cohomology ring in every
degree. The bounded small-monoid act theorem excludes it.
A one-row minimal ideal makes $BN$ rationally acyclic, and
Lemma~\ref{CP2:A-P5-four-state-acyclic-opposite} excludes it.

In the remaining case the minimal ideal has at least two
rows and two columns. Its maximal subgroup is trivial:
a nontrivial one would give at least eight ideal elements,
plus the identity outside the ideal. The only possible
rectangular shapes are therefore $2\times2$, $2\times3$
and $3\times2$. The complete Rees-coordinate enumeration
of Proposition~\ref{CP2:A-positive-corners-eight} produces
all tables of these shapes before its second-homology
filter; those full streams have the stated counts.
The new input checker
\path{monoid/code/p5_eight_corner_inputs.py} reads every
full-catalog table, recomputes its identity, associativity,
minimal ideal and two-sided shape, and verifies all $1205$
resulting records and their input/output digests.
\end{proof}


\subsection{Four further complete $P_5$-source profiles at seventeen arrows}

\begin{proposition}[Proved: the four new $P_5$-source scopes]
\label{CP2:A-P5-source-seventeen}
No two-object collage of profile
\[
(5,4,8),\quad(5,6,6),\quad(5,7,5),\quad(5,8,4)
\]
whose five-element corner has nonzero rational second
homology has classifying space of type $\CP^2$.
\end{proposition}
\begin{proof}
The positive five-element corner is $P_5$ by
Lemma~\ref{CP2:A-order-five-corner}.
The source acts are the complete rectangular-band normal
forms of Lemma~\ref{CP2:A-rectangular-band-act-normal-forms}:
there are respectively $1212,130,44,17$ state-isomorphism
classes at the four state sizes in the displayed order.
The eight-state count is checked independently by Burnside
counts in every core/transient stratum, and by extraction
of those strata from the literal act tables.

At opposite orders four, six and seven, all $35,2237,31559$
ordinary monoid classes are included before the proved
projective-right, fixed-state, and unique-core reductions
of Lemma~\ref{CP2:A-mixed-five-reductions}. At opposite
order eight, Proposition~\ref{CP2:A-P5-eight-corner-scope}
reduces the complete scope to all $1205$ two-sided minimal-band
tables, including those with vanishing second homology.
The omitted one-row and one-column cases are excluded by
that proposition. Trivial-right collapse is an actual
categorical retract, so the bounded-homology hypothesis of
the small-act theorem used there is inherited.

Every remaining commuting opposite action is enumerated
up to the exact automorphism group of the source act.
The separate checker independently constructs every commuting
transformation and uses fixed-order assignments, independently of the
producer's product propagation. It checks the complete
orbit/stabilizer counts and all stated analytic omissions.
The resulting counts are as follows; ``retained'' means
that the mod-two first and second Betti numbers are zero
and one, respectively.
\[
\begin{array}{c|r|r|r|r|r}
(m,n,s)&\text{labelled actions}&\text{state orbits}&
\text{retained}&\text{square zero}&H_3\ne0\\\hline
(5,4,8)&1688224&433652&110038&110038&0\\
(5,6,6)&1174129&471073&186364&186352&12\\
(5,7,5)&1141093&664300&316704&316630&74\\
(5,8,4)&952264&587016&200780&200626&154
\end{array}
\]
For every retained literal category, a further checker
verifies its identity laws, all typed associativity equations,
its first and second mod-two Betti numbers, and a nonzero
second cocycle. It then checks either every equation of
an explicit third cochain whose coboundary is the cup
square, or a positive third Betti number. Both alternatives
exclude $\CP^2$.
A third exact linker reconstructs each retained table from
the action record and matches its full multiplication word
and exclusion reason to the literal certificate stream,
requiring simultaneous end of file. All four complete
producer footers and all twelve independent receipts pass.
Their counts, hashes, programs and dependency paths are
recorded in
\path{monoid/certificates/p5_seventeen/manifest.json}.
\end{proof}


\subsection{The exact remaining two-object scope at seventeen arrows}

\begin{proposition}[Proved: the seventeen-arrow scope reduction]
\label{CP2:A-two-object-seventeen-scope}
To exclude all categories with at most two objects and at
most seventeen arrows from having type $\CP^2$, it remains
to exclude the following seventeen positive-source collage
scopes:
\[
\begin{gathered}
(5,4,8),\ (5,6,6),\ (5,7,5),\ (5,8,4),\\
(6,4,7),\ (6,5,6),\ (6,6,5),\ (6,7,4),\\
(7,3,7),\ (7,4,6),\ (7,5,5),\ (7,6,4),\\
(8,3,6),\ (8,4,5),\ (8,5,4),\\
(9,3,5),\ (9,4,4).
\end{gathered}
\]
Here the first corner is required to have nonzero rational
second homology; the entries specify the two corner orders
and connecting-biset cardinality. This proposition specifies the complete reduced scope; all
listed finite pieces are completed in
Theorem~\ref{CP2:A-two-object-seventeen}.
\end{proposition}
\begin{proof}
The one-object and at-most-sixteen-arrow cases are already
excluded. The wholly strong two-object support is covered
by the reviewed strong-support reduction. Thus a remaining
case is a collage. If both corners have zero rational
second cohomology, Lemma~\ref{CP2:A-collage-relative-cup}
annihilates every second square. By opposition choose a
positive-second-homology corner as the source. Its order
$m$ is at least five. Write $n$ for the other corner order
and $s$ for the number of connecting arrows, so
$m+n+s=17$ and $n,s\geq1$.

All cases with a corner of order at most two are excluded
by Corollary~\ref{CP2:A-act-cones-through-seventeen}.
A singleton connecting set gives the join of the two
connected corner spaces, whose second homology is zero.
Two connecting arrows are covered by
Lemma~\ref{CP2:A-high-two-cross}, and three by
Theorem~\ref{CP2:A-three-cross-reduction}.
We may therefore assume $n\geq3$ and $s\geq4$.
The elementary integer budget leaves exactly $21$ ordered
positive-source profiles. Three are excluded analytically:
$(5,3,9)$ and $(6,3,8)$ by the projective-right endpoint
reduction and Theorem~\ref{CP2:A-baseline-right-zero-three-six},
and $(10,3,4)$ by
Proposition~\ref{CP2:A-right-zero-four-states} together
with the projective-right endpoint reduction. The profile
$(5,5,7)$ was completed in
Proposition~\ref{CP2:A-five-five-seven}. Removing these
four gives precisely the displayed list.
The separately authored integer inventory
\path{monoid/code/two_object17_inventory.py} checks all
$66$ triples with $m\geq5$ and positive $n,s$, and the
exact $21$-profile remainder.

At source order five the only positive source is $P_5$.
At orders six through eight the source catalogs have
$5,43,510$ ordinary classes; at order nine,
Proposition~\ref{CP2:A-positive-corners-nine} gives all
$8046$ classes. Every opposite corner through order seven
ranges over its complete ordinary catalog. The only
order-eight opposite corner occurs in $(5,8,4)$ and is
covered by the full $1205$-table reduction of
Proposition~\ref{CP2:A-P5-eight-corner-scope}; that reduction
includes the zero-second-homology tables and proves the
omitted one-row and one-column cases separately.

In direct searches with opposite corner order five,
the $P_5$ opposite corner may be omitted because opposition
places that same case in the displayed $P_5$-source row.
This is an explicit dependency on the $(5,m,s)$ scope,
not an omission of all opposite corners with zero second
homology. The two nontrivial-minimal-group portions in
order nine are already completely excluded by
Proposition~\ref{CP2:A-C2-Rees-nine-bisets}; both rectangular-band
portions retain all $8044$ positive source tables.
Thus the stated finite and analytic dependencies cover
all two-object supports and all arrow budgets through
seventeen.
\end{proof}

\begin{remark*}[What each finished finite scope certifies]
All rectangular-band right acts are listed by
Lemma~\ref{CP2:A-rectangular-band-act-normal-forms}.
The additional state sizes use the same complete constructor,
with independent Burnside counts in every multiplicity
stratum. The seven-state $2\times3$ and $3\times2$ counts
are $1184$ and $9982$; the four/five-state $2\times4$ counts
are $93,349$, and the corresponding $4\times2$ counts
are $107,1187$.

Outside-source actors are assigned only from their exact
band compatibility equations. Every commuting unital
opposite action is then enumerated, modulo the exact
state automorphism group. The separate scope checker
independently constructs every commuting transformation and
uses fixed-order assignments
without the producer's product propagation. It checks every
orbit and stabilizer, and either checks the explicit proper
biset-retract witness of
Lemma~\ref{CP2:A-seventeen-central-actor-retract}, or
recomputes the first and second mod-two homology filters.
All surviving mod-two projective-plane rank candidates
receive a literal square primitive or a positive-third-homology
certificate, checked independently. A further exact linker
reconstructs each retained category from the action stream
and matches its full multiplication word and exclusion
reason, with simultaneous end of file.

If a right action is trivial, collapsing that corner gives
an actual categorical retract: the section includes its
identity element. The degree-two comparison is an
isomorphism, and boundedness passes to the retract. Thus
when the opposite corner has order eight, the bounded
small-act theorem is used with its boundedness hypothesis
fully justified. A mere degree-two isomorphism alone is
not used to infer boundedness.

The live inventory
\path{monoid/certificates/two_object17_inventory.json}
marks a finite piece complete only after a completed
producer footer and all three independent receipts. It
also checks dependencies on opposite $P_5$ scopes. A
producer prefix or an interrupted large-centralizer attempt
is explicitly insufficient.
\end{remark*}


\subsection{The complete two-object lower bound through seventeen arrows}

\begin{theorem}[Proved: the complete two-object seventeen-arrow exclusion]
\label{CP2:A-two-object-seventeen}
Every finite category with at most two objects and classifying
space of homotopy type $\CP^2$ has at least eighteen arrows.
\end{theorem}
\begin{proof}
Proposition~\ref{CP2:A-two-object-seventeen-scope} reduces the
new seventeen-arrow case to seventeen explicitly listed
positive-source collage profiles, with all smaller budgets,
one-object cases, wholly strong support, and other integer
profiles already excluded. The four new $P_5$-source profiles
are completed by Proposition~\ref{CP2:A-P5-source-seventeen}.
We now give the complete finite check of the thirteen
remaining profiles. The two nontrivial-minimal-group parts
at source order nine are separately completed by
Proposition~\ref{CP2:A-C2-Rees-nine-bisets} and are joined to
their rectangular-band parts below.

Use precisely the complete source catalogs, opposite-corner
catalogs, and right-act normal forms specified in the scope
proposition. Every outside-source actor is assigned subject
to all minimal-band product equations; every remaining
commuting opposite action is then assigned. Orbit reduction
uses the exact state automorphism group of the right action.
The independent scope checker constructs the transformation
sets by scanning all state maps and assigns source and
opposite actors in fixed order, without the producer's
product-propagation search. It verifies the full action
equations, orbit/stabilizer weights, and every declared
analytic omission, and requires simultaneous exhaustion of
the action records.

Four of the scopes also use the actual biset
retractions of Lemma~\ref{CP2:A-seventeen-central-actor-retract}.
The independent checker verifies the displayed map is a
proper idempotent and commutes with both actions. In a
whole-centralizer omission it independently constructs the
entire right-endomorphism monoid and checks commutation
with every member. These witnesses give a smaller actual
$\CP^2$ collage if the original were $\CP^2$, contradicting
the already proved two-object bound through sixteen arrows.
Thus these omissions have no dependence on an unfinished
seventeen-arrow case.

The complete rectangular-band-source results are:
\[
\begin{array}{c|r|r|r|r}
(m,n,s)&\text{biset orbits}&\text{retained}&
\text{square zero}&H_3\ne0\\\hline
(6,4,7)&415161&145399&145399&0\\
(6,5,6)&1007089&449128&448262&866\\
(6,6,5)&1350843&755573&755573&0\\
(6,7,4)&1773110&1209509&1183635&25874\\
(7,3,7)&617214&23295&23295&0\\
(7,4,6)&1750921&770642&767024&3618\\
(7,5,5)&1865215&1031261&1022501&8760\\
(7,6,4)&1983275&1349846&1329217&20629\\
(8,3,6)&3250246&1469720&1464623&5097\\
(8,4,5)&3335144&1767006&1757445&9561\\
(8,5,4)&2670565&1815569&1800165&15404\\
(9,3,5)&2957131&194126&194125&1\\
(9,4,4)&1619622&170706&170692&14\\
\end{array}
\]
The two order-nine rows in this table are the full $8044$
rectangular-band-source portions; each is joined to both
$C_2$-Rees source portions. Opposite $P_5$ cases omitted
from a direct row are joined to the already completed
opposite $P_5$-source row. Thus the table is used with the
explicit dependencies of the scope proposition, rather
than as an assertion about a smaller source catalog.

For every action orbit not removed by a proved gate or
an explicit proper retraction, the separate checker
recomputes the first and second mod-two homology filters.
Every retained table has first and second Betti numbers
$0,1$. The literal checker independently reconstructs its
typed composition table, tests all identity and associativity
equations, recomputes those two ranks, and verifies a
nonzero second cocycle. It then checks every equation
of the supplied square primitive or recomputes a positive
third Betti number. Each outcome excludes $\CP^2$.

A further exact linker reconstructs every retained
composition word from the full left/right-action streams
and matches it and its exclusion reason to the literal
certificate, with simultaneous end of file. This prevents
an unrelated valid certificate stream from substituting
for the categories actually enumerated.

The complete manifest
\path{monoid/certificates/two_object17_manifest.json}
joins all $21$ residual positive-source rows of the exact
integer scope, including all $21$ distinct new finite
pieces and the separately proved analytic and prior finite
rows. It checks every count equality and every finished
stream hash. Across the new finite pieces there are
26752578 biset state orbits and
11965998 retained literal tables;
11875860 have zero square and
90138 have positive third homology.
No unexcluded table remains. Immutable independent
receipts are in \path{monoid/certificates/seventeen_receipts/}.
The distinct stopped guard and superseded numerical
prefixes are explicitly marked unused and contribute no
scope assertion. All seventeen required profiles are now
excluded, completing the proof.
\end{proof}


\subsection{Three connecting arrows and constant minimal actions}

\begin{lemma}[Proved: a proper biset retract of a projective plane]
\label{CP2:A-biset-retract}
Let $T\subseteq S$ be nonempty finite left-$N$, right-$M$ bisets,
and let $r:S\to T$ be an equivariant retraction. Write $X_S,X_T$
for their two-object collage spaces, with the same corner monoids.
If $X_S\simeq\CP^2$, then $X_T\simeq\CP^2$. If $T\ne S$,
the latter category has strictly fewer arrows.
\end{lemma}
\begin{proof}
The inclusion and retraction are actual functors, identity on both
corners, and give a retraction $X_S\to X_T$. Use the two-sided
action-bar double mapping cylinders as CW models. Their vertex
spaces are identical and the edge inclusion $E_T\subseteq E_S$
is a simplicial subcomplex. Collapsing $X_T$ therefore gives
$X_S/X_T=\Sigma(E_S/E_T)$, with the evident pointed quotient.
All its positive-degree cohomology products vanish.

If $H^2(X_T;\mathbb Q)=0$, the generator of $H^2(X_S;\mathbb Q)$
lifts from this suspension quotient and has square zero, a
contradiction. Retraction makes every integral homology group
of $X_T$ a direct summand of the corresponding group of $X_S$.
Thus $H_2(X_T;\mathbb Z)=\mathbb Z$, and retraction induces
an integral isomorphism in degree two. Its generator has
primitive nonzero square, since its pullback is a generator
on $X_S$. Consequently $H_4(X_T;\mathbb Z)=\mathbb Z$ and
the degree-four pullback is an isomorphism as well. All other
positive homology groups vanish. The fundamental group of
$X_T$ is a retract of the trivial group, hence trivial.
The projective-plane recognition theorem used in the upper
construction now gives $X_T\simeq\CP^2$. Only the cross-hom
set changed, so the arrow difference is $|S|-|T|$.
\end{proof}

\begin{lemma}[Proved: the complete three-state image classification]
\label{CP2:A-three-state-images}
Two commuting transformation monoids on a three-element set
satisfy at least one of these alternatives: a common fixed
point; both images consist of permutations; one image is
trivial; a proper idempotent transformation commutes with
both images; or, after simultaneous state relabelling and
possibly exchanging the images, they are
\[
 \{1,t_0,t_1\},\qquad \{1,c_0,c_1\},
\]
where $t_i$ fixes $0,1$ and sends $2$ to $i$, while $c_i$
is constant at $i$.
\end{lemma}
\begin{proof}
We provide a fresh complete finite census and a separate
enumeration. Encode a map by $f(0)+3f(1)+9f(2)$, so every
one of the $27$ maps occurs. The producer
\path{code/three_state_images.py} starts with the identity,
adjoins each absent map in turn, and takes multiplication
closure. Every submonoid is reached by successively adjoining
its elements; equality of the complete bit sets removes only
duplicates. This gives exactly $699$ submonoids.

The independent C++ checker
\path{code/check_three_state_images.cpp} uses included and
excluded bit sets instead. It propagates the Horn implications
$a,b\in P\Rightarrow ab\in P$ and their two contrapositives
when the product is excluded, then branches on the least
undecided map. These are necessary closure conditions, and
both branches are tried, so induction on the decision tree
proves coverage. Its $1463$ nodes give exactly the same literal
set of $699$ submonoids; it does not import the producer.

Testing commutation for every ordered pair leaves $1854$
pairs. In the stated order of alternatives the disjoint counts
are $595,5,1158,90,6$. For the fourth alternative every map
is tested for idempotence, proper image and commutation with
every member of both monoids. The checker independently
constructs all displayed residual pairs by choosing the two
core states and exchanging the sides. Their literal set is
exactly the six remaining pairs. Full maps, monoid masks,
pair classifications, witnesses and independent receipts are
in \path{outputs/three_state_images/}. This classifies images
only; no original corner is replaced by its image.
\end{proof}

\begin{lemma}[Proved: a permutation biset has zero second products]
\label{CP2:A-permutation-biset}
If both actions on a nonempty finite biset are permutations,
all rational degree-two products of its collage vanish.
\end{lemma}
\begin{proof}
The two-sided edge bar $E$ projects to $BM\times BN$ as the
finite covering associated to the commuting permutation
actions. One may construct it from the action category of
the product monoid: every transition is a permutation, so
each simplex has exactly $|S|$ compatible lifts, and the
usual face and degeneracy identifications give a covering.
The bisimplicial two-sided bar and this diagonal action bar
have naturally equivalent realizations over $BM\times BN$.
Summing the lifts defines the chain transfer, whose composite
with pullback in rational cohomology is multiplication by
$|S|$. Thus pullback from the product is injective.

Finite monoids have zero rational first cohomology. Kunneth
therefore identifies $H^2(BM\times BN;\mathbb Q)$ with the
direct sum of the two corner second-cohomology groups.
The Mayer--Vietoris compatibility of the two restrictions
of a collage class says that their difference pulls back
to zero on $E$. Injectivity forces both restrictions to
be zero. The relative-cup lemma proves the assertion.
\end{proof}

\begin{theorem}[Proved: compression when the minimal ideal acts constantly]
\label{CP2:A-constant-act-compression}
Let $K$ be the proper minimal ideal of a finite monoid $M$,
and let $I$ be its set of minimal right ideals. Suppose
every element of $K$ acts constantly on a nonempty finite
right $M$-set $S$. There is a finite monoid $N$, with
minimal ideal the rectangular band $I\times S$, such that
\[
 |N|=|M|-|K|+|I||S|,
 \qquad BN\simeq B\mathcal C(M,S).
\]
Here the act cone has source corner $M$, trivial sink
corner, and cross-hom $S$. No triviality assumption is
made on the maximal group of $K$.
\end{theorem}
\begin{proof}
Write $c_k$ for the constant value of the actor $k$.
The right action law gives $c_{k\ell}=c_\ell$ for
$k,\ell\in K$. The minimal-ideal coordinates show that
$c_k$ depends only on its column: left multiplication
by $K$ traverses all row and group coordinates in a
fixed column. Let $\rho(k)$ be its row, and let
$\sigma_m:I\to I$ be the row map induced by left
multiplication by $m$. Then
\[
 \eta(k)=(\rho(k),c_k),\quad
 \eta(mk)=(\sigma_m\rho(k),c_k),\quad
 \eta(km)=(\rho(k),c_km).
\]
Define $N=(M\setminus K)\sqcup(I\times S)$. For old
elements retain their product when it is outside $K$
and otherwise take its $\eta$ image; set
\[
 m(i,s)=(\sigma_m(i),s),\qquad
 (i,s)m=(i,sm),\qquad
 (i,s)(j,t)=(i,t).
\]
The two displayed equivariance equations and the original
associativity check every product involving the new band;
products of three old elements reduce to the corresponding
old associative product and its image. The old identity
is the identity. Thus $q:M\to N$, equal to $\eta$ on
$K$ and the identity on its complement, is a unital
homomorphism. The new band is an ideal and is simple,
and every nonempty ideal meets it, so it is the minimal
ideal of $N$.

Fix $i_0,s_0$ and $e=(i_0,s_0)$. In the full subcategory
$D$ of the Karoubi envelope on $1,e$, the four hom-sets are
\[
 D(1,1)=N,\quad D(e,e)=\{e\},\quad
 D(1,e)=\{i_0\}\times S,\quad
 D(e,1)=I\times\{s_0\}.
\]
Composition is multiplication in $N$, in chronological
categorical order. The idempotent object is a retract
of $1$, hence $BD\simeq BN$. The functor
$F:\mathcal C(M,S)\to D$ sends the two objects to
$1,e$, the source corner by $q$, and $s$ to $(i_0,s)$.
The action equations verify that it preserves composition.
The comma category $(F\downarrow e)$ has terminal
object given by the sink and its identity.

We describe the other comma category, not merely its
homology. Its objects over the source are $n\in N$;
its objects over the sink are $*_i$, $i\in I$.
An arrow $n\to n'$ is labelled $m\in M$ with
$n'q(m)=n$. There is one arrow $(i,s)\to *_i$
labelled $s$, and no other nonidentity arrow involving
a star as source. For a right act $T$, write $A(T)$
for the category whose arrows $x\to y$ labelled $m$
satisfy $ym=x$. Its nerve in degree $d$ is naturally
$T\times M^d$, and so preserves set pushouts and
inclusions. The right-act pushout
$N=M\sqcup_K(I\times S)$ shows that the realization
of our comma category is
\[
 BA(M)\ \sqcup_{BA(K)}\
       \coprod_{i\in I}\operatorname{Cone}(BA(S)).
\]
Indeed each row is a sieve, so a chain ending at $*_i$
has every previous object in that row. There are no
additional mixed simplices. The inclusion $BA(K)\to
BA(M)$ is a simplicial cofibration. Each component
of $BA(K)$ is $BA(e_iM)$ for a minimal idempotent $e_i$.
It is contractible: the arrows $x\to e_i$ labelled
$x$ give a natural transformation to the constant
functor; its naturality is the equation $ym=x$.
The row projection $BA(K)\to I$ is therefore a
homotopy equivalence. Also $A(M)$ has terminal object
$1$. Contracting the cones identifies the displayed
homotopy pushout with a contractible space. Both comma
categories are contractible; Quillen's Theorem A gives
$BF$ a homotopy equivalence and proves the claim.
\end{proof}

\begin{lemma}[Proved: constant-action quotient detects second cohomology]
\label{CP2:A-constant-quotient-H2}
For a finite monoid $M$ with minimal ideal $K$ and nonempty
finite right act $S$, impose the equivalence relation generated
by $sk\sim s'k$ for $s,s'\in S$, $k\in K$. The quotient
$\bar S$ is a right act on which every $k$ is constant, and
the induced act-cone functor gives a surjection
\[
 H^2(B\mathcal C(M,\bar S);\mathbb Q)
 \longrightarrow H^2(B\mathcal C(M,S);\mathbb Q).
\]
If $|S|\leq3$ and some $K$ actor is nonconstant, every
rational degree-two product on the original cone is zero.
\end{lemma}
\begin{proof}
The relation is invariant since $km\in K$. Put
$R=\mathbb QM$, $L=\ker(\mathbb QS\to\mathbb Q)$.
Its relation subspace is exactly $LK$, because it is
spanned by $([s]-[s'])k$; differences along a chain of
generating identifications add, so taking equivalence
closure adds no further linear relation. Thus
$\bar L=\ker(\mathbb Q\bar S\to\mathbb Q)=L/LK$.

Average the maximal group of a minimal idempotent to
obtain the projective left row module $P=Rf\cong\mathbb QI$.
Its augmentation sequence is
$0\to V\to P\to\mathbb Q\to0$, with $KV=0$.
Write $D$ for rational duality. Every $R$-linear map
$V\to DL$ annihilates $LK$: for $v\in V$, $z\in L$,
$k\in K$, its value on $zk$ equals its value at $kv$
on $z$, hence is zero. Consequently
\[
 \operatorname{Hom}_R(V,D\bar L)
   =\operatorname{Hom}_R(V,DL).
\]
The reduced act-cone bar is $(L\otimes_R^{\mathbf L}\mathbb Q)[1]$.
Duality and the projective augmentation sequence identify
its second cohomology with the cokernel of
$\operatorname{Hom}_R(P,DL)\to\operatorname{Hom}_R(V,DL)$.
The quotient act has the same numerator and possibly a
smaller denominator. Naturality proves the stated surjection.
If a $K$ actor is nonconstant, at least two distinct states
are identified, so $|\bar S|\leq2$ when $|S|\leq3$.
Every action on two states extends affinely to the interval,
and its cone products vanish by Lemma~\ref{CP2:A-high-affine}.
The same is true for a singleton. Surjectivity and cup
naturality transfer this conclusion to the original cone.
\end{proof}

\begin{theorem}[Proved: the exact three-state act-cone minimum]
\label{CP2:A-three-state-minimum22}
An ordinary three-state act cone of type $\CP^2$ has at
least twenty-two arrows. The certified twenty-two-arrow
models attain this bound.
\end{theorem}
\begin{proof}
Suppose its monoid has order $m\leq17$. Nonzero degree-two
square requires $H^2(BM;\mathbb Q)\ne0$: otherwise all
second classes lift from the suspension quotient of $BM$.
Writing $|K|=rsh$ for the minimal row, column and group
counts, the rational minimal-group collapse and balanced
tensor formula force $r,s\geq2$. In particular $K$ is
proper, so $q=m-rsh\geq1$.
Lemma~\ref{CP2:A-constant-quotient-H2} forces constant
$K$ action. The compression theorem gives a homotopy-equivalent
monoid with order $m-rsh+3r$ and minimal band $r\times3$.
If $sh\geq3$, its order is at most $m<23$, impossible.
Thus $s=2,h=1$, and the compressed order is $m+r$, where
$r\leq\lfloor(m-1)/2\rfloor$.

For $m\leq15$ this is at most twenty-two. For $m=16$,
order at least twenty-three forces $r=7$, but the compressed
minimal band $7\times3$ differs from both $3\times5$ and
$5\times3$ of the complete minimum-monoid classification.
For $m=17$ the remaining rows are $r=6,7,8$. At $r=6$
the compressed monoid again has order twenty-three with
the wrong minimal-band shape. At $r=7,8$ its contracted
rational algebra $A$ has dimension $q=17-2r=3,1$, and
the row augmentation module has dimension $r-1=6,7$.
The target band--Tor dimensions would be $1,0,1,0,\ldots$.
If $A$ is local, tensor dimension one makes both modules
cyclic by Lemma~\ref{A:local-target-facts}; a cyclic module
cannot have dimension greater than $\dim A$, a contradiction.
If $A$ is nonlocal, Lemma~\ref{A:nonlocal-four} makes nonzero
degree-two Tor persist in arbitrarily high degrees, also
a contradiction. These exhaust all $m\leq17$. Adding
the three states and the sink identity proves the bound.
The already verified upper models have $m=18$ and three
states, proving equality.
\end{proof}

\begin{theorem}[Proved: three connecting arrows reduce to an act cone]
\label{CP2:A-three-cross-reduction}
If a two-object collage with exactly three cross arrows is
of type $\CP^2$, one of its actions is trivial, and its
ordinary three-state act-cone retract is also of type $\CP^2$.
Consequently it has at least twenty-two total arrows,
or at least twenty-three if both corners are nontrivial.
\end{theorem}
\begin{proof}
Apply the complete transformation-image classification.
A common fixed state supplies natural transformations
making both corner inclusions null-homotopic; the
relative-cup lemma excludes this case. Permutation images
are excluded by Lemma~\ref{CP2:A-permutation-biset}.
A proper commuting idempotent gives a biset retract
with at most two states. By Lemma~\ref{CP2:A-biset-retract}
it would still be a projective plane, contrary to the
universal two-cross lemma (or the singleton join).
The residual image pair is excluded by the universal
two-constant theorem, Theorem~\ref{CP2:A-two-constant-image}.

Only a trivial action side remains. Collapsing that
corner is an actual functor to the ordinary act cone,
and adjoining its identity defines a section. By
Lemma~\ref{CP2:A-trivial-biset-side-H2} the induced
map is a rational second-cohomology isomorphism. The
retract therefore retains the nonzero second class.
Its integral homology is a direct summand of that of
the projective plane; in degree two the split injection
is an isomorphism, and its primitive square forces
the same in degree four. Its fundamental group is
trivial by retraction. The recognition theorem gives
the asserted type. Apply the exact three-state act
minimum. Collapsing a corner of order greater than
one saves at least one arrow, proving the last clause.
\end{proof}


\subsection{Small act cones by quotient modules}

\begin{lemma}[Proved: at most two outside elements]
\label{CP2:A-two-outside-acts}
Let $M$ be a finite monoid with minimal ideal $K$ and
$|M\setminus K|\leq2$. For every finite-dimensional right
$R=\mathbb QM$-module $L$, eventual vanishing of
$\operatorname{Tor}^R_j(L,\mathbb Q)$ implies its vanishing
for every $j\geq3$. Thus a finite act cone over $M$ with
bounded rational homology has zero fourth rational homology.
\end{lemma}
\begin{proof}
If $K=M$, the monoid is a finite group and Maschke's
averaging makes its rational algebra semisimple. Otherwise
put $A=R/\mathbb QK$. At a minimal idempotent $e$, with
maximal group $G$, the ideal factorization is
$\mathbb QK=Re\otimes_{\mathbb QG}eR$.
The second factor is a free left $\mathbb QG$-module,
so this is projective as a left $R$-module. Hence
$\operatorname{pd}_R A\leq1$.
Averaging $G$ gives a projective row module $P$ and
$0\to V\to P\to\mathbb Q\to0$, with $KV=0$.
For $j\geq2$, dimension shifting identifies
$\operatorname{Tor}_j^R(L,\mathbb Q)$ with
$\operatorname{Tor}_{j-1}^R(L,V)$.

A unital rational algebra of dimension at most two is
semisimple or the dual numbers. In the first case $V$
is an $A$-projective module and has $R$-projective
dimension at most one. In the second case
$V=A^a\oplus k^b$, where $k=A/(\varepsilon)$.
The exact sequence $0\to k\to A\to k\to0$ and
$\operatorname{pd}_R A\leq1$ give
\[
 \operatorname{Tor}_j^R(L,k)
   \cong\operatorname{Tor}_{j-1}^R(L,k),\qquad j\geq3.
\]
If $b=0$ the assertion is immediate; otherwise the
eventually zero Tor for $V$ makes the displayed groups
eventually zero, hence zero already in degree two.
Shifting back proves the assertion. For an act cone
use $L=\ker(\mathbb QS\to\mathbb Q)$ and its reduced
Tor formula, with the one-degree suspension.
\end{proof}

\begin{lemma}[Proved: a connected graph with one cycle]
\label{CP2:A-one-cycle-graph-cofiber}
Let a finite monoid $M$ act cellularly on a connected
finite graph $\Gamma$, with $H_1(\Gamma;\mathbb Q)=\mathbb Q$.
For $X=\operatorname{hocofib}(\Gamma_{hM}\to BM)$,
\[
 \widetilde H_n(X;\mathbb Q)
  =\operatorname{Tor}_{n-2}^{\mathbb QM}
           (H_1(\Gamma;\mathbb Q),\mathbb Q),\qquad n\geq2.
\]
Nonzero $H_2(X;\mathbb Q)$ forces every actor to have
degree one on the graph cycle. If a commuting graph
self-map has degree $d\ne0,1$, all positive rational
cohomology products on $X$ vanish.
\end{lemma}
\begin{proof}
The augmented graph complex, with its actual right module
structure, has only its first homology as nonzero reduced
homology. Inclusion of that kernel is an equivariant
quasi-isomorphism. Tensoring a free left monoid resolution
and taking the cofiber shifts this kernel by two, proving
the formula. Its degree-two group is the coinvariant space
of the one-dimensional cycle module. Any actor degree
different from one annihilates that coinvariant space.
A commuting degree-$d$ map acts by $d$ on this kernel,
hence on all positive cofiber homology and cohomology.
For positive classes $u,v$, naturality gives $d(uv)=d^2uv$,
which proves the last assertion over $\mathbb Q$.
\end{proof}

\begin{proposition}[Proved: four states opposite a right-zero corner]
\label{CP2:A-right-zero-four-states}
Let $R_3=\{1,p,q\}$ with $ab=b$ for $a,b\in\{p,q\}$.
If a left-$R_3$, right-$M$ biset with at most four states
has collage of type $\CP^2$, an ordinary right-$M$ act
cone on at most four states also has that type.
Consequently no such collage exists when $|M|\leq12$.
For any three-element left corner, the profile
$(|M|,3,4)=(9,3,4)$ is excluded as well.
\end{proposition}
\begin{proof}
One or two states are excluded by the singleton join and
the universal two-cross theorem. Three states use
Theorem~\ref{CP2:A-three-cross-reduction}. Suppose there
are four. The two left actors $P,Q$ are retractions onto
the same nonempty subset $U$, by $PQ=Q$ and $QP=P$.
The natural graph model of the left action category has
vertices $U$ and one edge for each $s\notin U$, directed
from $Ps$ to $Qs$. Its construction and equivariance were
proved in Lemma~\ref{CP2:A-high-P5-graph}. A commuting
right actor maps an edge at its unit parameter to the
edge of the image state, or collapses it to a vertex.
Since $BR_3$ is contractible, the collage is the cofiber
of this graph's homotopy-orbit map to $BM$.

If $|U|=1$, its unique point is fixed by both sides,
and the relative-cup obstruction applies. If $P=Q$,
including $|U|=4$, the common map is a biset retraction
onto $U$. The retract is equivalent to the ordinary
$M$-act cone on $U$: the left side acts trivially and
its contractible classifying space can be collapsed
in the product edge bar and at its vertex. By
Lemma~\ref{CP2:A-biset-retract} this retract retains
the projective-plane type.

If $|U|=3$ and $P\ne Q$, the single graph edge joins
two distinct vertices. Its two components are
contractible. The equivariant map to its component
set is a homotopy equivalence, including after taking
the monoid bar. Its cofiber is therefore an ordinary
two-state act cone, already excluded.

The remaining graph has two vertices and two edges.
If it is disconnected, both edges are loops and
$P=Q$, already treated. If connected it has one
cycle. A putative projective plane forces every
actor to have degree one by the preceding lemma.
For one loop and one connecting edge, this forces
the unique loop to map to itself; its vertex is
a common fixed state. For two connecting edges
with the same direction, a degree-one actor fixes
both vertices and edges: swapping the edges has
degree minus one and folding or collapsing gives
degree zero. Again there is a common fixed state.
If the two edges have opposite directions, the
only degree-one actors are the identity and the
half-turn. Identify this graph with the circle
so that the half-turn is $z\mapsto-z$. The map
$z\mapsto z^3$ commutes with both actors. The
preceding degree lemma kills all positive products.
This exhausts the four-state graph types.

The resulting ordinary cones on one or two states
have zero products; on three states use the exact
minimum twenty-two; on four states use
Corollary~\ref{CP2:A-four-state-small}. This proves
the assertion for $|M|\leq12$.
Finally, every three-element monoid other than
$R_3$ has projective trivial right rational module.
Indeed its minimal ideal is a group, has one column,
or is the two-column one-row band with an identity;
the last case is exactly $R_3$. In the first two
cases finite-group averaging supplies the splitting
idempotent. Exact orbit comparison thus replaces
the left corner of $(9,3,4)$ by a trivial corner,
preserving rational homology and the cohomology ring.
The bounded-nine theorem excludes that orbit act
cone. This proves the last assertion without a
new multiplication census.
\end{proof}

\begin{lemma}[Proved: the two-column quotient has projective kernel]
\label{CP2:A-two-column-orbit-kernel}
Let the minimal ideal of $M$ be the band $I\times\{0,1\}$,
and suppose a finite nonempty right act $S$ has no point
fixed by the whole band. Let $o$ be the number of core
orbits in $SK$, and let $\bar S$ be the constant-action
quotient of Lemma~\ref{CP2:A-constant-quotient-H2}.
Over $R=\mathbb QM$ there is an exact sequence
\[
 0\to(eR)^{o-1}\to L_S\to L_{\bar S}\to0,
 \qquad L_T=\ker(\mathbb QT\to\mathbb Q).
\]
Every core orbit has two states, and
$|\bar S|=|S|-2(o-1)$.
Writing $X,Y$ for the two act cones, their actual quotient
map induces isomorphisms in rational homology in degrees
at least three and an exact sequence
\[
 0\to H_2(X)\to H_2(Y)\to\mathbb Q^{o-1}
       \to H_1(X)\to H_1(Y)\to0.
\]
If $X$ has the rational homology and ring of $\CP^2$,
then $H_2(Y)=\mathbb Q^o$, $H_4(Y)=\mathbb Q$, and
all its other positive rational homology vanishes.
Its degree-two cup pairing has rank exactly one.
\end{lemma}
\begin{proof}
A core orbit is a quotient of the two-point column act.
It cannot be a singleton by the hypothesis, so label it
$\{s_{a,0},s_{a,1}\}$, $0\leq a<o$. Every outside actor
has the same column map on every such orbit:
$s_{a,j}m=s_{a,\tau_m(j)}$. In particular the orbits
are individually $M$-invariant.

The kernel $L_SK$ of the quotient map has basis
$s_{a,j}-s_{0,j}$ for $1\leq a<o$, $j=0,1$.
Each is a generating relation after multiplication
by a band element in column $j$. Conversely every
generating relation is a difference of column-$j$
points, so this is the whole kernel. Its action formula
identifies each pair with $eR$, the projective two-point
row ideal. Noncore states are never identified, giving
the cardinality formula.

Apply the Tor exact sequence and
$eR\otimes_R\mathbb Q=\mathbb Q$ to the reduced
act-cone chain formula. This proves both the higher
isomorphisms and the stated low-degree sequence.
For a rational projective plane, $H_1(X)=0$ gives
the asserted dimensions, and the higher groups
transfer unchanged. Pullback is surjective in degree
two and an isomorphism in degree four. Naturality
therefore identifies the cup form on $Y$ with the
pullback of the nonzero rank-one form on $X$.
Finally compression gives $Y\simeq BN$ with minimal
band $I\times\bar S$ and the same outside elements
as $M$, so the band--Tor formula applies to this
actual quotient without any separate boundedness
assumption on $BM$.
\end{proof}

\begin{lemma}[Proved: a two-dimensional module over a small local algebra]
\label{CP2:A-local-three-small-module}
Let $A$ be a local rational algebra with $\dim A\leq3$.
If one of two finite modules $W,V$ has dimension at most
two, eventual vanishing of $\operatorname{Tor}^A_j(W,V)$
implies $\operatorname{Tor}^A_2(W,V)=0$.
\end{lemma}
\begin{proof}
By opposition suppose $\dim V\leq2$. If $V$ is
semisimple, it is a sum of copies of the residue
division algebra. A nonprojective $W$ then has
nonzero Tor in every positive degree by
Lemma~\ref{A:local-target-facts}; a projective $W$
has none. If $V$ is not semisimple, its dimension
is exactly two, its residue division algebra is
$\mathbb Q$, and its top has dimension one.
Thus it is cyclic. When $\dim A=2$ its cyclic
cover is an isomorphism. When $\dim A=3$ the
kernel of that cover is the one-dimensional simple
module $k$. For $j\geq2$, dimension shifting gives
$\operatorname{Tor}_j(W,V)=\operatorname{Tor}_{j-1}(W,k)$.
The same local-algebra fact excludes eventual
vanishing with a nonzero degree-two group. These
cases exhaust the possible dimensions.
\end{proof}

\begin{lemma}[Proved: the local rank-two additive obstruction]
\label{CP2:A-local-four-rank-two}
Let $A$ be local of dimension four, with
$\dim V=2$ and $\dim W\leq3$. There is no module
pair with bounded Tor, tensor dimension two,
$\operatorname{Tor}_1(W,V)=0$ and
$\operatorname{Tor}_2(W,V)\ne0$.
\end{lemma}
\begin{proof}
If $V$ is semisimple, boundedness and
Lemma~\ref{A:local-target-facts} force $W$ projective,
contrary to nonzero positive Tor. Otherwise the residue
field is $\mathbb Q$ and $V$ is cyclic of dimension
two. Its cyclic projective cover has nonzero kernel
$J_V$ of dimension two. The Tor exact sequence gives
\[
 0\to W\otimes_AJ_V\to W\to W\otimes_AV\to0,
 \qquad \dim(W\otimes_AJ_V)=\dim W-2.
\]
The tensor product of two nonzero modules over a local
finite algebra is nonzero, by its quotient onto the
tensor of their tops. Hence $\dim W=3$ and the
displayed tensor dimension is one. The same top
argument makes $W$ cyclic. Its cyclic cover then
has one-dimensional simple kernel. Consequently
$\operatorname{Tor}_j(W,V)=\operatorname{Tor}_{j-1}(k,V)$
for $j\geq2$, which is nonzero in arbitrarily high
degrees because $V$ is not projective. This contradicts
boundedness.
\end{proof}

\begin{theorem}[Proved: bounded act cones through monoid order nine]
\label{CP2:A-bounded-nine-acts}
For every finite monoid of order at most nine and every
nonempty finite right act, bounded rational homology
of its act cone forces every rational degree-two
product to vanish. There is no bound on the number
of states or on the second Betti number. The left-act
version holds by opposition.
\end{theorem}
\begin{proof}
If $H^2(BM;\mathbb Q)=0$, use the suspension quotient.
If there are at most two minimal rows, use
Lemma~\ref{CP2:A-high-two-row-act}. In the remaining
cases the minimal ideal has at least three rows and
two columns. Cardinality at most nine forces trivial
maximal group and either a $3\times2$ band with
at most three outside elements or a $4\times2$
band with one outside element. Two or fewer outside
elements are covered by Lemma~\ref{CP2:A-two-outside-acts}.

Only the $3\times2$ band with three outside elements
remains. A point fixed by all of $K$ is fixed by $M$:
if $ze=z$, then $zm=z(em)=z$, since $em\in K$.
It gives a section of the action bar and makes
all second classes relative, so their products vanish.
Otherwise the orbit-kernel lemma and compression give
a monoid $N$ with bounded rational homology and
$H_4(BN;\mathbb Q)=H_4(X;\mathbb Q)$. Its contracted
algebra $A$ has dimension three and its row
augmentation module has dimension two. For nonlocal
$A$, Lemma~\ref{A:nonlocal-four} forces degree-two
Tor to vanish under boundedness. For local $A$ use
Lemma~\ref{CP2:A-local-three-small-module}. The
band--Tor formula therefore gives $H_4(X;\mathbb Q)=0$.
This also annihilates every degree-two product.
\end{proof}

\begin{proposition}[Proved: ten-element corners and at most six states]
\label{CP2:A-ten-six-state-rational}
An act cone over a monoid of order at most ten, with
at most six states, cannot have the rational homology
and cohomology ring of $\CP^2$.
\end{proposition}
\begin{proof}
The preceding theorem handles orders through nine.
The relative and two-row obstructions leave, at order
ten, band shapes $3\times2$, $4\times2$, $3\times3$.
The last two have at most two outside elements and
are excluded by Lemma~\ref{CP2:A-two-outside-acts}.
Thus $K=3\times2$ and there are four outside elements.
A band-fixed state is a global fixed state and gives
zero products as above. Otherwise let $o$ count the
two-point core orbits. The act must be connected:
the cofiber exact sequence identifies $H_1(X;\mathbb Q)$
with the reduced zeroth homology of its action category,
since $H_1(BM;\mathbb Q)=0$. Because each core orbit
is invariant, $o=3$ at six or fewer states would have
$S=SK$ and a disconnected act. Hence $o=1$ or $2$.

The quotient compression has a four-dimensional
contracted algebra $A$, row module $V$ of dimension
two, and column module $W$ of dimension
$|S|-2(o-1)-1$. Its required Tor dimensions are
$(o,0,1,0,\ldots)$. Nonlocal $A$ is excluded by
Lemma~\ref{A:nonlocal-four}. If $A$ is local and
$o=1$, then $\dim W\leq5$, and
Lemma~\ref{A:local-four-small-target} excludes the
required four initial Tor dimensions. If $o=2$,
then $\dim W\leq3$, and
Lemma~\ref{CP2:A-local-four-rank-two} applies.
This excludes the stated rational target in all cases.
\end{proof}

\begin{lemma}[Proved: four-state band acts with zero first homology]
\label{CP2:A-four-state-band}
Let $M$ have a rectangular-band minimal ideal, and
let $S$ be a right act with at most four states.
If $H_1(B\mathcal C(M,S);\mathbb Q)=0$, either all
rational degree-two products vanish or every band
actor is constant on $S$.
\end{lemma}
\begin{proof}
A band-fixed point is globally fixed and gives zero
products. Otherwise every core orbit has at least
two points. Every such orbit is individually invariant
under $M$, since $yK$ is invariant and contains $y$
for a core point. Two or more core orbits at four
states force exactly two two-point orbits and $S=SK$.
The action category would be disconnected, contrary
to the cofiber first-homology sequence. Thus there
is just one core orbit.

On any right orbit of a rectangular band, each band
element acts constantly: writing a point as $yk$,
the value $(yk)\ell$ depends only on the fixed row
of the orbit and the column of $\ell$. For arbitrary
$s\in S$, idempotence gives $s\ell=(s\ell)\ell$,
and $s\ell$ lies in this single core orbit. Hence
the constant value is independent of $s$ as well.
\end{proof}

\begin{corollary}[Proved: four-state act cones through seventeen arrows]
\label{CP2:A-four-state-small}
An ordinary four-state act cone over a monoid of order
at most twelve cannot have type $\CP^2$.
\end{corollary}
\begin{proof}
Nonzero square requires nonzero corner $H^2$, at least
two minimal columns and at least three minimal rows,
by the relative and two-row results. A nontrivial
minimal group would cost at least $3\cdot2\cdot2+1=13$
elements. Thus the ideal is a band. The preceding
lemma forces constant action. Compression gives
a monoid of order
\[
 m-rs+4r\leq m+2r
 \leq m+2\lfloor(m-1)/2\rfloor\leq22
 \qquad(m\leq12).
\]
This contradicts the complete monoid minimum.
\end{proof}


\subsection{Eleven-element corners and the seventeen-arrow act gates}

\begin{lemma}[Proved: two-dimensional local module pairs]
\label{CP2:A-local-two-by-two}
Let $A$ be a finite-dimensional local rational algebra and let
$W,V$ be right and left modules of dimension two. They cannot have
\[
 (\dim\operatorname{Tor}_0^A(W,V),\ldots,
   \dim\operatorname{Tor}_3^A(W,V))=(1,0,1,0).
\]
Nor can they have bounded Tor with tensor dimension two,
zero first Tor, and nonzero second Tor.
\end{lemma}
\begin{proof}
If either module is semisimple, the local minimal-resolution
argument in Lemma~\ref{A:local-target-facts} makes its positive
Tor with a nonprojective opposite module nonzero in every
positive degree. A projective opposite module gives no positive
Tor. Thus both modules are nonsemisimple. Their residue division
algebra is then $\mathbb Q$, and both are cyclic. A cyclic
projective cover of $V$ has kernel $I$, of dimension $\dim A-2$.
We may suppose $I\ne0$, since otherwise $V$ is projective.
The first Tor vanishing gives
\[
 0\longrightarrow W\otimes_A I\longrightarrow W
   \longrightarrow W\otimes_A V\longrightarrow0.
\]
The tensor product of two nonzero modules over a finite local
algebra is nonzero: it surjects onto the tensor product of their
tops over the residue division algebra. This already excludes
tensor dimension two.

For tensor dimension one, $W\otimes I$ has dimension one, so
$I$ has one-dimensional top and is cyclic. Its cyclic cover
has kernel $L=\Omega^2 V$ of dimension two. If $L$ is semisimple,
then $\operatorname{Tor}_3(W,V)=\operatorname{Tor}_1(W,L)$
is nonzero, since $W$ is nonprojective. Otherwise $L$ is cyclic.
The first three projectives in a minimal resolution of $V$
are consequently $A$. After tensoring with $W$, the first two
differentials have ranks one and one. The required second Tor
dimension one forces the third differential to be zero.
The next projective is $A^t$ with $t>0$: otherwise $L$ is
projective, contrary to its dimension two and $\dim A>2$.
The fourth differential has image in $(WJ)^t$, a proper
subspace of $W^t$, by minimality. It cannot cover the whole
kernel of the zero third differential. Hence third Tor is
nonzero, a contradiction.
\end{proof}

\begin{lemma}[Proved: the five-dimensional two-by-two alternatives]
\label{CP2:A-five-algebra-two-by-two}
Let $A$ be a nonlocal rational algebra of dimension five,
$\dim W=\dim V=2$, and suppose its Tor is bounded,
$\operatorname{Tor}_1(W,V)=0$, and
$\dim\operatorname{Tor}_2(W,V)=1$.
If the tensor dimension is one, there is no such pair.
If it is two, the only remaining case in the classification
of Lemma~\ref{A:nonlocal-five-target} is the zero-composite
three-vertex line algebra, with
\[
 V=S_1\oplus S_3,\qquad W=T_1\oplus T_3.
\]
The two-dimensional space of balanced forms for this pair
has a basis consisting of matrices of rank one.
\end{lemma}
\begin{proof}
The persistence alternatives in that lemma contradict boundedness.
For $A=B\times\mathbb Q$, with $\dim B=4$ local, all positive
Tor is on the $B$ components. A one-dimensional nonzero
component forces persistence by the local minimal-resolution
argument. Thus both $B$ components have dimension two,
leaving no other components. Tensor dimension one is excluded
by Lemma~\ref{A:local-four-small-target}, and tensor dimension
two by Lemma~\ref{CP2:A-local-two-by-two}.

For the line algebra, use the five left and five right
indecomposables and their displayed resolutions in
Proposition~\ref{A:three-by-six-CP2}. Nonzero second Tor
requires a $T_1$ summand on the right and an $S_3$ summand
on the left. Each must occur exactly once. Neither remaining
one-dimensional summand can be the middle-vertex simple,
since that would give nonzero first Tor; repeating $T_1$
or $S_3$ would give second Tor dimension two. Thus the
pair is exactly the displayed pair. Its tensor dimension
is two, with one coordinate each at vertices1 and3.
The corresponding two coordinate functionals are rank-one
matrices and span the balanced-form space.

For the two-vertex exceptional algebra, the two-dimensional
left modules are
$S_1^{\oplus2},S_1\oplus S_2,S_2^{\oplus2},L,Ae_2$.
The first and last give no second Tor. For $S_2^{\oplus2}$
the second Tor dimension is even. For $S_1\oplus S_2$,
every right summand contributing to second Tor also gives
nonzero first Tor. With $V=L$, the only permissible
contributor is one copy of $T_2$: $T_1$ and $R$ give
nonzero first Tor. A right module of dimension two cannot
contain just one $T_2$ and a projective complement, since
the two indecomposable projectives have dimensions three
and two. Two copies of $T_2$ give second Tor dimension
two. This proves the exclusion. All these assertions also
follow by summing the independently checked literal
$5\times5$ Tor tables in
\path{monoid/certificates/five_dimensional_pairs.json}.
\end{proof}

\begin{proposition}[Proved: rational three-state cones through order eleven]
\label{CP2:A-eleven-three-state-rational}
An ordinary act cone over a monoid of order at most eleven,
with at most three states, cannot have the rational homology
and cohomology ring of $\CP^2$.
\end{proposition}
\begin{proof}
Orders through ten follow from
Proposition~\ref{CP2:A-ten-six-state-rational}.
Assume the corner order is eleven. The relative-cup and
two-row obstructions require at least three minimal rows
and two minimal columns. A nontrivial minimal group would
cost at least thirteen elements, so the ideal is a band.
If its action on the at most three states is not constant,
Lemma~\ref{CP2:A-constant-quotient-H2} pulls every second
class from a cone with at most two states, where every
second product is zero. We may therefore assume constant
band action. Fewer than three states are already excluded.

Compression gives a monoid with minimal band $I\times S$
and the same outside algebra $A$ as the original corner.
For at most two outside elements use
Lemma~\ref{CP2:A-two-outside-acts}. The only remaining
original band shapes and outside counts are
$3\times2$ with five outside elements and
$4\times2$ with three. In the second case the compressed
row and column augmentation dimensions are three and two;
Lemma~\ref{CP2:A-local-three-small-module} and
Lemma~\ref{A:nonlocal-four} exclude bounded nonzero
second Tor. In the first case $\dim A=5$ and the two
augmentation modules both have dimension two. The required
Tor dimensions are $(1,0,1,0,\ldots)$. The preceding two
lemmas exclude these dimensions for local and nonlocal
$A$, respectively. All comparisons are actual maps or
homotopy equivalences, so their cohomology pullbacks
preserve products.
\end{proof}

\begin{theorem}[Proved: eleven-element corners with five states]
\label{CP2:A-eleven-five-state-actual}
An ordinary act cone over a monoid of order at most eleven,
with at most five states, cannot have type $\CP^2$.
\end{theorem}
\begin{proof}
Orders through ten are excluded by the rational theorem;
through three states use the preceding proposition, and
four states use Corollary~\ref{CP2:A-four-state-small}.
Thus $|M|=11$ and $|S|=5$. The same row, column and
group budgets reduce to a rectangular minimal band.
A band-fixed state is globally fixed and kills second
products. Without such a state, the only shapes not
already excluded by at most two outside elements are
$3\times2$ with five outside elements and
$4\times2$ with three.

Let $o$ be the number of two-point core orbits. At five
states $o=1$ or2. For $o=1$, every band actor is constant:
its value on the sole core orbit is constant, and
idempotence determines its value on every transient state.
Compression therefore gives an actual monoid model of
order20 in the $3\times2$ case, contradicting the monoid
minimum. In the $4\times2$ case its order is23 and its
minimal band is $4\times5$, contrary to the complete
classification at order23.

For $o=2$, Lemma~\ref{CP2:A-two-column-orbit-kernel}
gives a compressed monoid $N$ with $H_2(BN;\mathbb Q)
=\mathbb Q^2$, $H_4(BN;\mathbb Q)=\mathbb Q$, no other
positive homology, and degree-two cup form of rank one.
Its column augmentation module has dimension two.
For the $4\times2$ shape, the outside algebra has
dimension three and the row module has dimension three;
the small-local and nonlocal-four lemmas again force
second Tor to vanish.

For the $3\times2$ shape, the outside algebra has
dimension five and both modules have dimension two.
Their Tor dimensions are $(2,0,1,0,\ldots)$.
The local case is excluded by
Lemma~\ref{CP2:A-local-two-by-two}; the nonlocal case
leaves only the line-algebra pair in
Lemma~\ref{CP2:A-five-algebra-two-by-two}.
Each of its two basis balanced forms has matrix rank one
and hence square zero by Lemma~\ref{A:rank-one-square}.
The symmetric cup matrix in this basis is therefore
$\left(\begin{smallmatrix}0&c\\c&0\end{smallmatrix}\right)$.
Over $\mathbb Q$ its rank is zero or two, never one.
This contradicts the quotient comparison's rank-one cup
form, and finishes every case.
\end{proof}

\begin{proposition}[Proved: the eleven-by-two four-arrow collage]
\label{CP2:A-eleven-two-four}
A two-object category with corner orders eleven and two
and four connecting arrows cannot have type $\CP^2$.
\end{proposition}
\begin{proof}
By opposition let the two-element corner act on the left
of a right-$M$ act $S$, where $|M|=11$ and $|S|=4$.
The two-element monoids are the group $C_2$ and
$\{1,e\}$ with $e^2=e$.
For the latter, put $T=eS$. Sending $S$ to $T$ and
both corner elements to the identity defines a functor
to the ordinary $M$-act cone on $T$. Its inclusion is a
section. The components $1_M$ and $e$ give a natural
transformation from the identity of the original category
to the inclusion followed by the quotient: on a cross
arrow the equation is precisely $es=es$.
Thus this is a homotopy equivalence, excluded by the
three-state rational or four-state actual theorem.

For the group $C_2$, exact rational orbit comparison
gives a rational homology and cohomology-ring isomorphism
to the ordinary $M$-act cone on $S/C_2$. If the action
is nontrivial, this orbit set has at most three elements,
contradicting Proposition~\ref{CP2:A-eleven-three-state-rational}.
If the action is trivial, the ordinary cone on all four
states is an actual categorical retract: quotient the
group corner to the identity, and include that identity
as a section. The quotient is still a rational homology
isomorphism. A retract of a simply connected space is
simply connected, and its integral homology groups are
direct summands of those of $\CP^2$. The rational
isomorphism consequently forces integral homology
isomorphisms in every degree. Naturality gives the
primitive square, so this retract has type $\CP^2$ by the same integral recognition argument as in
Theorem~\ref{A:monoid-upper23}. This contradicts
the four-state act theorem. Both monoids are excluded.
\end{proof}


\begin{corollary}[Proved: all act cones through seventeen arrows]
\label{CP2:A-act-cones-through-seventeen}
No ordinary finite-monoid act cone with at most seventeen
arrows has type $\CP^2$.
Every two-object category of at most seventeen arrows with one
corner of order two is excluded as well.
\end{corollary}
\begin{proof}
For an ordinary cone write $m$ for the corner order and $s$
for the positive number of states, so $m+s+1\le17$.
Orders through nine use the bounded-nine theorem; order ten
has $s\le6$ and uses the ten-state-budget rational theorem;
order eleven has $s\le5$ and uses
Theorem~\ref{CP2:A-eleven-five-state-actual}. For order twelve,
$s\le4$, use the four-state theorem or the three-state minimum.
For $m\ge13$, there are at most three states, so the universal
one- and two-state obstructions and the exact three-state
minimum twenty-two apply.

A disconnected category is excluded by connectedness. If both cross directions are nonempty, Proposition~\ref{CP2:A-baseline-strong-two} applies. In the remaining triangular case, take opposition if necessary. For a corner of order two, $m+s+2\le17$. One, two and three
connecting arrows are covered by the same universal cross-arrow
results. If $s\ge4$, then $m\le11$. Exact rational orbit
comparison at the order-two corner preserves the rational
homology and ring, with at most $s$ states. The bounded-nine
and ten-six-state theorems cover $m\le10$. The only remaining
budget is $(m,2,s)=(11,2,4)$, which is the preceding proposition.
\end{proof}


\subsection{Four-state act cones through corner order fourteen}

\begin{theorem}[Proved: four states and fourteen corner elements]
\label{CP2:A-four-state-fourteen}
Let $M$ be a finite monoid of order at most fourteen and $S$ a
finite right $M$-act with at most four states. The ordinary
two-object act-cone category $\mathcal C(M,S)$ cannot have
classifying space of type $\CP^2$. The left-act assertion holds
by opposition. In particular a four-state ordinary act cone
of type $\CP^2$ has at least twenty arrows, identities counted.
\end{theorem}
\begin{proof}
For at most three states use Theorem~\ref{CP2:A-three-cross-reduction}
and the three-state minimum; for four states and $|M|\leq12$
use Corollary~\ref{CP2:A-four-state-small}. It remains to consider
$|M|=m\in\{13,14\}$ and $|S|=4$.

Suppose the cone has projective-plane type. The relative-cup
obstruction requires nonzero rational second cohomology of $BM$,
and the two-row act obstruction requires at least three minimal
rows. Its minimal ideal therefore has $r\geq3$ rows and $c\geq2$
columns. If its maximal group has order at least two, the ideal
has at least twelve elements, leaving at most two outside
elements. Lemma~\ref{CP2:A-two-outside-acts} then gives zero
fourth rational homology from boundedness, a contradiction.
Thus the minimal ideal is the rectangular band $I\times J$,
of shape $r\times c$. Write $q=m-rc\geq1$ for the number of
outside elements; the identity is outside this proper ideal.

Lemma~\ref{CP2:A-four-state-band} says that either every rational
degree-two product is zero or the band acts constantly. The
first alternative is impossible. In the second, the actual
compression of Theorem~\ref{CP2:A-constant-act-compression}
gives a monoid $N$ of the same homotopy type, with minimal band
$I\times S$, the same $q$ outside elements, and
\[
 |N|=q+4r=m+r(4-c).
\]
If $c\geq3$, this is at most $m+\lfloor(m-1)/3\rfloor\leq18$
when $c=3$, and at most $m$ when $c\geq4$. The complete monoid
minimum excludes both. We may therefore assume $c=2$.

For $r=3$ or $4$, the compressed order is at most $22$, again
excluded. For $r\geq6$, the original outside count is at most
$14-12=2$, and the two-outside lemma applies directly to the
original act. Only $r=5$ remains. When $m=13$, the compressed
monoid has order $23$ and minimal band $5\times4$. This
contradicts the complete classification
Theorem~\ref{A:monoid-minimum23-classification}: the minimal-band
shapes of the two ordinary minimizing monoids are $3\times5$
and $5\times3$.

Finally let $m=14$, $r=5$, $c=2$, $q=4$. For the compressed
monoid put $A=\mathbb QN/\mathbb Q(I\times S)$. This is the
same contracted outside algebra as that of $M$, of dimension
four. Its row and column augmentation modules, respectively,
are the left $A$-module $V$ of dimension four and the right
$A$-module $W$ of dimension three. The band--Tor formula
would require
\[
 \dim(W\otimes_A V)=1,\qquad
 \dim\operatorname{Tor}_2^A(W,V)=1,\qquad
 \operatorname{Tor}_j^A(W,V)=0\quad(j\geq3).
\]
If $A$ is local, Lemma~\ref{A:local-target-facts} makes $V$
cyclic. Its cyclic cover $A\to V$ has equal source and target
dimensions and is therefore an isomorphism. This makes $V$
projective and contradicts the required second Tor. If $A$
is nonlocal, Lemma~\ref{A:nonlocal-four} says that nonzero
second Tor persists in arbitrarily high degrees, again a
contradiction. This exhausts the two remaining monoid orders.
Since a four-state ordinary act cone has $m+5$ arrows, the
last assertion follows. Every compression used above is an
actual homotopy equivalence; no rational comparison was
promoted to a homotopy equivalence.
\end{proof}


\subsection{A canonical biset census for the remaining linear profiles}

\begin{proposition}[Proved: exact biset and composition-map enumeration]
\label{CP2:A-linear-biset-enumeration}
Fix labelled corner monoids $M_0,M_1,M_2$ and positive
cross-hom sizes $x,y,z$ for a full-linear three-object category.
Enumerating the three connecting bisets separately up to their
state permutations, followed by all balanced and outer-equivariant
composition maps up to the target biset's automorphisms, covers
every category with those corner tables. No corner automorphism
is divided out. If the three state automorphism orders are
$A_S,A_T,A_U$ and a composition map has target stabilizer of
order $H$, its contribution to the completely labelled count is
\[
 \frac{x!}{A_S}\,\frac{y!}{A_T}\,\frac{z!}{A_U}\,
 \frac{A_U}{H}.
\]
The intermediate records are not claimed to be category
isomorphism classes.
\end{proposition}
\begin{proof}
Put $S=C(0,1)$, $T=C(1,2)$, and $U=C(0,2)$.
A biset is exactly a pair of commuting transformation actions,
with the source corner acting oppositely on the right.
For a single action, assign the first unknown actor every
transformation and propagate every product of assigned actors.
Reject only inconsistent equalities. This tries every action.
At each step choose the least representative under the state
permutations fixing all earlier actors, and then restrict that
group to the new actor's stabilizer. Every propagated actor is
a word in the earlier actors and the new actor, so this group
fixes all propagated values as well. The first unknown actor
index is invariant under state conjugation. Induction proves
exactly one representative of every action orbit and gives its
full stabilizer. For a biset, enumerate the second action in
the transformation centralizer of the first, using its state
stabilizer. The centralizer is closed under composition, so
propagation neither loses nor adds an allowed action.

The remaining data are exactly an equivariant function
\[
 \mu:T\otimes_{M_1}S\longrightarrow U.
\]
Construct the domain by imposing every literal relation
$(tm,s)\sim(t,ms)$. The outer actions descend by the commuting
biset identities. Assign the first unassigned domain class
every state of $U$ and propagate both outer actions, rejecting
only conflicting equalities. This produces every equivariant
function. At each step retain the least value under the current
pointwise stabilizer inside $\operatorname{Aut}(U)$ and then
fix that value. Since this group commutes with both target
actions, it fixes every propagated value too. The same induction
proves coverage and uniqueness up to $\operatorname{Aut}(U)$.
We do not divide composition maps by $\operatorname{Aut}(S)$
or $\operatorname{Aut}(T)$.

All associativity equations lie in a corner, involve one cross
arrow, or involve two successive cross arrows. These are,
respectively, the given corner axioms, the commuting biset
axioms, and the balancing and outer-equivariance of $\mu$.
There is no string of three cross arrows. Thus the construction
is both necessary and sufficient for a category.

Each labelled triple of bisets has a bijective set of possible
composition maps. Orbit--stabilizer gives $x!/A_S$, $y!/A_T$,
and $z!/A_U$ possible labellings of the three canonical actions;
a map orbit has $A_U/H$ members. Multiplying proves the formula.
Any invariant homology filter has the same weighted count.
\end{proof}

The executable implementation is
\path{code/linear17_bisets.cpp}; the subsequent explicit16/17
wrapper is \path{code/linear_bisets_16_17.cpp}.
The actual17-arrow input uses only corner orders at most seven,
for which the full ordinary catalogs have already been certified.
The implementation uses only the separately proved v2 corner
cuts. Its literal bar certificates are checked by the unchanged
independent checker. The separately written orbit comparator
\path{code/check_linear_biset_orbits.cpp} includes no producer code.

As a direct independent scope comparison, profile
$(1,5,4;3,3,1)$ gives exactly188608 fully labelled tables and
37584 weighted homology candidates, agreeing with the previous
typed-table DFS. Its7434 normalized table assignments give1062
retained categories. Exhausting every cross-hom relabelling in
the separate comparator proves exact equality of the two retained
canonical sets, not merely equality of their counts.
Profile$(3,1,7;1,2,3)$ independently gives7194510 labelled
tables and934200 weighted candidates in both implementations.
The old stream has102767 records and the new one102851;
the independent comparator proves equality of their102750
distinct canonical retained tables. This explicitly checks
that repeated records are not counted as distinct categories.
The new source and complete scope argument also received a
separate targeted proof check, recorded under
\path{acyclic/notes/linear_biset_targeted_review.md}.

\begin{theorem}[Proved: seven small-adjacent profiles at seventeen arrows]
\label{CP2:A-seven-small-adjacent-profiles}
No category has type $\CP^2$ with any of the following profiles
$(a,b,c;x,y,z)$, where $a,b,c$ are the three corner orders
and $x,y,z$ count $0\to1$, $1\to2$, $0\to2$, respectively:
\[
\begin{gathered}
 (3,1,3;2,2,6),\quad(3,1,3;2,3,5),\quad(3,1,3;2,4,4),\\
 (3,1,3;2,5,3),\quad(3,1,3;3,3,4),\quad
 (3,1,3;3,4,3),\quad(3,1,3;3,5,2).
\end{gathered}
\]
\end{theorem}
\begin{proof}
Use all seven ordinary monoids of order three at each endpoint,
including both orientations, and the trivial middle corner.
The preceding proposition proves complete category coverage;
the v2 cuts remove only the proved zero-product cases.
The exact results, in the displayed order, are
\[
\begin{array}{r|r|r}
 \text{labelled tables covered}&\text{normalized assignments}
    &\text{retained literal certificates}\\\hline
611330094&556346&92244\\
189135080&509966&90851\\
84287340&459119&95658\\
50074290&307277&46455\\
88224432&926459&117449\\
48945774&1347852&51904\\
20162124&302534&8539\\\hline
1092159134&4409553&503100
\end{array}
\]
Every retained table has binary first homology zero and second
homology of dimension one. For each, the independent checker
verifies the entire typed associative table, its nonzero
normalized two-cocycle, and a literal three-cochain whose
coboundary is its Alexander--Whitney square. All503100 squares
are zero; there is no survivor. All other tables fail the
necessary binary homology filter or a separately proved cut.
The run took139.30 seconds, including its literal checks.

The exact seven-profile input is
\path{outputs/linear17_bouquet_profiles.tsv}. Each complete
receipt binds that input, the producer source, the compressed
and uncompressed certificate hashes, its final count sentinel,
and checker completion. The complete interval and all literal
streams are archived under \path{outputs/linear17_bisets/bouquet/}.
This proves the finite seven-profile gate. We do not infer or
use the inherited unbounded bouquet theorem from this calculation.
\end{proof}


\subsection{Acyclic models: object count and an independently certified arrow bound}

We use ``acyclic'' here in the sense of Part A.7: every endomorphism
is an identity and the support relation is a partial order. Identities
are included in every arrow count.

\begin{proposition}[Proved]\label{A7:dimension-chain}
A finite acyclic category with $k$ objects has a nerve of dimension
at most $k-1$. Consequently $\nu(\mathbb{CP}^2)\ge5$. A five-object
acyclic model of $\mathbb{CP}^2$ must have totally ordered support.
\end{proposition}
\begin{proof}
Every nondegenerate simplex is a string of nonidentity arrows. Its
objects strictly increase in the support partial order and are
therefore distinct. Its dimension is at most $k-1$. Since
$H_4(\mathbb{CP}^2;\mathbb Z)=\mathbb Z$, at least five objects are
necessary. If there are exactly five objects, a nonzero fourth
homology class requires a four-simplex. Its vertices exhaust the
objects and order them $0<1<2<3<4$. Transitivity forces all ten
forward hom-sets to be nonempty. Acyclicity excludes reverse arrows
and nonidentity endomorphisms. This also independently confirms the
application of Lemma~1.1 of~\cite{CP2-Carcacia2026}.
\end{proof}

\begin{proposition}[Proved]\label{A7:adjacent-homs}
For a finite acyclic category on support $0<1<2<3<4$, if any one
of $C(0,1),C(1,2),C(2,3),C(3,4)$ is a singleton, every degree-two
square over any field is zero. Thus a five-object acyclic model
of $\mathbb{CP}^2$ has each of these four hom-sets of size at least two.
\end{proposition}
\begin{proof}
Suppose first that $C(0,1)$ is a singleton. In the full subcategory
on $012$, the boundary of a two-simplex, projected to the summand
of one-chains on $12$, determines that simplex uniquely. Thus its
second homology is zero. If $C(1,2)$ is a singleton, use instead
the projection to the $01$ one-chains. Over a field, second
cohomology is therefore zero in either case. Given a two-cocycle
$\alpha$ on $NC$, its restriction to $012$ is the coboundary of a
one-cochain $\lambda$. Extend $\lambda$ by zero on all other edges.
Subtracting its coboundary gives a representative of the same class
that vanishes on every $012$ triangle. On a four-simplex its square
is
\[
 (\alpha\smile\alpha)(a_{01},a_{12},a_{23},a_{34})
 =\alpha(a_{01},a_{12})\alpha(a_{23},a_{34}),
\]
which is now zero. The same argument on $234$ handles the other two
singletons. The square of the generator of the second cohomology
of $\mathbb{CP}^2$ is nonzero, giving the final assertion.
\end{proof}

\begin{theorem}[Proved by complete finite computation]
\label{A7:through24}
No finite acyclic category with five objects and at most twenty-four
arrows has classifying space of type $\mathbb{CP}^2$. If five objects
suffice, its least arrow count is therefore at least twenty-five.
\end{theorem}
\begin{proof}
By Propositions~\ref{A7:dimension-chain} and~\ref{A7:adjacent-homs}
we need only consider ten-tuples $h_{ij}=|C(i,j)|$, $i<j$, of
positive integers, with four adjacent entries at least two.
They have sum $n-5$. In particular $n\ge19$.
The number of nondegenerate $d$-simplices is
\[
 f_0=5,\qquad
 f_d=\sum_{0\le i_0<\cdots<i_d\le4}
       \prod_{j=1}^{d}h_{i_{j-1},i_j}\quad(1\le d\le4).
\]
This depends only on the profile, not on composition. We discard
exactly those profiles whose alternating sum is not $3$. At $n=19$
the only profile has $f=(5,14,25,28,16)$ and Euler characteristic
$4$, so it is excluded.

For $20\le n\le24$, our new program
\path{acyclic/code/search_five.cpp} generates every remaining
profile. It labels the identities $0,\ldots,4$ and then labels
arrows hom-set by hom-set in lexicographic order of $(i,j)$.
For each composable pair it permits every value in the required
hom-set. Identity products and products into singleton hom-sets
are fixed. For each composable triple it propagates
$(c\circ b)\circ a=c\circ(b\circ a)$ when the relevant inner
products are known, rejecting only unequal forced values. Each
propagated equality is a necessary associativity condition. It then
branches over the entire domain of an unassigned entry. Induction
on the search tree proves that every labelled associative table
in the profile occurs exactly once. There is no isomorphism,
opposition, relabelling, or cache discard.

The completed search counts are
\[
\begin{array}{r|r|r|r}
 n&\text{profiles}&\text{labelled associative tables}
   &\text{retained tables}\\\hline
19&0&0&0\\
20&2&18&0\\
21&8&224&64\\
22&16&13\,952&192\\
23&65&128\,868&19\,632\\
24&84&9\,250\,268&208\,896.
\end{array}
\]
The retained tables have mod-two Betti numbers exactly
$(1,0,1,0,1)$. All other tables are excluded by their computed
homology. The finite nerve has no cells above degree four, so
there is no uncomputed homology tail in this test. The search
checks that every chain-space basis fits its 128-bit linear
algebra representation and aborts if it does not; no truncation
is performed. For the retained cases the independent program
\path{acyclic/code/certify_five.py} reconstructs the typed table,
checks associativity including identity arrows, and reconstructs
the nerve and all boundary matrices without using generator code.
It records a two-cocycle $\alpha$, a two-cycle $z$ with
$\langle\alpha,z\rangle=1$, and a three-cochain $\beta$ satisfying
$\delta\beta=\alpha\smile\alpha$.

A further checker, \path{acyclic/code/verify_certificates.py},
imports neither producer nor search code. It reads every literal
table and cochain, rechecks typing and associativity, recomputes
the four boundary ranks, verifies $\partial z=0$ and the nonzero
pairing, verifies $\delta\alpha=0$, and verifies the square-primitive
equation on every four-simplex. All $228\,784$ retained tables
pass, with $53\,547\,280$ associativity checks. Their generators
therefore have zero mod-two square, excluding $\mathbb{CP}^2$.

Finally, \path{acyclic/code/audit_scope.py} independently generates
the complete required profile set and checks equality with each
completed profile log, not merely its cardinality. It checks
completion markers, all profile-sum counts, literal stream lengths,
and agreement of the producer and independent-checker digests.
The resulting \path{acyclic/outputs/scope_through_24.json} has status
\texttt{PASS}. The literal streams, cochains, full profile logs,
run times, and SHA-256 digests are supplied in
\path{acyclic/outputs/}. Large streams are stored as lossless gzip
files; the compression manifest records both names and the digest
of the uncompressed stream. This proves the asserted full scope.
\end{proof}

\begin{theorem}[Proved by complete finite computation]\label{A7:through25}
No finite acyclic category with five objects and at most twenty-five
arrows has classifying space of type $\mathbb{CP}^2$. If five objects
suffice, the least arrow count is at least twenty-six.
\end{theorem}
\begin{proof}
The preceding theorem handles at most twenty-four arrows. At
twenty-five, the same complete profile calculation gives $327$
Euler-compatible profiles. Exactly $40$ have all four adjacent
hom-sets of size two and are excluded by
Theorem~\ref{A7:four-double}. The program
\path{acyclic/code/search_five_v3.cpp} uses the same complete typed
table branching as before, with one additional exact homology
discard that we now justify.

Put $A_i=C(i,i+1)$ for $0\le i\le3$, and
$K_i=\ker(\mathbb F_2[A_i]\xrightarrow{\epsilon}\mathbb F_2)$,
where $\epsilon$ sends every arrow to $1$. The five face maps of
a four-simplex land in different object-set summands of $C_3$,
so a four-chain is a cycle precisely when each of its five faces
is zero. The intersection of the kernels of the two outside faces
is
\[
 K_0\otimes\mathbb F_2[A_1]\otimes
       \mathbb F_2[A_2]\otimes K_3.
\]
Choose one distinguished arrow in each of $A_0,A_3$. Differences
from these arrows give bases of $K_0,K_3$. The three remaining
face maps on this displayed space depend only on the three
adjacent composition tables
$A_i\times A_{i+1}\to C(i,i+2)$.
Their combined kernel is exactly $H_4(BC;\mathbb F_2)$, because
there are no five-chains. Thus, as soon as those three tables
are assigned, a branch whose kernel dimension differs from one
can be discarded independently of all unassigned longer products.
This discards $206\,618\,789$ partial branches. It neither
identifies tables nor drops an unseen case whose fourth Betti
number is one.

After these two discards the search completes $9\,064\,416$
labelled associative tables. Of these, $8\,877\,024$ have zero
first mod-two Betti number, and $7\,165\,632$ also have second
Betti number one. The Euler characteristic and fourth Betti
number force their third Betti number to be zero. These are the
retained cases; the table count here is explicitly after the
top-homology and structural discards.

The adapted producer \path{acyclic/code/certify_acyclic_cpp.cpp}
records each full typed composition table, a nonzero generator
$\alpha\in H^2(-;\mathbb F_2)$, and a three-cochain $\beta$.
The separate \path{acyclic/code/check_acyclic_cpp.cpp} imports
neither producer nor search code. It checks typing, identities,
associativity, connectedness, the first and second boundary ranks,
the cocycle condition, that $\alpha$ is not a coboundary, and
the literal equation $\delta\beta=\alpha\smile\alpha$ on every
four-simplex. Every retained case passes. The checker performs
$1\,940\,730\,432$ associativity checks,
$436\,573\,440$ cocycle checks, and $219\,424\,512$
square-primitive checks.

The independent \path{acyclic/code/link_acyclic_streams.cpp}
reconstructs the original profile and product JSON from every
full certificate table and compares each record byte for byte
with the corresponding search record, checking both streams
through their checked gzip ends. All $7\,165\,632$ records
match. Finally \path{acyclic/code/audit_scope_v2.py} checks the
entire required profile set, every profile count, the exact
per-profile join counts, completion markers, and the stream
digests. Its result is
\path{acyclic/outputs/scope_through_25.json}, with status
\texttt{PASS}. The complete literal tables and certificates are
the corresponding \path{five_25.tables.jsonl.gz} and
\path{five_25.certificates.tsv.gz}. The zero square excludes
$\mathbb{CP}^2$ in every remaining case.
\end{proof}

\begin{theorem}[Proved by structural reductions and complete computation]
\label{A7:through26}
No finite acyclic category with five objects and at most twenty-six
arrows has classifying space of type $\mathbb{CP}^2$. Thus a
five-object model, if one exists, has at least twenty-seven arrows.
\end{theorem}
\begin{proof}
At twenty-six arrows the complete Euler-profile formula gives
$333$ profiles. Applying Propositions~\ref{A7:short-composite-gates}
and~\ref{A7:long-composite-gates} and
Theorem~\ref{A7:two-double-graph} leaves exactly the following six,
in the hom-set order $01,02,03,04,12,13,14,23,24,34$:
\[
\begin{gathered}
(2,2,1,2,3,2,2,2,2,3),\quad(2,2,1,2,3,2,2,3,2,2),\\
(2,2,2,2,3,2,1,3,2,2),\quad(3,2,1,2,2,2,2,2,2,3),\\
(3,2,2,2,2,2,1,2,2,3),\quad(3,2,2,2,2,2,1,3,2,2).
\end{gathered}
\]
This is equality of explicit profile sets, checked by the
independent program \path{acyclic/code/profile_gates.py}; there
is no isomorphism or opposition discard.

In these profiles enumerate all three maps
\[
 C(0,1)\times C(1,2)\to C(0,2),\quad
 C(1,2)\times C(2,3)\to C(1,3),\quad
 C(2,3)\times C(3,4)\to C(2,4).
\]
The independent \path{acyclic/code/check_adjacent_four_cycles.cpp}
enumerates their entire Cartesian product, without imposing
associativity. The numbers of map triples, in the order above,
are $262\,144$, $2\,097\,152$, $2\,097\,152$, $65\,536$,
$65\,536$, and $262\,144$. Among all $4\,849\,664$ triples,
$4\,155\,200$ satisfy an explicit zero-product composition
condition of Proposition~\ref{A7:short-composite-gates}.
For each of the other $694\,464$, the program constructs the
full fourth boundary: every four-arrow string contributes its
five distinct faces, recorded in five separate object-subset
summands. Direct Gaussian elimination over $\mathbb F_2$ gives
kernel dimension at least two in every case. In particular it
never gives dimension one. This implementation imports no
generator or certificate code and uses the full boundary, rather
than the reduced-kernel formula used by the table search.

No actual associative category with one of these adjacent-map
triples can therefore have the cohomology ring of
$\mathbb{CP}^2$. This excludes even those triples admitting no
associative completion, which is harmless. The complete table
search \path{acyclic/code/search_five_v4.cpp} independently reaches
the same exclusion. Its profile log accounts for all $333$
profiles and $1\,614\,434$ search nodes, with no retained table.
The scope join \path{acyclic/code/audit_scope_v4.py} checks the
entire profile sets, every structural-skip flag, all counts, the
adjacent-map checker totals, and stream digests. Its output
\path{acyclic/outputs/five_26_v4.scope.json} is \texttt{PASS}.
The lower counts were already excluded by
Theorem~\ref{A7:through25}.
\end{proof}

\begin{theorem}[Proved by complete finite computation]\label{A7:through27}
No finite acyclic category with five objects and at most twenty-seven
arrows has classifying space of type $\mathbb{CP}^2$. A five-object
model, if one exists, therefore has at least twenty-eight arrows.
\end{theorem}
\begin{proof}
The preceding theorem handles at most twenty-six arrows. At
$27$ arrows the independent profile formula gives $1126$ profiles.
The same proved structural gates used there leave exactly eight.
The complete typed-table generator
\path{acyclic/code/search_five_v4.cpp} branches over every
unforced product value and uses only necessary associativity
propagation, the exact top-kernel test proved in
Theorem~\ref{A7:through25}, and the explicit zero-product
conditions of Proposition~\ref{A7:short-composite-gates}.
No isomorphism, opposition, or cache discard is used.
It completes $34\,810\,456$ nodes and $13\,305\,600$
labelled associative tables after these discards. All have zero
first mod-two Betti number; $2\,274\,048$ have second Betti
number one and are retained. Their third and fourth Betti numbers
are zero and one by the top test and Euler characteristic.

For every retained table the independent C++ checker described
in Theorem~\ref{A7:through25} verifies a nonzero two-cocycle
and a literal three-cochain primitive of its square. Every case
passes, with $724\,536\,576$ associativity checks,
$168\,058\,368$ cocycle checks and $81\,865\,728$
square-primitive checks. The literal table-to-certificate join
matches all $2\,274\,048$ records byte for byte. The independent
scope program verifies the entire profile set, every structural
skip, all profile counts and certificate counts, and both stream
digests. Its result is
\path{acyclic/outputs/five_27_v4.scope.json}, with status
\texttt{PASS}. Complete tables, cochains and full logs are
\path{acyclic/outputs/five_27_v4.*}. Every retained generator
has zero mod-two square, completing the exclusion.
\end{proof}

\begin{remark*}[Coverage of output-label normalization]
The next search version orders unassigned product positions by increasing
length of their target object interval. In a target hom-set $H=C(i,j)$,
it branches over the labels already used as outputs of nonidentity
products and the least unused label. Adjacent arrows stay labelled.
This restriction preserves coverage up to permutations of the hom-sets.
Indeed, when the current interval has length $L=j-i$, an arrow in $H$
cannot be an input to any nonidentity product of target length at most
$L$. No product into a longer non-singleton target has yet been
assigned: branching has not reached it, and associativity propagation
only copies an already assigned entry to an entry with the same target
hom-set. Longer singleton targets are constant. Thus any permutation
of $H$ fixing all previously used output labels preserves all assigned
information, and acts transitively on the unused labels. The selected
product position itself is fixed. Every completion using a new label
therefore has an isomorphic completion using the least unused label.
The rank and product-vanishing tests are invariant under these
permutations. Induction over the search proves coverage. Counts from
this version are counts of normalized tables, not isomorphism classes.

An independent program applies the corresponding completed-table
renaming to all $110592$ retained labelled twenty-seven-arrow tables
that survive the binary-act profile test. Their exact normalized set
is the set of $1152$ tables in the normalized run. The comparison,
\path{acyclic/outputs/five_27_v6_normalization_comparison.json},
has status \texttt{PASS}. This regression check supplements the
coverage proof. The complete argument and source inspection are
recorded in \path{acyclic/notes/output_label_normalization.md}.
\end{remark*}

\begin{theorem}[Proved: exclusion through twenty-eight arrows]
\label{A7:through28}
Every five-object acyclic category of type $\mathbb{CP}^2$ has
at least twenty-nine arrows, identities included.
\end{theorem}
\begin{proof}
Theorem~\ref{A7:through27} treats smaller counts. At twenty-eight
arrows the positive hom-set profiles with all four adjacent sizes
at least two and Euler characteristic three number $1198$.
The preceding structural gates, together with
Lemma~\ref{CP2:A-high-affine}, leave $15$ profiles.
The latter lemma applies when every hom-set out of object $0$
has at most two arrows, or every hom-set into object $4$ does:
the category is the corresponding objectwise binary act cone.

The generator \path{acyclic/code/search_five_v6.cpp} searches
all remaining associative tables with the normalization just proved.
It visits $19914994$ nodes and finishes $6919488$ normalized
associative tables after the exact fourth-homology and explicit
product-vanishing tests. Of these, $5220288$ pass the first-homology
test and $5069952$ also pass the second-homology test. These are
counts of normalized tables, with repetitions up to isomorphism
allowed. The complete literal table stream is
\path{acyclic/outputs/five_28_v6.tables.jsonl.gz}.

For each retained table the producer writes a nonzero degree-two
mod-two cocycle and an explicit degree-three primitive of its square.
The independent checker rebuilds all products and normalized chains,
checks the cocycle's nontriviality, and verifies the primitive equation.
All $5069952$ records pass, with $1819677312$ associativity equations,
$486466560$ cocycle equations and $273590784$ primitive equations.
No unexcluded record remains. The exact stream join checks, in order,
that the table in each certificate is the corresponding generator table.
The separate scope auditor reconstructs all $1198$ profiles, checks
all theorem flags, sums every search counter, and joins these records
to both independent certificate receipts and the exact stream hashes.
Its result, \path{acyclic/outputs/five_28_v6.scope.json}, is
\texttt{PASS}. Thus every remaining category has zero mod-two square,
which excludes the cohomology ring of $\mathbb{CP}^2$.
\end{proof}

\begin{remark*}[Coverage of successive adjacent-map reductions]
Write the three adjacent composition maps as
$u:A\times B\to U$, $v:B\times C\to V$ and $w:C\times D\to W$,
where $A,B,C,D$ are the four adjacent hom-sets.
The next generator completes $u$, then $v$, then $w$, before
branching on any longer target. We normalize each output table
by first use as above.

First retain the lexicographically least normalized $u$ under
all permutations of $A$ and $B$. Store exactly those permutations
$q$ of $B$ for which some permutation of $A$, followed by output
normalization in $U$, fixes $u$. This is the projection to
$\operatorname{Sym}(B)$ of the full stabilizer of $u$: the output
permutation is forced on used labels and extends arbitrarily on
unused ones. Next retain the least normalized $v$ under this
projected group and $\operatorname{Sym}(C)$. Store the permutations
of $C$ that can fix $v$ with one of the permitted $B$ permutations.
This is exactly the projection of the stabilizer of $(u,v)$:
the required $A$ and $U$ permutations exist by the first test and
do not affect $v$. Finally retain the least normalized $w$ under
that projected group and $\operatorname{Sym}(D)$.

Every triple of maps can be made canonical successively in this
way, first under the full group and then under the stabilizers
of the preceding maps. Each permutation extends to the whole
category. Longer non-singleton target tables remain unassigned;
longer singleton tables remain constant. Thus every associative
completion has an isomorphic completion surviving all prefix
tests. The earlier first-use argument covers the remaining
products. A cache miss performs the complete finite permutation
test; the cache stores only its answer and projected stabilizer,
and is cleared at every size profile. It drops no unseen prefix.
Residual automorphisms of the completed triple are not used to
identify longer tables, so the reported counts remain normalized
leaf counts, with possible isomorphic repetitions.

An independent program computes these prefix stabilizers directly
from all $1152$ retained twenty-seven-arrow tables of the previous
version. Exactly sixteen survive, and their complete table set
equals the sixteen leaves of the new version. The exact comparison
has status \texttt{PASS} in
\path{acyclic/outputs/five_27_v7_normalization_comparison.json}.
\end{remark*}

\begin{theorem}[Proved: exclusion through twenty-nine arrows]
\label{A7:through29}
Every five-object acyclic category of type $\mathbb{CP}^2$ has
at least thirty arrows, identities included.
\end{theorem}
\begin{proof}
Theorem~\ref{A7:through28} treats smaller counts. At twenty-nine
arrows there are $3214$ positive hom-set profiles satisfying the
adjacent-size and Euler-characteristic conditions. The proved
profile gates leave $83$. The generator with the preceding
stabilizer reduction visits $17860548$ nodes and finishes
$1017270$ normalized associative tables after the exact fourth-homology
and explicit product-vanishing tests. Of these, $919410$ pass the
first-homology test and $256856$ also pass the second-homology test.

Every retained table has a literal mod-two nonzero degree-two
cocycle and a degree-three primitive of its square. The independent
checker verifies all $256856$ certificates, including $96138292$
associativity equations, $24374556$ cocycle equations and $12696780$
primitive equations. The exact table-to-certificate stream join
also passes. The separate scope auditor reconstructs all $3214$
profiles, verifies every theorem flag, sums all counters including
the new prefix-pruning counter, and joins the profile totals to
the full certificate receipts and stream hashes. Its receipt is
\path{acyclic/outputs/five_29_v7.scope.json}, with status
\texttt{PASS}. The complete source, tables, certificates and
independent checkers are retained. Every possible category is
therefore excluded by a proved gate or a checked zero square.
\end{proof}

\begin{theorem}[Proved: exclusion through thirty arrows]
\label{A7:through30}
Every five-object acyclic category of type $\mathbb{CP}^2$ has
at least thirty-one arrows, identities included.
\end{theorem}
\begin{proof}
Theorem~\ref{A7:through29} treats smaller counts. At thirty
arrows there are $3086$ Euler-compatible profiles with the
required adjacent sizes. Using the strengthened longest-arrow
collapse gate, Proposition~\ref{A7:extreme-arrow-collapse},
together with the preceding gates, leaves $62$ profiles.
The fully checked first-use and adjacent-prefix reductions
remain unchanged.

The fourth-boundary calculation uses only its three internal
face components. Put $F=\mathbb F_2$. To see that this is exact, write its domain
as $F[A]\otimes F[B]\otimes F[C]\otimes F[D]$ for the four
adjacent hom-sets. The first and last face components have
supports $1234$ and $0123$. Their common kernel is
$K_A\otimes F[B]\otimes F[C]\otimes K_D$, where $K$ denotes
an augmentation kernel. On its usual tensor basis the exterior
components vanish identically. The remaining rows have supports
$0234,0134,0124$, hence are exactly the three-arrow chains from
$0$ to $4$. Their number is
\[
\begin{aligned}
 &|C(0,2)|\,|C(2,3)|\,|C(3,4)|\\
 &\quad+|C(0,1)|\,|C(1,3)|\,|C(3,4)|\\
 &\quad+|C(0,1)|\,|C(1,2)|\,|C(2,4)|.
\end{aligned}
\]
Deleting identically zero exterior rows and renumbering these
rows changes no kernel. The code checks the actual row capacities
and aborts rather than truncating. An independent full-matrix
comparison on the $59500$ distinct adjacent-map triples represented
by all $256856$ retained twenty-nine-arrow tables agrees exactly;
so do $8300$ additional arbitrary triples across all $83$ residual
profiles. This comparison supplements the algebraic argument.

The complete thirty-arrow run visits $18930744$ nodes and finishes
$7312590$ normalized associative tables after the exact top-homology
and product-vanishing tests. Of these, $5385994$ pass the first-homology
test and $2024190$ also pass the second-homology test. Every retained
table has a literal square-zero certificate. The independent checker
verifies all records, including $681990294$ associativity equations,
$153680124$ cocycle equations and $78111504$ primitive equations.
The exact table-to-certificate join passes. The separate scope
program reconstructs every size profile and theorem flag, verifies
all counters, and matches the complete certificate stream and its
hashes. The receipt \path{acyclic/outputs/five_30_v9.scope.json}
has status \texttt{PASS}. This excludes every thirty-arrow case.
\end{proof}

\begin{theorem}[Proved: exclusion through thirty-one arrows]
\label{A7:through31}
Every five-object acyclic category of type $\mathbb{CP}^2$ has
at least thirty-two arrows, identities included.
\end{theorem}
\begin{proof}
Theorem~\ref{A7:through30} treats all smaller counts. There are
$8567$ positive thirty-one-arrow hom-set profiles satisfying the
necessary Euler-characteristic and adjacent-size conditions.
The preceding profile gates, including
Proposition~\ref{A7:tree-stalk-profile}, leave exactly $110$.
The complete version-fifteen search uses the proved output-label
and adjacent-prefix normalizations without further identification
of completed categories. It visits $80058405$ nodes. No complete
associative table survives its obstruction tests.

The universal-short tests of
Theorem~\ref{A7:universal-short-model} account for $1067734$
literal records: $422168$ have positive third mod-two homology,
$165758$ have zero second mod-two homology, and $479808$ carry a
complete degree-two cohomology basis together with an explicit
degree-three primitive for the square of every basis element.
A separately written checker reconstructs each full component-graph
category and its normalized nerve, verifies associativity and
every boundary-square equation, and computes the actual mod-two
homology ranks. For the last class it also checks the basis modulo
all coboundaries and every Alexander--Whitney square equation.
All $1067734$ records pass.

The partial long-graph obstruction of
Proposition~\ref{A7:partial-long-graph} accounts for $13074446$
additional prefix records: $890746$ use the $03$ graph and
$12183700$ use the $14$ graph. A second independent checker
retains every graph vertex, checks the known-edge connectedness,
the global assignment prefix and first-use conditions, and the
pairwise prefix-free ordering. It verifies the exact unrestricted
short-map suffix counts using unbounded integer arithmetic. Their
sum is $53716847483$; this is a count of literal short-map
completions and is not reported as a count of categories. Every
record passes. The remaining branches are covered by the proved
profile, top-cycle and explicit product-vanishing tests described
above.

The independent scope auditor regenerates all $8567$ profiles in
their exact order, checks every theorem flag, joins all per-profile
counters to the complete run, and binds the closed universal and
partial streams to their independent receipts. It also verifies
that the partial-prefix cylinders are disjoint from every remaining
full record. Its receipt
\path{acyclic/outputs/five_31_v15.scope.json} has status
\texttt{PASS}. The separate receipt
\path{acyclic/outputs/five_31_v15.provenance.json}, with status
\texttt{PASS\_COMPLETE\_V15\_SCOPE\_AND\_PROVENANCE}, binds the
literal producer, the actual optimized binary and build manifests,
the proof files, the scope auditor and its imported helper sources,
and all complete streams and independent checkers by their hashes.
Thus every thirty-one-arrow case is excluded; no incomplete earlier
search is used as a complete census.
\end{proof}

\begin{remark*}[Exact unresolved state of A.7]
The construction below proves that five objects suffice, so
$\nu(\mathbb{CP}^2)=5$. The least arrow count at this object
count remains unresolved. Its lower bound is thirty-two,
and the explicit model below gives upper bound thirty-eight.
The next undecided arrow count is thirty-two; no claim about
that count is included in Theorem~\ref{A7:through31}.
\end{remark*}

\begin{remark*}[Constructions checked without success]
The full directed five-level subcategories of the Karoubi envelope
of the independently reconstructed twenty-three-element monoid were
checked for every choice of its four idempotent isomorphism classes.
Eight choices pass the Euler-characteristic and adjacent-size gates;
none has nonzero fourth mod-two homology. The corresponding
four-level corner-plus-act construction has no passing Euler profile.
These tests cover full hom-sets only.

We also constructed the categorical symmetric square of the
three-object sphere category with hom-set sizes $(2,2,1)$.
It has six objects and thirty arrows, but its mod-two Betti numbers
are $(1,0,0,0,3)$. In a further exact family of $164$ canonical
three-object sphere tables, with hom-set sizes $(2,q,q-1)$ for
$2\le q\le5$, every categorical symmetric square again has zero
second homology. Thus none provides the required model. In
particular we do not identify the classifying space of a categorical
orbit quotient with the topological symmetric square without proof.
The complete tested quotient tables and results are included.
\end{remark*}


\subsection{Four double adjacent hom-sets}
\begin{proposition}[Proved]\label{A7:four-double}
Let $C$ be a finite acyclic category on the totally ordered support
$0<1<2<3<4$. Suppose that each of $C(0,1),C(1,2),C(2,3),C(3,4)$
has exactly two elements. Then the product
$H^2(BC;\mathbb Q)\otimes H^2(BC;\mathbb Q)\longrightarrow H^4(BC;\mathbb Q)$
is zero, regardless of the sizes of the other six hom-sets.
\end{proposition}
\begin{proof}
Write the four adjacent sets as $A=\{a_0,a_1\}$,
$B=\{b_0,b_1\}$, $D=\{c_0,c_1\}$, and $E=\{d_0,d_1\}$.
We use $D,E$ for these sets only, not for categories. Let
\[
 u(a,b)=b\circ a\in C(0,2),\quad
 v(b,c)=c\circ b\in C(1,3),\quad
 w(c,d)=d\circ c\in C(2,4).
\]
Let $\alpha,\beta$ be normalized rational two-cocycles on $NC$.
Every four-simplex has vertices $01234$, and its cup value is
$\alpha(a,b)\beta(c,d)$. If the restriction of $\alpha$ to the
full subcategory on $012$ is zero in cohomology, extend a primitive
one-cochain by zero to $C$ and change $\alpha$ by its coboundary.
The new representative vanishes on all $012$ triangles, proving
$[\alpha][\beta]=0$. The same argument applies if $[\beta]$ restricts
to zero on $234$. We may therefore assume both restrictions nonzero.

The boundary of a two-chain in $C_{012}$ has separate components on
$A$, $B$, and $C(0,2)$. Vanishing of its first two components forces
the chain to be a multiple of
$(a_0-a_1)\otimes(b_0-b_1)$. Consequently
\[
 [u(a_0,b_0)]+[u(a_1,b_1)]
 =[u(a_0,b_1)]+[u(a_1,b_0)]
 \quad\text{in }\mathbb Q[C(0,2)],
\]
and $x=\alpha(a_0,b_0)-\alpha(a_1,b_0)
-\alpha(a_0,b_1)+\alpha(a_1,b_1)$ is nonzero.
Equality of these two multisets says that $u$ is constant on each
column or constant on each row. There are exactly two disjoint cases:
\begin{enumerate}
\item[(U0)] $u(a_0,b)=u(a_1,b)$ for both $b$;
\item[(U1)] $u(a,b_0)=u(a,b_1)$ for both $a$, and the two row values
are distinct.
\end{enumerate}
The analogous argument on $234$ shows that
$y=\beta(c_0,d_0)-\beta(c_0,d_1)
-\beta(c_1,d_0)+\beta(c_1,d_1)\ne0$, and gives two cases:
\begin{enumerate}
\item[(W0)] $w(c,d_0)=w(c,d_1)$ for both $c$;
\item[(W1)] $w(c_0,d)=w(c_1,d)$ for both $d$, with distinct column values.
\end{enumerate}

Let $z$ be any rational four-cycle. The faces omitting $0$ and $4$
land in separate three-chain summands. Their vanishing gives a unique
matrix $(t_{bc})_{b\in B,c\in D}$ such that
\[
 z=\sum_{b,c}t_{bc}
 (a_0-a_1)\otimes b\otimes c\otimes(d_0-d_1).
\]
The remaining three faces imply
\begin{align*}
 &\sum_b t_{bc}([u(a_0,b)]-[u(a_1,b)])=0 &&(c\in D),\\
 &\sum_c t_{bc}([w(c,d_0)]-[w(c,d_1)])=0 &&(b\in B),\\
 &\sum_{b,c}t_{bc}[v(b,c)]=0 &&\text{in }\mathbb Q[C(1,3)].
\end{align*}
Thus U1 forces every column sum of $t$ to vanish, while W1 forces
every row sum to vanish.

Put $\ell_b=\alpha(a_0,b)-\alpha(a_1,b)$ and
$r_c=\beta(c,d_0)-\beta(c,d_1)$, so $x=\ell_{b_0}-\ell_{b_1}$,
$y=r_{c_0}-r_{c_1}$, and
\[
 \langle\alpha\smile\beta,z\rangle=\sum_{b,c}t_{bc}\ell_b r_c.
\]
Define functions on $C(1,3)$ by
$F(s)=\alpha(a_0,s)-\alpha(a_1,s)$ and
$G(s)=\beta(s,d_0)-\beta(s,d_1)$. The cocycle equation
\[
 \gamma(g,h)-\gamma(g\circ f,h)+\gamma(f,h\circ g)-\gamma(f,g)=0
\]
on $0123$ gives $\ell_b=F(v(b,c))$ in case U0, and
$x=F(v(b_0,c))-F(v(b_1,c))$ in case U1. On $1234$ it gives
$r_c=G(v(b,c))$ in case W0, and
$y=G(v(b,c_0))-G(v(b,c_1))$ in case W1.

In case U0/W0 the cup evaluation is
$\sum t_{bc}F(v(b,c))G(v(b,c))=0$, because the coefficients sum to
zero separately over every fibre of $v$.

In case U1/W0 write $t_{b_0c}=q_c$ and $t_{b_1c}=-q_c$.
Applying $F$ to the last boundary equation gives
$x(q_{c_0}+q_{c_1})=0$. Since $x\ne0$, the matrix $t$ is a scalar
multiple of the checkerboard matrix. The same conclusion holds
in case U0/W1 by applying $G$, since $y\ne0$, and it holds directly
in case U1/W1. If that scalar is zero we are done. Otherwise the
last boundary equation says
\[
 [v(b_0,c_0)]+[v(b_1,c_1)]
 =[v(b_0,c_1)]+[v(b_1,c_0)].
\]
Hence $v$ is constant on rows or on columns. If it is constant on
columns, the formulas for $\ell$ (U0) or $x$ (U1) force $x=0$.
If it is constant on rows, the formulas for $r$ (W0) or $y$ (W1)
force $y=0$. Both contradict the nonzero restrictions. Thus the
cup evaluation vanishes on every rational four-cycle. Over a field,
cohomology pairs nondegenerately with homology, proving the claim.
\end{proof}


\subsection{Two double hom-sets force a nontrivial graded graph}
\begin{theorem}[Proved]\label{A7:two-double-graph}
Let $C$ be a finite acyclic category with support $0<1<2<3<4$.
Suppose $C(0,1)=\{a_0,a_1\}$ and $C(1,2)=\{b_0,b_1\}$.
Define a directed multigraph $\Gamma$ with vertex set $C(1,3)$,
and, for each $c\in C(2,3)$, one edge from $c\circ b_1$ to
$c\circ b_0$. If some product of rational degree-two classes on
$BC$ is nonzero, the following conditions are necessary:
\begin{enumerate}
\item There are distinct $u_0,u_1\in C(0,2)$ with
$b_j\circ a_i=u_i$ for every $i,j\in\{0,1\}$.
\item There is a function $h:C(1,3)\to\mathbb Q$ satisfying
$h(c\circ b_0)-h(c\circ b_1)=1$ for every $c\in C(2,3)$.
\item After identifying parallel edges, the undirected graph
underlying $\Gamma$ contains a cycle. Every such cycle has even
length, and hence length at least four.
\end{enumerate}
In particular, if $|C(2,3)|\le3$ or $|C(1,3)|\le3$, all rational
degree-two products vanish. The dual statement holds with the
roles of $012$ and $234$ reversed.
\end{theorem}
\begin{proof}
Write $A=C(0,1)$, $B=C(1,2)$, $D=C(2,3)$ and $E=C(3,4)$, and set
$u(a,b)=b\circ a$, $v(b,c)=c\circ b$, $w(c,d)=d\circ c$.
Let $\alpha,\beta$ be two-cocycles with nonzero product. The
restriction of $[\alpha]$ to the full subcategory on $012$ cannot
be zero: otherwise a primitive one-cochain extended by zero makes
$\alpha$ vanish on all $012$ triangles, annihilating its product
with $\beta$ by the Alexander--Whitney formula.

Every two-cycle on $012$ is a multiple of
$(a_0-a_1)\otimes(b_0-b_1)$. Thus
\[
 [u(a_0,b_0)]+[u(a_1,b_1)]
 =[u(a_0,b_1)]+[u(a_1,b_0)],
\]
and, with $\ell_b=\alpha(a_0,b)-\alpha(a_1,b)$, we have
$x=\ell_{b_0}-\ell_{b_1}\ne0$. Equality of these multisets gives
exactly two disjoint alternatives: either $u(a_0,b)=u(a_1,b)$
for each $b$ (case U0), or $u(a,b_0)=u(a,b_1)$ for each $a$
with distinct row values (case U1).

Define $F:C(1,3)\to\mathbb Q$ by
$F(s)=\alpha(a_0,s)-\alpha(a_1,s)$. The two cocycle equations
on $(a_0,b,c)$ and $(a_1,b,c)$, subtracted, show that
\[
 \ell_b=F(v(b,c))\quad\text{in U0},\qquad
 x=F(v(b_0,c))-F(v(b_1,c))\quad\text{in U1}.
\]
Let $z$ be any rational four-cycle. Its face omitting $0$ forces
\[
 z=(a_0-a_1)\otimes
   \sum_{b\in B,c\in D,d\in E}t_{bcd}\,b\otimes c\otimes d.
\]
The four remaining separate face summands give
\begin{align*}
 &\sum_d t_{bcd}=0 &&(b\in B,c\in D),\\
 &\sum_b t_{bcd}([u(a_0,b)]-[u(a_1,b)])=0 &&(c\in D,d\in E),\\
 &\sum_{b,c}t_{bcd}[v(b,c)]=0 &&(d\in E),\\
 &\sum_{c,d}t_{bcd}[w(c,d)]=0 &&(b\in B).
\end{align*}
All equalities involving brackets are in the free rational vector
space on the indicated hom-set.

In U0, the images $v(b_0,D)$ and $v(b_1,D)$ are disjoint: a common
value would give $\ell_{b_0}=\ell_{b_1}$, contrary to $x\ne0$.
The third boundary equality therefore holds separately for each
$b$. The cocycle equation on $(b,c,d)$ is
\[
 \beta(c,d)=\beta(v(b,c),d)-\beta(b,w(c,d))+\beta(b,c).
\]
Multiplying by $t_{bcd}$ and summing over $c,d$, all three terms
on the right vanish: respectively by the separated third boundary
equality, the fourth equality, and the first equality. Therefore
\[
 \langle\alpha\smile\beta,z\rangle
 =\sum_b\ell_b\sum_{c,d}t_{bcd}\beta(c,d)=0.
\]
Since this holds for every $z$, U0 contradicts the nonzero product.
Consequently U1 holds, proving condition (1).

In U1 the second boundary equality forces
$t_{b_0cd}=q_{cd}$ and $t_{b_1cd}=-q_{cd}$ for scalars $q_{cd}$.
The remaining equalities give
\[
 \sum_dq_{cd}=0,\qquad
 \sum_cq_{cd}([v(b_0,c)]-[v(b_1,c)])=0,\qquad
 \sum_{c,d}q_{cd}[w(c,d)]=0.
\]
The function $h=F/x$ satisfies condition (2). In each connected
component, its values differ from one fixed rational offset by
integers, with each edge increasing the integer by one. Hence the
underlying graph is bipartite; there are no loops and opposite
orientations on the same pair of vertices cannot both occur.

Suppose its underlying simple graph were a forest. The kernel of
its incidence matrix is then generated by differences of parallel
edges. More explicitly, orient each edge by increasing $h$.
The coefficients of a one-cycle sum to zero over each parallel
class: remove a leaf of each finite tree and use the incidence
equation at that leaf, then continue inductively. Applied to the
middle displayed equality, this says that for every $d$ and every
parallel class $P\subseteq D$, $\sum_{c\in P}q_{cd}=0$.

Using the cocycle equation for $b=b_0$ and the first and third
remaining boundary equalities gives
\[
 \sum_{c,d}q_{cd}\beta(c,d)
 =\sum_{c,d}q_{cd}\beta(v(b_0,c),d)=0.
\]
The last equality follows because $v(b_0,c)$ is constant on each
parallel class. But the cup evaluation is $x$ times the left side,
so again it vanishes on every four-cycle, a contradiction. This
proves condition (3). A bipartite simple cycle requires at least
four edges and four vertices, proving the numerical consequences.
Opposition reverses the five objects and proves the dual statement.
\end{proof}


\subsection{Short composite hom-sets and vanishing products}
\begin{proposition}[Proved]\label{A7:short-composite-gates}
Let $C$ be a finite acyclic category with support $0<1<2<3<4$,
and let $K$ be any field.
Write $A=C(0,1)$, $B=C(1,2)$, $D=C(2,3)$, $E=C(3,4)$, and
\[
 u(a,b)=b\circ a,\qquad v(b,c)=c\circ b,\qquad
 w(c,d)=d\circ c.
\]
Every product $H^2(BC;K)\otimes H^2(BC;K)\to H^4(BC;K)$ is zero
if any one of the following conditions holds:
\begin{enumerate}
\item for each $b\in B$, $u(a,b)$ is independent of $a\in A$;
\item for each $c\in D$, $w(c,d)$ is independent of $d\in E$;
\item there is $c\in D$ such that $v(b,c)$ is independent of $b\in B$;
\item there is $b\in B$ such that $v(b,c)$ is independent of $c\in D$.
\end{enumerate}
In particular, nonzero degree-two products require
$|C(0,2)|,|C(1,3)|,|C(2,4)|\ge2$.
\end{proposition}
\begin{proof}
Suppose (1) holds. Fix $a_*\in A$ and two cocycles $\alpha,\beta$.
For $a\ne a_*$ put
\[
 \ell_{ab}=\alpha(a,b)-\alpha(a_*,b),\qquad
 F_a(s)=\alpha(a,s)-\alpha(a_*,s),\quad s\in C(1,3).
\]
Subtracting the cocycle equations on $(a,b,c)$ and $(a_*,b,c)$
gives $\ell_{ab}=F_a(v(b,c))$ for every $b,c$.
For any four-cycle $z$, the face omitting object $0$ gives a unique
expression
\[
 z=\sum_{a\ne a_*}(a-a_*)\otimes
       \sum_{b,c,d}t^a_{bcd}\,b\otimes c\otimes d.
\]
Because the five faces of a four-simplex lie on different object
subsets, they vanish separately. The faces omitting $2$, $3$, and
$4$ give, for each $a\ne a_*$,
\[
 \sum_{b,c}t^a_{bcd}[v(b,c)]=0\quad(d\in E),\qquad
 \sum_{c,d}t^a_{bcd}[w(c,d)]=0\quad(b\in B),\qquad
 \sum_dt^a_{bcd}=0\quad(b,c).
\]
Here brackets denote the basis of the free $K$-vector space on the
indicated hom-set. The cocycle equation on $(b,c,d)$ gives
\[
 \beta(c,d)=\beta(v(b,c),d)-\beta(b,w(c,d))+\beta(b,c).
\]
Consequently the cup evaluation is
\[
 \langle\alpha\smile\beta,z\rangle
   =\sum_{a,b,c,d}t^a_{bcd}\ell_{ab}\beta(c,d)=0.
\]
Indeed, after the substitution, the first sum vanishes by the first
boundary equation weighted by $F_a(s)\beta(s,d)$; the second by
the second equation weighted by $\ell_{ab}\beta(b,r)$; and the
third by the last equation. Over a field, evaluation on all cycles
detects cohomology, proving (1). Reversing the objects proves (2);
the reversal interchanges the factors, whose degree is two.

For (3), fix such a $c$ and write $q=v(b,c)$, independent of $b$.
Define a one-cochain on the full subcategory on $012$ by
\[
 \lambda(a)=\alpha(a,q),\qquad
 \lambda(b)=\alpha(b,c),\qquad
 \lambda(p)=\alpha(p,c)\quad(p\in C(0,2)).
\]
The cocycle equation on $(a,b,c)$ says exactly
$\alpha(a,b)=\lambda(b)-\lambda(u(a,b))+\lambda(a)$.
Thus $\alpha$ restricts to a coboundary on $012$. Extending
$\lambda$ by zero to all other edges and subtracting its coboundary
gives a representative vanishing on every $012$ triangle. Its
Alexander--Whitney product with $\beta$ is zero on every
four-simplex. Reversal proves (4). The three stated cardinality
conditions follow from (1), (3), and (2), respectively.
\end{proof}

\begin{proposition}[Proved]\label{A7:long-composite-gates}
For a finite acyclic category with support $0<1<2<3<4$, nonzero
degree-two products over a field require that at least two of
$C(0,3),C(0,4),C(1,4)$ have size at least two.
\end{proposition}
\begin{proof}
If $C(0,4)$ and $C(1,4)$ are singletons, choose any arrow
$r:2\to4$. The unique arrows $0\to4$, $1\to4$, together with
$r$, are the components of a natural transformation from the
inclusion of the full subcategory on $012$ to the constant functor
at $4$. Naturality follows from the two singleton assumptions.
Thus the inclusion is null-homotopic on classifying spaces, so
every two-cocycle restricts to a coboundary on $012$.
Extending a primitive by zero shows that every degree-two product
is zero by the Alexander--Whitney formula. Opposition treats
singleton $C(0,3)$ and $C(0,4)$.

It remains that $C(0,3)=\{t\}$ and $C(1,4)=\{r\}$.
Fix arrows $c:2\to3$ and $d:3\to4$. For any two-cocycle $\alpha$,
the equation on $(a,s,d)$, where $a:0\to1$ and $s:1\to3$,
is
\[
 \alpha(a,s)=\alpha(s,d)-\alpha(t,d)+\alpha(a,r).
\]
Substitute $s=c\circ b$ in the cocycle equation on $(a,b,c)$.
This gives
\[
 \alpha(a,b)=\alpha(b,c)+\alpha(c\circ b,d)-\alpha(t,d)
             -\alpha(b\circ a,c)+\alpha(a,r).
\]
The right side is $\delta\lambda(a,b)$ for the one-cochain on
$012$ defined by
\[
 \lambda(a)=\alpha(a,r),\quad
 \lambda(b)=\alpha(b,c)+\alpha(c\circ b,d)-\alpha(t,d),\quad
 \lambda(p)=\alpha(p,c)\quad(p:0\to2).
\]
Its extension again kills the restriction to $012$ and therefore
every degree-two product. These are all pairs among the three
indicated hom-sets.
\end{proof}

\begin{proposition}[Proved: collapse of the longest hom-set]
\label{A7:extreme-arrow-collapse}
Let $C$ be a finite acyclic category on $0<1<2<3<4$ with total-chain
support, and let $C^{\flat}$ be obtained by replacing $C(0,4)$ by
one arrow and retaining every other arrow. The canonical quotient
functor $q:C\to C^{\flat}$ has
\[
 q^*:H^2(BC^{\flat};F)\twoheadrightarrow H^2(BC;F),
 \qquad
 q^*:H^j(BC^{\flat};F)\xrightarrow{\cong}H^j(BC;F)
 \quad(j\geq3)
\]
for every field $F$, and preserves cup products. Consequently, if
$H^2(BC;F)\otimes H^2(BC;F)\to H^4(BC;F)$ is nonzero, then
$|C(0,3)|\geq2$ and $|C(1,4)|\geq2$.
\end{proposition}
\begin{proof}
Keep all original compositions whose target hom-set is not $C(0,4)$,
and send every composition into the new singleton hom-set to its
unique arrow. This is an associative category: every associativity
equation with target $C(0,4)$ becomes an equality of that unique
arrow, and every other equation is unchanged. The map $q$ is the
identity on objects and all other arrows, and sends all $0\to4$
arrows to the new arrow. It is a functor and reduces the arrow count
by $|C(0,4)|-1$.

An arrow $0\to4$ cannot occur as a factor of a nondegenerate
simplex of degree at least two: there is no nonidentity arrow into
$0$ or out of $4$. Thus $q$ identifies the normalized chain bases
in every degree at least two. The boundary $d_2$ can change,
because a composite of two arrows may lie in $C(0,4)$, but
$d_j$ is identical for $j\geq3$. On cochains, therefore, $q^*$
is the identity in degrees at least two and intertwines identical
coboundaries there. The degree-two cocycles coincide, while
\[
 \operatorname{im}\bigl(\delta^1_{C^{\flat}}\bigr)
 \subseteq\operatorname{im}\bigl(\delta^1_C\bigr),
\]
where the inclusion follows from the cochain-map identity for
$q^*$. Taking quotients gives the asserted surjection in degree
two and isomorphisms in degrees at least three. Functoriality
of the Alexander--Whitney product gives compatibility with cups.
No assertion that $q$ is a homotopy equivalence or a degree-two
homology equivalence is used.

If either $C(0,3)$ or $C(1,4)$ is a singleton, then $C^{\flat}$
has that singleton hom-set and the singleton $C^{\flat}(0,4)$.
Proposition~\ref{A7:long-composite-gates} makes all its degree-two
products zero. Lift any two degree-two classes of $BC$ along the
proved surjection. Their product is the pullback of zero, hence
zero. This proves the final assertion.
\end{proof}


\subsection{A universal category determined by the three short products}

\begin{theorem}[Proved: a universal short-product obstruction]
\label{A7:universal-short-model}
Let $\mathcal C$ be a finite acyclic category with five objects and
full support $0<1<2<3<4$. Write its adjacent hom-sets as $A,B,C,D$
and its next three hom-sets as
\[
 U=\mathcal C(0,2),\qquad V=\mathcal C(1,3),
 \qquad W=\mathcal C(2,4).
\]
Let the chronological short-product maps be
$P:A\times B\to U$, $Q:B\times C\to V$ and $R:C\times D\to W$.
There is a finite acyclic category $F(P,Q,R)$, defined below solely
from these seven sets and three maps, such that
\[
 \dim_{\mathbb F_2}H_3(B\mathcal C;\mathbb F_2)
 \ \geq\ \dim_{\mathbb F_2}H_3(BF;\mathbb F_2).
\]
If every degree-two square of $BF$ over $\mathbb F_2$ is zero,
the same is true for $B\mathcal C$. Consequently either
$H_3(BF;\mathbb F_2)\ne0$ or $H^2(BF;\mathbb F_2)=0$
excludes $B\mathcal C\simeq\mathbb{CP}^2$.
\end{theorem}
\begin{proof}
Form two finite bipartite graphs. The first has tagged vertex set
$(A\times V)\sqcup(U\times C)$ and one edge for each $(a,b,c)$,
joining
\[
 (a,Q(b,c))\quad\text{to}\quad(P(a,b),c).
\]
Let $L$ be its set of connected components. The second has tagged
vertex set $(B\times W)\sqcup(V\times D)$, with an edge for each
$(b,c,d)$ joining
\[
 (b,R(c,d))\quad\text{to}\quad(Q(b,c),d).
\]
Let $K$ be its component set. Define $F$ to have the same five
objects, adjacent hom-sets $A,B,C,D$, short hom-sets $U,V,W$,
hom-sets $F(0,3)=L$, $F(1,4)=K$, and a singleton $F(0,4)$.
Its three short compositions are $P,Q,R$. A product in $A\times V$
or $U\times C$ is its graph component in $L$; the two analogous
products into $K$ are defined similarly. All products into $F(0,4)$
are the unique arrow there. Add the five identity arrows.

These rules define all products. Associativity on the object
strings $0123$ and $1234$ is exactly the relation imposed by the
two graphs. Every other triple of cross arrows starts at $0$ and
ends at $4$, so both associative composites are the unique arrow.
Thus $F$ is an actual finite acyclic category. The construction
includes every vertex of the two graphs, even if some element of
$U,V,W$ is not in the image of its short-product map.

First collapse $\mathcal C(0,4)$ to a singleton, obtaining
$\mathcal C_0$. By Proposition~\ref{A7:extreme-arrow-collapse},
the induced map is surjective on second cohomology and an
isomorphism on fourth cohomology. Its normalized chain groups in
degrees at least two, and all boundaries in degrees at least
three, are unchanged. In particular its third homology is
unchanged. We do not assert that this collapse preserves the
projective-plane homotopy type or the rank of second cohomology.

Call an arrow $0\to3$ generated if it is a product from
$A\times V$ or $U\times C$, and define generated arrows $1\to4$
using $B\times W$ or $V\times D$. Fold every nongenerated
$0\to3$ arrow onto a fixed generated one. This is an idempotent
endomorphism of the connecting profunctor across the cut
$\{0\}\mid\{1,2,3,4\}$. No two-arrow product hits a folded
arrow, and every further product lands in the singleton
$\mathcal C_0(0,4)$, so every naturality equation holds.
Likewise fold the nongenerated $1\to4$ arrows across the final
object cut. These operations give actual retractions onto the
subcategory $Y$ containing exactly the generated long arrows.

For either fold, write $i:Y'\hookrightarrow X$ and $r:X\to Y'$
for its inclusion and retraction. The relative mapping-cylinder
quotient is a suspension, as in the proof of
Theorem~\ref{CP2:A-cut-retract}, so all positive-degree relative
cup products vanish. For $\alpha\in H^2(BX;\mathbb F_2)$ put
$\lambda=\alpha-r^*i^*\alpha$. Since $i^*\lambda=0$,
$\lambda$ lifts to relative cohomology and $\lambda^2=0$.
Graded commutativity in characteristic two now gives
\[
 \alpha^2=r^*((i^*\alpha)^2).
\]
Thus zero degree-two squares transfer back across each fold
without a hypothesis on the dimension of second cohomology.

No arrow $0\to3$ or $1\to4$ can be a factor of a normalized
simplex of degree at least three. Moreover, a two-simplex
containing a nongenerated long arrow never occurs in a third
boundary: an inner face producing that long arrow would make
it generated, and an outer face would require such an arrow
already to be a factor of the three-simplex. Deleting these
arrows therefore deletes only zero rows from the third boundary;
the third and fourth chain groups and fourth boundary are
identical. Hence $H_3(BY;\mathbb F_2)=H_3(B\mathcal C_0;\mathbb F_2)$.

Associativity in $Y$ gives an actual functor $F\to Y$ which is
the identity on every adjacent and short hom-set. Its maps on
$L,K$ send graph components to their common long composite,
and are surjective because every long arrow of $Y$ is generated.
It is the identity on normalized chain groups in degrees three
and four, and gives a surjection $q_2:C_2(BF;\mathbb F_2)
\to C_2(BY;\mathbb F_2)$. The identity
\[
 \partial_3^Y=q_2\partial_3^F
\]
implies $\operatorname{rank}\partial_3^Y\leq
\operatorname{rank}\partial_3^F$. Since $\partial_4$ is identical,
this proves the asserted third-homology inequality.
The same identity on $C_3,C_4,\partial_4$ gives an isomorphism
on fourth cohomology. Naturality of the cup product then shows
that zero squares on $BF$ imply zero squares on $BY$. Transfer
back through the two folds, and then through the surjection
$H^2(B\mathcal C_0;\mathbb F_2)\to
H^2(B\mathcal C;\mathbb F_2)$, proves the square assertion.
The projective plane has zero third homology and a nonzero
degree-two square over $\mathbb F_2$, proving the obstruction.
\end{proof}

\begin{remark*}[Computational use]
The two component graphs can be computed before assigning any
long-product table. The theorem therefore allows a single check
of $F(P,Q,R)$ to exclude every longer-table completion of the
same three short maps. This is a proved reduction; it does not
by itself claim that any new finite search has been completed.
\end{remark*}


\subsection{Binary source and sink stalks after extreme-arrow collapse}

\begin{proposition}[Proved: a strengthened binary-stalk obstruction]
\label{A7:binary-extreme-collapse}
Let $\mathcal C$ be a finite acyclic category with full support
$0<1<2<3<4$, and write $h_{ij}=|\mathcal C(i,j)|$.
If
\[
 \max\{h_{01},h_{02},h_{03}\}\leq2
 \quad\text{or}\quad
 \max\{h_{14},h_{24},h_{34}\}\leq2,
\]
then every product of rational degree-two cohomology classes on
$B\mathcal C$ is zero.  No bound on $h_{04}$ is required.
\end{proposition}
\begin{proof}
Collapse $\mathcal C(0,4)$ to a singleton, obtaining an actual
quotient category $\mathcal C_0$.  Proposition~\ref{A7:extreme-arrow-collapse}
gives a surjection
$H^2(B\mathcal C_0;\mathbb Q)\to H^2(B\mathcal C;\mathbb Q)$
and an isomorphism in degree four.  If the first displayed
inequality holds, $\mathcal C_0$ is the cone obtained by adjoining
the trivial source $0$ to its full subcategory on $1,2,3,4$.
The associated categorical left act has one or two states at
every object, including its now singleton stalk at $4$.
Lemma~\ref{CP2:A-high-affine}, in its left-act form, says that
every positive rational product on $B\mathcal C_0$ vanishes.
Lift any two degree-two classes along the stated surjection
and apply naturality of the cup product.  Their product on
$B\mathcal C$ is therefore zero.  The second inequality is
the right-act version at the trivial sink $4$; equivalently,
apply the first case to the opposite category.
\end{proof}

\begin{remark*}
Earlier binary-stalk profile cuts required $h_{04}\leq2$ as
well.  They remain valid conservative cuts.  The proposition
strengthens only the profile reduction; it does not modify
any already completed search or claim completion of a new one.
\end{remark*}


\subsection{Connected partial long-product graphs}

\begin{corollary}[Proved: a partial-assignment obstruction]
\label{A7:partial-long-graph}
Fix the ten nonempty hom-sets of a five-object acyclic category
with full support $0<1<2<3<4$.  Use the notation $A,B,C,D,U,V,W$
of Theorem~\ref{A7:universal-short-model}.  Suppose only some
entries of its short-product maps $P,Q,R$ have been specified.
Form a graph on the entire tagged vertex set
$(A\times V)\sqcup(U\times C)$, inserting the edge
\[
 (a,Q(b,c))\;--\;(P(a,b),c)
\]
whenever both displayed product values have been specified.
If this graph is connected, every actual category completing
the partial assignment has zero degree-two squares over
$\mathbb F_2$.  The same conclusion holds if the corresponding
known-edge graph on $(B\times W)\sqcup(V\times D)$ is connected.
\end{corollary}
\begin{proof}
In any completion, the graph defining the universal hom-set
$F(0,3)$ has exactly the same vertex set and contains every
known edge.  Adding edges cannot disconnect a connected graph,
so this hom-set is a singleton.  The universal category already
has a singleton $F(0,4)$.  Proposition~\ref{A7:long-composite-gates}
therefore makes every degree-two product on $BF$ zero over
$\mathbb F_2$.  Theorem~\ref{A7:universal-short-model} transfers
zero squares from $BF$ to the completed category.  The proof
using the other graph is identical, with $F(1,4)$ a singleton.
All vertices are retained in this argument, including those
involving short arrows not yet used by any specified product.
\end{proof}

\begin{remark*}[Subtree scope]
The implemented search assigns the short cells in the fixed
row-major order $P$, then $Q$, then $R$.  A certificate records
all specified values and marks every remaining cell by $-1$.
Thus it represents an explicit initial prefix, not an
unspecified collection of completed maps.  If its numbers of
unassigned cells are $r_P,r_Q,r_R$, it excludes every one of
the $|U|^{r_P}|V|^{r_Q}|W|^{r_R}$ raw short-map completions of
that prefix, and every associative long-table completion of
each of those maps.  This number counts raw short maps, not
categories or isomorphism classes.  An independent prefix
check verifies legality, connectivity, the exact count and
prefix-freeness within each profile; the final scope join must
also account for every emitted branch and every remaining
full record.  No finite-search completion follows from this
corollary alone.
\end{remark*}


\subsection{Tree extensions of categorical acts}

\begin{lemma}[Proved: compatible trees]
\label{A7:tree-act}
Let $F$ be a nonempty finite-set-valued left or right act on a
finite category $S$.  Suppose that, on each set $F(s)$, one can
choose a tree $T(s)$ with exactly $F(s)$ as its vertex set, so
that every act map sends the endpoints of each tree edge to
equal or adjacent vertices.  Then every positive rational
cohomology product of the act-cone category is zero.
\end{lemma}
\begin{proof}
Give every edge length one. Extend each act map linearly on an
edge whose endpoints have distinct images, and constantly on
an edge whose endpoints have the same image. The hypothesis
makes these continuous maps of trees. They compose strictly:
an edge is sent either to a vertex or linearly to a single
edge, and its subsequent image is again determined by its
endpoints. Thus the extensions form an actual diagram $T$,
with a natural vertex inclusion $F\hookrightarrow T$.

Every $T(s)$ is contractible. Hence the natural map
$\operatorname{hocolim}T\to BS$ is a homotopy equivalence:
on each level of the simplicial bar it is a disjoint union of
products of a contractible tree with a simplex. The act cone
is therefore equivalent to the homotopy cofiber of
$\operatorname{hocolim}F\to\operatorname{hocolim}T$.
The objectwise inclusions are cellular cofibrations. This
cofiber is the pointed bar of the quotients $T(s)/F(s)$, with
the bar of their basepoints collapsed, exactly as in
Lemma~\ref{CP2:A-high-affine}.

Each $T(s)/F(s)$ is a wedge of one circle for each edge of
$T(s)$. Each actor sends any such circle to the basepoint or
linearly to one target circle, with degree $1$ or $-1$.
Choose an orientation on every tree edge. The degree-two maps
on all these circles commute strictly with every actor,
including orientation reversal, since $2(-t)=-2t$ modulo one.
They therefore induce a self-map of the pointed bar.
The reduced cellular chain complex of every quotient has
only degree one, and the self-map is multiplication by two
there. It is consequently multiplication by two on the
entire reduced bar total chain complex and on every positive
cohomology group. For positive classes $u,v$, naturality
gives $2uv=4uv$ over $\mathbb Q$, so $uv=0$.
The same construction with the opposite variance proves
both left- and right-act statements.
\end{proof}

\begin{proposition}[Proved: small trees at the source or sink]
\label{A7:tree-stalk-profile}
Let $\mathcal C$ be a finite acyclic category with full support
$0<1<2<3<4$, and put $h_{ij}=|\mathcal C(i,j)|$.
Every rational degree-two cohomology product of $B\mathcal C$
is zero if either of the following source conditions holds:
\[
\begin{split}
 &h_{02}\leq2,\qquad h_{03}\leq2;\qquad\text{(S1)}\\
 &h_{01}\leq2,\quad h_{03}\leq2,
       \quad h_{12}\leq3.\qquad\text{(S2)}
\end{split}
\]
The same conclusion holds under either dual sink condition
\[
\begin{split}
 &h_{24}\leq2,\qquad h_{14}\leq2;\qquad\text{(T1)}\\
 &h_{34}\leq2,\quad h_{14}\leq2,
       \quad h_{23}\leq3.\qquad\text{(T2)}
\end{split}
\]
There is no restriction on $h_{04}$.
\end{proposition}
\begin{proof}
Collapse the extreme hom-set $\mathcal C(0,4)$ to a singleton.
Proposition~\ref{A7:extreme-arrow-collapse} gives a surjection
on second cohomology from this quotient category to the
original category, preserving products. It suffices to work
in the quotient, which is the left-act cone of its full
subcategory on $1,2,3,4$. Write its source stalks as
$A=\mathcal C(0,1)$, $U=\mathcal C(0,2)$,
$L=\mathcal C(0,3)$ and the singleton at $4$.

Under (S1), choose any tree on $A$. Give each of $U,L$ its
unique tree with at most two vertices, and use the point at
$4$. Every act map carries an edge to an edge or a vertex:
all maps from $A$ have target of size at most two, as do the
remaining nonidentity maps. Lemma~\ref{A7:tree-act} proves
the conclusion, with no bound on $|A|$.

For (S2), write $B=\mathcal C(1,2)$. On the full vertex set
$U$, form a simple undirected graph $G$ by adding, for
each $b\in B$, the pair of distinct values of the map
$A\xrightarrow{\,b\circ-\,}U$, if that map has two values.
There are at most three edges. If $G$ is a forest, extend it
to a tree on $U$. Use the unique tree on $A$ and the unique
tree on $L$, both of which have at most two vertices.
Every map $A\to U$ sends the only possible edge to an edge
of the chosen tree or to a vertex. Every other act map has
target $L$ or the singleton at $4$, so it automatically has
the same property. Lemma~\ref{A7:tree-act} again applies.

The only remaining possibility is that $G$ is a triangle
together with isolated vertices: a simple graph with at
most three edges has no other cycle.
Then $|A|=2$, $|B|=3$, and the three maps $A\to U$ give its
three distinct edges. Let $A=\{a_0,a_1\}$ and write
$P(a,b)=b\circ a$. Any rational two-cycle on the full
subcategory with objects $0,1,2$ has, from its $B$-edge
boundary coefficients, the form
\[
 z=\sum_{b\in B}t_b\bigl[(a_0,b)-(a_1,b)\bigr].
\]
Its remaining boundary equations include
\[
 \sum_b t_b=0,
 \qquad
 \sum_b t_b\bigl([P(a_1,b)]-[P(a_0,b)]\bigr)=0.
\]
The second equation is the incidence-kernel equation of
the triangle with its three independently chosen edge
orientations; any isolated vertices of $U$ contribute
only zero rows. Its kernel is spanned by a vector whose
three entries are all $1$ or $-1$. Their sum is an odd
integer and cannot be zero over $\mathbb Q$. The first
equation therefore forces every $t_b=0$. The full
subcategory on $012$ has $H_2=0$, and hence $H^2=0$,
over $\mathbb Q$.

For any two global degree-two cocycles $\alpha,\beta$,
the restriction of $\alpha$ to $012$ is consequently a
coboundary. Extend a primitive from the edges of $012$
by zero and subtract its global coboundary. The resulting
representative of $\alpha$ vanishes on every $012$
two-simplex. Its Alexander--Whitney product with $\beta$
is zero on every four-simplex, since every such simplex
uses the object string $01234$. This proves (S2).
Opposition gives (T1) and (T2).
\end{proof}

\begin{remark*}
The restriction $h_{12}\leq3$ in (S2) is used only in
the triangle case. With more arrows in that hom-set,
repeated oriented edges can give incidence-kernel vectors
with coefficient sum zero, and the displayed argument
does not assert the same cardinality reduction.
\end{remark*}


\subsection{A five-object model and the exact object minimum}


\begin{remark*}[A singleton extreme hom-set and the fundamental group]
Suppose a category has objects $s,t$ such that every object $x$ admits
arrows $a_x:s\to x$ and $b_x:x\to t$, while $C(s,t)=\{e\}$.
Then its classifying space is simply connected. Indeed, choose
$a_s=1_s$ and $a_t=e$. In the groupoid obtained by inverting every
arrow, an arrow $f:x\to y$ satisfies
\[
 b_y f a_x=e=b_y a_y,
 \qquad\text{hence}\qquad f=a_y a_x^{-1}.
\]
Thus every edge path from $s$ to $x$ equals $a_x$, and every loop
at $s$ is trivial. This groupoid is the fundamental groupoid of
the nerve: its presentation has the category arrows as edges and
one relation $[g][f]=[g\circ f]$ for each two-simplex. All four
models below have a singleton $C(0,4)$ and full linear support,
so this gives a direct proof of their simple connectivity in
addition to the independently checked Tietze transcripts.
\end{remark*}

\begin{theorem}[Proved]\label{A7:model47}
There is a finite acyclic category with five objects and forty-seven
arrows, identities included, whose classifying space is homotopy
equivalent to $\mathbb{CP}^2$.
\end{theorem}
\begin{proof}
We give the entire category. Let $M$ be the first explicit monoid
of Theorem~\ref{A:monoid-upper23}, with its numeric element labels
$0=1$, $x_{ij}=1+5i+j$, and actors $16,\ldots,22$.
Our objects are $0<1<2<3<4$. Arrows $0,\ldots,4$ are their
identities. The nonidentity arrows and an auxiliary monoid label
$\ell$ are as follows; lists correspond in the displayed order.
\[
\begin{array}{c|c|l}
\text{hom-set}&\text{arrow labels}&\text{monoid labels }\ell\\\hline
01&5,6&1,3\\
02&7,8,9,10&1,2,5,3\\
03&11,12,13&1,4,5\\
04&14&1\\
12&15,\ldots,21&1,2,11,16,17,20,21\\
13&22,\ldots,34&1,1,1,4,4,5,11,11,11,14,14,19,21\\
14&35,36&1,11\\
23&37,\ldots,41&1,4,18,19,22\\
24&42,43,44&1,6,11\\
34&45,46&1,6
\end{array}
\]
Identities have their usual products and auxiliary label $0$.
For $b:1\to2$ and $c:2\to3$, the product $c\circ b$ is the
following entry, with rows $b$ and columns $c$:
\[
\begin{array}{c|rrrrr}
 &37&38&39&40&41\\\hline
15&22&25&23&22&24\\
16&24&26&27&25&22\\
17&28&31&29&30&29\\
18&24&26&24&33&33\\
19&23&25&27&33&34\\
20&28&31&34&31&33\\
21&29&32&34&29&34
\end{array}
\]
Every other composable pair of nonidentity arrows
$a:i\to j$, $b:j\to k$ has as product the unique arrow
in $C(i,k)$ with auxiliary label $\ell(a)\ell(b)$, multiplied
in $M$ in this chronological order. All such target hom-sets
have distinct auxiliary labels. Inspection of the monoid table
checks existence of the indicated target and that every entry
of the exceptional table has the prescribed product label.
These rules specify every composable pair and no untyped pair.

Any triple of nonidentity arrows has total object-distance at
least three. Its target hom-set has distinct auxiliary labels,
so the two bracketings agree by associativity in $M$. Thus these
rules define a category; its support and identity endomorphisms
make it acyclic. Its hom profile is
$(2,4,3,1,7,13,2,5,3,2)$, of sum $42$, hence the arrow count
is $47$. A literal $47\times47$ partial composition table is in
\path{acyclic/outputs/cp2_five_47_integral_certificate.json}.

Here are explicit cohomological data. On chronological
two-simplices $(a,b)$ define the integral cochain $\alpha$ to be
zero except in the following cases:
\[
\begin{array}{c|l|r}
a&b&\alpha(a,b)\\\hline
6&17,20,21,28,29,30,31,32,33,34,36&1\\
8,16,20&40,41,43,44&1\\
9,19,21&43&1\\
10&40,41,44&1\\
13,27,34&46&1\\
12,25,26,31,32,38&46&-1
\end{array}
\]
Every pair specified in this table is composable. Let
\[
 y=[5,15]-[5,17]-[6,15]+[6,17].
\]
Let $W$ be the integer matrix below, whose rows are indexed by
$15,\ldots,21$ and columns by $37,\ldots,41$:
\[
W=\begin{pmatrix}
-2&1&1&1&-1\\
1&0&-1&-1&1\\
1&-1&0&0&0\\
1&0&-1&-1&1\\
-1&0&1&1&-1\\
-1&0&1&1&-1\\
1&0&-1&-1&1
\end{pmatrix}.
\]
Define the four-chain
\[
 z=\sum_{b=15}^{21}\sum_{c=37}^{41}W_{bc}
 \bigl([5,b,c,45]-[5,b,c,46]
       -[6,b,c,45]+[6,b,c,46]\bigr).
\]
The indices on $W$ here mean the corresponding row and column.
Direct substitution in the five signed nerve faces verifies
\[
 \delta\alpha=0,\qquad \partial y=0,\qquad \partial z=0,
 \qquad \langle\alpha,y\rangle
 =\langle\alpha\smile\alpha,z\rangle=1.
\]
The program \path{acyclic/code/check_cp2_five_integral.py}
imports no constructor or search code, reconstructs the full
normalized nerve, verifies every displayed equation and all
$927$ associativity equations, and checks the $274$ cocycle
equations integrally.

The integral nerve has cell counts
\[
 (f_0,f_1,f_2,f_3,f_4)=(5,42,174,274,140).
\]
The four boundary ranks are $(4,38,135,139)$ and every nonzero
Smith invariant is $1$. These statements are checked by the
preceding program and independently by
\path{monoid/code/check_acyclic47.py}, with no acyclic
producer imported. The full latter certificate is recorded with
\path{monoid/certificates/acyclic47_integral_independent_verified.json}.
For clarity, Smith invariants of the full boundaries suffice:
$\ker\partial_d$ is saturated in the free group $C_d$, so
$C_d/\operatorname{im}\partial_{d+1}$ is an extension of $H_d$
by the free group $\operatorname{im}\partial_d$. The extension
splits, and the torsion of $H_d$ is exactly the torsion of this
boundary cokernel. The ranks therefore give
\[
 H_*(BC;\mathbb Z)=(\mathbb Z,0,\mathbb Z,0,\mathbb Z),
\]
with zero groups above degree four by acyclicity. The displayed
pairings make $[\alpha]$ a generator and its square a generator.

For simple connectivity, collapse the spanning tree formed by
arrows $5,15,37,45$. The edge-and-triangle presentation has
$38$ generators and $174$ relations. A transcript of $38$
single-occurrence Tietze eliminations reduces it to the empty
presentation. The independent
\path{code/check_acyclic_fundamental_group.py} reconstructs the
typed table and presentation and checks every substitution in
\path{outputs/five_object_47_pi1.json}; its verification output
is \path{outputs/five_object_47_pi1.json.verified.json}. Consequently
$\pi_1(BC)=1$.

We can now apply the homotopy classification in
\cite[Proposition~4C.3, pp.~430--431]{CP2-Hatcher2002}, whose
statement and proof were opened in this run. Its hypotheses are
simple connectivity, CW type, and reduced integral homology
$H_2=\mathbb Z^m$, $H_4=\mathbb Z$, with all other groups zero;
it says the integral cohomology ring determines the homotopy
type. Here $m=1$ and the primitive square gives precisely the
ring of $\mathbb{CP}^2$. Thus $BC\simeq\mathbb{CP}^2$.
\end{proof}

\begin{corollary}[Proved]\label{A7:nu}
The least number of objects of a finite acyclic category of type
$\mathbb{CP}^2$ is exactly $\nu(\mathbb{CP}^2)=5$.
\end{corollary}
\begin{proof}
Proposition~\ref{A7:dimension-chain} gives the lower bound, and
Theorem~\ref{A7:model47} gives a model with five objects.
\end{proof}

\begin{theorem}[Proved]\label{A7:model42}
There is a five-object acyclic model of $\mathbb{CP}^2$ with
forty-two arrows, identities included.
\end{theorem}
\begin{proof}
Use the same monoid and composition convention as in
Theorem~\ref{A7:model47}. Objects and identities are $0,\ldots,4$.
The complete nonidentity hom-sets and auxiliary monoid labels are
\[
\begin{array}{c|c|l}
\text{hom-set}&\text{arrow labels}&\ell\\\hline
01&5,6&1,3\\
02&7,8,9,10&1,2,5,3\\
03&11,12,13&1,4,5\\
04&14&1\\
12&15,\ldots,21&1,2,11,16,17,20,21\\
13&22,\ldots,30&4,4,1,5,14,14,11,19,21\\
14&31,32&1,11\\
23&33,34,35,36&4,18,19,22\\
24&37,38,39&1,6,11\\
34&40,41&1,6
\end{array}
\]
The exceptional composition table, again for $1\to2\to3$, is
\[
\begin{array}{c|rrrr}
&33&34&35&36\\\hline
15&23&24&24&24\\
16&22&25&22&24\\
17&26&28&28&28\\
18&23&24&29&29\\
19&22&25&29&30\\
20&26&30&26&29\\
21&27&30&28&30
\end{array}
\]
Every other composite is the unique arrow in its target hom-set
whose auxiliary label is the chronological monoid product.
As before this completely specifies a category and proves
associativity; the independent programs also check all $773$
associativity equations. Its hom profile is
$(2,4,3,1,7,9,2,4,3,2)$.

An integral two-cocycle is zero except in these cases:
\[
\begin{array}{c|l|r}
a&b&\alpha(a,b)\\\hline
6&17,20,21,26,27,28,29,30,32&1\\
8,16,20&35,36,38,39&1\\
9,19,21&38&1\\
10&35,36,39&1\\
13,25,30&41&1\\
12,22,23,26,27,33&41&-1
\end{array}
\]
Set $y=[5,16]-[5,17]-[6,16]+[6,17]$, and let
\[
W=\begin{pmatrix}
0&0&0&0\\
1&-1&-1&1\\
-1&1&1&-1\\
0&-1&-1&1\\
0&1&1&-1\\
0&1&1&-1\\
0&-1&-1&1
\end{pmatrix},
\]
with rows $15,\ldots,21$ and columns $33,\ldots,36$. Define
\[
 z=\sum_{b,c}W_{bc}
  \bigl([5,b,c,40]-[5,b,c,41]
        -[6,b,c,40]+[6,b,c,41]\bigr).
\]
The literal integral checks give $\delta\alpha=0$,
$\partial y=\partial z=0$, and
$\langle\alpha,y\rangle=\langle\alpha\smile\alpha,z\rangle=1$.
The normalized nerve has counts $(5,37,145,222,112)$ and boundary
ranks $(4,33,111,111)$; every nonzero Smith invariant is $1$.
Thus its integral homology is $(\mathbb Z,0,\mathbb Z,0,\mathbb Z)$,
with no higher cells. These facts are checked independently in
\path{acyclic/outputs/cp2_five_42_integral_certificate.json} and
\path{monoid/certificates/acyclic42_integral_independent_verified.json}.
Both checks reconstruct every chain and use exact integer matrices.

The spanning tree $5,15,33,40$ leaves $33$ generators and $145$
triangle relations. The same independent fundamental-group
checker verifies $33$ elementary Tietze eliminations to the empty
presentation, with transcript
\path{acyclic/outputs/cp2_five_42_pi1.json} and result
\path{acyclic/outputs/cp2_five_42_pi1.json.verified.json}.
Hence the space is simply connected. Proposition~4C.3 of
\cite{CP2-Hatcher2002} now applies exactly as in
Theorem~\ref{A7:model47}, proving its type.
The complete partial table and all literal cochains and cycles
are in the stated integral certificate.
\end{proof}

\begin{remark*}[The sink, join, and product tools]
The forty-two-arrow category is obtained from its full subcategory
$D$ on $0123$ by adding the sink $4$ along the presheaf
$F:D^{\mathrm{op}}\to\mathrm{FinSet}$, $F(i)=C(i,4)$,
where $F(f)(a)=a\circ f$. Its stalk sizes are $(1,2,3,2)$.
Thus the sink lemma, \cite[Lemma~2.1]{CP2-Carcacia2026}, identifies
its classifying space with the homotopy cofiber of the explicitly
specified projection $B\operatorname{El}(F)\to BD$. This is a
second, attachment-based description of the same fully given
category; the integral checks above compute the resulting type.

The join lemma, \cite[Lemma~1.2]{CP2-Carcacia2026}, also provides
a useful obstruction: no categorical join of two nonempty finite
categories can model $\mathbb{CP}^2$. Indeed $B(C\star D)$ is
$BC*BD$. After choosing basepoints, collapse the contractible
subcomplex consisting of the two basepoint cones. The quotient
is the reduced suspension of the smash product. Its positive
cohomology products vanish, whereas the degree-two square on
$\mathbb{CP}^2$ does not. This explains why the five-object
join constructions that suffice for spheres cannot give this
model. Finally, the same lemma gives
$B(C\times D)=BC\times BD$ and the product object bound.
Applied to two three-object sphere categories it gives a
nine-object model of $S^2\times S^2$, whose second homology
has rank two. Passing to a categorical symmetric quotient
requires a separate comparison; the explicit quotient checks
reported above show why the product bound alone does not
supply a projective-plane model.
\end{remark*}

\begin{theorem}[Proved: a thirty-nine-arrow model]
\label{A7:model39}
There is a finite acyclic category with five objects and thirty-nine
arrows whose classifying space is homotopy equivalent to
$\mathbb{CP}^2$.
\end{theorem}
\begin{proof}
Let $M$ be the explicit monoid of Theorem~\ref{A:monoid-upper23},
with elements numbered $0,\ldots,22$ and identity $0$.
Take objects $0<1<2<3<4$ and the following nonidentity arrows.
In each row, the arrow numbers and monoid labels correspond in order.
\[
\begin{array}{c|c|l}
(i,j)&\text{arrow numbers}&\text{monoid labels}\\\hline
(0,1)&5,6&1,3\\
(0,2)&7,8,9,10&1,2,3,5\\
(0,3)&11,12,13&1,4,5\\
(0,4)&14&1\\
(1,2)&15,16,17,18,19,20&1,11,16,17,20,21\\
(1,3)&21,22,23,24,25,26,27&1,5,11,14,15,19,21\\
(1,4)&28,29&1,11\\
(2,3)&30,31,32,33&18,19,21,22\\
(2,4)&34,35,36&1,6,11\\
(3,4)&37,38&1,6
\end{array}
\]
Add the five identity arrows, numbered by their objects and labelled
by $0$. There are no other arrows. To compose an arrow with monoid
label $x$ followed by one with label $y$, take the unique arrow in
the required hom-set with monoid label $xy$. The displayed hom-sets
are closed under these products: they are exactly the iterated
products of the four adjacent label sets. Thus this rule specifies
all products, and associativity follows from that of $M$.
The literal checker also verifies every one of the $683$ typed
associativity equations, including identities.

The normalized nerve has simplex counts $(5,34,128,192,96)$ and no
simplices in higher degrees. Its integral boundary ranks are
$(4,30,97,95)$ and every nonzero Smith invariant is one. Hence
its integral homology is $\mathbb Z$ in degrees $0,2,4$ and zero
otherwise. Here is explicit primitive cup-product data. Define
an integral degree-two cochain $\alpha$ by the following nonzero
values, with all unspecified values zero:
\[
\begin{array}{c|l|r}
a&b&\alpha([a,b])\\\hline
6&16,19,20,23,24,25,26,27,29&1\\
8,19&31,32,33,35,36&1\\
9&31,32,33,36&1\\
10,18,20&35&1\\
12,24&38&-1\\
13,22,25,27,32&38&1
\end{array}
\]
This is the pullback of the verified monoid cocycle and satisfies
all $192$ integral cocycle equations. The integral two-cycle
\[
 y=[5,15]-[5,16]-[6,15]+[6,16]
\]
satisfies $\alpha(y)=1$. Let $W$ be the matrix whose rows are
indexed by arrows $15,\ldots,20$ and columns by $30,\ldots,33$:
\[
 W=\begin{pmatrix}
 0&1&-1&0\\
 0&-1&1&0\\
 -1&0&-1&1\\
 1&0&1&-1\\
 1&0&1&-1\\
 -1&0&-1&1
 \end{pmatrix}.
\]
In the tensor notation for four-simplices, put
\[
 z=(5-6)\otimes W\otimes(37-38).
\]
Direct integral boundary and Alexander--Whitney calculations give
$\partial z=0$ and $\langle\alpha\smile\alpha,z\rangle=1$.
Thus $\alpha$ and its square are primitive in degrees two and four.
These assertions, including all integer matrices and chain vectors,
are certified in
\path{acyclic/outputs/cp2_five_39_integral_certificate.json}.
A second implementation, importing no construction or checker code
from this directory, independently computes saturated integral
kernels and the primitive pairing; its receipt is
\path{monoid/certificates/acyclic39_integral_independent_verified.json}.

For the fundamental group, take the spanning tree with arrows
$5,15,30,37$. The resulting presentation has thirty generators
and $128$ relations. The full transcript contains thirty
single-occurrence Tietze eliminations and ends with the empty
presentation. The independently reconstructed and checked transcript
is \path{acyclic/outputs/cp2_five_39_pi1.json}, with a
\texttt{PASS} receipt in the file of the same name followed by
\texttt{.verified.json}. Hence the nerve is simply connected.
Hatcher's Proposition~4C.3, in the form already checked above,
now gives the asserted homotopy equivalence.
\end{proof}

\begin{theorem}[Proved: a thirty-eight-arrow model]
\label{A7:model38}
There is a finite acyclic category with five objects and thirty-eight
arrows whose classifying space is homotopy equivalent to
$\mathbb{CP}^2$.
\end{theorem}
\begin{proof}
Use the same explicitly numbered monoid $M$ and five ordered objects
as above. The complete list of nonidentity arrows is
\[
\begin{array}{c|c|l}
(i,j)&\text{arrow numbers}&\text{monoid labels}\\\hline
(0,1)&5,6&3,5\\
(0,2)&7,8,9,10&1,2,3,5\\
(0,3)&11,12,13&1,4,5\\
(0,4)&14&1\\
(1,2)&15,16,17,18,19&5,16,17,18,20\\
(1,3)&20,21,22,23,24,25&1,5,11,18,19,21\\
(1,4)&26,27,28&1,6,11\\
(2,3)&29,30,31,32&1,18,21,22\\
(2,4)&33,34,35&1,6,11\\
(3,4)&36,37&6,19
\end{array}
\]
Add identities $0,\ldots,4$, all with monoid label $0$.
Composition is again the unique arrow in the required hom-set
labelled by the product of the two labels in chronological order.
Every longer hom-set is exactly the iterated product set of the
four adjacent ones. Closure and associativity therefore follow
directly from $M$; the separate literal checker also verifies all
$621$ typed associativity equations.

The normalized nerve has simplex counts $(5,33,117,166,80)$.
The integral boundary ranks are $(4,29,87,79)$ and every nonzero
Smith invariant is one. Thus the integral homology is
$\mathbb Z$ in degrees $0,2,4$ and zero in all other degrees.
Define the integral two-cochain $\alpha$ by these nonzero values:
\[
\begin{array}{c|l|r}
a&b&\alpha([a,b])\\\hline
5&19,22,24,25,28&1\\
6&27&1\\
8,19&31,32,34,35&1\\
9&31,32,35&1\\
10,15,17&34&1\\
12&36&-1\\
13,21,25,31&36&1
\end{array}
\]
All omitted values are zero. This is the pulled-back monoid
cocycle, and all $166$ integral cocycle equations hold.
In the formal tensor notation for composable arrow strings, put
\[
\begin{aligned}
y&=(5-6)\otimes(-16+17-18+19),\\
w&=(-29+30+31-32)\otimes(36-37),\\
z&=(5-6)\otimes(-16+17-18+19)
   \otimes(-29+30+31-32)\otimes(36-37).
\end{aligned}
\]
The first two expressions are integral two-cycles; the last is an
integral four-cycle with $64$ nonzero coefficients. Direct boundary
calculations and the Alexander--Whitney formula give
\[
 \alpha(y)=\alpha(w)=1,
 \qquad \langle\alpha\smile\alpha,z\rangle=1.
\]
This proves that the degree-two class and its square are primitive.
The literal category, every nonzero Smith invariant, and the
explicit cochain $\alpha$ and cycle $z$ are recorded in
\path{acyclic/outputs/cp2_five_38_integral_certificate.json}.
The second independent implementation gives the same integral
homology and primitive pairing in
\path{monoid/certificates/acyclic38_integral_independent_verified.json}.

The unique arrow $0\to4$ gives simple connectivity by the direct
groupoid argument at the start of this subsection. Independently,
the spanning tree $5,15,29,36$ gives $29$ generators and $117$
triangle relations. The transcript in
\path{outputs/five_object_38_pi1.json}, with its separate verified
receipt, consists of $29$ elementary Tietze eliminations to the
empty presentation. The integral homology and primitive square
now give the claimed homotopy type by
\cite[Proposition~4C.3]{CP2-Hatcher2002}.
\end{proof}

\begin{remark*}[The remaining arrow minimum]
Let $a_5(\mathbb{CP}^2)$ denote the least arrow count of a
five-object acyclic model. The proved interval is
$31\le a_5(\mathbb{CP}^2)\le38$. Its exact value is still
unresolved; this is the remaining part of A.7.
\end{remark*}


\begin{proposition}[Proved: a large leaf in a small three-object star]
\label{CP2:A-star-large-leaf}
Let $C$ be a connected three-object category whose acyclic support
is a star, with no chain of three distinct objects. If $C$ has at
most seventeen arrows, identities included, and a leaf corner has
order at least eight, then $BC$ is not homotopy equivalent to
$\mathbb{CP}^2$.
\end{proposition}
\begin{proof}
After taking the opposite category, write the two nonempty cross-homs
as $0\to1$ and $0\to2$. Let the corner orders be $a,b,c$, with
$b\geq8$, and let the stalk sizes be $s,t$. Thus
\[
 a+b+c+s+t\leq17,\qquad a,b,c,s,t\geq1.
\]
The center order satisfies $a\leq6$. Let $E_1,E_2$ be the full
edge subcategories on $\{0,1\}$ and $\{0,2\}$. Every nerve simplex
lies in one edge subcategory, so $BC=BE_1\cup_{BM}BE_2$, where
$M=C(0,0)$. These are CW subcomplexes and this union is their
homotopy pushout.

If $a=6$, the budget forces $b=8$ and $c=s=t=1$. All cross-homs
are singletons, so Lemma~\ref{CP2:A-singleton-support} identifies
rational second cohomology with that of the support tree. It is
zero. If $a=5$, both $s,t$ are at most two. A singleton cross-hom
is the join of the two corner classifying spaces, and has zero
rational second cohomology: both factors are connected and finite
monoids have zero rational first homology. A two-arrow cross-hom
has zero rational degree-two products by
Lemma~\ref{CP2:A-high-two-cross}. Also $H^3(BM;\mathbb Q)=0$ by
Corollary~\ref{CP2:A-five-types}. Mayer--Vietoris therefore makes
\[
 H^4(BC;\mathbb Q)\longrightarrow
 H^4(BE_1;\mathbb Q)\oplus H^4(BE_2;\mathbb Q)
\]
injective. Every global degree-two product restricts to zero in
both components, hence is zero.

It remains that $a\leq4$. Choose an idempotent $e$ in the minimal
ideal $I$ of $M$ and write $G=eMe$. The small-minimal-ideal proof
in Lemma~\ref{CP2:A-small-minimal-ideal} shows that at least one
of $Me=G$ and $eM=G$ holds, and that $BM\simeq BG$.
If $G$ is trivial, $BM$ is contractible. Collapsing this CW
subcomplex gives a homotopy equivalence
\[
 BC\simeq BE_1\vee BE_2.
\]
Indeed the common contractible subcomplex is collapsed in each
edge as well, and those quotient maps are homotopy equivalences.
If the wedge had type $\mathbb{CP}^2$, its free-product fundamental
group would force both edge spaces to be simply connected. Its
integral homology is the direct sum of theirs in positive degrees.
The nonzero primitive square forces the rank-one second and fourth
homology to lie in the same edge summand. That edge has exactly
the integral homology and primitive cup square of $\mathbb{CP}^2$,
and no other reduced homology. The recognition theorem makes it
an actual projective-plane model. Its arrow count is at most
$17-1-1=15$, contradicting the newly proved two-object baseline,
Theorem~\ref{CP2:A-two-object-baseline-sixteen}. This argument uses
that theorem for an actual homotopy type, not merely a rational
cohomology pattern.

Suppose finally that $G$ is nontrivial. In this small order both
$Me=G$ and $eM=G$ must hold. For example, if $a_0\in Me\setminus G$,
the map $G\to a_0G$, $g\mapsto a_0g$, is injective: multiplication
by $e$ reduces equality to equality in $G$, since $ea_0\in G$.
The sets $G$ and $a_0G$ are disjoint and lie in $I$, while the global
identity lies outside $I$ unless $M$ itself is a group. This would
require at least $2|G|+1\geq5$ elements. The other side is dual.
It follows that $I=MeM=G$. The averaging idempotent
\[
 f=|G|^{-1}\sum_{g\in G}g\in\mathbb QM
\]
satisfies $mf=f=fm$ for every $m\in M$: multiplication by $m$
acts on the ideal group by a group translation. Thus
$(\mathbb QM)f$ and $f(\mathbb QM)$ are the one-dimensional trivial
left and right modules, respectively, and both are projective.

Collapse the center to a trivial object, and replace each incident
stalk by its actual $M$-orbit set. Precomposition gives each stalk
its right $M$-action; the commuting leaf action descends to the
orbit set. These assignments define an actual quotient functor.
Lemma~\ref{CP2:A-high-orbits} makes it a rational homology and
cohomology-ring equivalence in all degrees. Its target $C'$ has
trivial center, hence
\[
 BC'=BX_1\vee BX_2,
\]
where $X_1,X_2$ are ordinary act cones, up to opposition, over the
leaf monoids. Their state counts are at most $s,t$. If $BC$ had
type $\mathbb{CP}^2$, this wedge would have its rational homology
and cohomology ring, and both summands would have bounded rational
homology.

Here $a\geq2$, so the smaller leaf has
$c\leq17-a-b-2\leq5$. Its act cone has zero rational degree-two
products by Theorem~\ref{CP2:A-bounded-nine-acts}. For the large
leaf, the number of states is at most
\[
 s\leq17-a-b-c-t\leq13-b.
\]
For $b=8$ or $9$, the same bounded-nine theorem applies. For
$b\geq11$, there are at most two states, so
Lemma~\ref{CP2:A-high-affine} gives zero products. The only remaining
case is $b=10$, with at most three states. Since the other wedge
summand has zero products, the nonzero square of the global
rank-one degree-two class would lie in $BX_1$. The direct-sum
homology then gives $BX_1$ exactly the rational homology and
cohomology ring of $\mathbb{CP}^2$. This contradicts
Proposition~\ref{CP2:A-ten-six-state-rational}. Every case is excluded.
\end{proof}


\subsection{Two further reductions for the four-object seventeen-arrow scope}

\begin{lemma}[Proved: a height-one support is a retract]
\label{CP2:A-height-one-support}
Let $C$ have acyclic support with no composable pair of cross-object
arrows. Its support poset $P$ is a categorical retract of $C$.
Consequently, if the undirected support graph has positive cycle rank,
then $H_1(BC;\mathbb Z)\ne0$.
\end{lemma}
\begin{proof}
The support functor $C\to P$ sends every corner to its identity and
every cross arrow to the corresponding unique poset arrow. Choose
one representative in each nonempty cross hom-set. These choices,
together with the corner identities, define a functor $P\to C$:
there are no composable pairs of nonidentity arrows in $P$ whose
composition would impose an additional condition. Its composite
with $C\to P$ is the identity. The nerve of $P$ is exactly its
undirected support graph. The induced retraction on first homology
proves the claim. In particular every four-object height-one crown
is excluded, regardless of its corner monoids or hom-set sizes.
\end{proof}

\begin{proposition}[Proved: the eight remaining large-corner profiles]
\label{CP2:A-four17-large-corner-profiles}
No four-object category of any of the following hom-cardinality
matrices has type $\CP^2$:
\[
\begin{gathered}
\begin{pmatrix}
1&0&0&1\\0&a&s&t\\0&0&b&0\\0&0&0&c
\end{pmatrix},
\quad 8\leq a\leq9,\quad b,c\leq2,\quad s,t\geq2,
\quad a+b+c+s+t+2=17;\\[3pt]
\begin{pmatrix}
1&0&0&2\\0&1&s&2\\0&0&b&0\\0&0&0&8
\end{pmatrix},
\quad (b,s)=(2,1)\text{ or }(1,2);\\[3pt]
\begin{pmatrix}
1&0&1&1\\0&8&0&3\\0&0&1&1\\0&0&0&1
\end{pmatrix}.
\end{gathered}
\]
There are respectively five, two, and one matrices in these families.
\end{proposition}
\begin{proof}
We use actual unions of full subcategory nerves. If every simplex
lies in $ND$ or $NE$, their intersection is a full subcategory $F$,
and $H_j(BF;\mathbb Q)=0$ for $j>0$, Mayer--Vietoris gives
\[
 H^j(BC;\mathbb Q)\cong
 H^j(BD;\mathbb Q)\oplus H^j(BE;\mathbb Q),\qquad j\geq2.
\]
Thus bounded rational homology of $BC$ implies the same for the two
pieces in degrees at least two. Moreover, products vanish on $BC$
if they vanish on both pieces, because restriction is injective
in degree four. All the intersections used below have corner
order at most two and hence are rationally acyclic.

For the first family take $D$ on objects $1,2,3$ and $E$ on $0,3$.
Every simplex belongs to one of these pieces and $F=\{3\}$.
The edge $E$ has a singleton cross hom-set and a trivial source,
so $BE$ is contractible. The category $D$ is a source star with
center order $a\leq9$ and leaves of orders at most two. Exact
orbit comparison at the leaves, Lemma~\ref{CP2:A-high-orbits},
makes the leaves trivial by a rational cohomology-ring
equivalence. Identifying these two trivial leaves gives an
ordinary act cone over the center, whose space is the reduced
star wedged with a circle, as in
Lemma~\ref{CP2:A-star-leaf-reductions}(3). It has bounded rational
homology, so Theorem~\ref{CP2:A-bounded-nine-acts} kills all its
degree-two products. Those of $BD$, and then of $BC$, vanish.

For the second family use $D$ on $0,1,3$ and $E$ on $1,2$,
intersecting at the trivial corner $1$. The first piece is the
opposite of the same bounded-nine act construction: its center
has order eight and its two leaves are trivial. The second piece
is either a singleton edge or a two-arrow edge with trivial
corners, so has zero second rational cohomology. The same
restriction argument applies.

For the third matrix use $D$ on $0,2,3$ and $E$ on $1,3$.
The first category is the three-element chain and is contractible;
the intersection is the trivial corner $3$. The second category
is an ordinary act cone over an eight-element monoid. Boundedness
and the bounded-nine act theorem again give zero products.
In every case this contradicts the nonzero rational square on
$\CP^2$.
\end{proof}


\subsection{New proofs of the local high-object reductions}
\label{CP2:A-high-local}
The following lemmas were reconstructed and checked at their actual
points of use in the excess arguments.  Their conclusions concern
rational products unless an actual homotopy equivalence is stated.
In particular, bounded homology of a whole category is never silently
assigned to one of its corners.  The subsequent theorem gives a
complete budget proof for acyclic support.  Explicit strong-component
collapses then remove the support restriction without using a small
category composition census.

\begin{lemma}[Proved: two flag-complex budgets]
\label{CP2:A-high-flag}
Let $G$ be a finite connected graph of cycle rank $R$, and let $X(G)$
be its clique complex.  If $R\leq6$, then
$H_2(X(G);\mathbb Q)=0$.  If $R\leq8$, then
$H_4(X(G);\mathbb Q)=0$; in this latter range $G$ contains at most
one clique of size five and none of size six.
\end{lemma}
\begin{proof}
Suppose first that second homology is nonzero.  Among induced
subgraphs having this property choose one with the fewest vertices,
and pass to its relevant connected component.  Deleting any vertex
$v$ kills second homology.  The deletion together with the star of
$v$ covers the clique complex, and their intersection is the link
of $v$.  Mayer--Vietoris therefore injects the nonzero second
homology into the first homology of that link.  A clique complex
on at most three vertices has zero first homology, so every vertex
has at least four neighbors.  There are at least six vertices:
with five vertices the degree condition forces a complete graph,
whose clique complex is contractible.  If their number is $a$ and
the edge number is $b$, then $b\geq2a$ and the cycle rank is at
least $a+1\geq7$.  Cycle rank is monotone under taking subgraphs,
which proves the first assertion.

Two distinct five-cliques sharing $t$ vertices have union with
cycle rank at least $12-\binom t2+t-1$ for $1\leq t\leq4$;
its minimum is nine, at $t=4$.  Disjoint five-cliques contribute
twelve independent cycles.  A six-clique already has cycle rank
ten.  Thus at most one four-simplex occurs when $R\leq8$, and
there is no five-simplex.  Its boundary, if it occurs, is a nonzero
sum of its five distinct faces.  The fourth chain boundary is
therefore injective, proving the second assertion.
\end{proof}

\begin{lemma}[Proved: a graph cofiber over the five-element band]
\label{CP2:A-high-P5-graph}
Let $P=\{1\}\sqcup\{x_{ij}:0\leq i,j\leq1\}$, with
$x_{ij}x_{kl}=x_{il}$.  Let a finite graph $\Gamma$ carry a
right $P$ action sending each oriented edge either to an
oriented edge by its linear parameter or to a vertex.
The cofiber of $\Gamma_{hP}\to BP$ has zero rational
degree-two products.  Consequently this conclusion holds
for every collage between $P$ and
$R_0=\{1,p,q\}$, $ab=b$ for $a,b\in\{p,q\}$,
with any nonempty finite $(R_0,P)$-biset.
\end{lemma}
\begin{proof}
Write $A=\mathbb QP$, $e=x_{00}$, $f=1-e$,
$a=x_{10}-x_{00}$ and $b=x_{01}-x_{00}$.  There is an exact
projective resolution of the trivial left module
\[
 0\longrightarrow Ae\xrightarrow{\cdot b}Af
   \xrightarrow{\cdot a}Ae\longrightarrow\mathbb Q
   \longrightarrow0.
\]
To check it directly, bases for the first and last projective
are $x_{00},x_{10}$; a basis for the middle projective is
$f,b,c=x_{11}-x_{10}$.  The first map sends the basis to
$b,c$, and the second sends $f$ to $a$ and kills $b,c$.
The last augmentation has kernel $\mathbb Qa$.

The images of the two actors in column $j$ are the same
retraction subgraph $\Gamma_j$.  The maps $x_{01}$ and
$x_{00}$ restrict to inverse isomorphisms between
$\Gamma_0$ and $\Gamma_1$.  If these subgraphs share a
vertex, it is fixed by every band element, hence by $P$.
It supplies a section of $\Gamma_{hP}\to BP$.
Mayer--Vietoris then makes all second classes of the
cofiber relative to $BP$.  They lift from its suspension
quotient and have zero products.

Otherwise the two subgraphs have disjoint vertex and edge
sets.  Put $W=H_1(\Gamma;\mathbb Q)$.  Tensoring the
displayed resolution with $W$ identifies
\[
 \operatorname{Tor}_2^A(W,\mathbb Q)
   =\ker\bigl(We\xrightarrow{\cdot b}Wf\bigr).
\]
Here $We=H_1(\Gamma_0;\mathbb Q)$, and $wb$ is the
difference between the image of the cycle $w$ in
$\Gamma_1$ and its original cycle in $\Gamma_0$.
A graph has no two-dimensional boundaries, and these
two cycles have disjoint edge supports.  Their difference
is zero only if $w=0$.  Thus this Tor group is zero.

This also proves the required fourth-homology vanishing
when $\Gamma$ is disconnected.  In fact, put
$L=\widetilde H_0(\Gamma;\mathbb Q)$ and let $K$ be the
two-term complex
$[C_1\Gamma\to\ker(C_0\Gamma\to\mathbb Q)]$.
The triangle $W[1]\to K\to L$ and the length-two
resolution show that $H_3(K\otimes_A^{\mathbf L}\mathbb Q)
=\operatorname{Tor}_2^A(W,\mathbb Q)$: the possible
outer terms are $\operatorname{Tor}_4^A(L,\mathbb Q)$
and $\operatorname{Tor}_3^A(L,\mathbb Q)$, both zero.
The augmented action-bar cofiber has this group as its
fourth homology.  Therefore its degree-two products are
zero in the disjoint-image case also.

For the biset conclusion, the left $R_0$-action consists
of two retractions $P_0,Q_0$ with a common image $U$.
Its action category has a natural graph model: vertices
are $U$, and each state $t\notin U$ gives an edge from
$P_0t$ to $Q_0t$.  Indeed, a state outside $U$ has only
its identity and its two outgoing arrows to $U$.  The
endomorphism monoid of each vertex of $U$ is $R_0$,
whose classifying space is contractible.  The two
outgoing-arrow cones attach along these contractible
spaces; collapsing them gives the graph with the outside
state as a bivalent vertex, then its stated edge.
On nerve simplices this collapse uses the sum of the
barycentric coordinates after the initial outgoing
arrow, so it commutes with every equivariant set map.
Such a map sends an edge to the edge of the image state
or collapses it when that state lies in $U$.  The
commuting right $P$ action therefore gives exactly the
graph action already treated.  Since $BR_0$ is
contractible, the collage is the stated cofiber.
\end{proof}

\begin{corollary}[Proved: two-object acyclic support through eleven arrows]
\label{CP2:A-high-two-objects}
A connected category with two objects, acyclic support, and at most
eleven arrows cannot have the homotopy type of $\mathbb{CP}^2$.
More strongly, its rational degree-two products are zero.
\end{corollary}
\begin{proof}
Write $a,b$ for the two corner orders and $t\geq1$ for
the cross-hom size.  If $t=1$, the collage is the join
of its two connected corner spaces and has zero second
cohomology.  If $t=2$, use
Lemma~\ref{CP2:A-high-two-cross}.  Suppose $t\geq3$.
If both corners have zero second cohomology, the
two-block lemma applies.  Otherwise one corner has
order at least five by Corollary~\ref{CP2:A-five-types}.
Its order is at most seven and the other corner has
order at most three.  With other corner of order at most
two, exact orbit comparison reduces to
Corollary~\ref{CP2:A-high-seven-act}.
The sole remaining budget is $(a,b,t)=(5,3,3)$ or its
opposite, and the positive corner is $P$ of the preceding
lemma.  If the order-three corner has a projective trivial
module on the required side, orbit comparison gives a
two-row act cone.  Otherwise its minimal ideal is the
two-element right-zero band, after choosing the direction
of the collage, and Lemma~\ref{CP2:A-high-P5-graph}
applies.  All connected cases have zero products.
\end{proof}

\begin{lemma}[Proved: a cut containing all parallel hom-sets]
\label{CP2:A-high-parallel-cut}
Let $C$ have acyclic support of cycle rank at most six.
Suppose no two hom-sets of size greater than one occur
serially: for such $C(x,y),C(z,w)$ one never has $y\leq z$.
If every global degree-two class restricts to zero at
each corner with nonzero second cohomology, all rational
degree-two products on $BC$ vanish.
\end{lemma}
\begin{proof}
We use the following join observation, including disconnected
predecessor categories.  If $D$ is nonempty and $M$ is a finite
monoid, then $BD*BM$ has zero rational second cohomology.
Indeed $BM$ is connected and $H_1(BM;\mathbb Q)=0$, since
its universal group is finite.  The reduced join Kunneth
formula has no nonzero summand in degree two.  The inclusion
of $BM$ in this join is null-homotopic, by choosing any point
of $BD$.  These facts also justify the join cuts below.

Let $D$ be the downset generated by the sources of the
parallel hom-sets.  Each such source lies in $D$, and
its target cannot lie there: otherwise that target is
at or below the source of another parallel hom-set,
contrary to the assumption.  Every parallel hom-set
therefore crosses from $D$ to its complementary upset.
The full categories on either side have singleton cross
hom-sets.  Each of their components with more than one
object has second cohomology equal to that of its support
by Lemma~\ref{CP2:A-singleton-support}, and the latter
vanishes by Lemma~\ref{CP2:A-high-flag}.  On a one-object
component, global classes restrict to zero by the
hypothesis.  All global classes therefore vanish on
both blocks.  Lemma~\ref{CP2:A-two-block} proves the claim.
When there are no parallel hom-sets, apply the same
singleton-support argument directly, component by component.
\end{proof}

\begin{lemma}[Proved: two parallel excesses]
\label{CP2:A-high-two-parallel}
Suppose $C$ is connected with acyclic support, every
corner has order at most four, its support cycle rank
is at most six, and
$\sum_{x<y}(|C(x,y)|-1)\leq2$.  If
$\dim H^2(BC;\mathbb Q)\leq1$, its degree-two squares vanish.
\end{lemma}
\begin{proof}
All corners are rationally acyclic.  Unless two doubled
hom-sets occur serially, the preceding lemma applies.
Write the two serial hom-sets as $a\to b$ and $c\to d$,
where $b\leq c$.  If $b<c$, the full downset of $c$ has
singleton hom-sets into $c$ and is the categorical join
of its predecessor category and its corner.  This join
has zero second rational cohomology: the corner is
rationally acyclic.  Its complement has only singleton
cross hom-sets and zero second cohomology by the support
budget.  The two-block lemma applies.

It remains that the two doubled sets are $a\to b$ and
$b\to c$.  If $b<x$ with $x\ne c$, the full downset
$D$ of $b$ has only singleton arrows into $x$.
Consequently $D\cup\{x\}$ is a join with the rationally
acyclic corner at $x$ and has zero second cohomology.
Every global second class vanishes on $D$.  The complement
has singleton cross hom-sets and zero second cohomology,
so the two-block lemma again applies.  An extra predecessor
of $b$ is handled by opposition.

Otherwise $b$ has just the two neighbors $a,c$.  The
category is the actual nerve union of the full triangle
on $a,b,c$ and the deletion of $b$.  Their intersection
is the singleton edge $a\to c$, a rationally acyclic
join.  The deletion has zero second cohomology by the
singleton-support lemma.  Mayer--Vietoris identifies
the triangle's second cohomology with the global group
and makes fourth-cohomology restriction to the two
pieces injective.  Its hom sizes are $(2,2,1)$, so
Lemma~\ref{CP2:A-high-small-binary} kills its rank-one
square.  Its restriction to the other piece is already
zero, proving the result.
\end{proof}

\begin{lemma}[Proved: finite affine acts and two states per object]
\label{CP2:A-high-affine}
Let $T\subset\mathbb Q$ be a nonempty finite right monoid act,
whose maps are affine of slopes $0,1,-1$.  Adjoin a trivial sink
to the monoid with outgoing arrows $T$.  Every positive rational
cohomology product of this act-cone category is zero.

More generally, the same conclusion holds for an act cone over
any finite category $S$ when its right act $F$ has one or two
states at each object.  Both statements have left-act duals.
\end{lemma}
\begin{proof}
A slope-one affine self-map of a finite set with at least two
elements is the identity.  A slope-minus-one map is its unique
reflection about the midpoint of its extreme elements.  Extend
the action to the interval $I=[\min T,\max T]$.  All affine
identities holding on $T$ hold on $I$, so this is a strict action.
The interval bar maps equivalently to the monoid classifying
space, level by level, because $I$ is contractible.  Consequently
the act cone is equivalent to the cofiber of the inclusion of
the $T$ bar in the $I$ bar.  This is the pointed bar with fiber
$I/T$, with its basepoint bar collapsed.  The fiber is a wedge
of circles indexed by consecutive elements of $T$.  Its actors
are identities, a reflection permuting and reversing circles,
or pointed constants.  On each circle take the map of degree two.
It commutes strictly with all these actors: in a circle coordinate,
$2(-t)=-2t$ modulo one.  It induces a self-map of the pointed bar.
The reduced cellular chain complex of the fiber is concentrated
in degree one, and this self-map is multiplication by two on
that whole complex.  The same is true in every degree of the
reduced bar total complex, hence on all its positive cohomology.
For positive classes $u,v$, functoriality gives $2uv=4uv$, and
thus $uv=0$ over $\mathbb Q$.  A singleton act has contractible
cone and presents no exception.

For the categorical assertion, replace each $F(s)$ by its convex
simplex $I(s)$ and each set map by its affine extension.  These
extensions compose exactly.  Each $I(s)$ is a point or interval,
and each quotient $I(s)/F(s)$ is a point or circle.  The induced
maps are pointed constants or circle maps of degree $1$ or $-1$.
Their degree-two self-maps form a strict natural transformation.
The cofiber of $\operatorname{hocolim}F\to\operatorname{hocolim}I$
is the pointed bar of these quotients, and
$\operatorname{hocolim}I\to BS$ is an equivalence.  It is thus
the required act cone.  The identical reduced-chain argument
proves the claim.  Reversing arrows proves the dual statements.
\end{proof}

\begin{lemma}[Proved: two minimal rows]
\label{CP2:A-high-two-row-act}
Suppose the minimal ideal of a finite monoid $M$ has at most
two rows.  For every nonempty finite right $M$-act $T$, all
rational degree-two products of its act cone vanish.  The
left-act statement holds with columns in place of rows.
\end{lemma}
\begin{proof}
Write $R=\mathbb QM$.  For a minimal idempotent $e$, the corner
$H=eMe$ is a finite group.  Average it to the algebra idempotent
$f=|H|^{-1}\sum_{h\in H}h$.  The projective left module $Rf$
has one basis vector per minimal row: the right $H$-orbits of
$Me$ are free, and multiplying by $f$ averages each such orbit.
Left multiplication acts by maps of the row set.  The row
augmentation $Rf\to\mathbb Q$ is onto.  With one row it is an
isomorphism, so the trivial left module is projective.
With two rows its kernel is a one-dimensional module
$\mathbb Q_\chi$, with $\chi(m)\in\{0,1,-1\}$ and
$\chi(mn)=\chi(m)\chi(n)$.

Put $L_T=\ker(\mathbb QT\to\mathbb Q)$ and
$D L_T=\operatorname{Hom}_{\mathbb Q}(L_T,\mathbb Q)$.
The augmented action bar gives, in degrees $j\geq2$,
\[
 H_j(\operatorname{cone}(T_{hM}\to BM);\mathbb Q)
       =\operatorname{Tor}_{j-1}^{R}(L_T,\mathbb Q).
\]
Indeed, the cone chain complex is the bar of the two-term
complex $[\mathbb QT\to\mathbb Q]$, which is
quasi-isomorphic to $L_T[1]$.  Dualizing a projective resolution
identifies degree-two cohomology with
$\operatorname{Ext}^1_R(\mathbb Q,D L_T)$.  The row augmentation
sequence makes the connecting map
\[
 \operatorname{Hom}_R(\mathbb Q_\chi,D L_T)
       \longrightarrow H^2(\operatorname{cone}(T_{hM}\to BM);\mathbb Q)
\]
surjective.  These constructions are natural in maps of acts.

Represent a class by a linear functional $a:L_T\to\mathbb Q$
satisfying $a(zm)=\chi(m)a(z)$.  Choose $t_0\in T$ and set
$q(t)=a(t-t_0)$.  Then
\[
 q(tm)=\chi(m)q(t)+q(t_0m).
\]
Thus $q$ is an equivariant map to the finite affine act $q(T)$.
On its augmentation kernel, the functional sending a formal
sum $\sum b_xx$ to $\sum b_xx\in\mathbb Q$ pulls back to $a$.
Naturality of the connecting map shows that the original class
pulls back from this affine act cone.  Its square vanishes by
Lemma~\ref{CP2:A-high-affine}.  Every square therefore vanishes,
and rational polarization kills all mixed products.  The one-row
case has zero second cohomology by projectivity.  Opposition
interchanges rows and columns.
\end{proof}

\begin{corollary}[Proved: act cones through order seven]
\label{CP2:A-high-seven-act}
Every act cone over a monoid of order at most seven has zero
rational degree-two products.  There is no bound on the number
of states and no bounded-homology hypothesis.
\end{corollary}
\begin{proof}
For a right act use the preceding lemma if the minimal row
number is at most two.  If the column number is one, the
group-corner comparison of Lemma~\ref{CP2:A-small-minimal-ideal}
gives $BM\simeq BG$ for a finite group, so $H^2(BM;\mathbb Q)=0$.
All second cohomology of the cone then lifts from its suspension
quotient, whose positive products vanish.  The only other
possibility at order at most seven has three rows, two columns,
trivial minimal group, and just the global identity outside the
minimal ideal.  This is the identity-adjoined $3\times2$ band.
The projective resolution in Lemma~\ref{CP2:A-band-spheres}
has length two over $\mathbb Q$ and resolves the trivial module.
Thus its act cone has
$H_4=\operatorname{Tor}_3^R(L_T,\mathbb Q)=0$.
This exhausts the row, column and group-size count.  The left-act
case follows by opposition.
\end{proof}

\begin{lemma}[Proved: exact orbit comparison]
\label{CP2:A-high-orbits}
In a collage with one monoid endpoint $N$, suppose its trivial
rational module is projective on the side used for the
coinvariants.  Replace $N$ by a trivial object and every
incident stalk by its actual $N$-orbit set.  The quotient functor
is a rational homology equivalence and induces an isomorphism
of rational cohomology rings.  The same operation can be made
at an extreme object of an acyclic-support category.
\end{lemma}
\begin{proof}
There is an actual quotient category: endpoint actions identify
only arrows incident with that endpoint, and compatibility of
composition with the actions defines all quotient products.
The action bar of an $N$-set maps to its discrete orbit set.
Its rational homology is the derived coinvariant functor on
the permutation module.  Projectivity of the trivial module
makes that functor exact, so this map is a rational homology
equivalence.  The vertex map $BN\to *$ is one also.
In each degree of every other bar direction the same argument
applies.  Filtering the resulting first-quadrant double bars
proves a rational homology equivalence on the cross-arrow
spaces.  The two-block homotopy pushout of
Lemma~\ref{CP2:A-two-block} completes the comparison.
For an extreme object with several incident stalks, the same
argument is natural in their remaining category variable.
As the comparison is induced by a functor, its cohomology
isomorphism respects products.
\end{proof}

\begin{lemma}[Proved: a small corner with binary stalks]
\label{CP2:A-high-small-binary}
Let a finite category be a collage of a source block $S$ and
one sink corner $N$ of order at most four.  Suppose every
stalk from an object of $S$ has one or two elements.  If the
second rational cohomology of the whole collage has dimension
at most one, every degree-two square is zero.  There is a
dual assertion for a small source corner.
\end{lemma}
\begin{proof}
If the trivial right $\mathbb QN$-module is projective,
Lemma~\ref{CP2:A-high-orbits} reduces the claim to
Lemma~\ref{CP2:A-high-affine}.
Otherwise the projectivity test
Lemma~\ref{CP2:A-projectivity-test} and the minimal-ideal
cardinality show that the ideal is a right-zero band of two
or three elements.  It has one row and trivial group: two
rows or a nontrivial group would require at least five
monoid elements when there are two columns.  Its elements
are right zeros of the whole monoid, so $BN$ is contractible.

On a set of at most two states, the left actions of this
right-zero band coincide.  Indeed they are idempotent maps
with a common image.  If that image has two elements all
are identities; if it has one element all are that same
constant.  Denote this retraction on each stalk by $E$.
The family is natural in $S$.  For $n\in N$, the equations
$nE=E=En$ follow because $nk=k$ and $kn$ remains in the
minimal ideal.  Its image is consequently a biset retract
with trivial $N$ action.  The core collage is equivalent
to an objectwise one- or two-state act cone, since $BN$
is contractible.  Its positive products vanish.

The inclusion of the core has an actual retraction.  Its
quotient in the double-mapping-cylinder model collapses
both vertex blocks and is a suspension of the quotient
correspondence.  If the core has zero second cohomology,
all second classes lift from this suspension.  Otherwise
its injective pullback spans the assumed space of dimension
at most one.  In either case every square vanishes.
\end{proof}

\begin{lemma}[Proved: two cross arrows, without corner restrictions]
\label{CP2:A-high-two-cross}
A two-object category with arrows in only one direction and
exactly two cross arrows has zero rational degree-two products.
\end{lemma}
\begin{proof}
Write its two-element biset as $T$, and consider the image of
the left corner in the four-element transformation monoid of
$T$.  If both constants occur, a commuting right transformation
fixes both points, so the right action is trivial.  If exactly
one constant occurs, no interchange can occur; both actions
fix its value.  That common fixed state supplies a section of
the correspondence projection to the product of the two corner
spaces.  Their second cohomology injects into the correspondence,
so Mayer--Vietoris makes all global second classes relative
to the vertex blocks.  Their products vanish by
Lemma~\ref{CP2:A-two-block}.

If no constant occurs, the left image is trivial or contains
the interchange.  In the latter case every commuting map is
a permutation.  The correspondence is the two-sheeted covering
of the product of the corner spaces determined by these two
permutation actions.  Its transfer makes pullback injective
over $\mathbb Q$, so the same relative argument applies.
The remaining case, and the first case after opposition, has
a trivial action at one corner.  Its correspondence is
$Z\times BN$, where $Z$ is the ordinary action bar of the
other corner.  A finite monoid has zero rational first
homology, hence
\[
 H^2(Z\times BN;\mathbb Q)
  =H^2(Z;\mathbb Q)\oplus
     H^0(Z;\mathbb Q)\otimes H^2(BN;\mathbb Q).
\]
The diagonal map from $H^2(BN)$ to the second summand is
injective.  In degrees one and two the Mayer--Vietoris
sequences therefore identify pullback from the ordinary
two-state act cone with an isomorphism onto the second
cohomology of the collage.  This is an actual functor map
and preserves products.  Lemma~\ref{CP2:A-high-affine}
annihilates all those products.
\end{proof}

\begin{lemma}[Proved: a nonlocal algebra of dimension at most four]
\label{CP2:A-high-small-Tor}
For a nonlocal finite-dimensional $\mathbb Q$-algebra $A$ with
$\dim A\leq4$ and finite-dimensional modules $W_A,{}_A V$,
nonzero $\operatorname{Tor}_2^A(W,V)$ forces nonzero Tor in
arbitrarily high degrees.  If a finite monoid $N$ has more
than one element, its algebra $\mathbb QN$ is nonlocal.
\end{lemma}
\begin{proof}
Put $J=\operatorname{rad}A$.  Semisimple algebras contribute
no positive Tor.  If the semisimple quotient has a nonrational
field factor and $J\ne0$, nonlocality and the dimension bound
leave only a quadratic field, a rational factor, and a
one-dimensional radical.  That radical cannot carry a unital
quadratic-field action and is supported at the rational
factor.  By lifting orthogonal idempotents, $A$ is a product
of a semisimple algebra and the dual numbers.  A matrix
simple factor of degree greater than one already uses all
four dimensions and is semisimple.

The remaining case has $A/J=\mathbb Q^r$, $r\geq2$, and
$\dim J\leq2$.  If $J^2\ne0$, necessarily $r=2$,
$\dim J=2$, $\dim J^2=1$, and $J^3=0$.  The unique
one-dimensional component of $J/J^2$ must be a loop:
an off-diagonal component cannot multiply with itself.
A lift $x$ has $x^2\ne0$, so
$A=\mathbb Q\times\mathbb Q[x]/(x^3)$.
For $\mathbb Q[x]/(x^m)$, every finite module splits into
modules $A/(x^a)$, and a nonprojective such module has the
two-periodic free resolution with alternating maps
$x^a,x^{m-a}$.  Nonzero degree-two Tor consequently recurs
in every subsequent even degree.  This treats these cases
and the dual numbers.

It remains that $J^2=0$.  Lift the $r$ orthogonal primitive
idempotents.  The algebra is the radical-square-zero quiver
algebra with $r$ vertices and $d=\dim J$ arrows, where
$r+d\leq4$.  First syzygies $\Omega W,\Omega V$ of minimal
projective covers are semisimple.  The inclusion
$\Omega W\to P_W$ tensors to zero with $\Omega V$,
because its image lies in $P_WJ$ and $J\Omega V=0$.
The two projective-cover sequences give
\[
 \operatorname{Tor}_n^A(W,V)
 =\operatorname{Tor}_{n-2}^A(\Omega W,\Omega V)
 \qquad(n\geq2),
\]
including ordinary tensor product at $n=2$.  A nonzero
tensor means that a vertex $v$ occurs in both syzygies.
It must have an entering and a leaving arrow.  With at
most two arrows and $r+d\leq4$, these arrows form a loop
or a two-cycle through $v$.  The resolution of a vertex
simple in a radical-square-zero algebra has one generator
for every directed path starting there.  Repeating that
cycle gives nonzero Tor from the right $v$-simple to the
left $v$-simple in arbitrarily high degrees.  These are
direct summands of the displayed groups, proving the claim.

Finally, a nonunit in a finite monoid has an idempotent
power different from the identity.  Its basis element is
a nontrivial algebra idempotent.  If there is no nonunit,
the monoid is a nontrivial finite group, whose averaging
idempotent is nontrivial.  A local algebra has no such
idempotent.  Thus $\mathbb QN$ is nonlocal.
\end{proof}

\begin{lemma}[Proved: augmentation kernels between two trivial extremes]
\label{CP2:A-high-linear-Tor}
Let $S$ be a finite category between a new trivial source $a$
and a new trivial sink $b$.  Suppose every incoming and
outgoing stalk at $S$ is nonempty, and put $R=\mathbb QS$.
Let $V$ be the left permutation module of the incoming
stalks and $W$ the right permutation module of the outgoing
stalks.  Their augmentations are onto the constant left
and right modules; write $\bar V,\bar W$ for the kernels.
Allow any nonempty finite set of direct arrows $a\to b$
and any compatible composition.  Then
\[
 H_j(BC;\mathbb Q)=\operatorname{Tor}_{j-2}^R(\bar W,\bar V)
 \qquad(j\geq3).
\]
For a three-object category whose extreme corner algebras
are semisimple and whose middle corner has order at most
four, bounded rational homology implies $H_4(BC;\mathbb Q)=0$.
\end{lemma}
\begin{proof}
Quotient the normalized category chains by the full subcategory
on $a,b$.  This is a graph, with no chains above degree one.
The quotient is the total two-sided categorical bar of the
two complexes $[W\to\mathbf1_r]$ and
$[V\to\mathbf1_\ell]$, with permutation terms in degree
one.  The four horizontal parts are respectively the bar
of $S$, its incoming and outgoing action bars, and its
two-sided action bar, shifted by zero, one, one, and two.
The face composing the two outer arrows lands in the
quotiented graph.  Every other face is the usual bar
face or an augmentation.  This gives an equality of
chain descriptions, including empty internal strings.

The two complexes are quasi-isomorphic to
$\bar W[1],\bar V[1]$.  This replacement preserves the
bar quasi-isomorphism because all tensor factors over
the ground field are exact; filtering by internal degree
is a first-quadrant comparison.  Its homology is therefore
the displayed Tor shifted by two.  The graph long exact
sequence proves the assertion in degrees at least three.
For semisimple extreme corners use the actual orbit
functors of Lemma~\ref{CP2:A-high-orbits}.  If the middle
monoid is trivial there is no positive Tor.  Otherwise
Lemma~\ref{CP2:A-high-small-Tor} applies to its algebra.
A nonzero fourth homology group would imply unbounded
nonzero homology, proving the last claim.
\end{proof}

\subsection{A complete proof of the acyclic-support excess bound}

\begin{theorem}[Proved: eight excess arrows with acyclic support]
\label{CP2:A-high-acyclic-eight}
Let $C$ be a finite category with $k$ objects, $n$ arrows including
identities, and acyclic support.  If $n\leq2k+7$, then
$BC\not\simeq\mathbb{CP}^2$.
\end{theorem}
\begin{proof}
Suppose otherwise.  The category is connected.  With one object
its monoid has order at most nine, excluded by
Corollary~\ref{A:monoids-nine}; with two objects use
Corollary~\ref{CP2:A-high-two-objects}.  Assume $k\geq3$.
Let $E$ be the number of comparable ordered pairs of distinct
objects in the support poset.  Define
\[
 D=\sum_x(|\operatorname{End}(x)|-1),\qquad
 P=\sum_{x<y}(|C(x,y)|-1),\qquad R=E-k+1.
\]
These are nonnegative integers and
$D+P+R=n-(2k-1)\leq8$.  The support functor is onto on
fundamental groups: represent a support edge loop by
comparable edges and lift each edge to any arrow, reversing
its path when necessary.  Thus the support order complex
has zero first homology over $\mathbb F_2$.

We first record how subpieces will be used.  A clique of the
comparability graph is a totally ordered set, so its full
subcategory contains every simplex with those distinct
vertices.  A graph decomposition into clique blocks therefore
gives an actual union of nerve subcomplexes.  Across an
intersection with $H^3(-;\mathbb Q)=0$, restriction on fourth
cohomology to the two pieces is injective.  Across a rationally
acyclic connected intersection, second cohomology splits and
bounded homology passes to the pieces.  If the intersection
has zero homology in degrees at least three, boundedness
still passes in sufficiently high degrees, which is all the
Tor argument needs.  These statements follow directly from
Mayer--Vietoris, and apply successively to the clique-tree
decompositions of Lemma~\ref{CP2:A-cycle-budget}.

\paragraph{Singleton hom-sets and semisimple corners.}
If $P=0$, Lemma~\ref{CP2:A-singleton-support} pulls every
second class back from the support poset.  Its fourth
cohomology is zero by Lemma~\ref{CP2:A-high-flag}, so every
square is zero.

Suppose all corner algebras are semisimple.  If $R\leq5$,
there is no five-object chain, and the relative-bar height
bound gives $H_4=0$.  If $R\geq6$, then $D+P\leq2$.
Since $P\geq1$, either $D=0$ and $P\leq2$, or $D=1$
and $P=1$.  There is at most one five-object chain and
none longer, by Lemma~\ref{CP2:A-high-flag}.  In that
chain choose a vertex which is neither an endpoint of
a parallel consecutive hom-set nor the one nontrivial
corner, if present.  At most four vertices are forbidden
when $D=0$, and at most three when $D=1$.
The component of the degree-four relative-bar boundary
deleting this vertex is injective.  It retains all varying
arrow labels.  A composed pair of singleton arrows maps
its single basis vector to one basis vector, even if the
target hom-set is larger.  A retained nontrivial corner
has semisimple algebra, so its balanced tensor products
preserve this injection.  Endpoint deletion is the same
argument with a singleton augmentation factor.  No other
five-object chain contributes to that face component.
The fourth boundary is injective, and $H_4=0$.

We may now assume $D\geq2$ and a nonsemisimple corner.
If $P=1$, then $R\leq5$ and there is at most one corner
$M$ with nonzero $H^2$, because each such corner costs
at least four units of $D$.  If there is none, use the
parallel cut.  A singleton hom-set incident with $M$
makes its inclusion null-homotopic inside the corresponding
join, so all global second classes restrict to zero there;
the same cut again applies.  Otherwise $M$ is a leaf on
the sole doubled hom-set.  Its other endpoint has order
at most four, hence is rationally acyclic.  The deletion
of $M$ has singleton cross hom-sets and zero second
cohomology.  The two-arrow edge collage has zero products
by Lemma~\ref{CP2:A-high-two-cross}, and fourth restriction
to it and the deletion is injective.  This closes $P=1$.
Henceforth $D,P\geq2$, and $R\leq4$.

\paragraph{Tree support: $R=0$.}
There is no chain of three distinct objects: its transitive
edge would create a cycle.  The nerve is a tree of
two-object edge collages meeting at corner spaces.
If every corner has order at most five, each intersection
has zero third rational homology by
Corollary~\ref{CP2:A-five-types}.  Every edge collage has
at most
\[
 (D_x+1)+(D_y+1)+(P_{xy}+1)\leq D+P+3\leq11
\]
arrows.  Corollary~\ref{CP2:A-high-two-objects} kills its
degree-two products.  Injectivity of fourth restriction
over the tree kills the global square.

Otherwise there is a unique corner of order six or seven,
and every other corner has order at most two.  This follows
from $D+P\leq8$, $P\geq2$: a corner of order six costs
five, and a further order-three corner would cost two more.
Group all edges incident with the large corner into its
star.  The other pieces have semisimple corners and only
one strict crossing, so have no homology in degree at
least two.  They attach at rationally acyclic corners.
The star has all its arrows oriented the same way.
Collapse its small leaf corners by exact orbit functors.
Identifying the resulting trivial leaf points gives an
ordinary act cone over the large corner.  At the level
of spaces this identifies finitely many points of a
connected space, adding only a wedge of circles; the
quotient map is an isomorphism on cohomology in degrees
at least two and respects products.  The order-seven
act lemma kills its products.  The attachments cannot
restore the fourth square.

\paragraph{One cycle: $R=1$.}
The support consists of one triangle and attached trees.
If all corners have order at most five, fourth restriction
to the triangle and its bridge collages is injective.
Each bridge has at most ten arrows and zero products.
If all corners have order at most four, the intersections
are rationally acyclic, so the triangle inherits bounded
homology and second-cohomology dimension at most one.

For $D\leq3$, at most one corner has order at least three.
When it is an endpoint of the triangle, cut that rationally
acyclic corner from the opposite semisimple two-object
block, whose second cohomology vanishes.  When it is in
the middle, use Lemma~\ref{CP2:A-high-linear-Tor}.
A triangle with only semisimple corners has no fourth
homology by its height.  These arguments also cover a
triangle not containing the larger corner.

For $D=4$ or $5$ with all corners of order at most four,
$P\leq3$ or $2$, respectively.  If an adjacent hom-set
is singleton, the corresponding two-object block is a
join and has zero second cohomology; cut it from the
remaining rationally acyclic corner.  Otherwise both
adjacent hom-sets have at least two elements.  The only
possible hom-size forms, allowing smaller excess, are
$(2,2,1)$, $(2,2,2)$, $(3,2,1)$ and $(2,3,1)$, with
the long hom-set listed last.  Some extreme has both
incident stalks of size at most two.  The rank-one
small-corner lemma applies to that extreme and kills
the triangle square.

If there is a corner of order five, all other corners
have order at most two and $P\leq3$.  With the order-five
corner in the middle, collapse the two semisimple extremes.
A singleton adjacent hom-set deletes the middle object:
the relevant comma category consists of its unique arrow
to, or from, a trivial extreme.  Otherwise the preceding
four hom-size forms have a trivial extreme with all its
stalks at most two.  The objectwise two-state lemma kills
all products.  This uses no triangle-compression theorem.

If the order-five corner is an endpoint and has zero
second cohomology, cut it from the opposite semisimple
block.  If it has nonzero second cohomology, it is $P$
of Lemma~\ref{CP2:A-high-P5-graph}.  A singleton incident
hom-set is a fixed arrow, and the opposite semisimple
block has zero second cohomology; apply
Corollary~\ref{CP2:A-fixed-arrow} or its dual.
Without a singleton, both incident hom-sets have size
at least two.  Unless the long hom-set has size three,
the opposite semisimple extreme has both its stalks at
most two; collapse it and use the categorical two-state
lemma.  The sole exceptional profile is $(2,1,3)$ or
its opposite.  It consumes $P=3$ and forces both other
corners to be trivial.  Delete the middle toward the
other trivial extreme, leaving a three-state act cone
over $P$, covered by the two-row act lemma.

The remaining possibility has one corner of order six,
$D=5$, $P=2$, and all other corners trivial.  Group all
simplices containing that corner into one full piece $X$:
take its incident bridges and the triangle if it lies
there.  All remaining pieces have trivial corners and
dimension at most two, and attach to $X$ and each other
at trivial points.  Thus fourth restriction is injective.
If the corner is off the triangle, $X$ is a star and
reduces, as above, to an act cone of order six.
If it is the middle of the triangle, there is no further
neighbor, since one would create another support cycle.
A singleton adjacent hom-set deletes it; otherwise the
profile is $(2,2,1)$ and a trivial extreme has binary
stalks.  If it is an extreme, its other neighbors are on
the same side.  A singleton opposite adjacent hom-set
allows deletion of the trivial middle, leaving a star.
If that opposite hom-set is parallel, the budget forces
one of the two hom-sets incident with the order-six corner
to be singleton.  Its complement in $X$ has trivial
corners and height at most one.  The fixed-arrow lemma
therefore kills the square.  This completes $R=1$.

\paragraph{At least two cycles and $P=2$.}
Here $D\leq4$.  If all corners have order at most four,
Lemma~\ref{CP2:A-high-two-parallel} applies.  Otherwise
the unique nontrivial corner has order five, $D=4$,
and $R=2$.  If its second cohomology vanishes, the proof
of that lemma still applies: all joins have zero second
cohomology, and in the final $(2,2,1)$ triangle a trivial
extreme replaces the small-corner step.  The same argument
works if global second classes restrict to zero at the
corner, in particular if a singleton hom-set is incident
with it.  A singleton edge intersection in this argument
has a trivial endpoint and is contractible, so the fourth
restriction step remains valid.

Otherwise the corner is $P$ and every incident hom-set
is parallel.  It has one or two neighbors.  Let $X$
consist of this corner and all its neighbors, and let
$D'$ delete only the corner.  These nerves cover $NC$.
Their intersection has at most two trivial objects and
zero third homology.  The category $D'$ has trivial
corners and cycle rank at most two, hence dimension at
most two.  With one neighbor, $X$ is a two- or three-state
act cone over $P$.  With two incomparable neighbors,
it is a star and identifying the trivial leaves reduces
to a four-state act cone.  With comparable neighbors,
it is a triangle with all hom-sets of size at most two
and a trivial extreme.  Every case has zero products by
the act lemmas.  Fourth restriction finishes this case.

\paragraph{Two cycles and at least three parallel excesses.}
Now $R=2$, $D\leq3$, and $P\leq4$.  For $D=2$ there
is one order-three corner and all other corners are
trivial.  Its trivial modules on both sides have projective
dimension at most one.  To see this without classification,
use its minimal ideal: if there is one row or one column,
the corresponding trivial module is projective; the sole
nonprojective case on the other side is the two-element
one-sided band with identity, whose resolution in
Lemma~\ref{CP2:A-band-spheres} has length one on that side.
Thus every edge intersection has zero homology in degrees
at least three.  Fourth restriction to triangle pieces is
injective and they inherit bounded high homology.  At an
endpoint, the order-three corner contributes at most one
internal bar degree, so the triangle has total homological
degree at most three.  In the middle apply the small-Tor
lemma.  Consequently every triangle has zero fourth
homology.  Bridges do also, and the case is closed.

For $D=3$, all corners are rationally acyclic and at most
one has order at least three; call it the larger corner.
Single triangles and bridges are handled by the preceding
one-cycle $D\leq3$ arguments, with inherited boundedness
across corner intersections.  It remains to treat two
triangles sharing an edge, the diamond.  Write its shared
edge as $u<v$ and its incomparable exclusive vertices as
$a,b$.  Transitivity allows, up to opposition, only
$u<a,b<v$ or $a,b<u<v$.

In the first orientation, cut either endpoint from the
opposite two-edge star.  If one star has zero second
cohomology, the cut proves product vanishing.  Otherwise
each star has an edge with potentially nonzero second
cohomology.  An edge can have such cohomology only if
it is parallel and touches the larger corner: singleton
edges are joins, and semisimple edges have height one.
Thus the larger corner must be one of $a,b$, say $a$.
The shared-edge endpoints are semisimple, so that edge
has zero homology in degrees at least two.  Fourth
restriction to the two triangles is injective and both
inherit bounded high homology.  The triangle through $a$
has middle order at most four and semisimple extremes,
so has zero fourth homology by the Tor lemma.  The other
triangle is semisimple and has height two.

In the second orientation, cut $\{a,b\}$ from the shared
edge.  If the shared edge has zero second cohomology,
the cut applies.  If one exclusive vertex has singleton
hom-sets to both $u,v$, its full triangle is a join and
has zero second cohomology.  Every global second class
then restricts to zero on the shared edge, and the same
cut applies.  Avoiding these alternatives requires the
shared edge and an incident hom-set at each exclusive
vertex to be parallel.  Since $P\leq3$, these three sets
are doubled and all others are singleton.  The shared
edge touches the larger corner.  If it is $u$, collapse
the semisimple final corner $v$; all its stalks are binary,
and the objectwise two-state lemma applies.  If it is
$v$, the remaining height-one source block is semisimple
and has zero second cohomology.  The corner at $v$ has
zero second cohomology also, so the two-block lemma
finishes.  This proves every two-cycle case.

\paragraph{Three or four cycles.}
The only remaining three-cycle budget is $D=2$, $P=3$.
There is one order-three corner and all others are trivial.
The small-cycle graph decomposition consists of triangle
pieces, or one four-clique, with bridge attachments.
All edge intersections have zero homology in degree at
least three, as just proved.  Triangle pieces inherit
boundedness and have zero fourth homology by the same
endpoint-degree or middle-Tor arguments.

Consider a four-clique.  If all its hom-sets have size at
most two, at least one extreme is trivial, so the whole
piece is an objectwise binary act cone.  One hom-set of
size four is covered by the parallel cut.  The only
other possibility is one hom-set of size three and one
of size two.  Nonserial sets use the same cut.  For
separated serial sets use the downset of the second
source, exactly as in the proof of
Lemma~\ref{CP2:A-high-two-parallel}.  For adjacent serial
sets their common vertex has an additional neighbor in
the four-chain; the join-downset argument in that proof
applies.  These arguments do not depend on the sizes
two versus three, only on which hom-sets are singleton.
Thus the four-clique has zero degree-two products, and
fourth restriction kills the global square.
Finally $R=4$ forces $D=P=2$, already covered by the
two-parallel lemma.  All nonnegative budgets of sum at
most eight have now been treated.  Each contradicts the
nonzero rational square of $\mathbb{CP}^2$.
\end{proof}

\begin{corollary}[Proved: the excess-six and excess-seven assertions]
\label{CP2:A-high-six-seven}
The preceding exclusion also holds under $n\leq2k+5$ or
$n\leq2k+6$.  In particular, all acyclic-support categories
with at least seven objects and at most twenty-one arrows
are excluded.
\end{corollary}
\begin{proof}
Each displayed bound implies $n\leq2k+7$; for the last
assertion, $k\geq7$ gives $21\leq2k+7$.
\end{proof}

\subsection{Strong components without a small-category census}

\begin{lemma}[Proved: the categorical minimal ideal]
\label{CP2:A-high-Rees}
Let $S$ be a finite strongly connected category.  There are a
finite group $G$, positive row and column numbers $r_i,c_i$
at each object, and a least nonempty ideal $K$, with
\[
 |K(i,j)|=r_j|G|c_i,\qquad
 |K|=|G|\Bigl(\sum_i r_i\Bigr)\Bigl(\sum_i c_i\Bigr).
\]
If every $r_i=1$, or every $c_i=1$, then $BS\simeq BG$.
If $1_i\in K$, object $i$ is a retract of every object.
When $G=1$, put $A=\mathbb QS/\mathbb QK$.  There are
left and right $A$-modules $V,W$, of dimensions
$\sum(r_i-1)$ and $\sum(c_i-1)$, such that
\[
 H_j(BS;\mathbb Q)=\operatorname{Tor}_{j-2}^{A}(W,V)
 \qquad(j\geq2).
\]
\end{lemma}
\begin{proof}
Fix an object $a$ and an idempotent $e$ in the least ideal
of its endomorphism monoid.  Its corner $G=eS(a,a)e$ is a
group by the finite-monoid argument in
Lemma~\ref{CP2:A-small-minimal-ideal}.  Set
$U_j=S(a,j)e$, $V_i=eS(i,a)$, and $K(i,j)=U_jV_i$.
Every nonempty ideal of $S$ meets $S(a,a)$ by composition
with arrows to and from $a$, and hence contains $e$.
It therefore contains all these $K(i,j)$.  They themselves
form an ideal, proving minimality.

The right $G$ action on $U_j$ and the left action on $V_i$
are free: composing an equation $ug=u$ with any
$v\in V_j$ gives $(vu)g=vu$ inside $G$.  Multiplication
induces a bijection $U_j\times_G V_i\to K(i,j)$.
For injectivity, if $uv=u'v'$, choose $w\in U_i$ and
adjust it by $G$ so that $vw=e$.  Then $u=u'g$ with
$g=v'w\in G$; composing the original equation with a
left inverse of $u'$ gives $gv=v'$.  This is exactly
the balanced equivalence relation.  Choose orbit
representatives.  The resulting coordinates are a row,
a group element, and a column, and composition multiplies
the group coordinates with the intervening fixed element
of $G$.  This gives both cardinality formulas.

Adjoin the retract object represented by $e$ in the
idempotent completion.  It is a retract of $a$, so its
addition preserves the classifying space by
Lemma~\ref{CP2:A-retract-deletion}.  If every row number
is one, the comma category from its group corner to each
object is the free transitive $G$ action on $U_i$.
It has exactly one arrow between any ordered pair of
objects and is contractible.  Quillen's Theorem~A gives
$BG\simeq BS$; the column case is dual.
If $1_i=uv\in K(i,i)$, then $vu$ is an idempotent in
$G$, hence $e$.  For any $j$ choose $u_j,v_j$ with
$v_ju_j=e$.  The maps $u_jv:i\to j$ and $uv_j:j\to i$
compose to $1_i$, proving the retract assertion.

Now suppose $G=1$, and write $R=\mathbb QS$.
The projectives $Re,eR$ augment onto the constant left
and right modules.  Their kernels, now denoted $V,W$,
have the stated dimensions.  The ideal acts constantly
on each row or column set and thus annihilates these
augmentation kernels.  Multiplication is a bimodule
isomorphism $Re\otimes_{\mathbb Q}eR\to\mathbb QK$,
as can be checked on the row-column basis.  In particular
$\mathbb QK$ is projective on both sides.
Dimension-shifting through the two augmentation sequences
gives $H_j(BS;\mathbb Q)=\operatorname{Tor}_{j-2}^R(W,V)$.
At $j=2$ the last map is into $eV=0$, so this includes
ordinary tensor product.  Finally the projective left
resolution $0\to\mathbb QK\to R\to A\to0$ gives
$W\otimes_R\mathbb QK=We\otimes eR=0$ and
$\operatorname{Tor}_{q}^R(W,A)=0$ for $q>0$.
The change-of-rings spectral sequence has only its
zero row and identifies Tor over $R$ with Tor over
$A$.  This proves the formula.
\end{proof}

\begin{lemma}[Proved: two explicit one-sided collapses]
\label{CP2:A-high-one-sided}
The following operations preserve the actual homotopy type.

\begin{enumerate}
\item Suppose $u:b\to a$, $v:a\to b$, and $vu=e\ne1_b$
is idempotent.  If $e$ and every unit of $\operatorname{End}(b)$
fix every incoming arrow from a different object, one may
delete $b$.  There is an outgoing dual.  The unit condition
holds automatically if $e$ is a zero of its corner.
\item Let a strong component $S$ have trivial minimal group
and one minimal row per object.  It may be replaced by one
trivial object if it is a source component.  The new outgoing
stalks are the actual right-act orbit sets.  More generally,
the replacement is allowed if all external incoming arrows
are fixed by local minimal idempotents.  In that case no
hom-set increases.  There are column/sink/exit duals.
\end{enumerate}
\end{lemma}
\begin{proof}
For the first assertion, let $D=C\setminus\{b\}$.
The nonunits of the finite corner at $b$ form an ideal:
if a product is a unit, each factor has a one-sided
inverse and finiteness makes it a unit.  Define
$F:C\to D$ by sending $b$ to $a$, fixing $D$, and
\[
 \begin{aligned}
 F(h)&=1_a\ (h\text{ a unit}),& F(h)&=uhv\ (h\text{ a nonunit}),\\
 F(g)&=ug\ (g:d\to b),& F(f)&=fv\ (f:b\to d).
 \end{aligned}
\]
Identities are preserved.  For nonunits $h,k$, the
equality $ekv=kv$ gives $uhvukv=uhkv$.  For incoming
$g$ it gives $uhvug=uhg$; for outgoing $f$ the equality
$ehv=hv$ gives $fvuhv=fhv$.  Unit factors use $hg=g$
and $hv=v$.  A passage through $b$ between retained
objects gives $fvug=feg=fg$.  In the other order,
$gf:b\to b$ is a nonunit, since $egf=gf$ and a unit
would force $e=1_b$; hence $F(g)F(f)=ugfv=F(gf)$.
These check every pair involving $b$; retained pairs
are unchanged.  Thus $F$ is a functor.  Its composite
with the inclusion is the identity on $D$.  The
components $v$ at $b$ and identities elsewhere give
a natural transformation from the other composite to
$\operatorname{id}_C$, by the same equations.  If $e$
is a zero, $hg=heg=eg=g$ for every unit $h$.

For the second assertion choose $a\in S$ and a minimal
idempotent $e$ there.  There is a unique arrow
$u_i\in S(a,i)e$ at each object, and for every
$f:i\to j$ one has $fu_i=u_j$.  Choose
$v_i\in eS(i,a)$; then $v_iu_i=e$, and put
$e_i=u_iv_i$.  The $u_i$ contract $BS$ by a natural
transformation from the constant functor at $a$.
For a right act $T$, the category of elements has
objects $(i,t)$ and arrows $f:(i,t)\to(j,t')$ with
$t=t'f$.  The map $(i,t)\mapsto tu_i\in T(a)e$
is constant on arrows, and its section sends $z$ to
$(a,z)$.  The $u_i$ give a natural transformation
from section followed by projection to the identity.
Thus the element category is equivalent on classifying
spaces to this discrete orbit set, naturally in $T$.
The two-block pushout proves the source-component
orbit-collapse equivalence.

For the general fixed-entrance assertion, keep the
complement $D$ and introduce one object $s$, with
\[
 \bar C(d,s)=eC(d,a),\qquad \bar C(s,d)=C(a,d)e,
 \qquad \bar C(s,s)=\{1_s\}.
\]
Use the old compositions between retained objects and
through $s$.  No path leaves $S$ and reenters it, since
it is a strong component.  Define $Q$ to send $S$ to
$s$, its arrows to $1_s$, entrances $p:d\to i$ to
$v_ip$, and exits $q:i\to d$ to $qu_i$.
The equations $p=e_ip$ and $fu_i=u_j$ give
$v_jfp=v_ip$, and passage through $S$ gives
$(qu_i)(v_ip)=qp$.  These and associativity verify
every composition.  A section chooses $a$ and the
literal displayed cross arrows.  Its identity at
$s$ is $1_a$.  The identities $v_ae=e$ and $u_a=e$
show that it is indeed a section.  The components
$u_i$ and identities on $D$ give the required natural
transformation to $\operatorname{id}_C$.
Every displayed cross-hom is an image of an old
cross-hom and the internal block becomes one arrow.
This proves the counting assertion as well.
\end{proof}

\begin{lemma}[Proved: nonzero second homology costs three excesses]
\label{CP2:A-high-H2-cost}
If a connected finite category has $k$ objects and at most
$2k+1$ arrows, then $H_2(BC;\mathbb Q)=0$.
\end{lemma}
\begin{proof}
Delete object retracts.  Each deletion removes at least
three arrows and one object and preserves the bound.
Assume none remains.  A proper strong component on
$t\geq2$ objects has at least $t^2+t$ internal arrows:
its minimal ideal has at least $t^2$ and all identities
lie outside it by Lemma~\ref{CP2:A-high-Rees}.
If there are $\ell$ outside objects, contracting the
component in the connected underlying graph leaves at
least $\ell$ edges, and one interface supplies at least
$t$ arrows.  With the outside identities,
\[
 n\geq t^2+t+2\ell+t-1=2k+t^2-1>2k+1.
\]
Thus no such proper component exists.  If the category
is wholly strongly connected, one row or one column
per object gives finite-group homotopy type.  Otherwise
its ideal has at least $(k+1)^2$ arrows, and for
$k\geq2$ its identities give
$n\geq(k+1)^2+k>2k+1$.  The one-object case has order
at most three and rationally acyclic classifying space.

The remaining support is acyclic, with $D+P+R\leq2$
in the notation of the preceding theorem.  If $P=0$,
singleton-support comparison and the flag budget give
$H_2=0$.  If $P>0$, then $D\leq1$ and every corner
algebra is semisimple.  With $R=0$, height one gives
$H_2=0$.  The only other budget is $R=P=1$, $D=0$.
A graph of cycle rank one has either no triangle or
one triangle with bridge attachments.  In the first
case there are no two-simplices.  In the second case
there is only one triangle, and at most one of its
hom-sets is doubled.  Its second boundary is injective:
if an adjacent set is doubled, the boundary component
retaining that set distinguishes the two basis
simplices; otherwise there is just one two-simplex
with nonzero boundary.  Bridges add no two-chains.
This proves the assertion without assuming the
support itself has zero first homology.
\end{proof}

\begin{theorem}[Proved: eight excesses for arbitrary support]
\label{CP2:A-high-all-eight}
Every finite category with $k$ objects and $n\leq2k+7$
arrows has classifying space different from
$\mathbb{CP}^2$.  There is no upper cap on $n$ in this statement.
\end{theorem}
\begin{proof}
Induct on $k$.  For $k=1$ use Corollary~\ref{A:monoids-nine}.
Suppose a connected counterexample has $k>1$ objects.
An object-retract deletion removes at least three
arrows and one object and preserves $n'\leq2k'+7$;
the inductive hypothesis excludes it.  Assume no
object is a retract of another.

A proper strong component on $t\geq3$ objects gives
$n\geq2k+t^2-1\geq2k+8$ by the count just proved.
If the whole category is strongly connected, the
one-row or one-column alternatives have finite-group
type.  Otherwise $n\geq(k+1)^2+k$, forcing $k=2$,
$n=11$, trivial minimal group, nine ideal arrows,
and exactly the two identities outside the ideal.
The quotient algebra in Lemma~\ref{CP2:A-high-Rees}
is $\mathbb Q^2$, so $H_4=0$ by its Tor formula.

It remains to exclude a proper strong pair $S$.
Write $s=|S|$ and $\ell=k-2\geq1$.
It has at least six arrows, and connectedness gives
\[
 s+2\ell+1\leq n\leq2\ell+11,
 \qquad 6\leq s\leq10.
\]
Here a support interface with $S$ has an arrow at each
of its two objects.  The row-column count shows that
$BS$ has finite-group homotopy type: otherwise both
total row and column numbers would be at least three,
requiring nine ideal arrows and two identities.
In particular its second rational cohomology vanishes.

First suppose $S$ is an extreme component; by opposition
make it a source.  If its complement $D$ has $q$ connected
components and nonzero second cohomology, the H2-cost
lemma gives
$|D|\geq2\ell-q+3$.  Every component meets $S$ in
at least two arrows.  Thus
\[
 n\geq s+2\ell+q+3\geq s+2\ell+4.
\]
For $s\geq8$ this contradicts the upper bound.  Both
blocks then have zero second cohomology, and the
two-block lemma excludes the square.  For $s=6,7$,
the minimal group is trivial and each object has one
minimal row and column.  Source orbit collapse in
Lemma~\ref{CP2:A-high-one-sided} saves $s-1$ internal
arrows and at least one interface arrow.  In fact its
orbit set is a surjective image of either nonempty
stalk: if $z\in T(a)e$, then
$(zv_i)u_i=z$.  Its size is at most the smaller of the
two stalk sizes.  The collapse saves at least $s$
arrows and one object, so induction applies.

Now suppose $S$ is nonextreme.  There are outside
objects $a,b$ with $a\to S\to b$, and composition
gives $a\to b$.  The contracted support has at least
$\ell+1$ edges, including two interfaces.  Let
$P_{\rm ext}$ be the sum of excesses $|\operatorname{Hom}|-1$
over all individual nonempty hom-sets between $S$ and
outside objects.  Counting gives
\[
 n\geq s+2\ell+3+P_{\rm ext},\qquad
 P_{\rm ext}\leq8-s,
\]
so $s\leq8$.  For $s=6,7$, the unique minimal
endomorphism $e_i$ at either object is a zero of its
corner.  Consider the four conditions that all incoming
cross arrows, or all outgoing cross arrows, are fixed
by $e_i$.  Each of the at most $s-6$ extra internal
arrows can violate at most two conditions.  An external
violation requires a nonsingleton hom-set and can
violate only the condition at its one incident object.
The number of failed conditions is therefore at most
\[
 2(s-6)+P_{\rm ext}\leq s-4\leq3.
\]
A condition holds.  The first one-sided collapse
deletes that object: its zero ensures the unit
hypothesis, and the minimal arrows to the other
object compose to $e_i$.  At least six arrows
disappear, namely two corner arrows, two internal
cross arrows, and an external entrance and exit.
Induction applies.  For $s=8$ one has
$P_{\rm ext}=0$.  Its minimal group is trivial and
the possible total row/column counts are
$(2,2),(2,3),(3,2)$.  Thus one row or one column occurs
at each object.  The external singleton hom-sets are
fixed, so the second one-sided collapse or its dual
replaces $S$ by a point, saving at least seven arrows
and one object.  Induction again applies.

Every strong-component alternative has been excluded.
The remaining support is acyclic and is excluded by
Theorem~\ref{CP2:A-high-acyclic-eight}.  This completes
the induction.
\end{proof}

\begin{proposition}[Proved: a proper strong pair at the ninth excess]
\label{CP2:A-high-pair-nine}
Suppose $C$ is connected, no object is a retract of
another, and $C$ has a proper strong component on two
objects.  If $n\leq2k+8$, then
$BC\not\simeq\mathbb{CP}^2$.
\end{proposition}
\begin{proof}
Use $S,s,\ell$ as in the preceding proof.  The count
is now $6\leq s\leq11$, with
$s+2\ell+1\leq n\leq2\ell+12$.  For $s\leq10$
the pair has finite-group homotopy type and zero
second rational cohomology.

For an extreme source pair, $s=6,7$ collapses with
at least $s$ arrows saved, reducing to the excess-eight
theorem.  For $8\leq s\leq10$, zero second cohomology
of the complement gives the two-block obstruction.
Otherwise the H2-cost count is
$n\geq s+2\ell+q+3$, forcing
$s=8$, $q=1$, $|D|=2\ell+2$, and exactly two
interface arrows.  All external hom-sets at $S$ are
singleton or empty.  Its one-row or one-column
fixed-interface collapse saves at least seven arrows,
again reducing to the preceding theorem.
For $s=11$, the elementary component count gives
\[
 n\geq11+(2\ell-q)+2q=2\ell+11+q.
\]
Equality forces $q=1$, $|D|=2\ell-1$ and two
interface arrows.  The complement is a trivial-corner
singleton-hom tree.  Both interface arrows end at
one object $w$; it has no outgoing cross arrow in
$D$, since that would create more interface arrows.
The full category on $S\cup\{w\}$ has terminal
object $w$.  Its nerve and $ND$ cover $NC$, meet
at $w$, and are both contractible.  The whole space
is contractible.  Opposites cover extreme sinks.

For a nonextreme pair, the same count becomes
$P_{\rm ext}\leq9-s$ and $s\leq9$.
For $s=6$, and for $s=7$ with $P_{\rm ext}\leq1$,
the four-condition bound is less than four, so
one-sided object deletion saves at least six arrows.
For $s=8,9$ with $P_{\rm ext}=0$, fixed-interface
block collapse saves at least seven.  Each reduces
to Theorem~\ref{CP2:A-high-all-eight}.

The remaining numerical pairs are
$(s,P_{\rm ext})=(7,2),(8,1)$.  Equality holds
in the support count.  All outside corners are
trivial, all outside hom-sets are singleton, the
contracted graph has one triangle $a,S,b$ with
tree attachments, and its only two interfaces with
$S$ are $a\to S$ and $S\to b$.  Delete each outside
tree leaf: it has trivial corner and one singleton
edge to a trivial neighboring corner, so its nerve
is a single interval attached at one endpoint.
This is a homotopy equivalence and removes two
arrows.  We reach a four-object core with sixteen
arrows and a singleton direct arrow $a\to b$.

For $(8,1)$ all parallel excess lies on one interface
side.  The other side has one arrow at each object
of $S$, making $a$ initial or $b$ terminal.  The
core is contractible.  The same holds for $(7,2)$
if both excesses lie on one side.  Otherwise there
is one doubled entrance and one doubled exit.
If any of the four fixed-arrow conditions holds,
one-sided deletion applies as before.  Assume all
four fail.  The unique internal arrow beyond the
four-arrow minimal band and two identities must
then be a cross arrow $t:i\to j$.  It accounts
for failure of the outgoing condition at $i$ and
the incoming condition at $j$.  Therefore the
doubled exterior stalks are $a\to i$ and $j\to b$;
the other two exterior stalks are singletons.

Write $R=\mathbb QS$ and $K$ for its four-arrow
minimal band.  The quotient $A=R/\mathbb QK$ is
the path algebra with two vertices and the one
arrow $t:i\to j$.  Let $V,W$ be the incoming and
outgoing augmentation kernels.  They are killed
by $K$.  At the doubled incoming stalk, its
minimal endomorphism factors through the other
singleton stalk and hence acts constantly; the
other augmentation space is zero.  This proves
$KV=0$, and the dual argument gives $WK=0$.
The projective ideal $\mathbb QK=Re\otimes eR$
and the change-of-rings argument in
Lemma~\ref{CP2:A-high-Rees} identify
$\operatorname{Tor}_2^R(W,V)$ with
$\operatorname{Tor}_2^A(W,V)$.
This last group is zero: a representation
$f:X_i\to X_j$ of the one-arrow quiver has the
explicit projective resolution
\[
 0\to P_j\otimes X_i\to
 (P_i\otimes X_i)\oplus(P_j\otimes X_j)
 \to X\to0,
\]
whose maps at $j$ are $x\mapsto(x,-f(x))$ and
$(x,y)\mapsto f(x)+y$.  The generalized middle-bar
formula in Lemma~\ref{CP2:A-high-linear-Tor}
now gives $H_4$ of the core equal to zero.
Thus even the exceptional cores are excluded
without any sixteen-arrow composition census.
\end{proof}

\begin{corollary}[Proved: every ninth-excess candidate has acyclic support]
\label{CP2:A-high-nine-support}
If $BC\simeq\mathbb{CP}^2$ and $n\leq2k+8$, then
$C$ has acyclic support, no object retract, and
$n=2k+8$.  Consequently a seventeen-arrow candidate
has at most four objects; eighteen- or nineteen-arrow
candidates have at most five; twenty- or twenty-one-arrow
candidates have at most six.  An eighteen-arrow
five-object candidate has acyclic support.
\end{corollary}
\begin{proof}
The equality follows from the excess-eight theorem.
A retract deletion would save at least three arrows
and one object, reducing to that theorem as well.
A proper component on $t\geq3$ objects has
$n\geq2k+t^2-1$.  The only possible case is $t=3$
with equality throughout.  Its ideal has one row
and one column per object and trivial group; it has
twelve internal arrows.  In the contracted support
it is a leaf, with just one interface containing
three singleton arrows.  It is an extreme component
and source-orbit collapse or its dual replaces it
by one point, saving at least thirteen arrows and
two objects.  This reduces to the excess-eight
theorem.

A wholly strong category with $k\geq2$ either has
finite-group type or has $n\geq(k+1)^2+k$.
The latter bound forces $k=2$, trivial minimal
group, nine ideal arrows, and at most three arrows
outside the ideal.  The quotient algebra contains
the two distinct nonzero object idempotents and
is nonlocal of dimension at most three.
Its degree-two Tor cannot be nonzero with bounded
higher Tor, by Lemma~\ref{CP2:A-high-small-Tor}.
The categorical minimal-ideal formula therefore
rules out the projective-plane homology.
A proper strong pair is ruled out by the preceding
proposition.  All nontrivial strong components are
thus excluded, proving acyclicity of the support.
The numerical object bounds are immediate from
$n>2k+7$, and the final assertion uses $18=2(5)+8$.
\end{proof}


\subsection{New reductions for the unrestricted sixteen-arrow baseline}

\begin{lemma}[Proved: bounded act cones through monoid order eight]
\label{CP2:A-baseline-eight-act}
Let $M$ be a finite monoid of order at most eight and $S$ a
nonempty finite right $M$-set.  If the act-cone category has
bounded rational homology, all its rational degree-two products
are zero.  The left-act assertion is dual.  The same holds
after adjoining an opposite endpoint with projective trivial
rational module on the appropriate side.
\end{lemma}
\begin{proof}
Put $R=\mathbb QM$ and
$L=\ker(\mathbb QS\xrightarrow{\epsilon}\mathbb Q)$,
with its actual right $R$-module structure.  To specify the
boundedness implication precisely, take a projective left
resolution $P_*\to\mathbb Q$.  The action-space map to
$BM$ is computed by
\[
 \mathbb QS\otimes_R P_*\longrightarrow
 \mathbb Q\otimes_R P_*.
\]
Projective left modules are flat, so
\[
 0\to L\otimes_RP_*\to\mathbb QS\otimes_RP_*
   \to\mathbb Q\otimes_RP_*\to0
\]
is exact in every degree.  The mapping cone of its last
map is quasi-isomorphic to $(L\otimes_RP_*)[1]$.
The two-block description of the act cone therefore gives
\[
 \widetilde H_j(B\mathcal C(M,S);\mathbb Q)
   =\operatorname{Tor}_{j-1}^R(L,\mathbb Q)
   \qquad(j\geq1).
\]
Thus boundedness means eventual vanishing of these particular
Tor groups.  It is not inferred separately for $BM$ or its
action space, whose homology can indeed cancel in a cofiber.

If $H^2(BM;\mathbb Q)=0$, second classes of the cone lift
from its suspension quotient and all their products vanish.
If there are at most two minimal rows, use
Lemma~\ref{CP2:A-high-two-row-act}.  Otherwise nonzero
vertex second cohomology requires at least two columns,
while there are at least three rows.  The minimal-ideal
size and the identity outside it force a trivial maximal
group, exactly a $3\times2$ band $K$, and at most two
outside elements.  Put $A=R/\mathbb QK$.

The ideal $\mathbb QK=Re\otimes_{\mathbb Q}eR$ is
projective on the left, so $\operatorname{pd}_R A\leq1$.
The row augmentation is an exact sequence
\[
 0\to V\to Re\to\mathbb Q\to0,
\]
where $KV=0$ and $\dim V=2$.  Hence $V$ is a left
$A$-module and, for $j\geq2$,
$\operatorname{Tor}_{j}^R(L,\mathbb Q)
=\operatorname{Tor}_{j-1}^R(L,V)$.

A unital rational algebra of dimension at most two is
semisimple or the dual-number algebra
$\mathbb Q[\varepsilon]/(\varepsilon^2)$.
In the semisimple case $V$ is a summand of a finite
sum of copies of $A$, so it has $R$-projective dimension
at most one and $\operatorname{Tor}_3^R(L,\mathbb Q)=0$.
In the dual-number case Jordan decomposition gives
$V=A^a\oplus k^b$, where $k=A/(\varepsilon)$.
Restrict the exact sequence
$0\to k\xrightarrow{1\mapsto\varepsilon}A\to k\to0$
to $R$.  Since $\operatorname{Tor}_{j}^R(L,A)=0$ for
$j\geq2$, its long exact sequence gives
\[
 \operatorname{Tor}_{j}^R(L,k)
   \cong\operatorname{Tor}_{j-1}^R(L,k)
   \qquad(j\geq3).
\]
If $b=0$, the projective-dimension argument already
applies.  If $b>0$, eventual vanishing of the cone
homology implies eventual vanishing for $V$, hence
for $k$.  The displayed recurrence then forces
$\operatorname{Tor}_2^R(L,k)=0$.  In either case
\[
 H_4(B\mathcal C(M,S);\mathbb Q)
   =\operatorname{Tor}_3^R(L,\mathbb Q)
   =\operatorname{Tor}_2^R(L,V)=0.
\]
This proves the obstruction in the remaining case.
Opposition gives left acts, and the actual exact-orbit
functor gives the projective opposite-endpoint version.
\end{proof}

\begin{proposition}[Proved: large corners in a linear sixteen-arrow profile]
\label{CP2:A-baseline-linear-large}
A three-object category with full linear acyclic support and
at most sixteen arrows cannot realize $\mathbb{CP}^2$ if any
corner has order at least eight.
\end{proposition}
\begin{proof}
Write the corner orders as $a,b,c$, and the hom-set sizes as
$x=|C(0,1)|$, $y=|C(1,2)|$, $z=|C(0,2)|$; all six
numbers are positive.  First suppose $c\geq8$.  Then
$a+b+x+y+z\leq8$ and $a\leq4$.  If $x,z\leq2$,
the small-source binary-stalk lemma excludes the square.
Otherwise $x\geq3$ or $z\geq3$.

If $x\geq3$, at least one of $y,z$ equals one: having
both at least two would already give
$a+b+x+y+z\geq9$.  If both are one, the sink cut is
a full singleton join.  If exactly one is one, the lower
two-object block has at most five arrows, so its rational
second cohomology is zero by the H2-cost lemma.  The
singleton incoming arrow to the sink then excludes all
second products by the fixed-arrow corollary.

If $z\geq3$ and $y=1$, the lower block has at most four
arrows and the same fixed-arrow argument applies.  If
$y\geq2$, the inequality forces
$a=b=x=1$, $y=2$, $z=3$, and $c=8$.  The comma category
for deleting the middle object has its one object given
by the unique arrow $0\to1$ and has trivial endomorphism
monoid.  Quillen's Theorem~A identifies the whole category
with the ordinary three-state act cone over the order-eight
sink.  Its boundedness follows from this equivalence, so
Lemma~\ref{CP2:A-baseline-eight-act} excludes its square.
Opposition handles a large source corner.

Finally suppose $b\geq8$.  Both endpoint corners have
order at most four.  Either $x,z\leq2$ or $y,z\leq2$
gives the binary-stalk obstruction.  Avoiding both would
require $z\geq3$ or both $x,y\geq3$.  The latter costs
at least seventeen arrows, so $z\geq3$ and
$a+c+x+y\leq5$.  In particular $a+c\leq3$, and at
least one of $x,y$ is one.  If $x=1$ but $a>1$,
then the same inequality forces $c=y=1$; the dual
statement holds if $y=1$ but $c>1$.  Thus there is
always a trivial endpoint with its adjacent hom-set
singleton.  Its singleton comma category permits
deleting the middle object.  The remaining two corners
have orders at most two, hence semisimple rational
algebras; the acyclic edge has rational homology only
through degree one.  Its fourth homology is zero.
All large-corner positions are excluded.
\end{proof}

\begin{proposition}[Proved: four-object large corners through sixteen arrows]
\label{CP2:A-baseline-four-large}
A four-object category with at most sixteen arrows and a
corner of order at least eight cannot realize $\mathbb{CP}^2$.
Consequently, to establish the four-object portion of the
unrestricted sixteen-arrow baseline it suffices to consider
acyclic support, exactly sixteen arrows, and corners of
order at most seven.
\end{proposition}
\begin{proof}
If the classifying space were $\mathbb{CP}^2$, the
excess-eight theorem would force $n=16$, and
Corollary~\ref{CP2:A-high-nine-support} would give
acyclic support and no object retract.  Let $v$ have
corner $M$ of order $q\geq8$.  With $d$ the corner
excess away from $v$, the budget becomes
\[
 d+P+R=10-q\leq2.
\]
The case $P=0$ is excluded by singleton-support
comparison, because a four-object poset has no
four-dimensional homology.

If $v$ is nonextreme, then $R\geq1$.  The remaining
budget forces $q=8$, $d=0$, $P=R=1$.  There is one
triangle through $v$ and one leaf edge; $v$ is its
middle object and has no other neighbor, since any
such comparison would create another cycle.
One adjacent hom-set is singleton.  Its other
endpoint is trivial and gives a singleton comma
category for deleting $v$, on the appropriate
side of Quillen's Theorem~A.  The result has three
trivial corners and acyclic support, hence nerve
dimension at most two.  This contradicts the
nonzero fourth homology.

Now take opposites so that $v$ is a sink, and let
$D=C\setminus\{v\}$.  If it has a singleton incoming
hom-set, then $|D|\leq7$.  At equality $q=8$ and
there is exactly one incoming arrow, so $D$ is
connected.  The H2-cost lemma gives $H^2(BD;\mathbb Q)=0$
in that case.  When $|D|\leq6$, each connected
component of size two has at most five arrows and
each isolated corner has order at most four;
the same lemma and the small-monoid result give
$H^2(BD;\mathbb Q)=0$, even if $D$ is disconnected.
The singleton is a fixed incoming arrow, and
Corollary~\ref{CP2:A-fixed-arrow} applies.

Suppose every incoming hom-set has at least two
elements.  There cannot be three neighbors, which
would require at least $8+3+6=17$ arrows.
With two neighbors the budget forces $q=8$, all
other corners trivial, two doubled incoming
hom-sets, and one additional singleton edge
connecting the fourth object $w$ to a neighbor
$x$.  This edge is $x\to w$: the reverse direction
would create a third hom-set into $v$.  Delete
the trivial leaf $w$.  The remaining star has
two trivial leaves and the order-eight center.
Identifying its two leaf points gives an ordinary
left-act cone with four states.  On spaces this
adds only a circle, so preserves bounded rational
homology and cohomology products in degrees at
least two.  Lemma~\ref{CP2:A-baseline-eight-act}
therefore kills its square.

With one neighbor $w$, write $s=|C(w,v)|\geq2$.
There is no incoming cross arrow to $w$ in $D$,
since its composite would create another neighbor
of $v$.  The nerves of $D$ and the edge collage
$Y$ on $w,v$ cover $NC$, meeting in the corner
at $w$.  The category $D$ is connected and needs
at least five arrows, so $s\leq3$.  If $s=2$,
$|D|\leq6$ and its second cohomology vanishes;
the shared corner has order at most two and is
rationally acyclic.  Fourth restriction is
injective, and the universal two-cross theorem
kills the square on $Y$.  If $s=3$, equality
forces $q=8$ and $|D|=5$: $D$ is a trivial-corner
singleton-hom tree and is contractible.  The
pushout identifies $BC$ with the three-state
ordinary act cone $BY$.  The order-eight act
lemma again excludes its square.  These cases
exhaust the budget.
\end{proof}

\begin{lemma}[Proved: a five-element extreme beyond two trivial corners]
\label{CP2:A-baseline-trivial-pair-five}
Suppose a three-object category has full linear acyclic support,
its first two corners are trivial, and its last corner has
order at most five.  Its rational degree-two products are
zero, with no bound on its hom-set sizes.  The dual statement
holds at the other end.
\end{lemma}
\begin{proof}
The full subcategory $D$ on the first two objects has graph
classifying space.  If the last corner $M$ is not $P_5$,
its second rational cohomology is zero by the small-monoid
classification, and the two-block lemma applies.

For $M=P_5$, the category of elements of the incoming
$D$-act is itself a graph $\Gamma$: it has one vertex for
each incoming arrow from either source object, and one edge
for each pair consisting of an arrow of $D$ and an incoming
arrow from its target.  There are no composable pairs of
nonidentity arrows inside $D$, so there are no higher
nondegenerate simplices.  The commuting $P_5$ action sends
vertices to vertices and each parametrized edge to the edge
of its image pair.  Thus it is precisely a finite graph
action of the kind in Lemma~\ref{CP2:A-high-P5-graph},
after taking opposites if needed; $P_5$ is isomorphic to
its opposite by transposing its band coordinates.

The two-block model identifies the quotient which collapses
$BD$ with the cofiber of $\Gamma_{hP_5}\to BP_5$.
Since $H^2(BD;\mathbb Q)=0$, the quotient map induces a
surjection onto $H^2(BC;\mathbb Q)$ by the relative
cohomology sequence.  Every second class is therefore the
pullback of a class whose products vanish by the graph
cofiber lemma.  Naturality of the cup product proves the
claim.  Neither boundedness nor a rank-one assumption was
used.
\end{proof}


\subsection{A fresh four-object census at sixteen arrows}

\begin{proposition}[Proved by complete finite computation]
\label{CP2:A-baseline-four-census}
No category on four objects with at most sixteen arrows has
classifying space of type $\mathbb{CP}^2$.
\end{proposition}
\begin{proof}
The excess-eight theorem excludes at most fifteen arrows.
At sixteen arrows the ninth-excess support theorem requires
acyclic support, and Proposition~\ref{CP2:A-baseline-four-large}
excludes every corner of order at least eight.  It remains
to check exactly sixteen arrows, four objects, acyclic
support, and corner orders at most seven.

\paragraph{Complete profile scope.}
The new program \path{converse/scope17/four16_scope.py}
enumerates all subsets of the six upper-triangular pairs on
four objects, retaining exactly the transitive connected
supports.  There are eighteen with these topological labels.
Every finite poset admits a linear extension, so no support
is lost.  It assigns positive integers to the four diagonal
positions and every supported cross-hom, with total sixteen,
and restricts the diagonal entries to at most seven.  This
gives $103370$ labelled decorations.  For each it tries every
object permutation and the opposite matrix, retaining the
least upper-triangular representative.  The resulting
$18217$ profiles cover every category in the stated scope;
the use of opposition preserves the homotopy type of the
nerve.

The following sufficient obstructions remove $17935$ of
these profiles.  Several can apply to the same profile;
the ledger records all reasons, and does not add overlapping
counts.  Singleton cross-homs use
Lemma~\ref{CP2:A-singleton-support}; four-object support has
no fourth homology.  Semisimple corners of orders one or
two use the support-height bound.  A full singleton cut is
a join.  A cut into two blocks with zero second cohomology
uses Lemma~\ref{CP2:A-two-block}; each proper block has at
most three objects.  The block test proves zero second
cohomology component by component using a corner of order
at most four, the H2-cost bound, singleton cross-homs, or
semisimple corners and support height at most one.

A singleton arrow into an extreme corner invokes the
fixed-arrow corollary whenever the complementary block
passes that test.  Extreme corners of order at most four
whose every stalk has size one or two use the binary-stalk
lemma, under the necessary rank-one cohomology hypothesis.
With at most two parallel excesses and corners at most
four, Lemma~\ref{CP2:A-high-two-parallel} applies.  More
generally, if no two parallel hom-sets are serial and every
corner of potentially nonzero second cohomology is incident
to a singleton edge, every global second class vanishes on
those corners: the singleton-edge join makes its endpoint
inclusions null-homotopic.  Thus
Lemma~\ref{CP2:A-high-parallel-cut} applies.  All support
cycle ranks here are at most three, within its hypothesis.

The scope also uses a literal singleton-comma deletion.
If an object has exactly one incoming cross arrow from a
trivial corner, its comma category from the deletion is a
point; the outgoing assertion is dual.  This reduction is
used only when the deletion saves at least three arrows.
The remaining three-object category then has at most
thirteen arrows and is excluded by the unrestricted
excess-eight theorem.  No unproved sixteen-arrow baseline
is invoked by this cut.

For a support tree, transitivity implies that there is no
chain of two strict comparisons.  Consequently its nerve
is the union of its edge-collage nerves along their corner
nerves.  The tree test is used only when every internal
corner has order at most five.  Such corner spaces have
zero rational homology in degrees at least three.  The
iterated Mayer--Vietoris sequences therefore pass bounded
homology to every edge and make fourth-cohomology
restriction to the edges injective.  Every edge is
required to satisfy one of the following already proved
tests: at most two cross arrows; both corners of order at
most four; at most eleven total arrows; one corner of order
at most two and the other of order at most eight; or corner
orders three and five.  The fourth case uses exact orbit
comparison and the bounded order-eight act theorem; the
last uses the $P_5$ graph theorem, after the small-corner
classification.  Thus every restriction of a global
degree-two square is zero, and injectivity makes the
global square zero.

Finally, a star with center order at most eight and all
leaves of order at most two admits exact orbit comparison
at every leaf.  This is a rational homology and cohomology
ring comparison to a star with trivial leaves.  Identifying
those leaf points adds only circles to its classifying
space and gives an ordinary act cone at the center.
Boundedness and all cohomology in degrees at least two
are preserved, so the order-eight act theorem applies.
These are all the profile cuts in the code.  An independent
targeted review checked every predicate against the named
lemma, including the transfer and deletion hypotheses.

\paragraph{Complete table search in the residual scope.}
There are $282$ residual profiles.  The largest normalized
two-chain basis has $66$ elements.  The program
\path{converse/scope17/four16_search.cpp} uses all ordinary
corner monoids from Theorem~\ref{CP2:A-monoid-census}, with
both isomorphism orientations retained.  It fixes these
corner multiplications and the identity laws.  Each other
composable pair has as its domain the complete cross-hom
of the required type.  Propagation uses only equations
$(ab)c=a(bc)$ with known inner products; it either rejects
unequal known values or forces an unknown outer product.
It then branches over every possible value of a remaining
entry.  Induction on this finite search tree proves that
every associative table with the chosen corners occurs.

The sole table-level structural cut is when all four
corner algebras are semisimple.  Nondegeneracy modulo
$32003$ of the regular trace pairing certifies
nondegeneracy over $\mathbb Q$ and hence semisimplicity.
The height-three bar bound then gives $H_4=0$.
No inherited act-cone or projective-endpoint table pruning
is enabled.  Retraction pruning in the generic engine
cannot fire, since the profiles have no arrows in both
directions between distinct objects.

The remaining tables undergo the necessary tests
$H_1(-;\mathbb F_2)=0$ and $\dim H_2(-;\mathbb F_2)=1$.
For tables passing both, exact cross-hom relabelling keys
remove only repeated typed composition tables.  Every
nonidentity composable product is in the key, while
identity products are already fixed.  Cache eviction can
retain duplicates and does not discard a new orbit.
The reported retained count is therefore a record count,
not an isomorphism-class count.

The producer supplies the full typed table, a nonboundary
two-cocycle, and either a literal three-cochain whose
coboundary is its Alexander--Whitney square, or a claim of
positive third mod-two homology.  The separate checker
\path{converse/scope17/check_four16.cpp} imports neither
the search nor the cup producer.  It reconstructs all
normalized chains, checks types, identities, every
associativity equation, connectedness, both preliminary
homology tests, the cocycle and its nonboundary class,
and every square-primitive equation.  For the other
reason it recomputes $H_3$.  The gzip end and CRC and the
generator/checker record totals are also checked.

All $282$ profile intervals completed.  The producer visited
$155765567$ corner-representative associative tables, of which
$15341745$ passed the first two homology tests before relabelling.
The independent checker accepted $1721983$ retained records:
$1719138$ have literal zero-square primitives and $2845$ have
positive third mod-two homology.  No record remains unexcluded.
The complete join \path{outputs/four16/scope_join.json} verifies
all profile indices, the exact residual ledger, every completed
interval, generator/checker count equality, source hashes and
all closed certificate-stream hashes.  The proof ledger and all
$282$ checked streams are retained under \path{outputs/four16/}.
The earlier interrupted FIFO-transport attempt is preserved
separately and contributes no completion claim.

This completes the residual exclusion by
Proposition~\ref{CP2:A-cup-certificate} and hence proves
the proposition in its full stated scope.
\end{proof}


\subsection{Further strong-component reductions for the baseline}

\begin{lemma}[Proved: a binary idempotent representable]
\label{CP2:A-baseline-binary-representable}
Let $S$ be a finite category, let $e\in S(a,a)$ be idempotent,
and suppose $1\leq |S(a,i)e|\leq2$ for every object $i$.
Every positive-degree rational cohomology product of $BS$
is zero.  There is a dual assertion for $eS(i,a)$.
\end{lemma}
\begin{proof}
Let $F(i)=S(a,i)e$, as a covariant finite act.  The category
of elements of the full representable $S(a,-)$ has initial
object $(a,1_a)$ and is contractible.  The maps $u\mapsto ue$
define a natural idempotent endomorphism of this act and
give a retraction of its category of elements onto the
category of elements of $F$.  Its classifying space is
therefore a retract of a contractible space and is itself
contractible: compose a contraction with the inclusion and
retraction.  The act cone of $F$ is the homotopy cofiber of
this contractible space mapping to $BS$, and the inclusion
of $BS$ in the cofiber is a homotopy equivalence.  The
left-act version of Lemma~\ref{CP2:A-high-affine} applies
to the one- or two-element stalks and annihilates all
positive rational products.  Taking opposites proves the
dual statement.
\end{proof}

\begin{lemma}[Proved: an objectwise binary projective cover]
\label{CP2:A-baseline-binary-projective-cover}
Let $S$ be a finite category.  Suppose the constant left
$\mathbb QS$-module $\mathbf1$ is the augmentation quotient
of a projective module $P$ which at each object has a basis
of one or two elements, and whose arrow maps send basis
elements to basis elements.  Every finite right $S$-act
with all stalks nonempty has an act cone with zero rational
degree-two products.  An opposite endpoint whose trivial
right rational module is projective may replace the trivial
sink.  There are left-right duals.
\end{lemma}
\begin{proof}
Put $R=\mathbb QS$ and $V=\ker(P\to\mathbf1)$.
Each $V_i$ has dimension zero or one.  Use the difference
of the two basis vectors when it has dimension one.
For an arrow $f:i\to j$ the induced difference map is
the scalar $\chi_f\in\{0,1,-1\}$, with zero whenever
one difference space is zero.  These typed maps compose
as module maps.

For the right act $F$, set
$L=\ker(\mathbb QF\to\mathbf1_{\rm right})$, using
objectwise augmentation.  All stalks are nonempty, so
this sequence is exact.  Tensoring a projective left
resolution with it, as in the relative calculation in
Lemma~\ref{CP2:A-baseline-eight-act}, identifies reduced
act-cone chains with $(L\otimes_R^{\mathbf L}\mathbf1)[1]$.
Finite-dimensional rational duality and the projectivity
of $P$ give a natural surjection
\[
 \operatorname{Hom}_R(V,DL)\longrightarrow
 \operatorname{Ext}^1_R(\mathbf1,DL)
  =H^2(B\mathcal C(S,F);\mathbb Q).
\]
Here $D$ denotes rational vector-space duality, object by
object, and the first map is the connecting homomorphism
of the projective-cover sequence.

Lift a second class through this map.  In the chosen
difference coordinates it is a family of functionals
$\ell_i:L_i\to\mathbb Q$, zero when $V_i=0$, satisfying
$\ell_i L(f)=\chi_f\ell_j$.  Choose $s_i\in F(i)$ and
put $q_i(s)=\ell_i([s]-[s_i])$.  Subtracting two such
values gives
\[
 q_i(F(f)s)=\chi_f q_j(s)+\beta_f,
 \qquad \beta_f=q_i(F(f)s_j).
\]
Thus the chosen class is encoded by a finite affine
diagram of rational values.

Take one rational grid spacing $\delta>0$ containing
all these values and translations.  Let $I_i$ be the
convex hull of $q_i(F(i))$ and let
$T_i=I_i\cap\delta\mathbb Z$.  The displayed affine
maps preserve the intervals and their grids.  They
compose exactly: on a nondegenerate source interval
their values at the two endpoints determine them, and
on a degenerate interval there is only one value to
check.  Therefore $T$ is an actual finite right act,
and $q:F\to T$ is an actual act map.

On $\ker(\mathbb QT_i\to\mathbb Q)$, the coordinate
functional $\sum a_t[t]\mapsto\sum a_t t$ satisfies
the same typed equation: translations disappear because
$\sum a_t=0$.  It is zero when $V_i=0$, and its pullback
is $\ell_i$.  Naturality of the connecting homomorphism
therefore lifts the original class from a second class
on the act cone of $T$.

For clarity, the common-grid extension of the interval
argument is strict.  Each $I_i/T_i$ is a wedge of circles
indexed by consecutive grid intervals of length $\delta$.
An affine map of slope $1$ or $-1$ takes one grid interval
onto another of the same length, preserving or reversing
its circle parameter; slope zero gives a pointed constant.
The degree-two circle maps commute with every such map.
They induce a self-map of the pointed category bar which
is multiplication by two on its reduced cellular chain
complex, in every total degree.  It is therefore
multiplication by two on every positive cohomology group.
Multiplicativity gives $2uv=4uv$ for positive classes,
so all those products vanish rationally.  In particular
the square of the lifted class, and hence of the original
class, is zero.  Rational polarization proves the mixed
degree-two assertion.

For a projective opposite endpoint, the exact orbit
comparison of Lemma~\ref{CP2:A-high-orbits} gives an
actual functor inducing a rational cohomology-ring
isomorphism to an act cone.  Its stalks remain nonempty,
so the preceding argument applies.  Opposition gives
the remaining statements.
\end{proof}

\begin{corollary}[Proved: two-row and two-orbit instances]
\label{CP2:A-baseline-two-row-instances}
The preceding lemma applies when $S$ is strongly connected
with trivial minimal group and at most two minimal rows
at every object.  It also applies whenever there are
$a,e,G$, with $e^2=e$ and $G$ a finite subgroup of the
units of $eS(a,a)e$, such that every orbit set
$S(a,i)e/G$ has one or two elements.
\end{corollary}
\begin{proof}
In the first case take $P=(\mathbb QS)e$ at a minimal
idempotent.  Its bases are the minimal rows at each object.
In the second case take
$P=(\mathbb QS)p$, where $p=|G|^{-1}\sum_{g\in G}g$.
It is projective because $p$ is idempotent.  The normalized
averages of the $G$-orbits give the required basis.  An
equivariant map sends a normalized orbit average to the
normalized average of its image orbit: all fibers have
equal size.  Thus every arrow acts by a basis map, and
objectwise augmentation onto the constant module is
surjective.  These are precisely the hypotheses needed.
\end{proof}

\begin{proposition}[Proved: wholly strong two-object categories through seventeen]
\label{CP2:A-baseline-strong-two}
A strongly connected category on two objects with at most
seventeen arrows cannot realize $\mathbb{CP}^2$.
\end{proposition}
\begin{proof}
If one object is a retract of the other, deletion saves
at least three arrows and leaves a monoid of order at most
fourteen.  Theorem~\ref{A:monoid-minimum23-classification}
excludes it.  Suppose no object retract is present.
Both identities lie outside the categorical minimal ideal.
Write $r=r_0+r_1$ and $c=c_0+c_1$ for its total minimal
row and column counts and $G$ for its maximal group.
If $r=2$ or $c=2$, all local row or column counts are one
and the category has finite-group homotopy type by
Lemma~\ref{CP2:A-high-Rees}; its positive rational
cohomology is zero.

Otherwise $r,c\geq3$ and $|G|rc+2\leq17$.
Consequently $G=1$ and $\min(r,c)=3$, since
$2\cdot3\cdot3+2=20$ and $4\cdot4+2=18$.
On the side with total three the two positive local
counts are one and two.  The corresponding idempotent
representable has one or two states at each object.
Lemma~\ref{CP2:A-baseline-binary-representable}, or its
dual, makes every positive rational product zero.
The nonzero projective-plane square is impossible.
\end{proof}

\begin{proposition}[Proved: cyclic-support three-object reduction through sixteen]
\label{CP2:A-baseline-strong-three}
A three-object category with cyclic support and at most
sixteen arrows either fails to realize $\mathbb{CP}^2$,
or admits an object-retract deletion to a two-object
acyclic-support category with at most thirteen arrows.
In particular, the cyclic-support part of the
three-object baseline requires only that explicit
two-object baseline dependency.
\end{proposition}
\begin{proof}
If the category is wholly strong, a retract deletion
leaves a wholly strong two-object category and the
preceding proposition applies.  Without an object
retract, its three identities are outside the minimal
ideal.  Unless it has finite-group type, the total
row and column counts are at least four, requiring
at least $4\cdot4+3=19$ arrows.  Thus the wholly
strong case is excluded.

The remaining connected case has a proper strong pair
$S$ and one object with corner $N$.  A retract inside
the pair has exactly the conclusion asserted: deletion
saves at least three arrows and the remaining support
is acyclic.  Assume that no object retract remains.
By opposition take $S$ as the source component.  There
are nonempty outgoing stalks at both its objects.

If $H^2(BS;\mathbb Q)$ and $H^2(BN;\mathbb Q)$ both
vanish, use the two-block lemma.  If the first is
nonzero, its total minimal rows and columns are at
least three.  Since $|S|\leq16-1-2=13$, the minimal
group is trivial, both totals are exactly three,
and $|S|\geq11$.  Every local row count is at most
two.  The other corner has order at most three.
If its order is at most two, its trivial right module
is projective and the two-row instance of the preceding
lemma kills all second products on the collage.  If
its order is three, the entire budget forces
$|S|=11$ and exactly one outgoing arrow at each object.
The collage is the full categorical join of $S$ and
$N$, whose positive cohomology products vanish.

It remains that $H^2(BS;\mathbb Q)=0$ and
$H^2(BN;\mathbb Q)\ne0$.  Then $5\leq|N|\leq8$
and $|S|\leq9$.  The two identities of $S$ are
outside its ideal.  Its minimal group is trivial,
since a group of order at least two would already
cost $2\cdot2\cdot2+2=10$ internal arrows.  One
total row or column count must be two, because
$3\cdot3+2=11>9$.

If the row total is two, the one-row source-component
collapse of Lemma~\ref{CP2:A-high-one-sided} gives an
ordinary left-act cone over $N$.  This is an actual
homotopy equivalence, so its rational homology is
bounded under the putative projective-plane hypothesis.
Lemma~\ref{CP2:A-baseline-eight-act} excludes it.
Otherwise the column total is two and the row total
is three, so $|S|\geq8$.  If its two outgoing stalks
are singleton, the join argument applies.  If not,
the only budget is $|S|=8$, $|N|=5$, and three
outgoing arrows.  The pair consists of its six-arrow
minimal band and its two identities; its local row
counts are two and one, and both local column counts
are one.  At the one-row object every incoming
internal arrow is fixed by its local minimal
idempotent, and its only unit is its identity.
There are no external entrances because $S$ is a
source.  The one-sided object deletion therefore
applies.  The retained collage has at most two
cross arrows, since either retained outgoing stalk
had size one or two.  It is excluded by the
singleton join or the universal two-cross lemma.
This exhausts the proper-pair cases without using
an inherited composition census.
\end{proof}


\begin{theorem}[Proved: cyclic three-object support through seventeen]
\label{CP2:A-strong-three-seventeen}
A three-object category with cyclic support and at most
seventeen arrows cannot have type $\CP^2$.
\end{theorem}
\begin{proof}
A disconnected category is excluded by connectedness of
the target. An object-retract deletion saves at least
three arrows. If its two remaining objects are strongly
connected, Proposition~\ref{CP2:A-baseline-strong-two}
applies; otherwise their arrow count is at most fourteen
and the complete two-object baseline applies. Assume
there is no object retract.

If all three objects are strongly connected, their
identities are outside the categorical minimal ideal.
Unless all local row counts or all local column counts
are one, both total counts are at least four and the
category has at least $4\cdot4+3=19$ arrows. The exceptional
one-row or one-column case has finite-group homotopy type
by Lemma~\ref{CP2:A-high-Rees}, and is also excluded.

It remains that a proper strong pair $S$ is connected
to one corner monoid $N$. By opposition take $S$ to be
the source component. Let $q=|S|$, $m=|N|$, and let $h$
be the total number of its outgoing arrows. Each of the
two stalks is nonempty, so
\[
 q+m+h\le17,\qquad h\ge2.
\]
If $h=2$, the collage is the full categorical join of
$S$ and $N$ and has zero positive rational products.
We may suppose $h\ge3$. If both $H^2(BS;\mathbb Q)$
and $H^2(BN;\mathbb Q)$ vanish, the two-block lemma
already excludes the target.

Suppose first $H^2(BS;\mathbb Q)\ne0$. Its total minimal
row and column counts $r,c$ are at least three; otherwise
it has finite-group type. Since $q\le13$, the minimal
group is trivial and $rc+2\le13$ forces $r=c=3$.
In particular $q\ge11$, every local row count is at most
two, and $m\le3$. If $m\le2$, the trivial right rational
$N$-module is projective. The objectwise binary projective
cover, Corollary~\ref{CP2:A-baseline-two-row-instances},
kills all second products on this collage.

If $m=3$, the budget forces $q=11$ and $h=3$. The strong
pair then consists exactly of its nine-element categorical
minimal band and its two identities. Choose the object
$b$ whose local row count is one. Every incoming arrow
from its other object lies in the minimal ideal and is
fixed by a local minimal idempotent $e_b$. The only unit
of its corner is its identity, since all its other elements
lie in that ideal. There are no external entrances into
$S$. Choose minimal cross arrows $u,v$ with $vu=e_b$.
All hypotheses of the one-sided object-deletion lemma
hold. Deleting $b$ gives an actual homotopy equivalence
to a two-object category with fewer than seventeen arrows,
contradicting the complete two-object baseline.

Finally suppose $H^2(BS;\mathbb Q)=0$ and
$H^2(BN;\mathbb Q)\ne0$. Then $m\ge5$. If either outgoing
stalk is a singleton, its arrow is fixed by every $N$
actor. Corollary~\ref{CP2:A-fixed-arrow}, with source
block $S$, kills all second products immediately.
Thus both stalks have at least two arrows: $h\ge4$ and
$q\le8$. The minimal group of $S$ is trivial, since a
nontrivial group would require at least
$2\cdot2\cdot2+2=10$ internal arrows. With two identities
outside the ideal, $rc+2\le8$ shows that one total row
or column count is two and the other is at most three.

If $r=2$, there is one minimal row per object. The actual
one-row source-component collapse replaces $S$ by a trivial
object and gives an ordinary left-act cone over $N$.
Here $m\le7$, since $q\ge6$ and $h\ge4$. Its rational
homology is bounded under the proposed target, and the
bounded-nine act theorem excludes it. If $r>2$, necessarily
$r=3$, $c=2$, $q=8$, $m=5$, and $h=4$. The pair consists
of its six-element minimal band and its two identities.
Its one-row object again has all incoming arrows fixed
by its minimal idempotent and no nontrivial units. There
are no external entrances. The same one-sided deletion
gives a two-object category with at most sixteen arrows,
which is excluded. This exhausts every cyclic support.
\end{proof}


\begin{theorem}[Proved: cyclic four-object support through seventeen]
\label{CP2:A-strong-four-seventeen}
A four-object category with cyclic support and at most
seventeen arrows cannot have type $\CP^2$.
\end{theorem}
\begin{proof}
Suppose $BC\simeq\CP^2$. In particular the category is
connected and $\dim H^2(BC;\mathbb Q)=1$, as required
by the binary-stalk obstruction below. If an object in a proper strong
component is a retract of another, its deletion removes
at least four arrows: its identity, arrows to and from
another object of the component, and an external incident
arrow. The last exists because every object of a strong
component inherits any external support adjacency. If the
whole category is strong, deletion removes at least seven
arrows. Thus an object retract always reduces to a
three-object category with at most thirteen arrows,
excluded by Theorem~\ref{CP2:A-high-all-eight}. Assume
there is no object retract.

A strong component on $r$ objects has its $r$ identities
outside its categorical minimal ideal, so costs at least
$r^2+r$ arrows. A wholly strong four-object category
therefore costs at least twenty. A proper three-object
strong component costs at least twelve and has at least
three interface arrows to the remaining corner. That
corner has order at most two. The strong component has
finite-group type: otherwise its total row and column
counts would both be at least four and its internal cost
would already be nineteen. Both blocks have zero second
rational cohomology, so the two-block lemma applies.
If there are two strong pairs, each costs at least six
and their interface at least four. Each internal pair
therefore has at most seven arrows and finite-group type.
The same two-block argument excludes this case.

There remains one strong pair $S$ and two singleton
components. Put $q=|S|$. The outside components have
corner monoids of orders $b,c$.

\emph{Extreme pair.} By opposition suppose $S$ is a
source. Let $D$ be its two-object complement, and let
$h$ count the interface arrows. Connectedness gives
$q\le12$. If $q=11$ or $12$ but $H^2(BS;\mathbb Q)=0$,
then $|D|\le4$ since $h\ge2$; the second-homology cost
lemma gives $H^2(BD;\mathbb Q)=0$, and the two-block
lemma applies. For $q\le10$, the row-column bound makes
$BS$ of finite-group type, so $H^2(BS;\mathbb Q)=0$.
If also $H^2(BD;\mathbb Q)=0$, the two-block lemma
finishes. If the latter group is nonzero, then $|D|\ge6$:
for connected $D$ this is the second-homology cost lemma,
and for disconnected $D$ one corner has order at least
five and the other at least one. If $q=6$ or $7$, the
minimal group is trivial and each object has one row
and one column. The actual source-component orbit
collapse saves at least $q$ arrows, as proved in
Theorem~\ref{CP2:A-high-all-eight}, reducing to the
three-object bound. The case $q=10$ is impossible here,
since $q+|D|+h\ge10+6+2>17$.

It remains in this branch that $q=8$ or $9$. The complement
must be connected: otherwise it costs at least six and
has at least four interface arrows, exceeding seventeen.
Moreover $S$ can meet only one of its objects, say $u$,
since meeting both would again require four interface
arrows. The outside edge is directed $v\to u$; an edge
in the other direction would force another interface.
The inequality is now $q+|D|+h\le17$, with $|D|\ge6$
and $h\ge2$. At $q=9$ this forces $|D|=6,h=2$; at
$q=8$ either $h=2$, or $|D|=6,h=3$.

The minimal group is trivial and one total row or column
count is two. For $h=2$, all external stalks are singleton,
so the one-row source collapse or the fixed-exit column
dual collapses the pair, saving at least seven arrows.
For $q=8,h=3$, a row total of two gives the source
collapse immediately. Otherwise the totals are three
rows and two columns, and $S$ consists exactly of its
six-element minimal band and two identities. At its
one-row object all incoming internal arrows are fixed
by the local minimal idempotent, and its only unit is
its identity. There are no external entrances. The
one-sided object-deletion lemma applies, removing at
least its two corner arrows, three internal cross arrows,
and one external arrow. This again leaves at most
thirteen arrows on three objects.

Now suppose $H^2(BS;\mathbb Q)\ne0$. Then $q=11$ or
$12$, with trivial minimal group and total rows and
columns both three. If $D$ is disconnected, the only
budget is $q=11$, two trivial outside corners, and four
singleton interface stalks. This is the full join of
$S$ with two discrete objects, hence has zero products.
If $D$ is connected, $S$ meets only its final object $u$,
as before. We have $|D|+h\le6$. If $h=2$, then
$|D|\le4$ and the corner of $u$ has order at most two.
The full category on $S\cup\{u\}$ is the join of $S$
and that corner and is rationally acyclic. Its nerve
and $ND$ cover the whole nerve, intersecting at the
rationally acyclic corner $u$; $H^2(BD;\mathbb Q)=0$
by its size. Hence the whole second cohomology is zero.
The only remaining case is $q=11$, $|D|=3$, $h=3$.
Here $D$ is a single interval between trivial corners,
so gluing it at $u$ does not change the homotopy type
of $S\cup\{u\}$. Each local row count of $S$ is at most
two and $u$ is trivial. The binary projective-cover
lemma kills every second product on this act cone.

\emph{Nonextreme pair.} The condensation is the full
chain $u<S<v$. Write $\alpha_i=|C(u,i)|$ and
$\beta_i=|C(i,v)|$ for its two entrance and exit
stalks, and $z=|C(u,v)|$. They are all positive. Put
\[
 P=\sum_i(\alpha_i-1)+\sum_i(\beta_i-1),\qquad
 \delta=(b-1)+(c-1)+(z-1).
\]
Then
\[
 q+7+P+\delta\le17,
 \qquad q\le10,\qquad P+\delta\le10-q.
\]
A nontrivial minimal group forces $q=10$, all five
external hom-sets singleton, and both outside corners
trivial. The initial object $u$ is then initial in
the whole category. Assume henceforth the minimal
group is trivial.

For $q\ge7$ and $z\le2$, both endpoint corners have
order at most four and $P\le3$. If both entrance
and exit sides had a stalk of size at least three,
they would cost at least four units of $P$. Thus at
least one endpoint has its two stalks at $S$ of size
at most two; its direct stalk $z$ is also binary.
Lemma~\ref{CP2:A-high-small-binary}, whose opposite
block may be any finite category, kills its square.

If $q\ge7$ and $z\ge3$, only $q=7$ or $8$ is possible.
For $q=8$ one has $P=0$; one row or one column at every
object and the fixed singleton interfaces permit the
block collapse. For $q=7$, $P\le1$. The pair consists
of its four-arrow minimal band, two identities, and
one extra arrow. Its local minimal idempotents are
zeros. Of the four conditions that one local idempotent
fix all incoming or all outgoing cross arrows, the
extra internal arrow can violate at most two, and
the external nonsingleton stalk at most one. At least
one condition holds. The zero guarantees the unit
condition in the one-sided deletion lemma. Deletion
removes at least six arrows, including one entrance
and one exit, so the three-object bound applies.

Finally let $q=6$. The pair is exactly its four-arrow
minimal band plus its two identities. All internal
cross arrows are fixed by the local zeros, so a failed
one-sided condition requires a nonsingleton external
stalk at that specific object and side. If an endpoint
has order five, the budget forces $P=0,z=1$ and the
pair collapses. Otherwise both endpoint corners have
order at most four. For $z\ge3$ one has $P\le2$,
so fewer than four fixed conditions can fail and
one-sided deletion applies. For $z\le2$, failure
of both small-endpoint binary tests requires a stalk
of size at least three on each side, using $P\ge4$.
Equality is forced, with $b=c=z=1$; each side has
one three-element stalk and one singleton stalk.
Thus two of the four fixed conditions still hold,
and one-sided deletion again applies. Every alternative
is excluded, completing the cyclic four-object scope.
\end{proof}


\subsection{Graph cofibers over corners of order at most six}

\begin{theorem}[Proved: the bounded small-corner graph obstruction]
\label{CP2:A-baseline-six-graph}
Let a finite monoid $M$ of order at most six act on the
right of a nonempty finite graph $\Gamma$.  Require each
actor to take a vertex to a vertex and each parametrized
edge to an edge, preserving or reversing its unit parameter,
or to a vertex.  If
\[
 X=\operatorname{hocofib}(\Gamma_{hM}\longrightarrow BM)
\]
has bounded rational homology, every rational degree-two
product on $X$ is zero.  The left-action statement is dual.
\end{theorem}
\begin{proof}
If $H^2(BM;\mathbb Q)=0$, each second class is relative
to $BM$ and lifts from the suspension quotient of the
cofiber.  Its products vanish.  Otherwise the minimal
ideal must have at least two rows and two columns;
the small cardinality forces a $2\times2$ rectangular
band $K$, with trivial maximal group and at most two
outside elements.  Put $R=\mathbb QM$ and
$A=R/\mathbb QK$.

Choose $e=x_{00}$ and $e_1=x_{01}$.  Their image graphs
$\Gamma_0=\Gamma e$ and $\Gamma_1=\Gamma e_1$ are
retract subgraphs.  An image under a band element depends
only on its column, and right multiplication by $e_1$
and $e$ gives mutually inverse graph isomorphisms between
these two images.  If they share a vertex $v$, then
$ve=ve_1=v$ and
$vx_{ij}=vex_{ij}=vx_{0j}=v$.  For every $m\in M$,
$vm=v(em)=v$ because $em\in K$.  This common fixed
vertex gives an action-bar section of
$\Gamma_{hM}\to BM$.  The cofiber sequence makes every
positive cohomology class relative to $BM$, and the
suspension quotient again kills its products.

Suppose instead that the two images have disjoint vertex
sets, hence disjoint edge sets.  Use the augmented graph
complex of right $R$-modules
\[
 C=[C_1(\Gamma;\mathbb Q)\longrightarrow
                  \ker(C_0(\Gamma;\mathbb Q)\to\mathbb Q)]
\]
in degrees one and zero.  The multiplication map
$Ce\otimes_{\mathbb Q}eR\to C$ is injective in both
degrees.  Indeed, $eR$ has the two column basis elements.
In degree one their images have disjoint edge supports;
orientation reversals change only signs.  In degree zero,
$Ce$ consists of the augmentation-zero vertex chains in
$\Gamma_0$, and its two column images have disjoint
vertex supports and each has augmentation zero.  Thus
there is no nonzero relation between the two copies.

Its image is exactly $C\mathbb QK$, using the literal
band factorization $\mathbb QK=Re\otimes_{\mathbb Q}eR$.
Consequently
\[
 D=C/(Ce\otimes_{\mathbb Q}eR)
\]
is an actual complex of right $A$-modules concentrated
in degrees zero and one.  The subcomplex being divided
out is a complex of projective right $R$-modules; after
tensoring with the trivial left module it is just $Ce$,
still in those two degrees.  Therefore
\[
 H_j(C\otimes_R^{\mathbf L}\mathbb Q)
  =H_j(D\otimes_R^{\mathbf L}\mathbb Q),\qquad j\geq3.
\]

Let $V$ be the one-dimensional row augmentation module.
The exact sequence $0\to V\to Re\to\mathbb Q\to0$
and the equality $De=0$ give
$D\otimes_R^{\mathbf L}\mathbb Q
 \simeq(D\otimes_R^{\mathbf L}V)[1]$.
There is no higher change-of-rings correction: the
left-projective resolution
$0\to\mathbb QK\to R\to A\to0$ gives
$U\otimes_R^{\mathbf L}A\simeq U$ for every right
$A$-module $U$, since $Ue=0$.  Applying this naturally
to the two terms of $D$ identifies
$D\otimes_R^{\mathbf L}V$ with
$D\otimes_A^{\mathbf L}V$.
Finally, the reduced graph-cofiber chains are
$(C\otimes_R^{\mathbf L}\mathbb Q)[1]$, by the
augmented graph sequence.  Combining the two shifts
gives the precise formula
\[
 H_n(X;\mathbb Q)
   =H_{n-2}(D\otimes_A^{\mathbf L}V),\qquad n\geq4.
\]
This calculation applies also when $\Gamma$ is disconnected.

The algebra $A$ has dimension one or two.  If it is
semisimple, $V$ is projective and its tensor with $D$
has no homology in degree two, so $H_4(X;\mathbb Q)=0$.
The only other case here is
$A=\mathbb Q[\varepsilon]/(\varepsilon^2)$: with one
outside element $a$ beyond the identity, its square
in the quotient is $0$, $a$, or $1$, the last two
giving semisimple algebras.  In the dual-number case
$V$ is its residue module $k=\mathbb Q$.

Only an elementary two-term recurrence is needed.
Tensor $0\to k\to A\to k\to0$ derivedly with $D$.
The middle tensor is $D$, with zero homology in
degrees at least two.  Its long exact sequence gives
\[
 H_j(D\otimes_A^{\mathbf L}k)
  \cong H_{j-1}(D\otimes_A^{\mathbf L}k),\qquad j\geq3.
\]
Boundedness of $X$, through the displayed localization
formula, makes these groups eventually zero.  The
recurrence forces the degree-two group to be zero,
again giving $H_4(X;\mathbb Q)=0$.  Thus either the
fixed-vertex relative argument or fourth-homology
vanishing annihilates every degree-two product.
\end{proof}

\begin{lemma}[Proved: an equivariant graph for a semisimple pair]
\label{CP2:A-baseline-semisimple-pair-graph}
Let $D$ have two objects with arrows only from the first
to the second and with semisimple rational corner
algebras $\mathbb QN,\mathbb QP$.  For any finite
right $D$-act $F$, its category of elements admits an
actual natural rational homology equivalence to a
finite graph.  Every commuting monoid action induces
strict unit-speed graph maps and makes this comparison
equivariant.  Also $H_j(BD;\mathbb Q)=0$ for $j\geq2$.
\end{lemma}
\begin{proof}
Put $U=F(0)$, $T=F(1)$ and $Z=D(0,1)$, viewed as
a left-$P$, right-$N$ biset.  The graph has vertex
sets $U/N$ and $T/P$ and edge set
$(T\otimes_P Z)/N$.  An edge $[t,z]$ joins $[tz]$
to $[t]$.  Here an orbit set is the quotient generated
by the action relations, including noninvertible actors.
The balanced relations $(tp,z)\sim(t,pz)$ and
$(t,zn)\sim(t,z)$ respect both endpoints.

The functor from the element category sends objects to
these vertex orbits, corner arrows to identities, and a
cross arrow $(t,z)$ to its displayed edge.  The balanced
relations are exactly what composition with corner
arrows requires.  No two cross arrows compose.  Thus
this is an actual functor, with no choice of orbit
representatives.  A commuting action preserves the
relations and sends each graph edge to an edge by
its parameter, so the comparison is equivariant.

Filter element-category bar chains by the number of
strict object changes, which is zero or one.  The
vertex complexes are
$\mathbb QU\otimes_{\mathbb QN}^{\mathbf L}\mathbb Q$
and
$\mathbb QT\otimes_{\mathbb QP}^{\mathbf L}\mathbb Q$.
The one-cross-arrow complex is
$\mathbb QT\otimes_{\mathbb QP}^{\mathbf L}\mathbb QZ
 \otimes_{\mathbb QN}^{\mathbf L}\mathbb Q$.
Semisimplicity puts their homology in degree zero,
where they are exactly the linearizations of the
two vertex orbit sets and the balanced edge set.
Their boundary maps are the actual displayed
endpoints.  The filtration therefore proves that
the functor induces the asserted rational homology
equivalence.  Apply the same construction to the
singleton act to obtain the last assertion about $BD$.
\end{proof}

\begin{corollary}[Proved: a semisimple pair opposite a corner through order six]
\label{CP2:A-baseline-semisimple-pair-six}
Let a three-object category have full linear acyclic
support.  If its first two corner algebras are semisimple
and its last corner has order at most six, bounded
rational homology implies that every rational degree-two
product is zero.  The opposite statement also holds.
In particular the size conditions $a,b\leq2$, $c\leq6$
exclude every projective-plane profile, regardless of
the three cross-hom sizes.
\end{corollary}
\begin{proof}
Use the first two objects as the source block $D$ and
the last corner as $M$.  The two-block model is
$BC\simeq\operatorname{hocolim}(BD\leftarrow E\to BM)$,
where $E$ is the homotopy-orbit space of the incoming
act's element category.  Collapse $BD$ and obtain
$X=\operatorname{hocofib}(E\to BM)$.
Since $H^2(BD;\mathbb Q)=0$, this actual quotient map
induces a surjection $H^2(X;\mathbb Q)\to H^2(BC;\mathbb Q)$.
Since all $H_j(BD;\mathbb Q)$ for $j\geq2$ vanish,
the cofiber sequence also passes bounded rational
homology from $BC$ to $X$.

The preceding lemma provides an equivariant rational
comparison of the element category to a graph $\Gamma$.
Its comparison remains a rational homology equivalence
after taking homotopy orbits: the nonnegative free
monoid bar preserves a coefficient quasi-isomorphism.
It gives an actual rational homology equivalence from
$X$ to $\operatorname{hocofib}(\Gamma_{hM}\to BM)$.
The bounded graph theorem, in its left-action version,
kills all second products there.  Pulling back through
the rational ring isomorphism and then through the
surjection in degree two proves the claim.  Corner
monoids of orders one and two have semisimple rational
algebras, giving the stated profile consequence.
\end{proof}

\begin{corollary}[Proved: the right-zero-three endpoint]
\label{CP2:A-baseline-right-zero-three-six}
Let $R_3$ be the identity-adjoined two-element right-zero
band.  A collage between $R_3$ and any monoid of order
at most six, with the $R_3$ acting on the left of the
connecting biset, has zero rational degree-two products
whenever its rational homology is bounded.  There is
no restriction on the number of connecting arrows.
\end{corollary}
\begin{proof}
The two band actors on the biset are retractions with
common image $U$.  Their action category has the
natural graph model used in the proof of
Lemma~\ref{CP2:A-high-P5-graph}: vertices are $U$,
and an outside state gives the edge between its two
retraction values.  Every commuting map sends such
an edge at unit speed to an edge or to a vertex.
Since $BR_3$ is contractible, the collage is the
graph cofiber of the other monoid's action.  Apply
Theorem~\ref{CP2:A-baseline-six-graph}.
\end{proof}


\subsection{Further complete linear-profile obstructions}

\begin{theorem}[Proved: a singleton long hom-set]
\label{CP2:A-baseline-singleton-long}
Let a three-object category have full linear acyclic support,
a singleton hom-set from the first to the last object, and
rationally acyclic corner classifying spaces.  Every positive
rational cohomology product of its classifying space is zero.
In particular this applies when all corner orders are at most
four.
\end{theorem}
\begin{proof}
Write $M,N,P$ for the corner monoids, $S,T$ for the two
adjacent bisets, and $U$ for the singleton long biset.
Decomposing nerve strings at their strict object changes
gives the triangular bar diagram with vertex spaces
$BM,BN,BP$, adjacent edge spaces
$A=B(N,S,M)$ and $B=B(P,T,N)$, long edge space
$K=B(P,U,M)=BP\times BM$, and face space
$F=B(P,T,N,S,M)$.

These are actual compatible bar maps.  A multi-simplex of
$F$ has the form $(p_*;t;n_*;s;m_*)$.  The maps to $A,B$
forget the opposite outer string and cross arrow.  The map
to $K$ retains $p_*,m_*$ and replaces the middle portion
by $t n_1\cdots n_q s$.  Middle faces preserve this
product; outer faces preserve it by the biset identities.
The composites to every vertex agree.  The homotopy
colimit of this diagram is $BC$: subdivision of the
products of these bar strings with their strict-object
simplex recovers the ordinary nerve.  In particular the
middle bar direction in $F$ is retained; no product
description $F=A\times B$ is assumed.

Replace the three vertex spaces and $K$ by points and
keep $A,B,F$ and both maps $F\to A,B$.  This is an actual
map of diagrams and an entrywise rational homology
equivalence.  For $K$ use rational Kunneth and the
endpoint acyclicity.  The finite homotopy-colimit
filtration makes the resulting map a rational homology
equivalence and hence a rational cohomology-ring
isomorphism.

Collapse the resulting long-edge interval, a contractible
CW subcomplex.  Its two endpoint vertices become one
suspension vertex; the middle vertex is the other.
Each triangular prism $F\times\Delta^2$ modulo the
long side is the suspension of an interval parametrized
by $F$, attached at its two ends through $F\to A,B$.
Thus the quotient is exactly the unreduced suspension
of the double mapping cylinder
$\operatorname{hocolim}(A\leftarrow F\to B)$.
All positive products vanish in that suspension, and
the collapse is a homotopy equivalence.  Pullback
through the rational comparison proves the theorem.
The final assertion uses the checked small-monoid
classification.
\end{proof}

\begin{corollary}[Proved: a projective initial singleton]
\label{CP2:A-baseline-projective-singleton}
Suppose a full-linear category has projective trivial left
rational module at its initial corner, a singleton first
adjacent hom-set, and final corner order at most nine.
Bounded rational
homology forces every rational degree-two product to
vanish.  In particular a profile $(a,b,c;x,y,z)$ with
$a\leq2$, $x=1$, and $c\leq9$ satisfies the hypothesis.
The opposite condition also suffices.
\end{corollary}
\begin{proof}
The initial corner's actual exact-orbit quotient replaces it
by a trivial object and replaces $x,z$ by their orbit
sets, inducing a rational homology and cohomology-ring
isomorphism.  The first adjacent hom-set remains a
singleton.  The comma category for deleting the middle
object is therefore a point: there is just the unique
incoming arrow from the now trivial first corner.
The deletion is an actual homotopy equivalence to
the ordinary left-act cone on the remaining long
hom-set over the final corner.  Its boundedness
follows from these comparisons.  The bounded-nine
act theorem, Theorem~\ref{CP2:A-bounded-nine-acts},
annihilates its second products.  Pull back and take
opposites for the other orientation.
\end{proof}

\begin{proposition}[Proved: two small incoming stalks]
\label{CP2:A-baseline-two-small-incoming}
Let a full-linear three-object category have trivial
first two corners, a last corner $M$ of order at most
eight, two arrows from the middle to the last object,
and at most three arrows from the first to the last.
It cannot realize $\mathbb{CP}^2$.  There is no bound
on the first adjacent hom-set.
\end{proposition}
\begin{proof}
Write $X=C(0,1)$, $Y=C(1,2)$ and $Z=C(0,2)$.
The first two objects form a graph category $D$, so
$H^2(BD;\mathbb Q)=0$.  If $H^2(BM;\mathbb Q)=0$,
the two-block argument applies.  Otherwise the small
cardinality and the minimal-ideal row and column
requirements force trivial minimal group.  The
universal group of $M$ is then trivial: for a minimal
idempotent $e$, every group image satisfies
$\eta(m)=\eta(eme)$, and $eMe=\{e\}$.

Each $x\in X$ gives an $M$-equivariant map
$\beta_x:Y\to Z$.  A constant such map supplies a
globally $M$-fixed point of $Z$, so the fixed-arrow
corollary excludes all second products.  Assume all
these maps are injective.  If $|Z|=3$ and two image
sets differ, their singleton intersection is invariant
under $M$ and again gives a fixed point.  Hence all
images are one common two-element invariant set $U$.
The same holds automatically when $|Z|=2$.

Choose one bijection $Y\to U$ to identify these two
sets.  If two of the maps $\beta_x$ differ, a
transposition commutes with every actor on $Y$.
Its transformation centralizer on two points consists
of permutations.  The resulting homomorphism from
$M$ to that permutation group is trivial because
$U(M)=1$.  Thus every point of $U$ is fixed, and
the fixed-arrow obstruction applies.  The same
conclusion holds if all maps agree and the $M$ action
on $Y$ consists of permutations.

In the only remaining case all $\beta_x$ agree and
some actor on $Y$ is constant.  The action category
of $Y$ is connected, since that constant joins both
states to its value.  Collapse $X$ to a singleton,
obtaining an actual category $C_0$ and functor
$q:C\to C_0$.  Choosing one member of $X$ gives a
section $j$; all compositions are preserved because
the maps $\beta_x$ coincide.

Here is the exact second-cohomology comparison.
Relative to the section category, a normalized chain
uses one extra arrow of $X$, optionally followed by
an arrow of $Y$ and a string in $M$.  For each extra
arrow the relative chain complex is the augmentation
cone of the chains of the $Y$ action category, shifted
once.  Its degree-two homology is
$\widetilde H_0(Y_{hM};\mathbb Q)=0$.  Consequently
$H^2(BC;\mathbb Q)\to H^2(BC_0;\mathbb Q)$ is
injective, and the retraction makes it an isomorphism
with inverse $q^*$.  This argument does not assume
separate boundedness of any corner or action space.

Under a putative projective-plane hypothesis, $BC_0$
has bounded rational homology because it is an actual
retract.  Its first adjacent hom-set is singleton
with trivial source, so deleting the middle object
is an actual homotopy equivalence to an ordinary
act cone over $M$ with $|Z|\leq3$ states.  The
bounded order-eight act theorem kills all its
degree-two products.  The displayed isomorphism
then kills the square on $BC$, a contradiction.
The case $|Z|=1$ was already the fixed-arrow case.
\end{proof}


\subsection{The remaining star-shaped support reductions}

\begin{lemma}[Proved: coning the source of a binary bridge]
\label{CP2:A-star-coned-binary}
Let $S$ be a nonempty left-$N$, right-$M$ biset with at most
two states, for finite monoids $M,N$. Put $E=B(N,S,M)$.
Every rational degree-two product on
$\operatorname{hocofib}(E\to BN)$ is zero.
Consequently the same is true for a source star if one stalk
is a singleton with rationally acyclic leaf corner, and the
other stalk has at most two states.
\end{lemma}
\begin{proof}
A single state gives a section of $E\to BN$. Suppose there
are two states and inspect the right $M$ transformation image.
If it contains a constant map with value $p$, commutation
forces every $N$ actor to fix $p$. Using this state and the
identity $M$ string gives a section $BN\to E$. If the image
contains no constant and contains the transposition, every
commuting $N$ actor is a permutation. In this case
$E\to BN\times BM$ is the finite covering associated to the
permutation action on $S$. Transfer, followed by the injection
from the $BN$ factor, makes $H^2(BN;\mathbb Q)\to H^2(E;\mathbb Q)$
injective. In either case the cofiber exact sequence makes
every second class vanish on $BN$. The suspension quotient
and relative-cup argument of Lemma~\ref{CP2:A-two-block} kill
all its degree-two products.

The remaining right image is the identity alone. Then
$E=BM\times Z$, where $Z=B(N,S)$ is the ordinary action space.
Projection to $Z$ induces a map of cofibers to the ordinary
two-state $N$-act cone $Y$. Since $H^1(BM;\mathbb Q)=0$,
it induces an isomorphism on $H^1$ of the bridge spaces.
The two kernels from $H^2(BN;\mathbb Q)$ agree: projection
$BM\times Z\to Z$ has a section. The cofiber exact sequences
therefore give an isomorphism $H^2(Y;\mathbb Q)\to H^2(X;\mathbb Q)$.
The affine two-state lemma, Lemma~\ref{CP2:A-high-affine},
kills products on $Y$, and naturality kills them on $X$.

For the stated star, its singleton edge is $BM*BN_0$ and
is rationally acyclic when $BN_0$ is. Collapse that edge.
The quotient map is a rational cohomology-ring isomorphism,
and the resulting space is exactly the cofiber just treated.
\end{proof}

\begin{corollary}[Proved: a center of order at most four]
\label{CP2:A-star-small-center}
A three-object star with at most seventeen arrows and center
of order at most four cannot have type $\CP^2$.
\end{corollary}
\begin{proof}
By Proposition~\ref{CP2:A-star-large-leaf}, assume both
leaves have order at most seven. The small-center argument
in that proposition applies without its large-leaf
assumption. If the center classifying space is contractible,
the star is the wedge of its two full edges. A projective
plane target would force one of these edges to be an actual
projective-plane model, with at most fifteen arrows,
contradicting the two-object baseline. If the center has
nontrivial minimal group, that proof shows that its minimal
ideal is a group and both trivial rational modules are
projective. The actual orbit quotient at the center gives
a rationally equivalent star with trivial center. Its two
ordinary act-cone wedge summands have bounded rational
homology and corner orders at most seven. The bounded-nine
act theorem kills their degree-two products, and hence
those of the whole wedge.
\end{proof}

\begin{lemma}[Proved: two small leaves and star retracts]
\label{CP2:A-star-leaf-reductions}
The following reductions hold for a source star.
\begin{enumerate}
\item A leaf corner whose trivial right rational module is
projective can be replaced by the trivial corner, with its
stalk replaced by the actual left orbit set. This gives a
rational homology and cohomology-ring isomorphism.
\item Suppose a leaf corner has order at most four and its
stalk has at most two states. Either the center fixes a
state of that stalk, or that leaf can be replaced as in (1)
by a trivial corner, with at most two states.
\item If both leaves are trivial, identifying their two
vertices gives the ordinary center-act cone on the disjoint
union of the stalks. Its classifying space is homotopy
equivalent to the original star space wedged with a circle.
\item Suppose a star of rational homology and cohomology
ring $\CP^2$ retracts onto a subcategory, and its quotient
by that subcategory is a suspension. Then the retract also
has the rational homology and cohomology ring of $\CP^2$.
This applies both to an equivariant retraction of one
stalk, leaving the other stalk unchanged, and to merging
one trivial leaf into the other along an equivariant map
of their center acts.
\end{enumerate}
\end{lemma}
\begin{proof}
Part (1) is the actual orbit quotient of
Lemma~\ref{CP2:A-high-orbits}, applied at that leaf. To prove
(2), only the nonprojective-right case remains. The small
minimal-ideal description shows that the minimal group is
trivial and the ideal is a one-row rectangular band. Indeed,
two rows and two columns already require five elements,
while one column makes the trivial right module projective;
a nontrivial minimal group with two columns requires at
least five elements. The left actors of this one-row band
are retractions with a common image $U$. On a set of at
most two states, either $|U|=1$ or all these retractions
are identities. In the first case the sole point of $U$
is fixed by the commuting center action. In the second,
every actor of the whole leaf monoid is the identity:
left multiplication fixes every element of a one-row band,
so its action fixes the image $U=S$. Its classifying space
is contractible by the small minimal-ideal lemma. The
quotient replacing it by a trivial corner is then a
homotopy equivalence in the double bar model.

For (3), the quotient identifying the leaf vertices is
the nerve realization of the ordinary act cone on the
disjoint union. Identifying two distinct points in a
connected CW complex has homotopy type $X\vee S^1$.
Equivalently, first attach an interval between the points
and then collapse that contractible interval.

For (4), if the second cohomology of the retract were zero,
all second classes would lift from the suspension quotient,
and their products would vanish. Thus the retract has
nonzero second cohomology. Its homology is a direct summand
of that of the original space, so it has one-dimensional
second homology and no other positive homology except
possibly in degree four. Pullback in second cohomology is
an isomorphism. The nonzero square forces a nonzero fourth
group in the retract, hence its entire rational target
pattern and ring. For a retraction of one stalk the quotient
is the suspension of the quotient of the corresponding
two-sided action spaces; all corner spaces and the other
edge are unchanged. For merging two trivial leaves the
retract is the retained edge, and the quotient is the
suspension of the removed stalk's ordinary action space.
These are literal quotients of the mapping-cylinder models.
\end{proof}

\begin{proposition}[Proved: a large center in a small star]
\label{CP2:A-star-large-center}
A connected three-object category with star-shaped acyclic
support, at most seventeen arrows, and center corner of
order at least eight cannot have type $\CP^2$.
\end{proposition}
\begin{proof}
Take opposition so the center is a source. Write its
order as $a$, its leaf orders as $b,c$, and its stalk sizes
as $s,t$. Thus $a+b+c+s+t\le17$. Large leaves are already
excluded by Proposition~\ref{CP2:A-star-large-leaf}; assume
$b,c\le7$.

Suppose first $a\ge10$. Both leaves have order at most
four, hence are rationally acyclic. A singleton stalk gives
a center-fixed state, and Corollary~\ref{CP2:A-fixed-arrow}
in its dual form kills every degree-two product. We may
assume $s,t\ge2$. The only possibilities are $a=11$ with
$b=c=1$ and $s=t=2$, or $a=10$ with leaf orders $(1,1)$
and stalk sizes at most $(2,3)$, or leaf orders $(1,2)$
and stalk sizes $(2,2)$, up to interchange. Orbit comparison
at the two-element leaf reduces the latter case to trivial
leaves. If an orbit set becomes a singleton, the same fixed
state argument applies. Otherwise the reduced star has the
rational homology and ring of $\CP^2$ and two trivial leaves,
one on two states and the other on at most three.

Merge the leaves as in Lemma~\ref{CP2:A-star-leaf-reductions}(3).
The resulting ordinary act cone has bounded rational
homology and the same positive-degree products in degrees
at least two. Relative cups require $H^2(BM;\mathbb Q)\ne0$;
the two-row obstruction requires at least three minimal
rows, and there are at least two columns. The cardinality
budget makes the minimal group trivial. If there are at
least three columns, the ideal is $3\times3$ and there
are at most two outside elements. Lemma~\ref{CP2:A-two-outside-acts}
then makes fourth homology of this bounded cone zero,
a contradiction.

Thus the minimal band has two columns. A band-fixed state
is center-fixed: if $xk=x$ for every $k$ in the ideal, then
$xm=xkm=x$ for every $m$. This again kills all products.
In its absence every core orbit has two points. On every
such orbit the action of $M$ is the same labelled column
action, because right multiplication sends $e_{ij}$ to
$e_{i,\rho_m(j)}$ independently of $i$. The two-state
stalk is one such orbit and therefore embeds equivariantly
into a core orbit of the other stalk. Merge its leaf into
the other leaf along this embedding. Part (4) of the
preceding lemma makes the retained ordinary act cone a
rational $\CP^2$ target. If $a=11$, it has two states and
the affine lemma excludes it. If $a=10$, it has at most
three states and Proposition~\ref{CP2:A-ten-six-state-rational}
excludes it. This finishes $a\ge10$.

Now let $a=8$ or $9$. If a leaf has order at least five,
the other has order at most two and $s+t\le3$. If the
large leaf's stalk is a singleton, that edge is a join
with zero second cohomology. It kills restrictions to
the center and large leaf; the remaining leaf is rationally
acyclic. The two-block relative-cup lemma applies. If the
small leaf has the singleton stalk, the other has at most
two states, and Lemma~\ref{CP2:A-star-coned-binary} applies.

We are left with $b,c\le4$. A singleton stalk again gives
a fixed state and zero products. Assume $s,t\ge2$. If
both leaf orders are at most two, orbit comparison makes
them trivial and Lemma~\ref{CP2:A-star-leaf-reductions}(3)
reduces to the bounded-nine act theorem. Every leaf on
two states can be treated in the same way by part (2),
since its fixed-state alternative kills products. The
only budget case not yet covered has
\[
 (a,b,c,s,t)=(8,3,1,3,2)
\]
up to interchange of the leaves. If the three-element
leaf has projective trivial right module, its orbit
quotient once again gives the bounded-nine act case.
Otherwise it is $R_0=\{1,p,q\}$, where $uv=v$ on the
two nonidentity elements.

Its two actors $P,Q$ on the three-state stalk are
retractions with the same image $U$. If $P=Q$, their
common retraction is equivariant for both corner actions.
If $|U|=1$, its point is center-fixed. Otherwise retract
this stalk to $U$. Part (4) transfers the rational target
to the retracted star. The $R_0$ action on $U$ is trivial
and $BR_0$ is contractible, so quotienting that leaf to
the trivial corner gives a rational equivalence (indeed,
a homotopy equivalence). The bounded-nine act theorem
after merging leaves again excludes the result.

If $P\ne Q$, necessarily $|U|=2$. The natural graph model
of the $R_0$ action category in
Lemma~\ref{CP2:A-high-P5-graph} has the two vertices $U$
and one edge from $P(x)$ to $Q(x)$ for the third state
$x$. It is an interval; its construction commutes with
the center action. Its homotopy-orbit space over $M$
therefore maps by a homotopy equivalence to $BM$.
Quotient the entire stalk to one state and the leaf to
the trivial monoid. The map on bridge spaces is this
equivalence, on the leaf spaces it is $BR_0\to *$, and
on the center it is the identity. Homotopy-pushout
invariance gives a homotopy equivalence of stars. Its
singleton stalk now supplies a center-fixed state, so
its degree-two products vanish. All cases are excluded.
\end{proof}


\begin{proposition}[Proved: the complete star cardinality reduction]
\label{CP2:A-star-seventeen-scope}
To exclude all three-object categories with star-shaped acyclic
support and at most seventeen arrows, it suffices to exclude
the $38$ exact seventeen-arrow profiles in
\path{converse/scope17/star17_profiles.tsv}.
Each has center order between five and seven, and both leaf
orders at most seven. No profile outside this explicit list
requires a composition-table search.
\end{proposition}
\begin{proof}
Identity padding, Proposition~\ref{CP2:A-identity-padding},
reduces to exactly seventeen arrows without changing the
support or classifying space. Taking opposition makes the
center a source. A profile is a positive tuple
$(a,b,c,s,t)$ with sum seventeen, in which $(b,s)$ and
$(c,t)$ may be interchanged. There are
$\binom{16}{4}=1820$ labelled tuples. The leaf-swap fixed
tuples satisfy $a+2b+2s=17$; their number is
$\sum_{d=2}^{8}(d-1)=28$. Thus there are exactly
$(1820+28)/2=924$ orbits. The fresh program
\path{converse/scope17/star17_scope.py} retains the representative
with $(b,s)\le(c,t)$ in lexicographic order and tests only
the following proved sufficient conditions.

The large-leaf, large-center and small-center results above
remove every corner configuration outside center orders
five through seven and leaf orders at most seven.
For completeness the program also records the overlapping
elementary cuts: all three corner second cohomology groups
zero; a center-fixed singleton stalk when both leaves have
zero second cohomology; a full edge with zero second
cohomology and a remaining leaf with the same property;
the coned-binary lemma; and the small-corner binary-stalk
lemma. The sufficient edge tests for zero second cohomology
are a singleton stalk, total edge size at most five, or
both corner orders at most two. They follow respectively
from the join model, the second-homology cost lemma, and
semisimplicity of the two corner algebras.

There are two further transfer tests. If the center order
is at most five, its rational homology above degree two
vanishes. The star is the union of its two full edges
along that corner. Mayer--Vietoris transfers boundedness
to both edges and makes restriction in fourth cohomology
injective. It is therefore enough that both edges have
zero second products. The program uses only: stalk size
at most two; both corner orders at most four; total edge
size at most eleven; a corner of order at most two opposite
a corner of order at most nine; or a corner of order at
most three opposite a corner of order at most six.
The first three are the binary, two-block and small-edge
lemmas. The fourth is exact orbit comparison followed by
the bounded-nine act theorem. For the fifth, the only
nonprojective three-element endpoint is the relevant
orientation of the identity-adjoined right-zero band;
its natural graph model and
Corollary~\ref{CP2:A-baseline-right-zero-three-six} apply.
All uses of boundedness occur after this Mayer--Vietoris
transfer.

Finally, if both leaves have orders at most four, each
nonprojective leaf is allowed only when its stalk has at
most two states. Lemma~\ref{CP2:A-star-leaf-reductions}
then either gives a center-fixed state or reduces both
leaves to trivial corners. Merging the leaves gives a
bounded ordinary act cone, excluded by the bounded-nine
theorem when the center order is at most nine.

The exact enumeration leaves $38$ profiles and removes
$886$. The full reason ledger, including every overlapping
predicate, is archived as
\path{outputs/star17/cardinality_scope.json}; its source
and ledger digests are in
\path{converse/scope17/star17_scope_summary.json}.
The maximum number of normalized two-simplices among the
residual profiles is $88$. The cardinality count and all
reduction predicates received a separate hostile proof
check. This proposition establishes the complete scope
reduction; it does not treat an unfinished composition
stream as a completed exclusion.
\end{proof}

\paragraph{Composition enumeration for the residual stars.}
The executable \path{converse/scope17/star17_canonical.cpp}
uses the action-orbit procedure proved in
Proposition~\ref{CP2:A-linear-biset-enumeration}. For a source
star there are exactly two independent connecting bisets,
and no composition of two cross arrows. Thus a pair of
canonical bisets supplies the complete category table.
If their state stabilizer orders are $A_S,A_T$, the
weight $s!t!/(A_SA_T)$ restores every labelled table.
All ordinary corner tables through order seven are used.

Only three types of corner-table cut precede this
enumeration. A projective trivial left module at the
center permits an actual orbit quotient to a trivial
center; the result is a wedge of bounded ordinary leaf
act cones of order at most seven. Projective trivial
right modules at both leaves permit their orbit quotients
to trivial leaves; merging them gives a bounded ordinary
center-act cone of order at most seven. The bounded-nine
theorem excludes both cases. The remaining cuts require
zero second cohomology on both blocks of a two-block
decomposition. Corner second-homology vanishing is checked
by a literal binary bar calculation, a sufficient condition
for rational vanishing by the universal coefficient theorem.
An edge is declared second-homology zero only for a singleton
stalk, two semisimple corner algebras, or total size at most
five. These tests received a separate targeted proof check.

The unchanged literal-certificate procedure checks the
necessary binary first and second homology pattern, then
produces either a square primitive or a positive third
homology calculation. The independent checker
\path{converse/scope17/check_star17.cpp} rechecks every
retained typed associative table and its bar certificate.
Only a closed, hashed stream with matching producer and
checker completion counts can certify its profile.


\begin{lemma}[Proved: an actual compression of a singleton-long triangle]
\label{CP2:A-singleton-triangle-monoid}
Let a full-linear three-object category have trivial endpoint
corners, middle monoid $M$, and hom-set sizes $(s,t,1)$.
Let $V=C(0,1)$ and $W=C(1,2)$. Its classifying space is
homotopy equivalent to that of the monoid
\[
 P=M\sqcup(V\times W),\qquad
 (v,w)(v',w')=(v,w'),\quad
 m(v,w)=(mv,w),\quad(v,w)m=(v,wm).
\]
In particular $|P|=|M|+st$, and its minimal ideal is the
rectangular band $V\times W$.
\end{lemma}
\begin{proof}
The given left and right actions, and the displayed band
rule, check every associative product; the identity of
$M$ is the identity of $P$. Fix $v_0\in V,w_0\in W$ and
put $e=(v_0,w_0)$. In the Karoubi envelope of $P$, retain
the two objects $x=1$ and $e$. Call the resulting category
$D$. Then
\[
 D(x,x)=P,\quad D(e,e)=\{e\},\quad
 D(e,x)=V\times\{w_0\},\quad
 D(x,e)=\{v_0\}\times W.
\]
The object $e$ is a retract of $x$, so deleting it gives
$BD\simeq BP$ by the object-retract lemma.

Define $F:C\to D$ by sending both endpoints to $e$ and
the middle object to $x$. Send $M$ by its inclusion,
$v\mapsto(v,w_0)$, $w\mapsto(v_0,w)$, and the unique
long arrow to $e$. The only two-cross-arrow equation is
$(v_0,w)(v,w_0)=e$, so $F$ is an actual functor.

For Quillen's Theorem~A consider $d\downarrow F$.
For $d=e$, its object at the initial endpoint is initial:
its arrows to the middle objects are exactly the states
$V$, and its arrow to the final endpoint is the unique
long arrow. This comma category is contractible.

For $d=x$, its objects over the initial endpoint are
indexed by $W$, those over the middle by $P$, and those
over the final endpoint by $W$. The full subcategory
on the middle $M$-objects and the final $W$-objects is
contractible, with initial object $(1,1_M)$. For each
$w\in W$, another full subcategory consists of the initial
object indexed by $w$, the middle objects $(v,w)$ for
$v\in V$, and the final object indexed by $w$. It has
the initial object indexed by $w$, hence is contractible.
It meets the first subcategory in exactly its final
object and meets no other such subcategory. Every nerve
simplex lies in one of these subcategories: multiplication
by $M$ preserves $M$ and each band column, and there is
no arrow out of a final object. Thus the comma nerve is
the union of a contractible CW complex and contractible
CW complexes attached individually at points. It is
contractible. Quillen's Theorem~A proves $BC\simeq BD$.

Finally every nonempty ideal of $P$ meets $V\times W$,
and multiplying any one of its elements on both sides
by band elements produces the entire band. This proves
the assertion about the minimal ideal and the cardinality.
\end{proof}

\begin{corollary}[Proved: the order-eight middle corner at seventeen]
\label{CP2:A-eight-middle-three-three}
The profile $(1,8,1;3,3,1)$ cannot have type $\CP^2$.
More generally the preceding triangle is excluded whenever
$|M|+st\le22$.
\end{corollary}
\begin{proof}
The actual compression would be a monoid model with at
most twenty-two elements, contrary to
Theorem~\ref{A:monoid-minimum23-classification}. In the
displayed profile its order is $8+3\cdot3=17$.
\end{proof}


\begin{theorem}[Proved: every full-linear three-object profile through sixteen]
\label{CP2:A-linear-sixteen-complete}
A three-object category with full-linear acyclic support
and at most sixteen arrows cannot have type $\CP^2$.
\end{theorem}
\begin{proof}
Identity padding reduces to exactly sixteen arrows without
changing support or homotopy type. Write the profile as
$(a,b,c;x,y,z)$, with all six entries positive. There are
$\binom{15}{5}=3003$ ordered tuples with sum sixteen.
Opposition sends this tuple to $(c,b,a;y,x,z)$ and fixes
$91$ tuples, so the exact scope has
$(3003+91)/2=1547$ opposition orbits.

The independently regenerated profile join has the disjoint
partition
\[
\begin{array}{l|r}
\text{proved small-corner profile reductions}&1399\\
\text{a corner of order at least eight}&88\\
\text{identity padding to a complete seventeen-arrow search}&54\\
\text{complete direct sixteen-arrow searches}&6\\\hline
\text{total}&1547.
\end{array}
\]
The large-corner cases are
Proposition~\ref{CP2:A-baseline-linear-large}. Every first-row
entry has an explicit direct proof label, or a specified
one-identity padding to a proved profile, in the delivered
join. Its predicates use the two-block and fixed-arrow
lemmas, binary stalks and binary projective covers, the
small middle Tor obstruction, the semisimple-pair graph
comparison, singleton-long and projective-singleton
reductions, the two-small-incoming lemma, and the actual
triangle compression. All are stated and proved above;
no inherited omitted-profile label serves as a premise.

For the last two rows, the join selects one complete
receipt for each profile or a stated one-identity padding.
There are $55$ distinct selected receipts, because one
seventeen-arrow profile can cover several sixteen-arrow
profiles. They contain $6535831$ independently checked
literal certificates, covering $2295552594$ labelled or
orbit-weighted tables under their respective proved
corner cuts. These totals are verification workload,
not category-isomorphism counts. Every retained table
has either a literal normalized primitive for the square
of its nonzero binary second class, or positive binary
third homology. Both contradict the projective-plane
target by Proposition~\ref{CP2:A-cup-certificate}.

The final join program
\path{converse/scope17/finalize_linear16.py} checks the
exact profile index and input digest of every selected
receipt, all its generator/checker count equalities,
every compressed certificate digest, and the producer
and independent-checker source digests. It rejects a
missing profile, incomplete stream, mismatched digest,
or unexcluded record. Its complete output is
\path{outputs/linear16_canonical/final_scope_join.json};
the compact receipt is
\path{converse/scope17/linear16_complete_summary.json}.
The join passes with no pending profile and no survivor.
Thus every exact sixteen-arrow profile, and by padding
every smaller full-linear profile, is excluded.
\end{proof}


\begin{theorem}[Proved: every three-object star through sixteen]
\label{CP2:A-star-sixteen-complete}
A three-object category with star-shaped acyclic support
and at most sixteen arrows cannot have type $\CP^2$.
\end{theorem}
\begin{proof}
Identity padding reduces to exactly sixteen arrows while
preserving the support and classifying space. Opposition
makes the center a source. Write its profile as
$(a,b,c,s,t)$, where the last two entries are the two
nonempty connecting hom-set sizes. There are
$\binom{15}{4}=1365$ positive tuples. The leaf swap
$(a,b,c,s,t)\mapsto(a,c,b,t,s)$ fixes $21$ tuples:
these satisfy $a+2b+2s=16$, with $b+s$ ranging from two
through seven. Thus there are exactly $693$ profile orbits.

The fresh scope join applies precisely the proved
predicates of Proposition~\ref{CP2:A-star-seventeen-scope}
at the actual sixteen-arrow budget. It gives the disjoint
partition
\[
\begin{array}{l|r}
\text{direct proved reductions}&679\\
\text{one identity padding to a proved profile}&1\\
\text{one identity padding to a complete search}&13\\\hline
\text{total}&693.
\end{array}
\]
The last row uses twelve distinct complete seventeen-arrow
receipts: indices $19$--$22$, $24$, and $28$--$34$
of \path{converse/scope17/star17_profiles.tsv}.
Their complete streams contain $533076$ independently
checked literal certificates and cover $28163325$
orbit-weighted tables after the proved corner cuts.
These figures count verification work, rather than
category-isomorphism classes. Every retained table has
either zero square on its nonzero binary second class
or positive binary third homology, contradicting
Proposition~\ref{CP2:A-cup-certificate}.

The join program \path{converse/scope17/finalize_star16.py}
checks every selected profile index, the frozen input
digest, producer and independent-checker source digests,
compressed certificate digests, and all producer/checker
completion counts. It also checks that every padding
changes only one corner order by one. All $693$ profiles
are covered, with no incomplete stream and no survivor.
The full join is
\path{outputs/star17_canonical/star16_final_scope_join.json};
its digest and compact totals are in
\path{converse/scope17/star16_complete_summary.json}.
Padding then excludes every smaller star profile.
\end{proof}


\begin{theorem}[Proved: every three-object star through seventeen]
\label{CP2:A-star-seventeen-complete}
No three-object category with star-shaped acyclic support
and at most seventeen arrows has classifying space of type
$\CP^2$.
\end{theorem}
\begin{proof}
Proposition~\ref{CP2:A-star-seventeen-scope} supplies the
complete $924$-profile scope: $886$ profiles meet explicit
proved reductions and $38$ require the stated finite searches.
The canonical action enumeration described there is now
complete for all $38$ frozen indices. Its totals are
\[
\begin{array}{l|r}
\text{labelled tables restored by action-orbit weights}&113338106\\
\text{retained literal certificates}&1564645\\
\text{explicit zero-square certificates}&1538408\\
\text{positive third-homology certificates}&26237\\
\text{unexcluded records}&0.
\end{array}
\]
The last two exclusion rows sum to the retained-record
count. Every complete stream passes the independent
literal checker \path{converse/scope17/check_star17.cpp}.
The table counts are verification workload, not counts
of category-isomorphism classes. The exact generic-DFS
versus canonical-action comparisons for indices1 and2
also agree on the complete retained orbit sets, as
recorded in \path{converse/scope17/star_canonical_comparison.json}.

The finalizer \path{converse/scope17/finalize_star17.py}
regenerates the positive-tuple scope by cut positions,
independently of the nested loops of the original scope
program. It verifies all $1820$ labelled tuples, the $28$
leaf-swap fixed tuples, and the resulting $924$ orbits.
It matches the complete frozen reason ledger and every
one of the $38$ residual matrices. It then verifies each
closed receipt's index, range, input and source digests,
compressed certificate digest, independent outcome counts,
and absence of a survivor. All these checks pass.
The complete join is
\path{outputs/star17_canonical/star17_final_scope_join.json},
with digest and compact totals in
\path{converse/scope17/star17_complete_summary.json}.
The literal obstructions exclude the projective-plane
target by Proposition~\ref{CP2:A-cup-certificate}.
Identity padding handles every smaller arrow count.
\end{proof}


\subsection{Affine act reductions at the eighteen-arrow frontier}

The following proofs reconstruct the affine reductions at their
points of use. They require no inherited finite-search output.

\begin{lemma}[Proved: a common rational affine row origin]
\label{CP2:A-act-common-affine-origin}
Let $M$ have rectangular minimal ideal $K=I\times\{0,1\}$.
Write $\sigma_m$ and $\tau_m$ for its row and column maps.
Suppose there is $p\in\mathbb QI$, of augmentation one,
fixed by every $\sigma_m$ for which $\tau_m$ is nonconstant.
For every finite nonempty right $M$-act $S$, all products
of rational degree-two classes of $B\mathcal C(M,S)$ vanish.
Neither bounded homology nor a bound on $|S|$ is required.
\end{lemma}
\begin{proof}
First take the actual constant-action quotient of
Lemma~\ref{CP2:A-constant-quotient-H2}. Its degree-two
cohomology surjects onto that of the original cone, so it
suffices to treat constant $K$ action. Compress this cone
by Theorem~\ref{CP2:A-constant-act-compression} to the
monoid $N$ with minimal band $I\times S$. Write $s_j$
for the constant value of an old column $j$; these values
may coincide. They satisfy $s_jm=s_{\tau_m(j)}$.

Choose $i_0\in I$ and use base column $s_0$. Put
$V=\ker(\mathbb QI\to\mathbb Q)$ and
$W=\ker(\mathbb QS\to\mathbb Q)$. The band--Tor
formula identifies $H^2(BN;\mathbb Q)$ with all balanced
forms $L:W\times V\to\mathbb Q$. For any element $f$
of $N$, define
\[
 a_f=[\sigma_f(i_0)]-[i_0],\qquad
 b_f^L(v)=L([\tau_f(s_0)]-[s_0],v),
\]
and let $A_f$ be its row augmentation operator. Balance
and the row/column composition laws give
\[
 a_{fg}=a_f+A_fa_g,\qquad
 b_{fg}^L=b_g^L+b_f^LA_g.
\]
Thus $u_L(f,g)=b_f^La_g$ is the cocycle realizing $L$,
as checked on the band cycles in Lemma~\ref{A:rank-one-square}.

This cocycle transfers through actual functors as follows.
Set $e=(i_0,s_0)$ and use the full Karoubi subcategory
on $1,e$. Extend $u_L$ by the same formula on its
underlying $N$ labels. Both object identities have
$a=b^L=0$, so the extension is normalized. The displayed
identities prove its cocycle equation on every typed pair.
The inclusion of object $1$ is a classifying-space
equivalence, and the compression functor is an equivalence
on classifying spaces. Therefore these pulled-back
cocycles represent all degree-two classes of the act cone.
The compression functor sends a cross state $s$ to
$(i_0,s)$; on that arrow $a_s=A_s=0$ and
$b_s^L=L([s]-[s_0],-)$.

Set $t=p-[i_0]\in V$. Change the row translations on
old monoid arrows to
\[
 a'_m=a_m+(A_m-1)t=\sigma_mp-p.
\]
Set $a'_s=0$ on cross arrows and zero on the sink identity.
These are typed affine coordinates: on any normalized
composable pair, the second arrow is an old monoid arrow
unless the first is the sink identity. Hence the same
affine equation $a'_{fg}=a'_f+A_fa'_g$ holds.
The new cocycle $u'_L(f,g)=b_f^La'_g$ represents the
same class, since
\[
 u'_L-u_L=-\delta h_L,\qquad h_L(f)=b_f^Lt.
\]
For instance the coboundary on a cross/old pair is
$b_f^L(1-A_g)t$, by the displayed identity for $b_{fg}^L$.
All identity cases give zero.

Define, only on the old monoid,
\[
\begin{aligned}
 \chi(m)&=\tau_m(1)-\tau_m(0),&
 \beta(m)&=\tau_m(0),\\
 \ell_L(v)&=L([s_1]-[s_0],v),&
 \alpha_L(m)&=\ell_L(a'_m).
\end{aligned}
\]
Equivariance of the old column line and balance imply
$\ell_LA_m=\chi(m)\ell_L$. Consequently
\[
\begin{aligned}
 \beta(mn)&=\beta(n)+\chi(n)\beta(m),\\
 \alpha_L(mn)&=\alpha_L(m)+\chi(m)\alpha_L(n),\\
 b_m^L&=\beta(m)\ell_L.
\end{aligned}
\]
The common-origin hypothesis is precisely
$\chi(m)a'_m=0$ for every old actor.

Use algebra-ordered normalized chains: a cross arrow,
when present, is first, and every later arrow is old.
For two balanced forms $L,D$, the normalized three-cochain
\[
 P(f,g,h)=u'_L(f,g)\,\beta(g)\,\alpha_D(h)
\]
is therefore defined on every nondegenerate typed triple.
It is zero on degenerate triples. Expanding its five-term
coboundary and substituting the preceding identities gives
\begin{align*}
 (\delta P)(f,g,h,k)
 &=u'_L(f,g)\beta(h)\alpha_D(k)\\
 &\quad+\chi(h)\beta(g)\alpha_D(k)
       \bigl(u'_L(f,gh)-u'_L(f,g)\bigr)\\
 &=u'_L(f,g)u'_D(h,k)
   +\chi(h)\beta(g)\alpha_D(k)b_f^LA_ga'_h\\
 &=u'_L(f,g)u'_D(h,k).
\end{align*}
The last equality uses $\chi(h)a'_h=0$. Thus every
mixed product has a literal rational primitive. Pullback
through the constant-action quotient proves the assertion.
\end{proof}

\begin{lemma}[Proved: projective row dimension at most one]
\label{CP2:A-act-row-pd-one}
Let $M$ have rectangular minimal ideal $I\times J$, put
$A=\mathbb QM/\mathbb QK$ and
$V=\ker(\mathbb QI\to\mathbb Q)$.
If $\operatorname{pd}_A V\leq1$, all rational degree-two
products of every finite right-act cone vanish.
\end{lemma}
\begin{proof}
Take the constant-action quotient and its compression $N$.
The compressed outside algebra is literally $A$, and
its row module is $V$. If $\bar W$ is its column
augmentation module, Lemma~\ref{A:band-Tor} gives
$H_4(BN;\mathbb Q)=\operatorname{Tor}_2^A(\bar W,V)=0$.
The degree-two quotient surjection and naturality of
products give the result, with no boundedness inference.
\end{proof}

\begin{corollary}[Proved: local two-column contracted algebras]
\label{CP2:A-local-two-column-acts}
If $M$ has rectangular minimal ideal with two columns
and its contracted rational algebra is local, all
rational degree-two products of every finite right-act
cone vanish, at any row and state counts.
\end{corollary}
\begin{proof}
A high power of an outside basis element in the pointed
quotient monoid $M/K$ is zero or an idempotent basis
element. A local algebra has only the idempotents zero
and one. In the second case this element is a finite
unit. A nontrivial cyclic group of such units has an
average which is an idempotent different from zero and
one, since its group elements are distinct basis vectors.
Thus every outside element other than the identity is
nilpotent in $A$. The column augmentation character
$\chi:A\to\mathbb Q$ vanishes on them, so their column
maps are constant. Every band actor also has constant
columns. The identity is the only remaining actor, and
it fixes every row. Apply the common-origin lemma.
\end{proof}

\begin{proposition}[Proved: three rows and five outside elements]
\label{CP2:A-three-row-five-outside-acts}
If the rectangular minimal ideal of $M$ has three rows,
two columns, and at most five outside elements, every
finite right-act cone has zero rational degree-two products.
\end{proposition}
\begin{proof}
If $H^2(BM;\mathbb Q)=0$, the suspension quotient of
the monoid vertex already kills all such products.
Otherwise the balanced description of $H^2(BM)$ gives
a nonzero covector $\ell$ on the two-dimensional row
augmentation module satisfying
$\ell A_m=\chi(m)\ell$. Regard its three row values
as a finite rational set, up to a common translation.

If the values are distinct, a slope-one row actor fixes
every row: a nonzero translation cannot preserve a
finite set. Any slope-minus-one actor is its unique
reflection and fixes the middle row. A common row
origin therefore exists. Otherwise translate and rescale
the row values to $(0,0,1)$. A nonconstant column actor
with slope one fixes row2 and preserves the pair of rows0,1.
An actor of slope minus one has one of the two row maps
\[
 r_0=(2,2,0),\qquad r_1=(2,2,1),
\]
and its column map is the swap. If neither occurs, row2
is a common origin. If only $r_0$ occurs, every slope-one
actor fixes row0: composing on the left with $r_0$
would otherwise produce $r_1$. The affine midpoint of
rows0 and2 is then fixed by every nonconstant-column
actor. The other one-flip case is symmetric.
All these cases have zero products by the first lemma.

A possible nonzero product therefore forces both flips,
their squares
$p_0=(0,0,2)$ and $p_1=(1,1,2)$ with identity columns,
and the actual identity. These five distinct profiles
are all outside the minimal ideal. The outside budget
is therefore exactly five, with one element of each
profile. Products among them stay outside, since column
maps remain bijective. Their row and column profiles
force the multiplication
\[
 (i,\epsilon)(j,\delta)=(i,\epsilon+\delta),
 \qquad (i,0)=p_i,\quad(i,1)=r_i,
 \quad\epsilon,\delta\in\mathbb Z/2.
\]
The contracted algebra is consequently the unitalization
of $\mathbb Q[L_2\times C_2]$.

Put $e_-=(p_0-r_0)/2$ and
$n_-=(p_1-p_0-r_1+r_0)/2$. Direct multiplication gives
$e_-^2=e_-$ and $Ae_-=\mathbb Qe_-\oplus\mathbb Qn_-$.
In the row basis $u=[1]-[0]$, $v=[2]-[0]$ we have
$e_-v=v$ and $n_-v=-u$. The $A$-linear map
$Ae_-\to V$, $ae_-\mapsto av$, is therefore an
isomorphism. Thus $V$ is projective and the preceding
row-dimension lemma excludes the last possibility.
This proves the assertion without a source catalog.
\end{proof}

\begin{proposition}[Proved: two columns and three outside elements]
\label{CP2:A-three-outside-two-column-acts}
If the rectangular minimal ideal has two columns and at
most three outside elements, every finite right-act cone
has zero rational degree-two products, at any row count.
\end{proposition}
\begin{proof}
Let $P$ be the outside actors with nonconstant column
map. Their slopes multiply in $\{1,-1\}$, so $P$ is
an actual monoid containing the identity, of order at
most three. If $P$ is a group, average a row orbit.
If its minimal ideal is a singleton zero $z$, any row
in the image of $\sigma_z$ is fixed by $P$. The only
other possible proper minimal ideal has two elements:
it is a group of order two, a left-zero band, or a
right-zero band. A group ideal has a $P$-invariant
affine average, since the only outside element is the
identity. For a right-zero ideal, $xe=e$ shows that
every row in the image of $\sigma_e$ is fixed by $P$.
The common-origin lemma applies in all these cases.

The remaining monoid is the identity-adjoined left-zero
two-band. It uses all three outside elements, so
$A=\mathbb QP$. If its two left zeros are $p,q$, the
orthogonal idempotents $p,1-p$ and the vector $q-p$
give the three-dimensional one-arrow path algebra:
the radical squares to zero and lies in $(1-p)Ap$.
The one-arrow path algebra is hereditary. For example
its representation consists of a single linear map;
splitting that map into kernel, image and cokernel
decomposes every module into the two simples and the
single arrow module, whose projective resolutions have
length at most one. Hence $\operatorname{pd}_A V\leq1$,
and the row-dimension lemma applies.
\end{proof}

\begin{theorem}[Proved: bounded act cones through order eleven]
\label{CP2:A-bounded-eleven-acts}
For any finite monoid $M$ of order at most eleven and
any finite nonempty right act $S$, bounded rational
homology of $B\mathcal C(M,S)$ forces all rational
degree-two products to vanish. There is no bound on
the state count or the second Betti number. The left-act
version and the exact-projective-opposite-corner version
follow by opposition and actual orbit comparison.
\end{theorem}
\begin{proof}
If the monoid vertex has zero rational second cohomology,
use its suspension quotient. At most two minimal rows
are excluded by Lemma~\ref{CP2:A-high-two-row-act}.
Otherwise there are at least three rows and two columns.
A nontrivial maximal group would require at least
$3\cdot2\cdot2+1=13$ elements, so the ideal is a band.
Writing $q=|M\setminus K|$, the complete size list is
\[
\begin{gathered}
 (r,c,q)=(3,2,q\leq5),\quad (4,2,q\leq3),\\
 (5,2,q\leq1),\quad (3,3,q\leq2).
\end{gathered}
\]
The first case is the three-row proposition. The next
two are the three-outside proposition. In the last case
Lemma~\ref{CP2:A-two-outside-acts}, applied to the reduced
augmentation module of the actual act, gives zero fourth
rational homology from boundedness. This proves the theorem.
For an opposite projective endpoint, its exact coinvariant
functor is the actual finite orbit quotient and preserves
the rational homology and cohomology ring, including
boundedness. The same conclusion therefore transfers.
\end{proof}

\begin{corollary}[Proved: projective initial singleton through order eleven]
\label{CP2:A-eleven-projective-singleton}
In a full-linear three-object category, suppose the initial
corner has projective trivial left rational module, the
initial adjacent right act has one orbit, and the final
corner has order at most eleven. Bounded rational homology
forces all rational degree-two products to vanish.
The opposite assertion holds at the final endpoint.
\end{corollary}
\begin{proof}
The actual exact-orbit quotient makes the initial corner
trivial and its adjacent hom-set singleton. Its singleton
comma category deletes the middle object by an actual
homotopy equivalence. The remaining ordinary act cone
over the final corner has bounded rational homology.
Apply the preceding theorem and pull back its product
vanishing along the specified comparisons.
\end{proof}


\subsection{The last ordinary act profile at eighteen arrows}

\begin{proposition}[Proved: six outside elements over three rows]
\label{CP2:A-three-row-six-outside-acts}
If a finite monoid has minimal band $3\times2$ and
at most six outside elements, every finite right-act
cone has zero rational degree-two products.
\end{proposition}
\begin{proof}
The case of at most five outside elements is
Proposition~\ref{CP2:A-three-row-five-outside-acts}.
Its preceding affine argument applies at any outside
count: a possible nonzero product forces row values
$(0,0,1)$ and the four distinct outside profiles
\[
 \begin{aligned}
 p_0&=((0,0,2),(0,1)),& p_1&=((1,1,2),(0,1)),\\
 r_0&=((2,2,0),(1,0)),& r_1&=((2,2,1),(1,0)).
 \end{aligned}
\]
Together with the identity these occupy five of the
six outside labels. The sixth has one of exactly $24$
balanced profiles: four row maps over identity columns,
two over swapped columns, and nine over each of the
two constant column maps. For a constant column map,
all three row images must lie in one of the two value
levels, giving $2^3+1=9$ choices.

We give a complete small literal calculation for this
last outside label. Fix identity0, the four forced labels
1 through4, the additional label5, and six row-major
band labels6 through11. All products involving the
identity or band are forced. For every other ordered
pair, its allowable outputs are exactly all twelve
labels with the composite row and column profile.
The new producer \path{converse/scope18/forced_three_row12.py}
tries every such output, propagating only required
associativity equalities. Every monoid under consideration
is obtained after the indicated relabelling: no ordinary
isomorphism or opposition identification removes a table.
There are exactly $27$ associative tables, reached in
$44$ search nodes across the $24$ extra profiles.

The independent program
\path{converse/scope18/check_forced_three_row12.py}
does not import that producer. It reconstructs allowable
products from associativity with the six band elements,
assigns the undetermined cells in reverse fixed order,
and never propagates an unknown outside product. It
checks the entire literal set, with $258$ search nodes
and $300612$ evaluated associativity equations, and
finds exactly the same $27$ tables. It also independently
checks their identities, full associativity and exact
minimal ideals.

Every table has a rational row-projector certificate.
Put $A=\mathbb QM/\mathbb QK$ and let $V$ be its
two-dimensional row augmentation module. The certificate
specifies a literal pair $e,g$ satisfying
$e^2=e$, $g^2=e$, $eg=ge=g$, the idempotent
$\varepsilon=(e-g)/2$, a basis for the full left ideal
$A\varepsilon$, and a vector $v\in V$. The independent
checker verifies that the ideal has dimension two and
that $x\mapsto xv$ sends its basis to a basis of $V$.
It verifies this map's equivariance on every algebra
basis element. Thus $V\cong A\varepsilon$ is projective
in all $27$ cases. The $324$ explicit basis-equivariance
equations and all rational ranks pass.
Lemma~\ref{CP2:A-act-row-pd-one} now kills the products.

All tables, idempotents and bases are in
\path{converse/scope18/forced_three_row12.json}; its
digest and the independent completion receipt are in
\path{converse/scope18/forced_three_row12_checked.json}.
The calculation is newly generated from the finite
profiles and uses no inherited table or module certificate.
\end{proof}

\begin{lemma}[Proved: four outside elements and two columns]
\label{CP2:A-two-column-four-outside-bounded}
Let $M$ have a rectangular minimal ideal with two
columns and at most four outside elements. Every finite
right-act cone with bounded rational homology has zero
rational degree-two products, at any row or state count.
\end{lemma}
\begin{proof}
If its contracted algebra $A$ is local, use
Corollary~\ref{CP2:A-local-two-column-acts}.
Suppose $A$ is nonlocal. A point fixed by the whole
minimal band is fixed by $M$, since $zm=z(em)=z$
for a band idempotent $e$ and $em\in K$. Such a
point makes every second class relative to the monoid
vertex and gives zero products.
Without such a point, the projective orbit-kernel lemma
and constant-action compression give a monoid $N$
with the same homology as the act cone in every degree
at least three. In particular its homology is bounded.
Its contracted algebra is the same $A$, of dimension
at most four. The band--Tor formula and
Lemma~\ref{A:nonlocal-four} show that $H_4(BN;\mathbb Q)$
must vanish: otherwise it would recur in arbitrarily
high degrees. Thus the original cone also has zero
fourth rational homology, proving the claim.
\end{proof}

\begin{lemma}[Proved: a column quotient preserves act-cone $H^2$]
\label{CP2:A-column-quotient-act-H2}
Let $q:M\to M'$ collapse columns of a rectangular
minimal ideal, preserve its row set bijectively, and
leave every outside element distinct. Suppose the
right $M$-action on $S$ factors through this actual
monoid quotient. The induced functor between the
ordinary act cones induces an isomorphism on rational
second cohomology. No assertion about their higher
homology is needed here.
\end{lemma}
\begin{proof}
Write $R=\mathbb QM$, $R'=\mathbb QM'$ and
$L=\ker(\mathbb QS\to\mathbb Q)$. The projective
row modules $P=Re$ and $P'=R'q(e)$ identify under
$q$ with the same permutation module $\mathbb QI$;
their augmentation kernels identify with the same
module $V$. All of $P,V,L$ therefore factor through
$R'$, and their $R$-linear maps are precisely their
$R'$-linear maps. Dualizing the reduced act-cone
Tor description and using $0\to V\to P\to\mathbb Q\to0$
identifies second cohomology with
\[
 \operatorname{coker}\bigl(
 \operatorname{Hom}_R(P,DL)\longrightarrow
 \operatorname{Hom}_R(V,DL)\bigr),
 \qquad D=\operatorname{Hom}_{\mathbb Q}(-,\mathbb Q).
\]
The identical presentation over $R'$ proves the stated
isomorphism. Naturality of the augmentation sequence
identifies it with the pullback of the actual quotient
functor, so it preserves cup products.
\end{proof}

\begin{lemma}[Proved: three columns, three outside elements, five states]
\label{CP2:A-three-by-three-five-state}
A five-state ordinary right-act cone over a monoid
with minimal band $3\times3$ and at most three outside
elements cannot have type $\CP^2$.
\end{lemma}
\begin{proof}
A band-fixed state is globally fixed and gives zero
products. Otherwise every core orbit in $SK$ has
two or three states. Each such orbit is individually
$M$-invariant: if $x=y(i,j)$, then
$xm=y(i,\tau_m(j))$ lies in the same orbit.
A projective-plane cone has connected action category,
since its first rational homology is the reduced
zeroth homology of that category. Thus two core orbits
cannot exhaust $S$.

With one core orbit every band element acts constantly
on all states: it is constant on that orbit, and
$sk=(sk)k$ determines the same value for every $s$.
Compression then gives an actually equivalent monoid
of order at most $3+3\cdot5=18$, contradicting the
monoid minimum.

The only other possibility is two two-point core
orbits and one transient state $x$. The column maps
of the two orbits give surjections
$\pi_0,\pi_1:\{0,1,2\}\to\{0,1\}$, up to labelling
each orbit. If their two-part partitions differ,
the three identifications
$\pi_0(j)\sim\pi_1(j)$ join all four core states.
Indeed two distinct two-part partitions of a three-set
give a connected bipartite graph on their four parts.
The constant-action quotient therefore has just two
states, including $x$. Its degree-two products vanish,
and its $H^2$ surjection kills those on the original cone.

Suppose the partitions agree, and label both orbits
by a common quotient $\pi$ of the old column set.
Every right actor preserves each core orbit and obeys
$s_jm=s_{\tau_m(j)}$, so $\pi\tau_m$ depends only on
$\pi(j)$. Collapsing the band columns by $\pi$ is
therefore an actual congruence of $M$. Let $M'$ be
the quotient, with minimal band $3\times2$ and the
same at most three outside elements.

The action on $S$ factors through $M'$. This is clear
on the cores. For the remaining state $x$, the map
$j\mapsto x(i,j)$ lands, for each fixed row $i$, in
one core orbit and is its column map. Thus equal
$\pi$-values give the same result on $x$ as well.
Lemma~\ref{CP2:A-column-quotient-act-H2} makes the
quotient functor an $H^2$ isomorphism.
Every degree-two product of the $M'$ act cone vanishes
unconditionally by
Proposition~\ref{CP2:A-three-row-five-outside-acts}.
Pullback gives the contradiction. No boundedness
claim was transferred across this column quotient.
\end{proof}

\begin{theorem}[Proved: twelve corner elements and five states]
\label{CP2:A-twelve-five-state}
An ordinary act cone over a monoid of order at most
twelve with at most five states cannot have type
$\CP^2$.
\end{theorem}
\begin{proof}
Orders through eleven use
Theorem~\ref{CP2:A-bounded-eleven-acts}; at order twelve,
at most four states are already excluded. A possible
nonzero square requires at least three minimal rows
and two columns. A nontrivial maximal group would
cost at least thirteen elements, so the ideal is a band.
At order twelve the only shapes are $3\times2$ with
six outside elements, $4\times2$ with four,
$5\times2$ with two, and $3\times3$ with three.
The first is the fresh projective-row calculation,
the next two are the four-outside bounded lemma,
and the last is the preceding five-state lemma.
Every possibility is excluded.
\end{proof}

\begin{corollary}[Proved: ordinary act cones through eighteen arrows]
\label{CP2:A-ordinary-act-eighteen}
No ordinary finite-monoid act cone with at most eighteen
total arrows, identities counted, has type $\CP^2$.
\end{corollary}
\begin{proof}
Write its size as $m+s+1\leq18$. For $m\leq11$
use the bounded-eleven theorem. If $m=12$, then
$s\leq5$ and the preceding theorem applies. For
$m=13$, $s\leq4$ and
Theorem~\ref{CP2:A-four-state-fourteen} applies.
For $m\geq14$ there are at most three states, excluded
by the universal one- and two-state obstructions and
the exact three-state minimum of twenty-two arrows.
\end{proof}


\subsection{Further large-corner reductions at eighteen arrows}

\begin{lemma}[Proved: a trivial action at one end of a cut]
\label{CP2:A-trivial-cut-H2}
Let a category be a collage of a full source block $D$ and a
finite monoid $M$. Suppose $M$ acts trivially on every connecting
hom-set. Replacing $M$ by a trivial corner is an actual category
retract, and its pullback is an isomorphism on rational $H^2$.
The opposite statement also holds.
\end{lemma}
\begin{proof}
Let $T$ be the classifying space of the connecting $D$-act.
The two actual bar pushouts are
\[
 BD\longleftarrow T\times BM\longrightarrow BM,
 \qquad BD\longleftarrow T\longrightarrow *.
\]
The quotient functor has the section that includes the identity
of $M$. Every finite monoid has finite universal group, and
therefore $H^1(BM;\mathbb Q)=0$. Kunneth identifies the first
cohomology of $T\times BM$ with $H^1(T)$ and its second
cohomology with
\[
 H^2(T)\ \oplus\ H^0(T)\otimes H^2(BM).
\]
In the Mayer--Vietoris map from $H^2(BD)\oplus H^2(BM)$,
the two components are respectively the pullback to $T$ and
the injective map $b\mapsto-1_T\otimes b$.
Thus both the kernel in degree two and the preceding cokernel
in degree one are unchanged by the quotient. Naturality of
the exact sequence gives the stated $H^2$ isomorphism.
Here $T$ need only be nonempty; connectedness is not assumed.
\end{proof}

\begin{lemma}[Proved: a permutation orbit opposite a zero second group]
\label{CP2:A-permutation-fiber-cups}
In a collage $D\to M$, suppose $H^2(BD;\mathbb Q)=0$ and
one connecting stalk contains a nonempty finite $M$-invariant
subset on which every actor is a permutation. All rational
degree-two products of the collage vanish.
\end{lemma}
\begin{proof}
The homotopy-orbit space of that subset maps to the two-sided
bridge. Its projection to $BM$ is a finite covering, since
the action factors through a finite permutation group. The
covering transfer makes pullback from $H^2(BM;\mathbb Q)$
injective. Compatibility of the two restrictions of a collage
class therefore forces its restriction to $BM$ to be zero.
The restriction to $BD$ is zero by hypothesis. The two-block
relative suspension argument kills all second products.
\end{proof}

\begin{proposition}[Proved: two small incoming stalks through order eleven]
\label{CP2:A-two-small-incoming-eleven}
Suppose a full-linear category has first two corner orders
at most two, last corner order at most eleven, two arrows
from its middle to its last object, and at most three arrows
from its first to its last. If its rational homology is
bounded, every rational degree-two product vanishes. There
is no bound on its first adjacent hom-set.
\end{proposition}
\begin{proof}
First let the first two corners be trivial. Write $X,Y,Z$
for the three forward hom-sets, so $|Y|=2$, $|Z|\leq3$.
Every $x\in X$ supplies an $M$-equivariant map
$\beta_x:Y\to Z$. A constant map gives a fixed incoming
arrow. Otherwise every map is injective. Two different
two-point images in a three-set have a singleton invariant
intersection, again giving a fixed point. Hence all images
are one common invariant two-set $U$.

If the action on $Y$ consists of permutations, the preceding
lemma applies because the first two objects form a graph.
If two maps $Y\to U$ differ, their quotient is the
transposition of $Y$, commuting with every actor. A map of
two points commuting with their transposition is a
permutation, so this is the same case. Otherwise all
$\beta_x$ agree and some actor on $Y$ is constant. Its
action category is connected.

Fold $X$ to one chosen member, leaving $Y,Z,M$ fixed.
This is an actual category retract $C\to C_0$ because the
composition maps agree. The relative normalized chains,
one summand for each removed member of $X$, are the
augmentation cone of the $Y$ action-category chains,
shifted once. Their second homology is
$\widetilde H_0(Y_{hM};\mathbb Q)=0$. The retraction
therefore induces an isomorphism on $H^2$. The retract
has bounded homology. Deleting its middle object gives
the ordinary $M$-act cone on $Z$, and
Theorem~\ref{CP2:A-bounded-eleven-acts} kills its products.

For an initial corner of order at most two, take its
actual exact-orbit quotient first. This preserves the
rational ring, does not increase $|X|,|Z|$, and leaves
$Y$ unchanged. The initial corner is now trivial.
Let $N$ be the remaining middle corner, of order at most
two. Its rational algebra is semisimple.

If the right $N$-action on $Y$ is nontrivial, its orbit
set is a singleton. The equivariant graph model of
Lemma~\ref{CP2:A-baseline-semisimple-pair-graph} has
$Y/N$ as one vertex set. Its single vertex is fixed by
$M$, giving a section to the graph's homotopy-orbit
projection. The source block has zero rational $H^2$,
and the same two-block argument applies.

If the right action on $Y$ is trivial, replace $N$ by
a point and $X$ by $X/N$. Composition descends since
$\beta(y,nx)=\beta(yn,x)=\beta(y,x)$. This is an
actual rational homology equivalence. Indeed, the
strict-object bar filtration has only its corner,
edge and triangle terms. At $N$, exact semisimple
coinvariants identify respectively the corner bar
with a point, the $X$ bar with $X/N$, and the middle
triangle tensor with
$\mathbb QY\otimes\mathbb Q(X/N)$, since $N$ acts
trivially on $Y$. The other terms are unchanged.
The finite filtration proves the comparison and
passes boundedness to the quotient. The already
proved trivial-corner case now applies.
\end{proof}

\begin{proposition}[Proved: two first arrows and four long arrows]
\label{CP2:A-two-two-four-incoming}
A full-linear category with trivial first two corners,
last corner of order at most eleven, $|X|=|Y|=2$ and
$|Z|\leq4$ cannot have type $\CP^2$.
\end{proposition}
\begin{proof}
Consider its two maps $\beta_0,\beta_1:Y\to Z$.
Constant maps, distinct images meeting in one point,
or a permutation action on $Y$ are excluded by the
fixed-arrow or permutation-fiber lemmas. If the images
agree and the maps differ, the commuting transposition
again forces the permutation case. If the maps agree,
folding $X$ to a singleton is a cut-profunctor retraction.
Lemma~\ref{CP2:A-cut-retract} would retain projective-plane
type on its image. Deleting the middle object there
gives a bounded ordinary act cone over a monoid of
order at most eleven, a contradiction.

The only remaining case has disjoint images. They
exhaust $Z$, since both have size two. Fold $X$ onto
$x_0$ and fold the second image onto the first by
$\beta_0\beta_1^{-1}$. The latter map is $M$-equivariant.
These two maps form an actual idempotent natural
endomorphism of the cut from object0 to objects1,2:
its composition equation is exactly
$\gamma_Z\beta_i=\beta_0$.
Its image has $|X|=1$, $|Z|=2$. The same cut-retract
lemma and middle deletion give an ordinary binary act
cone, whose second products vanish universally.
\end{proof}

\begin{lemma}[Proved: a trivial middle corner and a singleton long hom-set]
\label{CP2:A-trivial-middle-singleton-long}
If a full-linear category has a trivial middle corner
and a singleton long hom-set, every rational
degree-two product vanishes, with no restriction on
its two other corners or adjacent hom-sets.
\end{lemma}
\begin{proof}
Let $D$ be the first two objects and $M$ the last.
Choose an arrow $y$ from the middle to the last object.
The unique long arrow and $y$ give a natural
transformation from the inclusion of $D$ to the
constant functor at the last object: every two-cross
composite is the unique long arrow, the first corner
fixes it, and the middle corner is trivial.
Thus $BD\to BC$ is null-homotopic. The unique long
arrow also gives a natural transformation from the
constant functor at the first object to the inclusion
of $M$, so $BM\to BC$ is null-homotopic.
Both restrictions of every positive class vanish.
Apply the two-block relative suspension argument.
\end{proof}

\begin{lemma}[Proved: a three-element middle opposite a binary stalk]
\label{CP2:A-small-middle-binary-singleton-long}
Suppose the first corner is trivial, the middle corner
$N$ has at most three elements, $|Y|=2$ and $|Z|=1$.
Every rational degree-two product vanishes.
\end{lemma}
\begin{proof}
The singleton $Z$ makes the last corner's inclusion
null-homotopic. If the right $N$-action on $Y$ has a
common fixed point, that point and $Z$ give a natural
transformation from the source block to the constant
last object, as above, and both restrictions vanish.

Suppose that action has no common fixed point. Its
transformation image has at most three elements. On
a two-set the possibilities are the permutation
group $C_2$ and the identity together with both
constants. In the former case $N$ surjects onto
$C_2$. Its minimal group has at least two elements;
the order bound forces $N=C_2$ or $C_2$ with a new
identity. Both rational algebras are semisimple, so
the source two-object block has zero rational $H^2$.
Together with the already null last-corner inclusion
this kills the products.

In the constants case the commuting last-corner
action on $Y$ is the identity: a map commuting with
both constants must fix both points. That corner
also acts trivially on $Z$. Apply
Lemma~\ref{CP2:A-trivial-cut-H2} to replace it by a
point. The resulting two trivial endpoint corners
and binary adjacent stalk satisfy the binary
projective-cover lemma, so all its second products
vanish. The $H^2$ isomorphism pulls back the result.
\end{proof}

\begin{proposition}[Proved: a three-element initial singleton]
\label{CP2:A-three-initial-singleton}
A full-linear category with initial corner order at
most three, trivial middle corner, last corner order
at most eleven, $|X|=1$, $|Y|=2$ and $|Z|\leq3$
cannot have type $\CP^2$.
\end{proposition}
\begin{proof}
If the initial trivial left module is projective,
use Proposition~\ref{CP2:A-eleven-projective-singleton}.
The small-monoid classification leaves only the
identity-adjoined two-element left-zero band.
Its two right actors $P,Q$ on $Z$ are retractions
onto a common image $U$, and the composition map
$\beta:Y\to Z$ lands in $U$. A constant $\beta$,
or $|U|=1$, supplies a fixed incoming arrow. The
source block has a terminal middle object, so its
second homology is zero and this case is excluded.
Otherwise $\beta$ is injective and $|U|\geq2$.

If $P=Q$, their common map on $Z$, together with
the identity on $X$, is an idempotent natural
endomorphism of the cut from the first object.
Its image retains a putative projective-plane type
by the cut-retract lemma. The initial corner acts
trivially on both image stalks, so the opposite
version of Lemma~\ref{CP2:A-trivial-cut-H2} replaces
it by a point without changing $H^2$. The replacement
is itself a category retract and hence bounded.
Middle deletion and the bounded-eleven act theorem
give the contradiction.

If $P\ne Q$, the size bound forces $|Z|=3$,
$|U|=2$, and one state outside $U$. The equivariant
graph of the right left-zero action has vertices
$U$ and the one edge from $Pz$ to $Qz$. Its two
endpoints are distinct. The element category of
the full incoming source-block act adds the two
vertices $Y$ and the edges from $\beta(y)$ to $y$.
This follows from the natural action-category graph
comparison: its fixed core vertices are unchanged,
and the unique first-adjacent arrow attaches exactly
the indicated two edges. All comparisons commute
with the last-corner action. The resulting graph
is a tree, hence contractible. The source block is
also contractible, so the two-block bar pushout is
the cofiber of an equivariant contractible graph's
homotopy-orbit map to the last corner. That map is
a homotopy equivalence, giving a contradiction.
\end{proof}

\begin{proposition}[Proved: three middle-to-last arrows and two long arrows]
\label{CP2:A-three-two-incoming}
A full-linear category with trivial first two corners,
$|Y|=3$, and $|Z|=2$ cannot have type $\CP^2$,
for any last monoid and any nonempty $X$.
\end{proposition}
\begin{proof}
Any constant composition map $\beta_x:Y\to Z$ gives
a fixed point. Assume they are all surjective. If two
have the same partition and differ, they differ by
the transposition of $Z$, so every last-corner actor
on $Z$ is a permutation. The permutation-fiber lemma
excludes this case. If all maps agree, fold $X$ to
a singleton by a cut retraction. Its image would
still be a projective plane, but middle deletion
gives an ordinary two-state cone, a contradiction.

Thus choose two maps with different two-part
partitions. Their joint map embeds the three-set
$Y$ into $Z^2$. The diagonal transposition of $Z^2$
has no fixed point and cannot preserve a three-set.
No last-corner actor can therefore transpose $Z$.
Unless there is a globally fixed point, the actor
image on $Z$ contains both constants. The joint
image consequently consists of $(0,0),(1,1)$ and
one off-diagonal pair. Name the corresponding states
$y_0,y_1,w$, and interchange the two chosen maps
if needed so that the last pair is $(0,1)$.
An actor is the identity on $Y$ when it is the
identity on $Z$, and is constant to $y_i$ when
it is constant to $i$. Every other $\beta_x$ must
take $y_i$ to $i$ and $w$ to either0 or1. Thus
there are exactly two composition maps in use.

Fold all repetitions in $X$, retaining one copy of
each map. This is a cut retraction fixing $Y,Z$.
Its image would still have projective-plane type.
There $|X|=2$ and the first two objects form a
circle $BD$. The incoming act graph has vertex
sets $Y,Z$ and the six edges of the two maps.
Contract the three edges of the first map. They
form an equivariant forest. The quotient graph
$\Gamma$ has two vertices, one loop at each,
and a connecting edge. It is connected and its
first homology is the permutation module
$\mathbb Q\{0,1\}$: the two constant actors send
both loop generators to the selected loop.

Put $X_0=\operatorname{hocofib}(\Gamma_{hM}\to BM)$.
For a connected graph its augmented chain complex
is equivariantly quasi-isomorphic to $H_1(\Gamma)[1]$.
Hence
\[
 H_2(X_0;\mathbb Q)=H_1(\Gamma;\mathbb Q)_M=\mathbb Q.
\]
Collapsing $BD$ in the collage gives $X_0$.
Its cofiber boundary $H_2(X_0)\to H_1(BD)$ sends
each loop generator to the generator of $BD$:
before contraction it follows the second first-arrow
and returns along the first. It is therefore the
augmentation on this permutation module, an
isomorphism after coinvariants. Since $H_2(BD)=0$,
the cofiber exact sequence forces $H_2(BC)=0$.
The retract cannot be a projective plane.
\end{proof}

\begin{proposition}[Proved: a two-element endpoint in the three-by-three triangle]
\label{CP2:A-three-three-two-endpoint}
Every full-linear category with profile
$(1,m,2;3,3,1)$ and $m+9\leq22$ is excluded.
\end{proposition}
\begin{proof}
Take the exact-orbit quotient at the last corner,
whose rational algebra is semisimple. It preserves
the rational cohomology ring. If its three-element
incoming stalk shrinks, it becomes binary and the
binary projective-cover obstruction applies. If
it does not shrink, the last corner acts trivially
on that stalk and on the singleton long stalk.
The quotient then has an actual section. Its integral
homology groups are direct summands of those of a
putative projective plane. The rational $H^2$
isomorphism forces its integral $H_2$ to be $\mathbb Z$;
naturality of the primitive square then forces
$H_4=\mathbb Z$ with the same unit square. Its
fundamental group is a retract of the trivial group.
The projective-plane recognition argument of the
cut-retract lemma therefore shows that this quotient
also has type $\CP^2$.
It has trivial endpoints and singleton long hom-set.
Triangle compression gives a monoid of order $m+9$,
contradicting the monoid minimum of twenty-three.
\end{proof}

\begin{remark*}[Exact remaining large-corner scope]
The fresh eighteen-arrow profile ledger initially has fourteen
full-linear profiles with a corner of order at least eight.
The preceding lemmas remove twelve; the full three-by-three
triangle theorem removes $(1,8,1;3,3,2)$.
The single remaining profile is $(1,1,8;2,3,3)$, together
with its opposite. Its separate complete reduction and fresh
literal calculation are given in
Theorem~\ref{CP2:A-extreme-eight-two-three-three} below.
\end{remark*}


\subsection{The last large endpoint at eighteen arrows}

\begin{theorem}[Proved: the profile $(1,1,8;2,3,3)$]
\label{CP2:A-extreme-eight-two-three-three}
No full-linear category with profile $(1,1,8;2,3,3)$,
or its opposite, has type $\CP^2$.
\end{theorem}
\begin{proof}
Write $M$ for the last corner, $X$ for its two first
adjacent arrows, and $Y,Z$ for the two three-point
incoming stalks. The two compositions are equivariant
maps $\beta_0,\beta_1:Y\to Z$. The first two objects
form a graph $D$. If $H^2(BM;\mathbb Q)=0$, the
two-block lemma excludes a nonzero square. We may
therefore use the complete $510$-table positive-corner
catalog of Proposition~\ref{CP2:A-positive-corners-eight}.
It contains $492$ minimal bands of shape $2\times2$
and nine each of shapes $2\times3$ and $3\times2$.

A constant $\beta_i$ supplies a fixed incoming arrow.
If one map is bijective, identify $Y$ with $Z$ by it
and write the other as the commuting endomorphism $T$.
If $T$ is nonbijective, an idempotent power of $T$,
on both stalks, is a proper natural endomorphism of
the cut from $D$ to $M$. Its image has both stalks
of size at most two. The cut-retract lemma transfers
a putative projective plane to that image, excluded
by the fixed-arrow or two-small-incoming obstruction.
If $T$ is a transposition, its unique fixed point
is fixed by every commuting $M$ actor. If $T$ is a
three-cycle, its transformation centralizer consists
of permutations, and the permutation-fiber lemma
applies. If $T$ is the identity, the two composition
maps agree. Folding $X$ to a singleton is a cut
retraction, whose image is an ordinary three-state
act cone after middle deletion. It is excluded by
the bounded-eleven act theorem.

Both maps consequently have rank two. If their images
differ, their singleton intersection is invariant and
gives a fixed point. Let their common image be $U$.
If their two-part partitions agree, the maps either
agree, already treated, or differ by the transposition
of $U$. Equivariance then forces every actor on $U$
to be a permutation, another excluded case.

For different partitions the joint map embeds $Y$
as a three-point subset of $U^2$. No diagonal
transposition preserves such a subset. An actor on
$U$ is therefore the identity or a constant. Unless
there is a common fixed point, both constants occur.
The joint image contains $(0,0),(1,1)$ and one
off-diagonal point. Relabel $Y,U$ and the two members
of $X$ to obtain
\[
 \beta_0=(0,1,0),\qquad \beta_1=(0,1,1),
 \qquad U=\{0,1\}\subset Z=\{0,1,2\}.
\]
An actor on $Y$ is now uniquely determined by its
restriction to $U$: it is the identity, constant0,
or constant1. The extra point of $Z$ may have any
of three images. Thus there are exactly nine paired
actors. This finite reduction is also checked directly
by trying all $27^2$ pairs of transformations of $Y,Z$
against both equivariance equations.

The minimal ideal $K$ of $M$ is a band. Its image
on $U$ is the two-constant minimal ideal, because a
surjective monoid map sends the minimal ideal onto
the minimal ideal of its image. Each band element
actually acts constantly on all of $Z$. To see this,
suppose some $k\in K$ fixed the extra point $z$.
Since $z=kz$, the orbit $Kz$ is a row quotient of
the principal left ideal $Kk$; every band actor is
constant on this orbit. In particular $k$ is constant
there with value $z$. The orbit cannot meet $U$,
where $k$ takes a value in $U$. Since there is only
one point outside $U$, it follows that $Kz=\{z\}$.
Then $mz=m(kz)=(mk)z=z$ for every $m$, a forbidden
global fixed point. Hence $kz\in U$ for every $k$;
idempotency makes it the same constant value as on $U$.
These values are constant on each minimal row and
take both values0 and1.

This gives the following complete fresh finite scope.
For each of the $510$ literal monoid representatives,
choose every surjective assignment of its minimal
rows to $\{0,1\}$. Its band actions are now forced.
Choose every one of the nine paired actors at each
remaining nonidentity label, subject to multiplication.
The producer
\path{converse/scope18/extreme8_actions.py} uses only
the forced mixed band equations to shorten these
domains, then checks every full multiplication equation.
Its $26244$ outside assignments give exactly $3858$
literal eighteen-arrow categories. It does not divide
out any action or category isomorphisms.

The separate program
\path{converse/scope18/check_extreme8_scope.py}
imports no constructor. It obtains the nine actors
from the $27^2$ direct tests, recovers band rows as
principal right ideals, and assigns all nine possibilities
to outside labels in reverse order. It propagates no
unassigned value. Its $61032$ nodes and $2896504$
evaluated action equations recover exactly the same
action sets. It independently constructs every literal
category from named arrows and verifies the exact
complete table-stream correspondence.

Finally, every one of these $3858$ tables is checked
through the normalized bar complex. Exactly $2272$
have mod-two second Betti number different from one.
Each of the other $1586$ has an explicit nonboundary
two-cocycle and a literal three-cochain whose coboundary
is its square. The independent scalar checker
\path{converse/scope18/check_extreme8.cpp} checks all
tables, including the low-homology rejections. It verifies
$3927444$ typed associativity equations, $1076894$
cocycle equations and $7538258$ square-primitive
equations. There is no unexcluded record.
All table, input, code, action-scope and checker hashes
are joined in
\path{converse/scope18/extreme8_complete.json}.
Thus every branch of the stated profile is excluded.
\end{proof}


\subsection{Further star reductions at eighteen arrows}

\begin{lemma}[Proved: merging a trivial leaf into an arbitrary leaf]
\label{CP2:A-star-general-leaf-merge}
Consider a source star with center $M$, one leaf $N$ on a
left-$N$, right-$M$ biset $S$, and a trivial leaf on a right
$M$-act $T$. If there is an $M$-equivariant map $T\to S$,
the star retracts onto its full $(M,N,S)$ edge. If the star
has the rational homology and cohomology ring of $\CP^2$,
so does that edge.
\end{lemma}
\begin{proof}
Map the trivial leaf to the object with corner $N$, its
identity to $1_N$, and its stalk by the specified map.
Keep the retained edge unchanged. Equivariance is exactly
the remaining functoriality condition, so this is an actual
category retraction.

In the double mapping-cylinder model its quotient by the
retained edge is the unreduced suspension of $B(T,M)$:
the center, the retained leaf and its entire bridge are
collapsed, leaving the other bridge with its two ends
collapsed separately. If the retract had zero second
cohomology, every second class on the star would lift from
this suspension and all second products would be zero.
Thus its second group is nonzero. The retraction makes its
homology a direct summand of the rational projective-plane
homology. Pullback in degree two is consequently an
isomorphism; the nonzero square forces its fourth group
to be nonzero as well. These facts give exactly the stated
rational homology and ring. No condition on $BN$ is used.
\end{proof}

\begin{lemma}[Proved: a constant actor on a connected leaf action space]
\label{CP2:A-connected-leaf-constant}
Let $S$ be a left-$N$, right-$M$ biset, where $BN$ is
rationally acyclic. Suppose the ordinary left action space
$Z=B(N,S)$ is connected and some right $M$ actor on $S$
is constant. The full two-object collage $(M,N,S)$ has
zero rational second homology.
\end{lemma}
\begin{proof}
Collapsing $BN$ in the double-bar pushout is a rational
homology equivalence onto
$Y=\operatorname{hocofib}(Z_{hM}\to BM)$.
The cone of the augmented right $\mathbb QM$ chain complex
$C_*(Z;\mathbb Q)\to\mathbb Q$ has zero homology in
degrees zero and one, since $Z$ is nonempty and connected.
Its degree-two homology is $H_1(Z;\mathbb Q)$.
Tensoring a projective left monoid resolution gives the
relative bar complex for $Y$. Its lowest-degree homology,
or equivalently the first nonzero row of the hyperhomology
spectral sequence, therefore gives
\[
 H_2(Y;\mathbb Q)=H_1(Z;\mathbb Q)_M.
\]
If $m$ acts constantly with value $u$, commutation makes
$u$ fixed by every $N$ actor. The induced map on $Z$
factors as the forgetful map $Z\to BN$ followed by the
fixed-state inclusion $BN\to Z$. It acts as zero on
$H_1(Z;\mathbb Q)$. The coinvariant relation $v m=v$
then kills every class $v$, proving the claim.
\end{proof}

\begin{proposition}[Proved: a binary leaf beside a trivial leaf]
\label{CP2:A-star-binary-leaf-eleven}
A source star with center of order at most eleven, a
trivial leaf on any nonempty finite stalk, and another
leaf on at most two states cannot have the rational
homology and cohomology ring of $\CP^2$. There is no
restriction on the other leaf corner.
\end{proposition}
\begin{proof}
Write $S$ for the binary stalk and $N$ for its leaf corner.
A singleton stalk has a center-fixed state; more generally,
if the left $N$ image contains a constant map, its value is
fixed by the commuting center action. The constant map
from the other stalk to that point is center-equivariant.
Lemma~\ref{CP2:A-star-general-leaf-merge} transfers the
rational projective-plane target to the retained binary
edge, contradicting the universal two-cross obstruction.

If the left image contains a transposition but no constant,
both corner actions on $S$ are permutations. Its bridge
maps as a finite covering to $BN\times BM$. Transfer
makes rational cohomology pullback injective. Compatibility
of the restrictions of a global second class to the two
corners, together with $H^1(BN;\mathbb Q)=H^1(BM;\mathbb Q)=0$,
therefore forces both restrictions to be zero. The other
leaf is trivial, so the two-block relative-cup lemma kills
all second products.

The remaining left image is the identity. Apply
Lemma~\ref{CP2:A-trivial-cut-H2} to replace $N$ by a
trivial corner. This is an actual retract and an $H^2$
isomorphism, so its homology is bounded. Merging its two
trivial leaves gives the ordinary center-act cone on the
disjoint union of the stalks, homotopy equivalent to the
retracted star wedged with a circle. The bounded-eleven
act theorem kills every second product there. Pullback
along the $H^2$ isomorphism gives the contradiction.
\end{proof}

\begin{proposition}[Proved: a small three-state leaf]
\label{CP2:A-star-small-three-leaf}
A source star with center of order at most eleven, one
trivial leaf, and another leaf of order at most four on
at most three states cannot have the rational homology
and cohomology ring of $\CP^2$. The other stalk may have
any positive size. The same conclusion holds if the
otherwise trivial leaf has order at most two.
\end{proposition}
\begin{proof}
At a leaf of order at most two, exact orbit comparison
replaces it by a trivial leaf without increasing its
stalk. We may therefore assume that leaf trivial.
The binary case is the preceding proposition. Write
$N$ for the small leaf corner and $S$ for its three-set.
If its trivial right rational module is projective,
exact orbit comparison makes this leaf trivial as well;
merging the leaves and applying the bounded-eleven theorem
gives the result.

Otherwise the small minimal-ideal classification makes
the minimal ideal of $N$ a one-row band with at least
two columns and trivial maximal group. Its left actors
on $S$ are retractions onto a common nonempty set $U$,
fixed pointwise by every actor of $N$. In particular
$BN$ is contractible. If $|U|=1$, its point is fixed by
the center. Restrictions of every global second class
to the center are zero, as are those to both leaves;
the two-block lemma applies. If $|U|=3$, the entire
$N$ action is trivial, so the leaf can be collapsed by
a homotopy equivalence and the bounded-eleven argument
applies.

Let $|U|=2$. If all band actors are the same retraction
$P$, then $P$ commutes with every actor of $N$: left
multiplication fixes each band element, and right
multiplication remains in the band. It also commutes
with the center. Retracting the stalk to $U$ transfers
the rational target by Lemma~\ref{CP2:A-star-leaf-reductions}(4).
On $U$ the leaf action is trivial. Collapse its contractible
corner, merge the now trivial leaves, and use bounded-eleven.

Finally suppose two band actors $P,Q$ differ. For the
unique state $w\notin U$, their values are the two distinct
points of $U$. A transformation commuting with $P,Q$ is
either the identity on $S$ or a constant map to one of
these two points: the image of $w$ determines all three
values. If the center image consists only of the identity,
there is a center-fixed state and the relative-cup argument
already applies. Otherwise a center actor is constant.
The left $N$ action category is connected, since its arrows
from $w$ reach both points of $U$. The preceding connected
leaf lemma makes the full $(M,N,S)$ edge have $H^2=0$.
Its union with the remaining trivial leaf is then a
two-block collage with both second restrictions zero.
This finishes all cases.
\end{proof}

\begin{proposition}[Proved: four states opposite an order-three leaf]
\label{CP2:A-star-three-leaf-four-states}
A source star with center of order at most eight, one
trivial leaf, and another leaf of order at most three on
at most four states cannot have the rational homology
and cohomology ring of $\CP^2$.
\end{proposition}
\begin{proof}
The preceding proposition treats at most three states.
An order-three leaf with projective trivial right module
is replaced by its exact orbit set, and the merged
ordinary act is excluded by bounded-eleven. Otherwise
the leaf is $R_3=\{1,p,q\}$, with $ab=b$ on its two
nonidentity elements. Its actors $P,Q$ are retractions
onto the same subset $U$ of the four-state stalk.

If $|U|=1$, a center-fixed state kills the products. If
$P=Q$, their common biset retraction transfers the rational
target to the stalk $U$. The leaf now acts trivially;
collapse its contractible corner and use bounded-eleven.
If $|U|=3$ and $P\ne Q$, the natural graph model of
Lemma~\ref{CP2:A-high-P5-graph} has a single connecting
edge and one isolated vertex. Equivariant collapse to
its two contractible components replaces the entire
leaf and stalk by a trivial two-state leaf, by homotopy
pushout invariance. Again merge and use bounded-eleven.

It remains that $|U|=2$, $P\ne Q$. The natural graph
has two vertices and two edges; it is connected and
has one cycle. By Lemma~\ref{CP2:A-one-cycle-graph-cofiber},
if any center actor has cycle degree different from one,
the retained full edge has zero $H^2$, and its union
with the trivial leaf has zero second products.
Assume every actor has degree one.

If the graph has one loop and one connecting edge, all
actors fix the loop vertex. If its two connecting edges
have the same direction, degree one forces both vertices
fixed. Either case gives a center-fixed state. In the
last case the two edges have opposite directions and
the only possible actors are the identity and half-turn,
as in the four-state graph proof cited above.

If the center has $H^2(BM;\mathbb Q)=0$, both sides of
the cut from the center to its two leaves have zero
second restriction, and the two-block lemma applies.
Otherwise the center's minimal ideal has at least two
rows and two columns. A nontrivial maximal group would
require at least eight ideal elements and an additional
identity, contrary to $|M|\leq8$. Its maximal group is
therefore trivial. That local group surjects onto the
universal group of $M$, so every homomorphism from $M$
to the half-turn group is trivial. All actors fix both
vertices, giving the remaining center-fixed state.
This proves the assertion.
\end{proof}

\begin{proposition}[Proved: a twelve-element center with two binary leaves]
\label{CP2:A-star-twelve-binary-leaves}
A source star with a center of order twelve, two trivial
leaves and two arrows to each leaf cannot have type
$\CP^2$; indeed it cannot have its rational homology and ring.
\end{proposition}
\begin{proof}
Merge the two leaves. The resulting ordinary act cone
on the disjoint union of two invariant two-sets has
bounded rational homology and the same products in
degrees at least two. A possible nonzero product requires
at least three minimal rows and two columns, by the
two-row act obstruction; its maximal group is trivial
by the order budget. The possible minimal bands are
$3\times2$, $4\times2$, $5\times2$, and $3\times3$,
with respectively six, four, two, and three outside
elements. The first three are excluded by
Propositions~\ref{CP2:A-three-row-six-outside-acts} and
\ref{CP2:A-two-column-four-outside-bounded}.

For the $3\times3$ band, a band-fixed state is fixed by
the whole monoid and kills products. In its absence
each invariant two-set is a full core orbit. Their
column maps give two two-part partitions of the three
band columns. If the partitions differ, the constant
action quotient identifies all four states, since
the bipartite graph on the four parts is connected.
Its cone has zero second products; the natural $H^2$
surjection from the constant-action quotient kills
those on the original cone.

If the partitions agree, collapsing columns by their
common two-part partition is a monoid congruence:
every outside column actor preserves that partition,
as witnessed on either core orbit. The entire four-state
action factors through the quotient monoid, which has
minimal band $3\times2$ and three outside elements.
Lemma~\ref{CP2:A-column-quotient-act-H2} gives an $H^2$
isomorphism. The three-row act theorem kills products
on the quotient without any boundedness assumption,
and naturality finishes the proof.
\end{proof}

\begin{remark*}[Exact scope of the new large-corner reductions]
In the notation $(a,b,c,s,t)$ for center, leaf orders and
the corresponding stalk sizes, these propositions exclude
the following seven eighteen-arrow profiles and their
leaf interchanges:
\[
\begin{gathered}
 (8,1,3,2,4),\quad(8,1,3,3,3),\quad(8,1,4,2,3),\\
 (8,1,5,2,2),\quad(8,2,3,2,3),\quad(9,1,3,2,3),\\
 (12,1,1,2,2).
\end{gathered}
\]
The two coned-three-state profiles $(5,1,8,1,3)$ and
$(8,1,5,1,3)$ remain separate obligations here.
\end{remark*}


\subsection{A first exact star reduction at eighteen arrows}

\begin{proposition}[Proved: a small center through eighteen arrows]
\label{CP2:A-star-small-center-eighteen}
A three-object star with at most eighteen arrows and center
of order at most four cannot have type $\CP^2$.
\end{proposition}
\begin{proof}
Take opposition so the center is a source. Its classifying
space has finite-group homotopy type by the small-monoid
lemma. If it is contractible, the union of the two full
edges along that corner is a homotopy wedge of those edges.
A projective-plane wedge target forces one edge to have
the projective-plane type: the integral homology summands,
the primitive square and the free-product fundamental
group give the usual simply connected recognition.
That edge has at most sixteen arrows, contradicting the
complete two-object baseline.

Suppose instead that the center has a nontrivial minimal
group. Its minimal ideal must itself be a group. Indeed,
a nontrivial group together with at least two rows or
two columns would require at least four ideal elements
and an additional identity. Both trivial rational modules
of the center are therefore projective, by averaging
over this group ideal. Exact orbit comparison at the
center gives a rationally equivalent star with trivial
center, hence a wedge of two ordinary leaf act cones.
Bounded rational homology passes to each wedge summand.
A leaf of order at most eleven is covered by the
bounded-eleven act theorem. For a leaf of order $b\ge12$,
the original center has at least two elements and the
other leaf and its stalk have at least one each. Its
own stalk has at most $18-2-b-1-1=14-b\le2$ states.
The binary act obstruction applies. Thus both wedge
summands have zero second products, a contradiction.
\end{proof}

\begin{corollary}[Proved: two small leaf reductions through center order eleven]
\label{CP2:A-star-leaves-bounded-eleven}
For a source star with bounded rational homology, suppose
the center order is at most eleven, both leaf orders
are at most four, and each leaf either has order at
most two or is joined by at most two arrows.
Every rational degree-two product vanishes.
\end{corollary}
\begin{proof}
Apply Lemma~\ref{CP2:A-star-leaf-reductions}(1)--(2)
at each leaf. A center-fixed state kills restrictions
to the center; the leaf corners have zero rational
second cohomology, so the relative two-block argument
kills products. Otherwise both leaves become trivial
by rational homology and ring equivalences. Merging
the two leaves gives the ordinary center-act cone on
their disjoint union, homotopy equivalent to the star
wedged with a circle. Its rational homology is bounded,
and the bounded-eleven theorem kills its second products.
\end{proof}

\begin{remark*}[Scope only]
The fresh program \path{converse/scope18/star18_scope.py}
uses these two statements and the previously checked
budget-independent star cuts. It does not apply the
seventeen-arrow large-center or large-leaf theorems at
eighteen arrows. Its residual list remains a separate
proof or finite-enumeration obligation.
\end{remark*}


\subsection{The exact full-linear scope at eighteen arrows}

\begin{proposition}[Proved: the remaining full-linear eighteen-arrow profiles]
\label{CP2:A-linear-eighteen-scope}
A three-object category with full acyclic support, eighteen
arrows and classifying space of type $\CP^2$ must have one
of the $373$ cardinality profiles in
\path{converse/scope18/linear18_residual_profiles_v3.tsv},
up to opposition. Every remaining corner has order at
most seven. This statement is a scope reduction and does
not exclude those remaining profiles.
\end{proposition}
\begin{proof}
Write the six positive cardinalities as
$(a,b,c;x,y,z)$, where $a,b,c$ are the successive corner
orders and $x,y,z$ are the $01,12,02$ stalk sizes.
There are $\binom{17}{5}=6188$ positive tuples of sum
eighteen. Opposition sends this tuple to
$(c,b,a;y,x,z)$. Its fixed tuples satisfy $a=c$ and
$x=y$; their number is
\[
 \sum_{r=2}^{8}(r-1)(17-2r)=140.
\]
Thus there are $(6188+140)/2=3164$ opposition orbits.

The fresh source \path{converse/scope18/linear18_scope.py}
generates every tuple by its five partial sums, canonically
chooses the lesser of the two orientations, and tests the
named, already proved two-block, fixed-arrow, binary,
projective-cover, semisimple-graph, singleton-long,
singleton-comma and triangle-compression hypotheses.
These are applied with their actual bounds; no
seventeen-arrow-only exclusion is substituted. In
particular singleton-comma deletion leaves an ordinary
act cone with at most fifteen arrows, and the initial
projective-singleton cut uses the bounded-eleven theorem.

The separate v2 source adds exactly the hypotheses of
the preceding large-corner propositions, the complete
$(1,1,8;2,3,3)$ literal gate, and
Corollary~\ref{CP2:A-three-three-triangles}, in either
orientation. This gives $2790$ proved profiles and
$374$ residual profiles, all of whose corners have
order at most seven. The independent program
\path{acyclic/code/check_linear18_scope_predicates.py}
imports neither scope source. It regenerates all $6188$
tuples, all $3164$ canonical profiles, every proof reason
and the entire residual ledger. Its per-predicate hostile
review checks the actual theorem hypotheses at each use;
its complete arithmetic and digest join passes.

Finally Corollary~\ref{CP2:A-three-four-two-triangles}
removes exactly one additional residual tuple,
$(1,7,1;3,4,2)$. The v3 source leaves both earlier
sources and ledgers unchanged and binds the new proof
to its complete independent $23$-orbit certificate join.
The separate extension checker confirms that this is the
only changed row and proof reason. The final partition is
\[
 3164=2791+373.
\]
The maximum-corner distribution of the residual ledger is
\[
\begin{array}{c|rrrrr}
\text{maximum corner order}&3&4&5&6&7\\\hline
\text{profiles}&47&85&140&69&32.
\end{array}
\]
The exact v3 ledger digest is recorded in
\path{converse/scope18/linear18_scope_v3.json} and checked
again in
\path{acyclic/outputs/linear18_scope_v3_independent_audit.json}.
This proves the stated complete cardinality reduction.
\end{proof}

\begin{remark*}[The remaining composition search]
The fresh canonical-biset producer
\path{converse/scope18/linear18_bisets.cpp} uses only
the general corner-table cuts audited in
\path{converse/scope18/linear18_engine_scope.md}.
The whole residual input has at most $100$ normalized
two-simplices and $600$ normalized three-simplices;
all representation bounds are checked before search.
The independent literal checker imports no producer
code. Each completed profile receives its own closed
stream and source-bound verification receipt. A partial
profile or unresolved literal record supplies no
exclusion in this proposition.
\end{remark*}


\subsection{Cyclic four-object support at eighteen arrows}

\begin{lemma}[Proved: a one-row source component reduces to its minimal group]
\label{CP2:A-source-minimal-group-collapse}
Let $S$ be a strong source component on at least two objects in
a finite category $C$. If each local minimal row count of $S$
is one, there is an actual homotopy equivalence from a category
with strictly fewer arrows to $C$, replacing $S$ by its finite
minimal group. There is a one-column sink dual.
\end{lemma}
\begin{proof}
Choose a minimal idempotent $e\in S(a,a)$ and adjoin its retract
object to obtain the full one-object extension $C^e$ in the
idempotent completion. Its group corner is $G=eS(a,a)e$.
The added object is a retract of $a$, so
$BC\to BC^e$ is a homotopy equivalence. Let $D$ be the full
complement of $S$ and let $C'$ be the full category on $D$ and
the new object $e$.

For each retained object the comma category of the inclusion
$C'\to C^e$ has its usual terminal identity. For $x\in S$,
there are no arrows from $D$ into $x$, because $S$ is a source
component. The comma category therefore consists of
$C^e(e,x)=S(a,x)e$, with its right $G$ action. The categorical
minimal-ideal lemma makes this action free, and row count one
makes it transitive. There is exactly one morphism between
each ordered pair of its objects, so the comma category is
contractible. Quillen's Theorem~A proves $BC'\simeq BC$.

For $d\in D$, the new outgoing stalk $C^e(e,d)=C(a,d)e$ is a
subset of the old stalk $C(a,d)$. No other outside hom-set grows,
and there are no incoming stalks from $D$. The internal part
$S$ has at least $|G|\,|\operatorname{Ob}S|^2>|G|$ arrows,
whereas its replacement has exactly $|G|$. The total arrow
count strictly decreases. Taking opposites proves the dual.
\end{proof}

\begin{lemma}[Proved: two supported augmentation modules at a small strong pair]
\label{CP2:A-small-pair-augmentation-Tor}
Let $S$ be a strongly connected two-object category with at most
seven arrows. Suppose its categorical minimal ideal $K$ has
exactly one arrow in each ordered hom-set and contains neither
object identity. Adjoin a trivial initial corner $u$ and a trivial
final corner $v$, with every entrance and exit stalk nonempty
and any nonempty direct hom-set $u\to v$. If each of the entrance
and exit augmentation modules is supported on at most one object
of $S$, bounded rational homology of the resulting category
implies that its fourth rational homology is zero.
\end{lemma}
\begin{proof}
Write $R=\mathbb QS$ and let $V,W$ be the left entrance and
right exit augmentation modules. If either is zero, the
middle-bar formula of Lemma~\ref{CP2:A-high-linear-Tor} already
gives the conclusion. Otherwise, at either object the local
minimal idempotent factors through the other object, using
the two minimal cross arrows. It therefore acts by zero on
a module supported at that object. The other minimal arrows
also act by zero, since one of their module spaces is zero.
Thus $K$ annihilates both $V$ and $W$.

Put $A=R/\mathbb QK$. The bimodule description
$\mathbb QK\cong Re\otimes_{\mathbb Q}eR$ in
Lemma~\ref{CP2:A-high-Rees} makes the ideal projective on both
sides. Since $We=0$, the resolution
$0\to\mathbb QK\to R\to A\to0$ and change of rings give
\[
 \operatorname{Tor}_j^R(W,V)\cong
 \operatorname{Tor}_j^A(W,V)\quad(j\ge0).
\]
The two object identities give distinct nonzero orthogonal
idempotents in $A$, so $A$ is nonlocal; its dimension is at
most $7-4=3$. By Lemma~\ref{CP2:A-high-small-Tor}, a nonzero
second Tor group would force nonzero Tor in arbitrarily high
degrees. The exact middle-bar shift identifies these with
homology in degrees two higher. Boundedness therefore forces
$H_4=\operatorname{Tor}_2^A(W,V)=0$.
\end{proof}

\begin{theorem}[Proved: cyclic four-object support through eighteen]
\label{CP2:A-strong-four-eighteen}
A finite category on four objects with cyclic support and at
most eighteen arrows cannot have classifying space of type $\CP^2$.
\end{theorem}
\begin{proof}
Assume a connected counterexample. It has bounded rational
homology and a one-dimensional second cohomology group with
nonzero square. Every actual arrow-reducing homotopy equivalence
would contradict the already proved global lower bound
Theorem~\ref{CP2:A-global-lower-eighteen}; in particular no
object is a retract of another.

A strong component on $t$ objects has its identities outside
the categorical minimal ideal and costs at least $t^2+t$ arrows.
A wholly strong four-object category therefore costs at least
twenty. A proper three-object component costs at least twelve
and its interface at least three, leaving a corner of order
at most three. The strong component has finite-group type:
otherwise its total row and column counts are at least four
and its internal cost is already nineteen. Both blocks have
zero second rational cohomology, and the two-block lemma applies.
If there are two strong pairs, the interface costs at least
four. Each pair then has at most eight internal arrows and
finite-group type, giving the same obstruction.

It remains to consider one strong pair $S$ and two singleton
components. Put $q=|S|\ge6$.

\emph{An extreme pair.} By opposition suppose $S$ is a source.
Let $D$ be its two-object complement and let $h$ count all
interface arrows. Connectedness gives $q\le13$ and $h\ge2$.
If both blocks have zero second rational cohomology, the
two-block lemma finishes.

First suppose $H^2(BD;\mathbb Q)\ne0$. Then $|D|\ge6$:
for connected $D$ this is the second-homology cost bound,
and for disconnected $D$ one corner has order at least five.
Consequently $q\le10$. The internal pair has finite-group
type, since total rows and columns both at least three
would cost at least eleven. If its total row count is two,
Lemma~\ref{CP2:A-source-minimal-group-collapse} is an actual
arrow-reducing equivalence, excluding the case even when its
minimal group is nontrivial.

Otherwise its total column count is two, its minimal group
is trivial, and its total row count is three or four. In
particular $q\ge8$. At $q=8$ it consists exactly of its
six-element minimal band and its two identities. The object
with one local row has all incoming internal arrows fixed by
its minimal idempotent and no units other than its identity.
There are no external entrances, so the one-sided deletion
of Lemma~\ref{CP2:A-high-one-sided} applies. If $h=2$ at
any of the remaining orders, every external stalk is singleton
or empty. The fixed-exit column version of the same lemma
collapses the whole pair and strictly decreases arrows.

The only remaining budget is $q=9$, $|D|=6$, $h=3$.
The complement is connected: two separate components would
require at least four interface arrows. The pair meets only
one object $u$ of $D$, and its other object $v$ points to $u$;
the opposite edge direction would force another interface.
Nonzero second cohomology on a connected six-arrow two-object
category forces corner orders $3,1$ and stalk size two.
Indeed a singleton stalk is a join; if both corners have
orders at most two the semisimple height-one bar has zero
second homology; and a corner of order at least four leaves
only a singleton stalk in this budget.

If the corner at $u$ has order three, all stalks into $u$
from the other three objects have size one or two: the two
from $S$ total three and the one from $v$ has size two.
Lemma~\ref{CP2:A-high-small-binary} applies with the arbitrary
source block $S\sqcup\{v\}$ and kills the square. If the
corner at $u$ is trivial, the full category on $S,u$ is an
ordinary act cone with one or two states at each object,
so Lemma~\ref{CP2:A-high-affine} kills its products. The
full category $D$ also has zero degree-two products by the
two-block lemma, since its corners have orders three and one.
Their nerves cover the whole category and meet at the point
$u$. Fourth-cohomology restriction is injective, so the
global products vanish too.

Now suppose $H^2(BS;\mathbb Q)\ne0$. Then $11\le q\le13$,
the minimal group is trivial, and both total row and column
counts are three. Thus each local row count is at most two,
and Corollary~\ref{CP2:A-baseline-two-row-instances} applies
to ordinary acts over $S$.
If $D$ is disconnected, connectedness requires four interface
stalks, giving $|D|\le3$. Both outside corners have order at
most two. Their exact orbit functors replace them by trivial
leaves. Merging those two vertices is an isomorphism on
cohomology in degrees at least two, respecting products: it
adds one circle to the connected realization. The result is
an ordinary act cone over $S$, whose degree-two products are
zero by the binary projective-cover lemma. Pullback excludes
the original category.

If $D$ is connected and $S$ meets both objects, the budget
forces $q=11$, $|D|=3$, $h=4$. All cross arrows between
$S$ and $D$ are singleton, so this is a categorical join
and has zero positive products. Otherwise $S$ meets only
the final object $u$ of the outside edge $v\to u$. Let $E$
be the full category on $S,u$. The corner at $u$ has order
at most three, since
\[
 |\operatorname{End}(u)|\le18-q-h-2\le3.
\]
The nerves of $E$ and $D$ cover the whole nerve and meet
at this rationally acyclic corner. Fourth restriction is
injective. Also $|D|\le5$, so $H^2(BD;\mathbb Q)=0$.
If the corner at $u$ has order at most two, exact orbits
and the binary projective-cover lemma kill the products
of $BE$. If its order is three, equality in the budget
forces $q=11,h=2,|D|=5$, and $E$ is a singleton-stalk
join with zero products. Restriction therefore kills the
global square in every case. The two extreme branches
exhaust the possibilities, since nonzero second cohomology
on both blocks would require at least $11+6+2=19$ arrows.

\emph{A nonextreme pair.} The condensation is the full chain
$u<S<v$. Let the extreme corner orders be $b,c$, put
$z=|C(u,v)|$, and write the two entrance and exit sizes
as $\alpha_i,\beta_i$ for $i=0,1$. All are positive. Set
\[
 P=\sum_i(\alpha_i-1)+\sum_i(\beta_i-1),\qquad
 \delta=(b-1)+(c-1)+(z-1).
\]
The complete arrow count is
\[
 q+7+P+\delta\le18,\qquad q\le11,
 \qquad P+\delta\le11-q.
\]

For $q\ge8$ and $z\le2$, both endpoint corners have order
at most four. Failure of the binary-stalk test at each end
would require one entrance of size at least three and one
exit of size at least three, costing $P\ge4$. This is
impossible because $P\le3$. Thus one endpoint satisfies
Lemma~\ref{CP2:A-high-small-binary}. For $q\ge8$ and
$z\ge3$, the budget gives $P\le1$ and both endpoint corners
have order at most two. Apply their actual exact-orbit
functors to make them trivial. At least one whole entrance
or exit side originally consisted of singletons, and still
does. Its augmentation module is zero. The middle-bar
formula gives $H_4=0$, again a contradiction.

It remains that $q=6$ or $7$. The minimal group is trivial,
there is one minimal row and column at each object, and
the local minimal idempotents $e_i$ are corner zeros.
Consider the four conditions that $e_i$ fix all incoming
cross arrows, or all outgoing cross arrows. Any one of
these conditions permits the one-sided object deletion;
the zero guarantees its unit condition.

At $q=6$, all internal cross arrows are fixed. If no
condition holds, all four external stalks must be
nonsingleton, costing $P\ge4$. Then $P+\delta\le5$ implies
both endpoint corners and the direct stalk have size at
most two. An entrance and an exit both of size at least
three, in addition to the other two nonsingleton stalks,
would cost $P\ge6$. Hence at least one endpoint has all
three incident stalks binary, and the small-binary lemma
excludes its square.

At $q=7$, the single extra internal arrow can violate at
most two of the four fixed conditions, and a nonsingleton
external stalk can violate only its one incident condition.
Thus $P\le1$ gives a one-sided deletion. Assume $P\ge2$.
If $z\le2$, both endpoint corners have order at most four.
Unless a binary endpoint already applies, one entrance
and one exit have size at least three, so the budget
forces $P=4$, $\delta=0$. Thus $b=c=z=1$, and both the
entrance pair and the exit pair have sizes $(3,1)$ in
some order. Their augmentation modules are each supported
on one object of $S$. Lemma~\ref{CP2:A-small-pair-augmentation-Tor}
gives $H_4=0$ under boundedness.

Finally $q=7$, $z\ge3$, $P\ge2$ forces
$P=2$, $b=c=1$, $z=3$. If one whole side is singleton,
its augmentation module is zero and the middle-bar
formula finishes. Otherwise each side has excess one,
so each augmentation module is one-dimensional and
supported at one object. The same small-pair Tor lemma
applies. This exhausts every nonextreme budget, and
therefore every cyclic four-object category in the
eighteen-arrow range.
\end{proof}


\subsection{An actual reduction at extreme corners}

\begin{lemma}[Proved: an extreme one-row corner reduces to its minimal group]
\label{CP2:A-extreme-minimal-group-reduction}
Let $x$ be a source object in a finite category $C$: no other
object has an arrow to $x$. Write $M=C(x,x)$, choose an
idempotent $e$ in its minimal ideal, and put $G=eMe$.
If $Me=G$, there is a finite category $C'$ of the same
homotopy type obtained by replacing the corner $M$ by $G$
and every outgoing stalk $C(x,y)$ by $C(x,y)e$. In particular,
\[
 |\operatorname{Mor}C'|\leq
 |\operatorname{Mor}C|-\bigl(|M|-|G|\bigr).
\]
The hypothesis $Me=G$ is equivalent to the minimal ideal
having one row. There is a dual statement for a sink object
with $eM=G$, equivalently one minimal column.
\end{lemma}
\begin{proof}
Adjoin the retract object $e$ to $C$, inside its idempotent
completion, and call the resulting category $C^e$. The old
object $x$ retracts onto the added object, so the inclusion
$BC\to BC^e$ is a homotopy equivalence by the retract-object
lemma. Let $C'$ be the full category obtained from $C^e$
by deleting $x$.

Apply Quillen's Theorem~A to the inclusion $i:C'\to C^e$,
using the comma categories $(i\downarrow z)$. If $z$ is
retained, its identity is terminal in the comma category.
For $z=x$, no old retained object has an arrow into $x$.
Thus the comma category is precisely the right action
category of $G$ on
\[
 C^e(e,x)=Me=G.
\]
Right multiplication is free and transitive. Consequently
this comma category has exactly one morphism between each
ordered pair of objects and is contractible. The theorem
gives $BC'\simeq BC^e\simeq BC$.

The retained old hom-sets are unchanged, the new corner is
$G$, and the new outgoing hom-set to $y$ is $C(x,y)e$.
These are subsets of the corresponding old hom-sets, and
there are no new incoming arrows. This proves the arrow
inequality. The finite minimal-ideal description identifies
$Me$ with a disjoint union of one free transitive $G$-set
for each minimal row, proving the stated equivalence of
hypotheses. Opposition proves the sink assertion.
\end{proof}

\begin{corollary}[Proved: the extreme-corner cut at the next arrow budget]
\label{CP2:A-extreme-minimal-group-eighteen}
An eighteen-arrow category of type $\CP^2$ cannot have
a source corner which is a nongroup monoid with one minimal
row, or a sink corner which is a nongroup monoid with one
minimal column.
\end{corollary}
\begin{proof}
For a nongroup monoid, its minimal group has strictly fewer
elements: equality would imply $M=G$ as a set and hence as
a monoid. The preceding actual homotopy equivalence would
give a category with at most seventeen arrows, contrary
to Theorem~\ref{CP2:A-global-lower-eighteen}.
\end{proof}

The literal table predicate uses only minimal-ideal data.
For an idempotent $e$ in a corner's minimal ideal it computes
the sets $Me$, $eM$, and $eMe$ directly. At a source, the
test is $|Me|=|eMe|<|M|$; at a sink it is
$|eM|=|eMe|<|M|$. Since $eMe$ is contained in both other
sets, equality of cardinalities is equality of sets.
This test involves neither a homology approximation nor
an assumption that the chosen corner action is trivial.


\subsection{A rational pendant reduction for a three-object triangle}

\begin{lemma}[Proved: a singleton pendant at an acyclic corner]
\label{CP2:A-rational-pendant-deletion}
Let $p$ be an object of a finite category $C$ whose only
strictly adjacent object is $q$. Suppose their unique nonempty
strict hom-set is a singleton and $BC(q,q)$ is rationally
acyclic. Let $E$ be the full subcategory obtained by deleting
$p$. Then the inclusion $BE\to BC$ is a rational homology
equivalence and induces a rational cohomology-ring isomorphism.
There is no restriction on the monoid $C(p,p)$.
\end{lemma}
\begin{proof}
There are no arrows between $p$ and any object other than $q$,
and arrows between $p$ and $q$ go in one direction. A nerve
simplex that contains $p$ can therefore use no object besides
$p$ and $q$: any further strict object change would force an
additional arrow incident to $p$ by composition. Thus
\[
 BC=BE\cup B C\{p,q\},\qquad
 BE\cap B C\{p,q\}=BC(q,q),
\]
as a union of nerve subcomplexes. The singleton cross hom-set
makes the two-object category the categorical join of its
two corner monoids. Consequently
\[
 B C\{p,q\}\simeq BC(p,p)*BC(q,q).
\]
The join with a nonempty rationally acyclic space is rationally
acyclic, by the reduced homology formula for a join over the
field $\mathbb Q$. Both the intersection and the attached
piece therefore have zero positive rational homology and are
connected. The Mayer--Vietoris sequence proves that
$BE\to BC$ is a rational homology equivalence. The induced
cohomology map is multiplicative, proving the ring assertion.
\end{proof}

\begin{corollary}[Proved: a singleton pendant over a singleton-long triangle]
\label{CP2:A-pendant-singleton-long-triangle}
Suppose a finite four-object category has a singleton pendant
as in the preceding lemma. If its three retained objects
form full linear acyclic support, their long hom-set is a
singleton, and their three corner classifying spaces are
rationally acyclic, every positive rational cohomology product
of the whole category is zero. In particular it suffices
that the retained three corner orders are at most four.
\end{corollary}
\begin{proof}
The pendant's neighbor is retained, so its corner is rationally
acyclic. The lemma identifies the rational cohomology ring
with that of the retained triangle. Theorem
\ref{CP2:A-baseline-singleton-long} annihilates every positive
product in that ring. The corner-order consequence uses
Lemma~\ref{CP2:A-five-types}.
\end{proof}

For example, the frozen four-object eighteen-arrow profile
with index $175$ has hom-set cardinality matrix
\[
 \begin{pmatrix}
 1&0&0&1\\
 0&1&2&1\\
 0&0&3&7\\
 0&0&0&2
 \end{pmatrix}.
\]
Delete the pendant object $0$. The retained triangle has
corner orders $1,3,2$ and singleton long hom-set $1\to3$,
so the corollary excludes the entire profile without any
composition enumeration or any assumption on the sink's
two-element monoid table.


\begin{corollary}[Proved: bounded triangle gates after rational pendant deletion]
\label{CP2:A-pendant-bounded-triangle-gates}
Use the hypotheses and notation of
Lemma~\ref{CP2:A-rational-pendant-deletion}, and suppose the
three retained objects form full linear acyclic support.
Write their corner orders as $a,b,c$, and their first
adjacent, second adjacent and long hom-set sizes as $x,y,z$.
Assume $BC$ has bounded rational homology and
$\dim H^2(BC;\mathbb Q)\leq1$. Every rational degree-two
product is zero if any one of the following holds:
\begin{enumerate}
\item $a,c\leq2$ and $b\leq4$;
\item $a,b\leq2$, $c\leq6$, or $b,c\leq2$, $a\leq6$;
\item $a\leq2$, $x=1$, $c\leq11$, or
      $c\leq2$, $y=1$, $a\leq11$;
\item $a\leq4$ and $x,z\leq2$, or
      $c\leq4$ and $y,z\leq2$.
\end{enumerate}
\end{corollary}
\begin{proof}
The rational pendant comparison transfers bounded homology,
the stated second-cohomology dimension bound, and the entire
cohomology ring to the retained triangle $E$.

In the first case, the two extreme corner algebras are
semisimple. The exact orbit comparisons of
Lemma~\ref{CP2:A-high-orbits} replace them by trivial
corners, inducing rational homology equivalences. All
three strict hom-sets remain nonempty. Lemma
\ref{CP2:A-high-linear-Tor}, including its small-middle
consequence, gives $H_4(BE;\mathbb Q)=0$ under boundedness.
This annihilates every degree-two product.

The second case is precisely Corollary
\ref{CP2:A-baseline-semisimple-pair-six}, which has no
restriction on the strict hom-set sizes and uses the
transferred boundedness. The third is Corollary
\ref{CP2:A-eleven-projective-singleton}: the indicated
corner of order at most two has projective trivial module
on the required side. The other orientation follows by
opposition.

For the fourth case, take the extreme corner as the
one-object block and the other two objects as its
complement. Both incident hom-sets are nonempty and
have at most two elements. The extreme corner has order
at most four, and the transferred dimension bound is
exactly the premise of Lemma~\ref{CP2:A-high-small-binary}.
It gives zero square for every second class; rational
polarization gives zero for every pairwise degree-two
product. Pullback along the pendant comparison finishes
all four cases.
\end{proof}


\subsection{Gluing small triangles along an acyclic corner or edge}

\begin{lemma}[Proved: boundedness and products through a low-homology intersection]
\label{CP2:A-low-homology-gluing}
Suppose full subcategories $E,F$ of a finite category $C$
cover its nerve, and their intersection $I$ satisfies
$H_j(BI;\mathbb Q)=0$ for every $j\geq2$.
Then restriction gives an isomorphism
\[
 H^n(BC;\mathbb Q)\longrightarrow
 H^n(BE;\mathbb Q)\oplus H^n(BF;\mathbb Q)
 \quad(n\geq3)
\]
and a surjection in degree two. Consequently bounded
rational homology of $BC$ passes to both pieces, and
each piece's second Betti number is at most that of $BC$.
If both pieces have zero rational degree-two products,
then so does $BC$.
\end{lemma}
\begin{proof}
The realizations are CW subcomplexes and form a pushout
along their intersection. Apply the rational cohomology
Mayer--Vietoris sequence. For $n\geq3$ both adjacent
groups $H^{n-1}(BI;\mathbb Q)$ and $H^n(BI;\mathbb Q)$
are zero, giving the isomorphism. The group after
$H^2(BE)\oplus H^2(BF)$ is $H^2(BI)=0$, giving the
surjection. Over a field, the corresponding homology
and cohomology vanishing statements agree. These facts
give boundedness and the asserted dimension inequalities.
Finally restriction is multiplicative and is injective
in degree four. A degree-two product whose restrictions
to both pieces vanish must therefore vanish on $BC$.
\end{proof}

In the next two statements a \emph{permitted triangle}
means a full-linear three-object category with profile
$(a,b,c;x,y,z)$ satisfying at least one of these
cardinality conditions:
\[
\begin{gathered}
 z=1,\quad \max(a,b,c)\leq4;\\
 a,c\leq2,\quad b\leq4;\\
 a,b\leq2,\ c\leq6,\quad\hbox{or}\quad b,c\leq2,\ a\leq6;\\
 a\leq2,\ x=1,\ c\leq11,\quad\hbox{or}\quad
 c\leq2,\ y=1,\ a\leq11;\\
 a\leq4,\ x,z\leq2,\quad\hbox{or}\quad
 c\leq4,\ y,z\leq2.
\end{gathered}
\]
Every such triangle with bounded rational homology and
second Betti number at most one has zero degree-two
products. For the five lines, respectively, apply
Theorem~\ref{CP2:A-baseline-singleton-long},
Lemma~\ref{CP2:A-high-linear-Tor},
Corollary~\ref{CP2:A-baseline-semisimple-pair-six},
Corollary~\ref{CP2:A-eleven-projective-singleton}, and
Lemma~\ref{CP2:A-high-small-binary}. The last use takes
the indicated endpoint as one block and the other two
objects as the arbitrary opposite block. Its second-rank
premise is the stated Betti bound; polarization over
$\mathbb Q$ converts zero squares to zero pairwise products.
These are the same underlying triangle statements used
in Corollary~\ref{CP2:A-pendant-bounded-triangle-gates}.

\begin{theorem}[Proved: a pendant edge attached to a permitted triangle]
\label{CP2:A-pendant-edge-permitted-triangle}
Let $C$ have four objects and acyclic support. Suppose
one object $p$ is strictly adjacent only to $q$, the
corner at $q$ has order at most four, and the retained
three objects form a permitted triangle. Let $s$ be the
size of the nonempty hom-set between $p$ and $q$.
If $s\leq2$ or the corner at $p$ also has order at most
four, bounded rational homology and second Betti number
at most one force every rational degree-two product of
$BC$ to be zero.
\end{theorem}
\begin{proof}
Let $E=C\setminus\{p\}$ and $F=C|_{\{p,q\}}$. The absence
of any other strict arrow incident to $p$ ensures that
every nerve simplex lies in $E$ or $F$: a strict chain
containing $p$ and a third object would compose to such
an additional arrow. The intersection is the corner at
$q$, whose classifying space is rationally acyclic by
the small-monoid lemma. Therefore
Lemma~\ref{CP2:A-low-homology-gluing} applies. In fact
here restriction is an isomorphism also in degree two,
since the intersection has zero first rational cohomology.
The permitted triangle $E$ inherits bounded homology
and second Betti number at most one, so its degree-two
products vanish by the five triangle statements above.

For the attached two-object piece $F$, a singleton stalk
is a categorical join and has zero positive products.
A two-element stalk has zero degree-two products by
Lemma~\ref{CP2:A-high-two-cross}, with no corner assumption.
In the remaining case both corner orders are at most
four, so both corner second rational cohomology groups
are zero. The two-block lemma again kills every
degree-two product on $BF$, regardless of its stalk size.
The degree-four restriction injection in the gluing
lemma now kills every such product on $BC$.
\end{proof}

\begin{theorem}[Proved: two permitted triangles meeting in a small edge]
\label{CP2:A-two-permitted-triangles}
Let $C$ have four objects and acyclic support. Suppose
two objects $p,r$ are incomparable, both full subcategories
$E=C\setminus\{p\}$ and $F=C\setminus\{r\}$ are
permitted triangles, and their two-object intersection
$I$ has a nonempty strict hom-set. Assume either that
both corner orders in $I$ are at most two, or that its
strict hom-set is singleton and at least one of its
corner orders is at most four. Bounded rational homology
and second Betti number at most one then force every
rational degree-two product on $BC$ to be zero.
\end{theorem}
\begin{proof}
A nerve simplex cannot contain both $p$ and $r$, since
composing its strict arrows would make these objects
comparable. Thus $BE$ and $BF$ cover $BC$ with intersection
$BI$. If both intersection corners have order at most
two, their rational algebras are semisimple. The relative
bar complex of the two-object edge has strict height
one, hence $H_j(BI;\mathbb Q)=0$ for $j\geq2$.
In the other case $BI$ is the join of its corner spaces;
at least one is rationally acyclic, so their join is
rationally acyclic as well.

The gluing lemma passes bounded homology and the
second-rank bound to each triangle. The five triangle
statements give zero degree-two products on both pieces.
Injectivity of degree-four restriction gives the result.
\end{proof}

\begin{corollary}[Proved: actual deletion of a contractible-neighbor pendant]
\label{CP2:A-contractible-pendant-deletion}
If a singleton pendant $p$ has neighbor $q$ with contractible
corner classifying space, the inclusion
$B(C\setminus\{p\})\to BC$ is an actual homotopy equivalence.
In particular an eighteen-arrow projective-plane model
cannot have such a pendant when the neighbor's minimal
group is trivial and its minimal ideal has one row or
one column.
\end{corollary}
\begin{proof}
The same nerve union has a contractible intersection and
contractible attached join, so the inclusion of the
retained piece is a homotopy equivalence. A monoid with
one minimal row or one minimal column is homotopy
equivalent to its minimal group, by
Lemma~\ref{CP2:A-extreme-minimal-group-reduction} applied
to its one-object category or its opposite. A trivial
minimal group therefore makes the corner contractible.
Deleting $p$ removes its identity and its singleton
stalk, contradicting the global bound through seventeen.
\end{proof}


\begin{proposition}[Proved: the additional gluing partition at four objects]
\label{CP2:A-four-eighteen-gluing-scope}
Among the $2724$ four-object eighteen-arrow profiles left
by Proposition~\ref{CP2:A-four-five-eighteen-scopes}, exactly
$1176$ satisfy the gluing criteria of
Theorems~\ref{CP2:A-pendant-edge-permitted-triangle} and
\ref{CP2:A-two-permitted-triangles}. They consist of $1001$
pendant-edge profiles and $175$ two-triangle profiles.
The original $61420$-profile scope therefore has the
further partition
\[
 61420=59872\text{ structurally excluded}
       +1548\text{ remaining finite profiles}.
\]
This proposition does not assert that those remaining
composition searches have completed.
\end{proposition}
\begin{proof}
The fresh program
\path{converse/scope18/four18_gluing_supplement.py}
reads the unchanged $2724$-row ledger and its complete
independent scope receipt. It tests every possible pendant
object and every possible pair of incomparable exclusive
objects. For a pendant it separately requires the neighbor
order to be at most four, and checks either the stalk bound
two or the other corner bound four. For two triangles it
tests the intersection's semisimple-height-one or acyclic-join
condition. Every retained triangle is then tested against
the five displayed cardinality gates in the preceding
proofs, in the appropriate topological order.

The output records the original index, the full cardinality
matrix, and every literal witness and applicable reason.
The two support shapes are disjoint and give $1001+175=1176$
profiles. The frozen earlier $282$ singleton-pendant profiles
are explicitly checked to be a subset, with all their old
matrix and proof-source bindings preserved. The remaining
$1548$ rows are written in their original order and their
original indices are retained. Adding the $1176$ new proofs
to the original $58696$ structural exclusions gives the
displayed partition. All previous ledgers, producer sources,
partial attempts and numerical receipts remain unchanged.
The full supplement is
\path{converse/scope18/four18_gluing_supplement.json}.
A separately written root checker,
\path{code/check_gluing1176_scope.py}, independently
regenerates every witness and reason without importing
the supplement program. It matches all $1176$ exclusions,
the prior $282$-profile inclusion, and the exact residual
ledger bytes. Its complete passing receipt is
\path{outputs/gluing1176_root_independent.json}.
\end{proof}


\subsection{Small-leaf refinements of the eighteen-arrow star scope}

\begin{lemma}[Proved: extending the binary and four-state leaf gates]
\label{CP2:A-star-small-leaf-orbit-extension}
Consider a source star with center $M$ and two nonempty stalks.
Suppose one leaf corner has order at most two. The star cannot
have the rational homology and cohomology ring of $\CP^2$ in
either of the following cases:
\begin{enumerate}
\item $|M|\leq11$ and the other stalk has at most two states,
with no restriction on that other leaf corner;
\item $|M|\leq8$, the other leaf corner has order at most three,
and its stalk has at most four states.
\end{enumerate}
The stalk at the first leaf can have any positive size.
\end{lemma}
\begin{proof}
The rational algebra of a monoid $L$ of order at most two is
semisimple: besides the trivial monoid the possibilities are
$C_2$ and $\{1,e\}$ with $e^2=e$, and in both cases the
algebra is $\mathbb Q\times\mathbb Q$. In particular its
trivial right module is projective. Replace this leaf by a
trivial corner and its left-$L$ stalk $T$ by the orbit set
$L\backslash T$. This is an actual functor, and exact orbit
comparison gives a rational cohomology-ring and homology
isomorphism. More explicitly, the bridge bar
$\mathbb Q\otimes^{\mathbf L}_{\mathbb QL}\mathbb Q[T]
\otimes^{\mathbf L}_{\mathbb QM}\mathbb Q$ is naturally
identified with the ordinary $M$ action bar on $L\backslash T$;
the $L$ corner is rationally acyclic. The other edge and its
intersection $BM$ are unchanged, so the mapping-cylinder
comparison gives the same isomorphism for the whole star.

The orbit stalk remains nonempty and the other stalk does not
change. Apply Proposition~\ref{CP2:A-star-binary-leaf-eleven}
in the first case and
Proposition~\ref{CP2:A-star-three-leaf-four-states} in the
second. Both propositions exclude the rational target, which
is exactly what the orbit comparison transfers.
\end{proof}

\begin{lemma}[Proved: a singleton leaf with a rationally acyclic corner]
\label{CP2:A-star-singleton-small-leaf-retract}
Let a source star have a leaf on a singleton stalk, with
nontrivial leaf monoid $L$ such that $BL$ is rationally acyclic.
Replacing $L$ by the trivial monoid is an actual categorical
retract and a rational homology and cohomology-ring equivalence.
If the original star has type $\CP^2$, so does the retract.
Consequently an eighteen-arrow projective-plane star cannot
have a singleton stalk at a leaf corner of order two, three,
or four.
\end{lemma}
\begin{proof}
All corner actions on the singleton stalk are trivial. Send
every element of $L$ to the identity of the replacement corner,
and keep every other arrow unchanged. This defines a functor;
its section sends the new identity to $1_L$ and fixes all
other arrows. Their composite on the replacement category is
the identity.

Write $E$ for the other full edge. The nerve of the star is
the union of $BE$ and the singleton-edge nerve, meeting in
$BM$, where $M$ is the center. The singleton-edge space is
$BM*BL$, and its replacement is $BM*\{*\}$, a cone. Both
are rationally acyclic: for the first this follows from the
join homology formula and the rational acyclicity of $BL$.
The replacement map is the identity on $BE$ and $BM$.
Mayer--Vietoris therefore makes the actual retraction a
rational homology equivalence in every degree. Its induced
cohomology isomorphism respects products.

Suppose the original space has type $\CP^2$. The categorical
retraction makes the retract's integral homology a direct
summand of the original integral homology, and its fundamental
group a retract of the trivial group. The rational equivalence
forces the degree-two and degree-four integral summands to be
$\mathbb Z$, and all other reduced integral groups vanish.
The split pullbacks in those two degrees are isomorphisms;
naturality consequently transfers the primitive square.
The integral recognition argument in
Theorem~\ref{A:monoid-upper23} gives type $\CP^2$ for the
retract. The replacement removes $|L|-1\geq1$ arrows, so at
eighteen arrows this contradicts
Theorem~\ref{CP2:A-global-lower-eighteen}.
Finally all monoids of orders at most four have rationally
acyclic classifying spaces by Lemma~\ref{CP2:A-five-types}.
\end{proof}

\begin{proposition}[Proved: fifteen further frozen star profiles]
\label{CP2:A-star-eighteen-small-leaf-scope}
In the frozen sixty-one-profile star scope at eighteen arrows,
the original indices
\[
14,17,19,22,24,26,38,40,42,43,44,49,57,58,59
\]
are excluded by the preceding two lemmas. This gives a separate
partition of the $1204$ canonical cardinality profiles into
$1158$ proved profiles and $46$ profiles requiring finite
composition checks. It does not assert that those forty-six
finite checks have all finished.
\end{proposition}
\begin{proof}
The separately frozen program
\path{converse/scope18/star18_small_leaf_supplement.py}
tests each original matrix in both leaf orders. It records
the actual small leaf, the relevant stalk sizes, and the
specific lemma clause. The singleton-corner predicate tests
orders two through four and singleton stalk size. The orbit
predicates test a leaf order at most two and, respectively,
the exact bounds $(|M|,|S|)\leq(11,2)$ or
$(|M|,|N|,|S|)\leq(8,3,4)$ on the retained center, leaf and
stalk. These are precisely the hypotheses proved above.
Exactly the fifteen displayed original rows qualify. The
old ledger and its $1143$ proved rows are unchanged, and
$1143+15=1158$, $61-15=46$.
\end{proof}



\subsection{The independently certified unrestricted lower bound}
\begin{theorem}[Proved: the global lower bound of seventeen arrows]
\label{CP2:A-global-lower-seventeen}
Every finite category whose classifying space is homotopy equivalent
to $\CP^2$ has at least seventeen arrows, including identities.
Consequently its least possible arrow count $n$ satisfies
\[
 17\le n\le22.
\]
\end{theorem}
\begin{proof}
Suppose there are at most sixteen arrows. The category is connected.
For one object use the complete monoid minimum
Theorem~\ref{A:monoid-minimum23-classification}; for two objects use
Theorem~\ref{CP2:A-two-object-baseline-sixteen}.
For three objects, cyclic support is excluded by
Theorem~\ref{CP2:A-strong-three-seventeen}. A connected acyclic
support on three objects is either a full linear order or a star:
a directed two-step path forces its composite edge, and otherwise
there are precisely two support edges sharing their source or sink.
Theorems~\ref{CP2:A-linear-sixteen-complete} and
\ref{CP2:A-star-sixteen-complete} exclude these alternatives.
For four objects use the full acyclic-support census
Theorem~\ref{CP2:A-baseline-four-census} and the separate cyclic-support
Theorem~\ref{CP2:A-strong-four-seventeen}.
For $k\ge5$ objects, $16\le2k+7$, so
Theorem~\ref{CP2:A-high-all-eight} applies. These results explicitly
include object-retract reductions and the small corner orders.
They exhaust all finite categories under the bound.
The upper bound is the independently checked twenty-two-arrow
category in Theorem~\ref{A:category-upper22}.

The two final three-object joins cover respectively all $1547$
linear and $693$ star profiles at exactly sixteen arrows.
A separate root audit reconstructs both complete sets directly
from all positive integer tuples, checks equality with their
joined sets, checks every padding changes exactly one corner
order by one, and binds both joins by SHA-256. Its receipt is
\path{outputs/global16_scope_independent.json}. Their literal
finite gates and the four-object gate have separately verified
cup or third-homology obstructions. Thus this lower bound does
not presume the missing inherited certificate files.
\end{proof}


\subsection{All degree-two products for a three-by-three minimal ideal}

\begin{theorem}[Proved: the full finite-type three-by-three obstruction]
\label{CP2:A-three-by-three-all-products}
Let $M$ be a finite monoid whose minimal ideal has three rows and
three columns, with any finite maximal group. If $BM$ has finite
CW homotopy type, then every rational product
\[
 H^2(BM;\mathbb Q)\otimes H^2(BM;\mathbb Q)
 \longrightarrow H^4(BM;\mathbb Q)
\]
is zero. No hypothesis on the dimension of second homology is imposed.
\end{theorem}
\begin{proof}
We extend the one-dimensional balanced-form calculation used in
Proposition~\ref{A:three-row-classification}; the higher-dimensional
cases are included explicitly here. First suppose the minimal ideal
is the band $K=I\times J$, $|I|=|J|=3$. For a transformation $f$
of a three-point set its augmentation matrix is
\[
 A_f(i,j)=\mathbf1_{f(j)=i}-\mathbf1_{f(0)=i},
 \qquad 1\le i,j\le2.
\]
Write an actor as $(\sigma,\tau)$ with multiplication
$(\sigma,\tau)(\alpha,\beta)=(\sigma\alpha,\beta\tau)$.
Let $R$ be the rational span of all coefficient relations in
$B_\tau^tL=LA_\sigma$ for the actors of $M$. If $R=\mathbb Q^4$
there is no second rational cohomology. Otherwise form the monoid
$H(R)$ of all actors whose relation space is contained in $R$.
The actual action homomorphism $M\to H(R)$ fixes the band.
Its relation space is exactly $R$: it contains all original actors
and introduces no additional relation. The balanced-tensor
identification of Lemma~\ref{A:band-Tor} consequently makes its
map on rational second homology an isomorphism. Every degree-two
class on $BM$ is therefore pulled back from this specified hull.

The previously independently checked relation-lattice computation
tests all $729$ actors, closes their $81$ distinct proper relation
spaces under every proper sum, and obtains $111$ spaces and $13$
orbits under row and column permutations. Opposition is not used
in this numbering. The integer relation bases and complete actor
sets are in \path{monoid/certificates/relation_lattice_3x3.json};
its separate checker establishes both closure and generation by
the displayed actors. Thus it covers every possible $R$, including
second-homology dimensions greater than one.

The six one-dimensional cases have orbit numbers $7$ through $12$.
For $7,8,9$ the balanced matrix has rank one, so
Lemma~\ref{A:rank-one-square} gives zero square. For $10,12$ the
already checked integral projective prefixes give zero rational
fourth homology of the hull. Orbit $11$ has $30$ actors and one
integral balanced cocycle. Exactly eight actors have a nonzero
restriction of this class to their eventual cyclic group. None
can occur in the image of the original finite-CW monoid: its
eventual cyclic group surjects onto that image group, so inflation
preserves the nonzero character, contradicting the line-bundle
nilpotence argument proved in Proposition~\ref{A:three-row-classification}.
This argument is applied to the original $BM$, not to the hull.
The independently checked closure enumeration gives $154$ allowed
closed subsets, each contained in one of nine maximal submonoids
of order $16$. All nine have verified zero rational fourth
homology. Factoring the action through such a submonoid kills
the pulled-back square. These arguments use the one-dimensional
balanced-form condition and finite CW type only; they do not
require the target to have the integral homology of $\CP^2$.

For clarity, the seven remaining hulls in the same numbering are
as follows. A four-tuple lists a $2\times2$ matrix by rows.
\[
\begin{array}{c|r|r|l}
\text{orbit}&\dim H^2&|H|&\text{balanced-form basis}\\\hline
0&4&10&\text{the four matrix units}\\
1&2&31&(1,0,0,0),\ (0,1,0,0)\\
2&2&12&(1,0,0,0),\ (0,1,1,0)\\
3&2&14&(1,0,0,0),\ (0,1,1,1)\\
4&2&31&(1,0,0,0),\ (0,0,1,0)\\
5&2&11&(1,0,0,0),\ (0,1,1,2)\\
6&2&12&(1,1,1,0),\ (-1,0,0,1).
\end{array}
\]
In orbits $1,4$ every matrix in the entire displayed space has
rank at most one. Every square is zero by
Lemma~\ref{A:rank-one-square}, and rational polarization kills
all products. For $0,5,6$ the actors outside the minimal band
form respectively the groups $1,C_2,C_3$. Their contracted
algebras $\mathbb QH/\mathbb QK$ are the corresponding rational
group algebras. Averaging splits every module surjection, so
these algebras are semisimple and have zero positive Tor.
Lemma~\ref{A:band-Tor} gives $H_4(BH;\mathbb Q)=0$.

For orbit $3$ we construct a fresh integral projective prefix
through degree five. Its projective abelian ranks and
coinvariant dimensions are respectively
\[
 (3,10,16,16,16,16),\qquad(1,0,2,2,2,2).
\]
The separate literal checker reconstructs every principal
idempotent summand and its action, verifies every module-map
equation, verifies zero composites, and checks all image Smith
invariants and complementary kernel ranks. It proves exactness
through degree four. The induced coinvariant differential gives
Betti numbers $(1,0,2,0,0)$ through degree four. In particular
this hull has zero rational fourth homology, without any
assertion about an uncomputed resolution tail.

Orbit $2$ has nonzero higher homology, so we check its products
directly. Let $u,v$ be the affine cocycles for its two displayed
balanced forms. Three fresh integral projective/bar certificates
use the forms for $u,v,u+v$. Their common projective ranks and
coinvariant dimensions are
\[
 (3,9,15,21,39),\qquad(1,0,2,1,1).
\]
Each prefix is exact through degree three. On degree-two
coinvariants the three cocycles are $(1,0),(0,1),(1,1)$;
on the entire degree-four coinvariant module each square is
the zero cochain. The checker verifies the comparison with
the normalized integral bar resolution on all $13$ projective
generators and their $87$ displayed translates, together with
all $1728$ cocycle and associativity equations per certificate.
The comparison extends to a projective resolution because
the prefix is exact in the preceding degrees. Thus these are
the actual cup classes and
$[u]^2=[v]^2=([u]+[v])^2=0$ on the hull. Polarization over
$\mathbb Q$ gives $[u][v]=0$ and hence all products vanish.

The constructor is \path{code/extend_three_by_three_products.py}.
The four new literal certificates have prefix
\path{monoid/certificates/full_three_by_three_}.
The independent projective checker is the one already used
for the monoid upper and lower bounds; it imports no constructor.
The additional join \path{code/check_three_by_three_extension.py}
checks all seven higher-dimensional cases, matches every new
table and form to the complete relation lattice, checks the
three outside group tables and the two rank-one form spaces,
and reruns every new literal prefix check. Its complete receipt
is \path{outputs/full_three_by_three_products_independent.json}.

Finally suppose the original minimal ideal has a nontrivial
maximal group. The actual minimal-group quotient in
Lemma~\ref{A:minimal-group-collapse} induces a rational
cohomology-ring isomorphism and leaves the three row and column
sets unchanged. Apply the same hull classification to its
action. All zero-product and zero-fourth-homology cases pull
back to the original monoid. In the exceptional orbit $11$,
pull the integral hull cocycle directly to the original $M$;
the eventual local group there still surjects onto the forbidden
cyclic image group. The same nilpotence contradiction applies
to the original finite-CW space. No finite-type assumption on
the intermediate minimal-group quotient is needed.
\end{proof}

\subsection{Collapsing the long arrows of a triangle}

\begin{lemma}[Proved: the actual long-arrow collapse]
\label{CP2:A-triangle-long-collapse}
Let $C$ be a three-object category with full linear support
$0<1<2$ and trivial corners at $0,2$. Write
$z=|C(0,2)|\ge1$. Let $\overline C$ be the category obtained
by replacing $C(0,2)$ by a singleton and sending every
two-cross-arrow composite to that singleton. The identity
on all other arrows and this collapse define a functor
$\pi:C\to\overline C$. It induces a surjection
\[
 \pi^*:H^2(B\overline C;\mathbb Q)\longrightarrow H^2(BC;\mathbb Q).
\]
If $BC$ has finite CW homotopy type, so does $B\overline C$.
If in addition $BC$ is simply connected, then
\[
 B\overline C\simeq BC\vee\bigvee^{z-1}S^2.
\]
\end{lemma}
\begin{proof}
The composition and functor laws follow directly: every equation
with a long-arrow value becomes the unique such arrow, and
every other equation is unchanged. No nonidentity chain of
length at least two contains a long arrow, because the endpoint
corners are trivial. Thus the normalized nerve has exactly the
same simplices in every degree at least two; only the long-edge
faces of two-simplices are identified. The degree-two cocycle
spaces are identical and pullback identifies the smaller
degree-two coboundary space with a subspace of the old one.
This proves the cohomological surjection, with its actual
cup-preserving map.

Let $A\subset BC$ be the finite graph consisting of the two
endpoint vertices and the $z$ long edges. It is a CW subcomplex.
The nerve description identifies $B\overline C$ with the
pushout of $BC\leftarrow A\to[0,1]$, where each long edge
maps linearly onto the interval. Since the first map is a
cofibration, this is a homotopy pushout. Replacing the
contractible interval by a point shows that it is the homotopy
cofiber of the specified map $A\to BC$. A homotopy cofiber
of a map from a finite graph into a space of finite CW
homotopy type again has finite CW homotopy type.

The connected graph $A$ has homotopy type
$\bigvee^{z-1}S^1$. When $BC$ is simply connected, its map
to $BC$ is null-homotopic. The cofiber is then the stated
wedge with its reduced suspension. This also includes $z=1$,
when the graph is contractible and the collapse changes nothing.
\end{proof}

\begin{corollary}[Proved: all three-by-three triangles]
\label{CP2:A-three-three-triangles}
Let $C$ have full linear support on three objects, trivial
endpoint corners, and $|C(0,1)|=|C(1,2)|=3$. If $BC$ has
finite CW homotopy type, every rational product of two
degree-two classes is zero. In particular, no profile
$(1,m,1;3,3,z)$, for any positive integers $m,z$, has type
$\CP^2$.
\end{corollary}
\begin{proof}
Apply Lemma~\ref{CP2:A-triangle-long-collapse}. The actual
singleton-long compression of
Lemma~\ref{CP2:A-singleton-triangle-monoid} identifies
$B\overline C$ up to homotopy with $BP$ for the finite
monoid $P=M\sqcup(C(0,1)\times C(1,2))$, whose minimal
ideal is a three-by-three band. Its finite CW type was
just proved. Theorem~\ref{CP2:A-three-by-three-all-products}
kills all its rational degree-two products. Pullback along
$\pi$, surjective in degree two, kills every such product
on $BC$. The generator of $H^2(\CP^2;\mathbb Q)$ has
nonzero square, giving the exclusion.
\end{proof}

\begin{remark*}[Scope and source comparison]
This newly checked extension recovers the arbitrary-second-rank
three-by-three assertion in the supplied section
\path{FINMON/sections/B_rational_three_by_three.tex}, whose
proof was opened and read in full before this comparison.
Its missing certificates were not used. We reconstructed the
needed additional projective/bar certificates from our own
complete relation lattice. The broader inherited assertion
excluding $\CP^2\vee S^2$ for every monoid of order at most
$22$ has additional three-by-four, three-by-five and
four-by-four obligations; it is not a premise of the present
triangle result.
\end{remark*}


\subsection{The complete linear seventeen-arrow exclusion}

\begin{theorem}[Proved: all linear three-object categories through seventeen]
\label{CP2:A-linear-seventeen-complete}
No finite three-object category with full linear acyclic support
and at most seventeen arrows, identities counted, has classifying
space of homotopy type $\CP^2$.
\end{theorem}
\begin{proof}
The at-most-sixteen-arrow case is already settled by
Theorem~\ref{CP2:A-linear-sixteen-complete}. At exactly seventeen
arrows write a profile as $(a,b,c;x,y,z)$, with positive corner
orders $a,b,c$ and positive hom-set sizes $x,y,z$ on $01,12,02$.
Every such profile is one of the $\binom{16}{5}=4368$ positive
ordered six-tuples of sum seventeen. Opposition acts by
\[
 (a,b,c;x,y,z)\longmapsto(c,b,a;y,x,z).
\]
There are $112$ fixed tuples and hence $(4368+112)/2=2240$
orbits. The explicit profile map and the independent final
positive-tuple reconstruction give the exact partition
\[
\begin{array}{l|r}
\text{original proved profile reductions}&2065\\
\text{additional three-by-three triangle theorem}&1\\
\text{complete independent finite-search receipts}&174\\\hline
\text{all exact seventeen-arrow profiles}&2240.
\end{array}
\]

The first line is the full proved scope map in
\path{converse/scope17/linear17_proof_coverage.json}.
Its hypotheses include the separately proved large-corner,
small-module, semisimple, graph, singleton-long,
projective-singleton and actual singleton-triangle reductions.
Each profile is attached to its explicit theorem conditions;
no corner catalog beyond its proved range is silently used.
The second line is exactly $(1,7,1;3,3,2)$, excluded by
Corollary~\ref{CP2:A-three-three-triangles}. It is attached as a
new proof with its own exact source and certificate hashes;
the frozen original profile input and proof map are unchanged.

The third line comprises all $150$ profiles in the inherited
residual list, now independently computed, and $24$ further
profiles found in our full scope audit. The latter audit had
found $25$ additional profiles; the one just specified has
been removed by the new theorem. Every finite profile uses
the complete ordinary corner catalogs through order seven.
Proposition~\ref{CP2:A-linear-biset-enumeration} proves the
canonical action-orbit and equivariant-composition enumeration.
The earlier general table producer instead assigns every
unforced typed product and propagates required associativity.
Both methods have complete scope proofs. They retain exact
orbit weights where used; a finite cache may repeat a record
after eviction but never removes an unseen case.

The completed selected finite receipts account for
$18114854855$ labelled or orbit-weighted tables and
$30884728$ retained literal records. Of those records,
$30702087$ have a checked zero square and $182641$ have
positive third mod-two homology. These are computational
workload counts, not counts of isomorphism classes.
For every retained record the independent checker reconstructs
the typed table, verifies every identity and associativity
equation, recomputes the first two mod-two homology ranks,
and verifies either every equation of a degree-three square
primitive or the nonzero third homology. No unexcluded
record remains.

The finalizer \path{code/finalize_linear17_scope.py}
reconstructs all $2240$ profiles independently of either finite
producer. It distinguishes the two frozen input lists,
requires exactly the right profile and index in each receipt,
checks completed range footers, all table/certificate count
equalities, all current source hashes and every finished
certificate-stream digest, and binds the extra triangle proof
to its two full proof readings and new independent hull checks.
The command with \texttt{--require-all} rejects any pending
profile. The complete join is
\path{outputs/linear17_bisets/linear17_final_scope_join.json};
its compact hash-bound record is
\path{outputs/linear17_complete_summary.json}.

Five earlier receipts quote one exact historical version of
the projective-singleton proof. Its corollary was broadened
from a small corner-order condition to the actual sided
projectivity hypothesis already used in that proof. The
finalizer checks the exact old source inside checkpoint four,
the exact current source, and the independently reviewed
comparison. Only this specified digest pair is accepted.
All five actual producer predicates are covered, including
the one receipt that used this pruning rule. This explicitly
resolved version difference supplies no arbitrary exception
to the source checks.

The disjoint exhaustive profile partition and the complete
obstruction receipts now exclude every exact seventeen-arrow
linear category. Together with the earlier bound this proves
the stated theorem.
\end{proof}


\subsection{Rank-two products for a three-by-four minimal ideal}

\begin{lemma}[Proved: a relative three-by-three obstruction]
\label{CP2:A-relative-three-by-three}
Let $P,N$ be finite monoids and let $f:P\to N$ be a unital
monoid homomorphism. Suppose $BP$ has finite CW homotopy type
and the minimal ideal of $N$ has three rows and at most three
columns, with any finite maximal group. Then
\[
 f^*(x)f^*(y)=0\quad\text{in }H^4(BP;\mathbb Q)
 \qquad(x,y\in H^2(BN;\mathbb Q)).
\]
The same assertion holds with rows and columns interchanged.
No finite CW hypothesis is imposed on $BN$.
\end{lemma}
\begin{proof}
The actual minimal-group quotient of
Lemma~\ref{A:minimal-group-collapse} preserves the rational
cohomology ring. Composing $f$ with this quotient therefore
reduces to a minimal band. With only one column there is no
second rational cohomology. With two columns every balanced
matrix has rank at most one. Lemma~\ref{A:rank-one-square}
and rational polarization kill all products already on $BN$.

For three columns form the specified relation hull $H(R)$
of the action of $N$, as in the proof of
Theorem~\ref{CP2:A-three-by-three-all-products}. The actual
homomorphism $N\to H(R)$ induces an isomorphism on rational
second cohomology, so $x,y$ lift to that hull. The complete
thirteen-orbit proof there kills all products on the hull
itself, except for the character-pruned orbit $11$.
In that orbit second rational cohomology has dimension one,
represented by the displayed integral affine cocycle $u$.

Let $g:P\to H(R)$ be the composite. If a forbidden actor
$a$ belonged to $g(P)$, choose $p\in P$ with $g(p)=a$.
The eventual cyclic group of $p$ maps surjectively onto the
eventual cyclic group of $a$: the map sends sufficiently
large powers to the corresponding powers, and their periodic
tails are precisely these cyclic groups. The nonzero integral
character of $u$ remains nonzero under this surjection.
The local group is an endomorphism group of its idempotent
in the Karoubi envelope, giving the specified map of its
classifying space into $BP$. The line-bundle nilpotence
argument in Proposition~\ref{A:three-row-classification}
applies to the pullback integral class on the original
finite-CW space $BP$ and contradicts that nonzero character.
Thus $g(P)$ contains no forbidden actor.

To use the already enumerated maximal targets, adjoin the
whole minimal band and the identity to $g(P)$. This still
avoids forbidden actors: multiplying by any band element
lands in the band, and all band actors are allowed.
The resulting submonoid is contained in one of the nine
checked maximal allowed submonoids of order sixteen.
Each has zero rational fourth homology. Factoring $g$
through that target kills every pulled-back degree-two
product. The argument uses finite CW type only for $BP$;
it never assigns that property to $N$ or its hull.
Opposition, which preserves the classifying-space homotopy
type and rational degree-two products, gives the dual case.
\end{proof}

\begin{theorem}[Proved: rank two with a three-by-four minimal ideal]
\label{CP2:A-three-by-four-rank-two}
Let $M$ be a finite monoid whose minimal ideal has three rows
and four columns, or four rows and three columns, with any
finite maximal group. If $BM$ has finite CW homotopy type
and $\dim_{\mathbb Q}H^2(BM;\mathbb Q)=2$, then every product
of rational degree-two classes is zero.
\end{theorem}
\begin{proof}
By opposition it suffices to treat three rows and four columns.
Use the actual minimal-group quotient, retaining the original
$M$ as the source whenever finite CW type is needed. The
quotient has a rectangular minimal band $I\times J$, with
$|I|=3$, $|J|=4$, and the same rational cohomology ring.
The action on that band gives a homomorphism to the specified
relation hull $H(R)$. The balanced-tensor calculation makes
this homomorphism an isomorphism on rational second homology.
Consequently all degree-two classes of the original $M$ are
pulled back from the two-dimensional balanced-form space of
$H(R)$.

The independently checked complete relation lattice
\path{monoid/certificates/relation_lattice_3x4.json}
has $2767$ proper relation spaces in $53$ row/column
permutation orbits. Exactly $23$ orbits have second rank two.
The following disjoint partition covers all of them in the
unchanged numbering of that complete lattice:
\[
\begin{array}{l|l}
\text{obstruction}&\text{orbit numbers}\\\hline
\text{common column quotient}&9,10,11,14,15,16,18,19,20,27\\
\text{all forms have rank at most one}&17\\
H_4(BH;\mathbb Q)=0&12,13,21,30,32,41,42,43,45\\
\text{character exclusion and an }H_4\text{-zero target}&35\\
u^2=v^2=(u+v)^2=0&29,47.
\end{array}
\]

We give the exact quotient argument for the first line.
Write the two integral balanced forms as matrices $L_1,L_2$
with zero initial row and column, indexed by $J\times I$.
For $j,k\in J$ put $j\sim k$ if
$L_t(j,i)=L_t(k,i)$ for every $i\in I$ and $t=1,2$.
There is a repeated column value in each listed orbit, so
$|J/{\sim}|\leq3$. Every actor $(\sigma,\tau)$ preserves
this equivalence relation. Indeed balance says
\[
 L_t(\tau(j),i)-L_t(\tau(0),i)
   =L_t(j,\sigma(i))-L_t(j,\sigma(0)),
\]
and subtraction for $j,k$ proves the assertion.
Thus the map on actors
$(\sigma,\tau)\mapsto(\sigma,\overline\tau)$ is an actual
unital monoid homomorphism, with the usual opposite
composition on the column transformations. Its image $N$
contains all constant row/column pairs and has minimal band
$I\times(J/{\sim})$.

Both forms descend to balanced forms on that band. The
affine-cocycle formula commutes with this quotient, since the
class of the initial column is the initial class. Hence their
degree-two classes lift from $BN$. They span all second
cohomology of the original $M$. Apply
Lemma~\ref{CP2:A-relative-three-by-three} to the actual
composite $M\to N$. It kills every product of those lifts.
Neither $BN$ nor the intermediate hull is assumed to have
finite CW type. For orbit $17$, every $2\times2$ minor of
every rational linear combination of the two basis forms
is zero. The rank-one lemma and polarization therefore
kill all products on that hull itself.

For the nine hulls on the third line we construct integral
projective prefixes through degree five. The independent
literal checker verifies module actions and maps, zero
composites, all image Smith invariants, and complementary
kernel ranks. Each prefix is exact through degree four,
and its coinvariant differential has zero fourth rational
homology. This establishes the stated obstruction without
any assertion about an uncomputed tail.

Orbit $35$ has $21$ actors and two integral affine basis
cocycles. Direct cyclic restrictions identify exactly
actors $17$ and $20$ as forbidden: each has an eventual
group of order two and residues $(1,0)$ on the two forms.
All other actors have zero residues on both. If either
forbidden actor occurred in the original action image,
its eventual cyclic preimage in $M$ would contradict the
same finite-CW line-bundle argument used in the preceding
lemma. The other nineteen actors form a closed submonoid
containing the identity and the entire minimal band.
Its integral projective resolution is finite, with
underlying abelian ranks and coinvariant dimensions
\[
 (3,10,12,4),\qquad (1,0,2,0).
\]
The literal checker proves exactness including the zero
final kernel. Its fourth rational homology is therefore
zero. Factoring the original action through this allowed
target kills every product.

The two remaining hulls have order seventeen. For each,
let $u,v$ be its two specified integral affine cocycles.
We construct projective/bar prefixes for $u,v,u+v$ with
common projective abelian ranks and coinvariant dimensions
\[
 (3,14,48,174,690),\qquad (1,0,2,2,6).
\]
All are exact through degree three. The literal bar
comparison has degree-two cochain vectors
\[
\begin{array}{c|ccc}
\text{orbit}&u&v&u+v\\\hline
29&(1,0)&(0,-1)&(1,-1)\\
47&(1,1)&(1,0)&(2,1).
\end{array}
\]
Each pair has determinant $-1$. Every one of the six
degree-four square cochains is zero on the entire
six-dimensional coinvariant module. The independently
checked comparisons extend to projective resolutions
because of exactness in the preceding degrees. These
are therefore the actual cup classes. The three zero
squares and rational polarization kill all products.

The constructors are
\path{code/explore_three_by_four_rank_two.py} and
\path{code/finish_three_by_four_rank_two.py}.
The separate source/scope checker
\path{monoid/code/check_three_by_four_rank_two_scope.py}
imports neither constructor. It reconstructs the exact
partition of all $23$ lattice orbits; verifies both full
forms, every actor congruence and quotient product in
the ten quotient cases; recomputes every cyclic restriction
for orbit $35$; and binds each certificate's actors,
multiplication and affine cocycle to the source hull.
It reruns all seventeen successful projective-prefix
checks, including one diagnostic prefix which is explicitly
not an exclusion. The complete receipt is
\path{monoid/certificates/three_by_four_rank_two_scope_verified.json}.
The earlier failed degree-five attempts for orbits $29,47$
remain recorded and are not premises; their successful
shorter bar prefixes supply the proof just given.

Finally all maps used above start at the original $M$,
possibly through its minimal-group quotient. The rational
ring isomorphism of that quotient transfers the balanced
classes, and every character obstruction is applied to
the original finite-CW source. This covers nontrivial
maximal groups as well and completes the proof.
\end{proof}

\begin{corollary}[Proved: three-by-four triangles with two long arrows]
\label{CP2:A-three-four-two-triangles}
For every positive integer $m$, neither full-linear profile
$(1,m,1;3,4,2)$ nor its opposite has type $\CP^2$.
\end{corollary}
\begin{proof}
Suppose such a category $C$ had projective-plane type.
The actual long-arrow collapse of
Lemma~\ref{CP2:A-triangle-long-collapse} gives
$B\overline C\simeq\CP^2\vee S^2$ and a pullback
surjective in degree two. The specified singleton-long
compression in Lemma~\ref{CP2:A-singleton-triangle-monoid}
identifies $B\overline C$ up to homotopy with the classifying
space of the monoid $P=M\sqcup(3\times4)$, whose minimal
ideal is that rectangular band. It has finite CW type
and second rational rank two. The preceding theorem
kills all its degree-two products. Pullback kills the
square of the generator on $BC$, a contradiction.
Opposition gives the other profile.
\end{proof}


\subsection{A tree with one possibly large corner}

\begin{proposition}[Proved: a tree reduces to one bounded act cone]
\label{CP2:A-tree-one-large-corner-eleven}
Let $C$ be a finite connected category with at least two objects and acyclic support,
whose undirected comparability graph is a tree. Suppose
there is an object $v$ with $|\operatorname{End}(v)|\leq11$
and every other corner has order at most two. If $BC$ has
bounded rational homology, then all rational degree-two
cohomology products on $BC$ are zero.
\end{proposition}
\begin{proof}
A tree support has no chain
of three distinct objects, because transitivity would supply
the third edge of a triangle. All arrows incident with $v$
therefore point in the same direction.

Let $X$ be the full subcategory on $v$ and all its neighbors.
It is a nonempty source or sink star. Each component of
the graph outside this vertex set attaches to exactly one
vertex $w$ of $X$, and $w\ne v$. Let $Y$ be the full
subcategory on that component and its attachment vertex.
Every corner of $Y$ has order at most two. Its rational
corner algebras are semisimple, and its support has strict
height at most one. The relative bar complex therefore
gives $H_j(BY;\mathbb Q)=0$ for every $j\geq2$.
The same corners have rationally acyclic classifying spaces.

The nerves of $X$ and these $Y$ pieces cover $NC$ exactly:
each simplex uses at most two distinct objects and hence
lies over one vertex or one tree edge. Attach the pieces
one at a time. Each intersection is the corner category at
its attachment vertex $w$, which is rationally acyclic.
Mayer--Vietoris gives, in every degree $j\geq2$, the
restriction isomorphism
\[
 H^j(BC;\mathbb Q)\xrightarrow{\ \cong\ }H^j(BX;\mathbb Q).
\]
It is induced by the actual inclusion and respects products.
The corresponding homology sequences also show that $BX$
has bounded rational homology.

At every leaf of $X$, the trivial rational module is
projective on both sides, since its corner has order at
most two. Apply the actual orbit functors of
Lemma~\ref{CP2:A-high-orbits}, replacing these leaves by
trivial corners. They give rational cohomology-ring
isomorphisms and preserve bounded rational homology.
There is at least one leaf. Merge all the resulting trivial
leaves into one, retaining the disjoint union of their
stalks as the actual finite act of $\operatorname{End}(v)$.
There are no compositions between different leaves, so this
is a functor and its nerve realization identifies exactly
those finitely many leaf vertices. On a connected CW space,
identifying $r$ distinct points adds a wedge of $r-1$ circles
up to homotopy: this follows by attaching a contractible
tree joining them and then collapsing that tree. Thus the
quotient map is an isomorphism in cohomological degrees at
least two, preserving products, and the quotient still has
bounded rational homology.

The quotient is the ordinary act cone over the corner at
$v$, or its opposite. Theorem~\ref{CP2:A-bounded-eleven-acts}
kills every rational degree-two product there. Pullback
through the specified quotient and orbit functors kills
all such products on $BX$. The degree-four restriction
isomorphism then kills all products on $BC$.
\end{proof}


\subsection{Universal short products over arbitrary fields}

\begin{lemma}[Proved: unused long arrows split off two-spheres]
\label{A7:unused-long-wedges}
Let $X$ be a finite acyclic category with full support
$0<1<2<3<4$ and a singleton $X(0,4)$.
Delete every arrow in $X(0,3)$ and $X(1,4)$ which is not
a composite of two nonidentity arrows, obtaining a
subcategory $Y$ on the same objects. Put
$A=X(0,1)$, $D=X(3,4)$, and let $r,s$ be the respective
numbers of deleted $0\to3$ and $1\to4$ arrows. Then
\[
 BX\simeq BY\vee\bigvee^{\,r(|D|-1)+s(|A|-1)}S^2.
\]
This equivalence can be chosen under the specified
inclusion $BY\hookrightarrow BX$. In particular, vanishing
of all degree-two products on $BY$ implies the same
vanishing on $BX$, over any field.
\end{lemma}
\begin{proof}
First, every full-support subcategory retaining the
singleton extreme hom-set is simply connected. Indeed,
choose arrows $a_i:0\to i$, with $a_0=1_0$ and $a_4$
the unique extreme arrow. For $f:i\to j$, choose
$b_j:j\to4$, using $b_4=1_4$. In the universal groupoid,
\[
 b_j f a_i=a_4=b_j a_j.
\]
Cancellation gives $f=a_j a_i^{-1}$. Thus all arrows are
determined by a tree of the $a_i$, and the fundamental
group at $0$ is trivial. The assertion applies to $Y$:
each long hom-set retains a composite and hence remains
nonempty, and deletion of arrows not hit by composition
is closed under composition.

Adjoin one deleted arrow $e:0\to3$ to any such
subcategory. The only new nondegenerate simplices are
the edge $e$ and the triangles $(e,d)$ for $d\in D$.
No new simplex has dimension at least three, since
there is no nonidentity arrow into $0$ and only one
further object after $3$. Nor is a new triangle a face
of an old three-simplex: such a face would express
$e$ as a two-arrow composite, contrary to its definition.
Every composite $de$ is the unique extreme arrow.

Choose $d_0\in D$. The subcomplex obtained by adjoining
only $e$ and $(e,d_0)$ is an elementary expansion of
the old nerve realization: the new edge is a free
face of that triangle, and its other two edges already
belong to the old complex. It deformation retracts
onto the old complex while fixing it. Each of the
remaining $|D|-1$ triangles is a two-cell attached
along a loop in this subcomplex. That subcomplex is
simply connected by the first paragraph, so every
attaching map is null-homotopic. Attaching these
two-cells therefore gives a wedge of $|D|-1$ copies
of $S^2$, up to homotopy under the old complex.

Repeat for every unused $0\to3$ arrow. The dual
argument for an unused $1\to4$ arrow adjoins one
edge and $|A|$ triangles, giving $|A|-1$ two-spheres.
The two kinds of new arrows cannot be consecutive
factors, and neither changes $A$ or $D$. The asserted
formula follows. In a wedge with two-spheres every
positive product involving a sphere summand is zero,
while products of classes from $BY$ are their original
products. This proves the ring assertion over any field.
\end{proof}

\begin{theorem}[Proved: universal product transfer over any field]
\label{A7:universal-short-all-fields}
Use the three short maps and the universal category
$F=F(P,Q,R)$ of Theorem~\ref{A7:universal-short-model}.
For every field $k$, if all products
$H^2(BF;k)\otimes H^2(BF;k)\to H^4(BF;k)$ are zero,
then all such products on $B\mathcal C$ are zero.
\end{theorem}
\begin{proof}
Collapse $\mathcal C(0,4)$ to a singleton, obtaining
$\mathcal C_0$. Proposition~\ref{A7:extreme-arrow-collapse}
gives the actual product-preserving surjection from
its second cohomology to that of $\mathcal C$.
Delete its unused $03$ and $14$ arrows to obtain $Y$.
Lemma~\ref{A7:unused-long-wedges} shows that vanishing
of products transfers from $BY$ to $B\mathcal C_0$.

The same explicit functor $F\to Y$ used in
Theorem~\ref{A7:universal-short-model} is the identity
on the normalized chain groups in degrees three and
four and on the fourth boundary. Its map on fourth
cohomology is therefore an isomorphism over $k$.
For $u,v\in H^2(BY;k)$, naturality sends $uv$ to
the product of their two pullbacks on $BF$. The
hypothesis makes this zero, and the fourth-cohomology
isomorphism forces $uv=0$. Transfer successively
through the wedge decomposition and the degree-two
surjection proves the claim. No bound on second
Betti number and no characteristic-two argument
is required.
\end{proof}

\begin{corollary}[Proved: two components suffice for tree-stalk prefixes]
\label{A7:partial-long-tree}
In a partial initial assignment of the three short
maps $P,Q,R$, retain all vertices of the two known-edge
graphs from Proposition~\ref{A7:partial-long-graph}.
Every completion has zero rational degree-two products
if either of the following conditions holds:
\begin{enumerate}
\item the $03$ graph has at most two connected components,
and either $|U|\leq2$ or $|A|\leq2$, $|B|\leq3$;
\item the $14$ graph has at most two connected components,
and either $|W|\leq2$ or $|D|\leq2$, $|C|\leq3$.
\end{enumerate}
\end{corollary}
\begin{proof}
Adding the still unknown edges can only decrease the
number of connected components. In the first case,
every completed universal category has $|F(0,3)|\leq2$.
Its source stalks satisfy (S1) or (S2) of
Proposition~\ref{A7:tree-stalk-profile}, respectively.
Thus all its rational degree-two products are zero.
Theorem~\ref{A7:universal-short-all-fields} transfers
this to every category with those completed short
maps. The second case uses (T1) or (T2) at the sink.
All vertices, including unused short-map values,
are essential to the component count.
\end{proof}


\subsection{The five-object scope at eighteen arrows}

\begin{lemma}[Proved: semisimple corners in the five-object budget]
\label{CP2:A-five18-semisimple-small}
Let $C$ have five objects, acyclic support, eighteen arrows,
and corner orders at most two. Then $BC$ does not have the
homotopy type of $\CP^2$.
\end{lemma}
\begin{proof}
The rational algebra of every corner is semisimple. If the
support is not a full five-object chain, its strict height
is at most three. The relative bar resolution of
Lemma~\ref{CP2:A-semisimple-height} then has no homology in
degree four.

For the full chain there are ten strict hom-sets. Put
$D=\sum_x(|\operatorname{End}(x)|-1)$ and
$P=\sum_{x<y}(|C(x,y)|-1)$. The arrow count gives
$D+P=18-5-10=3$. If $D=0$, all corners are trivial and
at least one of the four consecutive hom-sets is a
singleton: otherwise $P\geq4$. Proposition~\ref{A7:adjacent-homs}
annihilates all degree-two products. If $D\geq1$, then
$P\leq2$, and the support cycle rank is $10-5+1=6$.
Lemma~\ref{CP2:A-high-two-parallel} applies; under the
assumption of projective-plane type its second Betti number
hypothesis is exactly one. Its zero-square conclusion again
contradicts the projective-plane cohomology ring.
\end{proof}

\begin{lemma}[Proved: deletion through a unique singleton comma]
\label{CP2:A-five18-singleton-deletion}
Suppose a category $C$ with acyclic support has an object $x$
such that its incoming arrows from other objects consist
of a single arrow $y\to x$, and $\operatorname{End}(y)$ is
trivial. The inclusion of the full deletion
$i:C\setminus\{x\}\longrightarrow C$ is a homotopy
equivalence on classifying spaces. The dual assertion
holds for a unique outgoing arrow to a trivial corner.
Consequently such a five-object eighteen-arrow category
does not have type $\CP^2$.
\end{lemma}
\begin{proof}
For every retained object $z$, the comma category
$(i\downarrow z)$ has terminal object $(z,1_z)$.
For $z=x$, that comma category consists of the one object
$(y,y\to x)$ and its identity, by the two stated hypotheses.
Quillen's Theorem~A therefore applies to this specified
inclusion. Opposition proves the dual assertion.
Deleting $x$ removes its identity and the specified cross
arrow, and possibly more arrows, so the deletion has four
objects and at most sixteen arrows. The independently
proved complete sixteen-arrow exclusion rules it out.
\end{proof}

\begin{lemma}[Proved: the tree-edge scope test]
\label{CP2:A-five18-tree-edges}
Let $C$ have connected acyclic support whose undirected
comparability graph is a tree. Suppose every nonleaf corner
has order at most five, and every edge with endpoint corner
orders $a,b$ and stalk size $s$ satisfies at least one of
\[
\begin{gathered}
s\leq2,\qquad \max(a,b)\leq4,\qquad a+b+s\leq11,\\
\min(a,b)\leq2\ \text{ and }\ \max(a,b)\leq11,
\qquad \{a,b\}=\{3,5\}.
\end{gathered}
\]
If $BC$ has bounded rational homology, all its rational
degree-two products vanish.
\end{lemma}
\begin{proof}
The tree condition forbids a strict three-object chain.
Thus the nerve is the union of its full edge subcategories,
attached successively along corner categories. Every
intersection corner is a nonleaf corner. By
Corollary~\ref{CP2:A-five-types}, its rational homology
vanishes in degrees at least three. The Mayer--Vietoris
sequences consequently give injectivity of restriction in
degree four, and show that every edge piece inherits
bounded rational homology. These statements also apply
when the intersection corner has the type of $S^2$;
rational acyclicity of that corner is not required.

For $s=1$ the edge space is the join of its two corner
spaces and has zero second rational cohomology. For $s=2$
use Lemma~\ref{CP2:A-high-two-cross}. When both corners
have order at most four, both restrictions in degree two
are zero, so Lemma~\ref{CP2:A-two-block} applies. For
$a+b+s\leq11$ use Corollary~\ref{CP2:A-high-two-objects}.

If one corner has order at most two, its trivial rational
modules are projective on both sides. The appropriate
actual orbit functor of Lemma~\ref{CP2:A-high-orbits}
is a rational cohomology-ring equivalence to an ordinary
act cone over the other corner. Boundedness passes through
this equivalence, so Theorem~\ref{CP2:A-bounded-eleven-acts}
applies when the other corner has order at most eleven.

Finally consider orders three and five. If the five-element
corner is not the rectangular-band monoid $P_5$, both
corners have zero rational second cohomology and the
two-block lemma applies. Otherwise the three-element
corner either has the required projective trivial module,
giving an ordinary $P_5$ act cone by the same orbit
comparison, or is the exceptional oriented right-zero
band monoid. The ordinary act is covered by
Lemma~\ref{CP2:A-high-two-row-act}, and the exceptional
biset by Lemma~\ref{CP2:A-high-P5-graph}. These arguments
allow every nonempty finite stalk. In all five cases the
edge products vanish. Injectivity of fourth restriction
therefore kills the global products.
\end{proof}

\paragraph{Complete interpretation of the cardinality predicates.}
The source \path{code/five18_scope.py} considers only connected
acyclic supports on exactly five objects and positive
decorations of total eighteen. For each support its cycle
rank is at most six. Its support-first canonization chooses
one representative under object relabelling and opposition,
and then divides decorations by the complete automorphism
and anti-automorphism group of that fixed support.
No corner multiplication or composition table is assumed
to exist merely because its cardinalities have been listed.

Here is the exact mathematical meaning of every enabled
proof predicate. Nonzero first mod-two homology of the
support is excluded by the fundamental-group surjection:
lift each traversed support edge to any category arrow,
using the reversed path when necessary. Singleton cross
hom-sets use Lemma~\ref{CP2:A-singleton-support}; every
second class pulls back from the support, whose fourth
cohomology is zero by Lemma~\ref{CP2:A-high-flag}.
Corners at most two use
Lemma~\ref{CP2:A-five18-semisimple-small}. Parallel excess
at most two with corners at most four uses
Lemma~\ref{CP2:A-high-two-parallel} with the candidate's
second Betti number one. Nonserial parallel hom-sets use
Lemma~\ref{CP2:A-high-parallel-cut}: a small corner has
zero second cohomology, and a corner incident with a
singleton cross hom-set has nullhomotopic inclusion into
that edge join, so every global second class restricts
to zero there.

For every proper downset and complementary upset, the
source tests the two-block lemma and the fixed-arrow
corollary. Its sufficient zero-second-cohomology test is
applied component by component. A one-object component
has corner order at most four. A larger component either
has at most $2k+1$ arrows, covered by
Lemma~\ref{CP2:A-high-H2-cost}; has singleton cross
hom-sets, covered by the singleton-support lemma and the
flag budget on at most four vertices; or has semisimple
corners of orders at most two and strict height at most
one. This last case follows from the relative bar bound.
These are sufficient tests only. A failed test removes
nothing. A full singleton join across a cut has zero
positive products, since the join is a suspension after
collapsing its two contractible basepoint cones. A small
extreme corner with every incident stalk nonempty and of
size at most two uses Lemma~\ref{CP2:A-high-small-binary},
again with the candidate's second Betti number one.

Unique singleton comma deletions use
Lemma~\ref{CP2:A-five18-singleton-deletion}. The tree-edge
predicate is exactly Lemma~\ref{CP2:A-five18-tree-edges}.
A full star with center at most eleven and every other
corner at most two is a special case of
Proposition~\ref{CP2:A-tree-one-large-corner-eleven}.
Thus every recorded proof reason is a sufficient exclusion
under explicitly stated candidate hypotheses. The numerical
profile and residual joins are separate finite obligations.

\paragraph{Reproducible orbit counts and the separate tree supplement.}
For a fixed support with $q$ nonempty hom-sets, including
the five diagonal positions, all positive decorations
number $\binom{17}{q-1}$. The source lists every positive
composition of eighteen into these $q$ positions, retains
the least decoration in each support symmetry orbit, and
records its orbit size as group order divided by its
stabilizer order. The orbit-size sum is checked against
that binomial coefficient for every support. The full
stream preserves each profile and every applicable named
reason. An independent program must still regenerate the
full list and its predicate partition before the numerical
scope is adopted in a completed lower bound.

The later source \path{code/five18_tree_supplement.py}
leaves this first stream and its residual ledger unchanged.
It checks the exact hypothesis of
Proposition~\ref{CP2:A-tree-one-large-corner-eleven} on
each residual row and supplies a separately hashed proof
partition. Rows not covered by it remain finite-composition
obligations; no running or incomplete search excludes one.


\subsection{The last four profiles in the seventeen-arrow census}

\begin{proposition}[Proved: three pendant small-corner profiles]
\label{CP2:A-four17-final-pendant}
None of the following four-object hom-cardinality matrices
has a classifying space of type $\CP^2$:
\[
\begin{gathered}
\begin{pmatrix}1&0&0&1\\0&1&0&1\\0&0&7&3\\0&0&0&3\end{pmatrix},\qquad
\begin{pmatrix}1&0&0&1\\0&6&2&3\\0&0&1&0\\0&0&0&3\end{pmatrix},\\
\begin{pmatrix}1&0&0&1\\0&7&2&2\\0&0&1&0\\0&0&0&3\end{pmatrix}.
\end{gathered}
\]
These are rows $5,237,255$ of the frozen four-object
seventeen-arrow residual ledger, respectively.
\end{proposition}
\begin{proof}
Let $N$ be the corner at object $3$, of order three.
It is a sink corner. Object $0$ is trivial and is joined
only to $3$, by a singleton stalk. Let $D$ be the full
deletion of object $0$. The full subcategory on $0,3$
has initial object $0$, hence contractible classifying
space. Its nerve and $ND$ cover $NC$, and their intersection
is the nerve of $N$: no strict chain can pass through $0$
and a further object other than $3$.

First suppose $BN$ is contractible. The inclusions of
these nerves are CW cofibrations, so their union is their
homotopy pushout. The actual inclusion $BD\to BC$ is
therefore a homotopy equivalence. Deleting $0$ removes
exactly its identity and its singleton arrow, so $D$ has
fifteen arrows. The complete sixteen-arrow exclusion
rules out $BD\simeq\CP^2$.

Otherwise the small-corner classification
Corollary~\ref{CP2:A-five-types} identifies $BN$ with the
classifying space of a nontrivial finite minimal group.
The minimal ideal of $N$ is that group itself. Indeed,
a nontrivial group with two rows or two columns would
already require at least four ideal elements; if the
ideal is not a group, an identity outside it is additionally
required. This is impossible at order three. Averaging
over the group ideal makes both trivial rational $N$
modules projective.

Apply the actual sink-orbit functor
$C\to\overline C$ from Lemma~\ref{CP2:A-high-orbits}:
replace $N$ by the trivial corner and every incoming
stalk by its $N$-orbit set. It induces a rational
cohomology-ring isomorphism. All the stalks remain
nonempty, so the comparability support is unchanged.
In each of the three displayed profiles that support is
a tree. After the orbit functor there is just one corner
of order greater than two, and its order is six or seven.
Under the assumption $BC\simeq\CP^2$, the space
$B\overline C$ has bounded rational homology.
Proposition~\ref{CP2:A-tree-one-large-corner-eleven}
annihilates all its degree-two products. Pullback through
the specified functor contradicts the nonzero rational
square on $BC$. This proves all three exclusions.
\end{proof}

\begin{proposition}[Proved: the final triangle with a small sink]
\label{CP2:A-four17-final-triangle}
The hom-cardinality matrix
\[
\begin{pmatrix}
1&0&0&1\\
0&6&2&3\\
0&0&1&1\\
0&0&0&2
\end{pmatrix}
\]
cannot have classifying space of type $\CP^2$.
It is row $238$ of the frozen residual ledger.
\end{proposition}
\begin{proof}
The sink corner at object $3$ has order two, so its
trivial rational modules are projective on both sides.
The actual sink-orbit functor $C\to\overline C$ replaces
that corner by the trivial monoid and gives a rational
cohomology-ring isomorphism. Its singleton stalks from
$0$ and $2$ remain singletons. The stalk $1\to2$ still
has two arrows, and the stalk $1\to3$ becomes a nonempty
act of the six-element corner at $1$.

Delete $0$ from $\overline C$. Its only outgoing cross
arrow is the singleton arrow to the trivial corner $3$.
Lemma~\ref{CP2:A-five18-singleton-deletion}, in its dual
form, proves that the full deletion inclusion is a
homotopy equivalence. In the resulting category on
$1,2,3$, delete $2$ by the same argument: its only
outgoing cross arrow is again the singleton to $3$.
The resulting full category $E$ on $1,3$ is an ordinary
act cone over the six-element corner, and the two actual
inclusions give $BE\simeq B\overline C$.

If $BC$ were a projective plane, this comparison would
give bounded rational homology of $BE$.
Theorem~\ref{CP2:A-bounded-eleven-acts} annihilates its
degree-two products. The two homotopy equivalences and
the orbit functor transfer this vanishing to $BC$, a
contradiction. The argument uses rational boundedness
after the orbit functor, not an asserted integral
homotopy equivalence for that functor.
\end{proof}


\subsection{Completion of the four-object seventeen-arrow census}

\begin{theorem}[Proved: all four-object acyclic-support profiles at seventeen]
\label{CP2:A-four-seventeen-complete}
A finite category with four objects, acyclic support, and exactly
seventeen arrows, including identities, cannot have classifying
space homotopy equivalent to $\CP^2$.
\end{theorem}
\begin{proof}
A candidate is connected. Every acyclic support admits a
topological ordering, so write its nonempty strict hom-sets above
the diagonal. Independently enumerate all subsets of the six
possible strict positions, imposing transitivity and connectedness.
There are eighteen topologically labelled supports. Positive
decorations of their diagonal and supported strict positions with
total seventeen give $186\,186$ labelled decorations. Quotienting
by every object permutation preserving an upper triangular
presentation and by opposition gives exactly $34\,218$ profiles.
The independent program
\path{converse/scope17/finalize_four17.py} reconstructs this entire
set without importing the producer or cardinality-scope program.
It checks equality with every literal matrix in the frozen scope
stream, including each normalized two-simplex count.

The original proved cardinality reductions leave $1081$ profiles,
all with corner orders at most seven. The later proofs are joined
to this unchanged ledger, with their exact source digests and
literal profile witnesses. Their disjoint contribution is
\[
\begin{array}{lr}
\text{original general proof rows}&33\,137\\
\text{one-large-corner tree rows}&56\\
\text{additional singleton-comma deletion rows}&224\\
\text{four final displayed profiles}&4\\
\text{required complete finite profiles}&797\\\hline
\text{total}&34\,218.
\end{array}
\]
The tree rows use
Proposition~\ref{CP2:A-tree-one-large-corner-eleven}.
There are $244$ rows with a deletion covered by the general
statement of Lemma~\ref{CP2:A-five18-singleton-deletion}; exactly
twenty were already covered by the tree proof. For each of the
$244$ rows the finalizer records the deleted object, the unique
incoming or outgoing arrow, its trivial neighboring corner, and
the number of removed arrows. Each deletion leaves three objects
and at most fifteen arrows, so the complete sixteen-arrow
baseline, Theorem~\ref{CP2:A-global-lower-seventeen}, applies.
Propositions~\ref{CP2:A-four17-final-pendant} and
\ref{CP2:A-four17-final-triangle} cover the four specified final
rows. Thus none of these proof rows depends on an unfinished
composition search.

For the remaining $797$ profiles, choose all corner monoids from
the independently complete ordinary-isomorphism catalog through
order seven. For each strict hom-set, enumerate every commuting
left and right action of its endpoint corners, up to permutations
of that hom-set. These permutations are independent for distinct
strict hom-sets. Consequently every category table can be
transported to one with all these actions canonical. With the
actions installed, the producer enumerates every typed product
of two strict arrows and imposes every typed associativity
equation. It takes no further quotient of these composition maps.
The product of the action-orbit sizes recovers the count of
labelled action tables. This proves coverage of the full category
scope, including all compatible composition maps.

One additional corner-table test precedes this enumeration.
If all four rational corner algebras are semisimple, the relative
bar complex has strict height at most three, so fourth homology
vanishes. This is the only enabled corner-table exclusion at four
objects: the shared two-block test explicitly returns false when
the object count is not three, and the inherited projective-module
and act-cone table cuts are disabled. The complete profile receipts
retain the exact source versions for this test and the enumeration.

The resulting complete gates contain $650\,840\,107$ weighted
labelled tables. After the explicit first- and second-homology
filters, their literal retained certificate streams contain
$7\,267\,931$ records. The separately written normalized-nerve
checker verifies $7\,247\,034$ zero-square certificates and
$20\,897$ positive-third-homology certificates, with no survivors.
In the former case it checks a nonzero degree-two cohomology
class $\alpha$ over $\mathbb F_2$ and a literal cochain $w$ with
$\delta w=\alpha\smile\alpha$; this is incompatible with the
projective-plane ring. In the latter case the nonzero third
homology is incompatible with $\CP^2$.

The finalizer requires a closed producer footer and an independent
checker pass for every one of these $797$ profiles. It binds each
profile index and matrix, source and input versions, compressed
and uncompressed stream digests, stream length, and both
certificate counts. The flag \texttt{--require-all} completed
successfully with an empty missing-profile list. Its compact
receipt is
\path{converse/scope17/four17_checked_scope_summary.json}; the
separately mapped full join retains every scope row and adopted
receipt.

Four additional, proof-redundant historical receipts had stale
failure-file path bindings after a later publication failure.
They are explicitly excluded from every adopted finite count;
their profiles lie in the independently proved tree set. Their
bytes and the exact mismatch are preserved separately. No
required finite profile is exempted from a verification check.
Together the disjoint proof partition and the $797$ completed
finite gates exhaust all $34\,218$ profiles, proving the theorem.
\end{proof}


\subsection{The complete exclusion of seventeen arrows}

\begin{theorem}[Proved: the unrestricted lower bound is eighteen arrows]
\label{CP2:A-global-lower-eighteen}
Every finite category whose classifying space is homotopy
equivalent to $\CP^2$ has at least eighteen arrows, counting
identities. In particular, the seventeen-arrow alternative
in Part~A.2 is false, and the least arrow number satisfies
\[
\boxed{18\leq n\leq22}.
\]
\end{theorem}
\begin{proof}
The complete bound through sixteen arrows is
Theorem~\ref{CP2:A-global-lower-seventeen}. It therefore
suffices to exclude exactly seventeen arrows. A category
of projective-plane type is connected; write $k$ for its
number of objects.

For $k=1$, the complete minimum-monoid classification
Theorem~\ref{A:monoid-minimum23-classification} excludes
all orders less than twenty-three. For $k=2$, use
Theorem~\ref{CP2:A-two-object-seventeen}; it includes
both support types and all object-retract reductions.

For $k=3$, cyclic support is excluded by
Theorem~\ref{CP2:A-strong-three-seventeen}. A connected
acyclic support is either a source or sink star, or a
full linear order: a directed two-step path forces its
third edge by composition. The star case is
Theorem~\ref{CP2:A-star-seventeen-complete}, and the
linear case is Theorem~\ref{CP2:A-linear-seventeen-complete}.

For $k=4$, use the cyclic-support theorem
\ref{CP2:A-strong-four-seventeen} and the complete
acyclic-support theorem \ref{CP2:A-four-seventeen-complete}.
Finally, if $k\geq5$, then $17\leq2k+7$ and $17\leq21$,
so the all-category excess theorem
\ref{CP2:A-high-all-eight} excludes the category. This
exhausts every finite category with seventeen arrows,
including empty cross hom-sets and nontrivial corners.

The independently checked construction of
Theorem~\ref{A:category-upper22} supplies the upper bound.
No inherited certificate stream absent from this run is
a premise of any finite gate used in this proof.
\end{proof}

\paragraph{Exact scope joins.}
The two-object gate has its separate root arithmetic and
receipt join in \path{outputs/two_object17_scope_independent.json}.
The star and linear gates cover all $924$ and $2240$
canonical cardinality profiles, respectively, with their
separate root joins in
\path{outputs/star17_scope_independent.json} and
\path{outputs/linear17_scope_independent.json}.

The final four-object scope has $34218$ canonical profiles.
Its original $33137$ structural cases are supplemented
by $56$ tree cases and $244$ actual singleton-comma
deletions, with $20$ overlapping cases, and the four
explicit final profiles. Thus its disjoint partition is
\[
34218=(33137+56+244-20+4)+797=33421+797.
\]
The required finite part contains $650840107$ weighted
corner-representative tables and $7267931$ literal
obstruction records. Of those records, $7247034$ have
a checked zero-square primitive and $20897$ have positive
third mod-two homology. There is no missing required
profile and no unexcluded record. These workload counts
are not counts of category isomorphism classes.

The separate root program
\path{code/check_four17_scope_join.py} imports no profile
enumerator or finalizer. It regenerates every canonical
cardinality row, reconstructs the three supplemental proof
sets, and checks all $797$ required receipts, input and
source hashes, compressed-stream digests and independent
literal-obstruction counts. Its complete receipt is
\path{outputs/four17_scope_independent.json}. Four historical
proof-redundant receipts whose failure-history filenames
changed are explicitly excluded from every adopted finite
count; their profiles are covered by the separately bound
tree proof. Their preserved records and subsequent history
recovery have no role in the required finite partition.

The unrestricted eighteen-arrow decision is the next
open case. No theorem in this subsection excludes an
eighteen-arrow category.


\subsection{A pendant singleton at a small tree corner}

\begin{proposition}[Proved: a small-corner pendant in the eighteen-arrow budget]
\label{CP2:A-pendant-small-tree-eighteen}
Let $C$ be a finite connected category with at least three
objects, acyclic support, and at most eighteen arrows.
Suppose its undirected comparability graph is a tree.
Assume that a leaf $p$ has trivial corner and a singleton
stalk to its neighbor $v$, whose corner $N$ has order at
most four. Suppose that, apart from $v$, at most one corner
has order greater than two, and that corner has order at
most eleven. Then $BC$ does not have type $\CP^2$.
\end{proposition}
\begin{proof}
The full edge on $p,v$ has initial or terminal object $p$,
according to its orientation, so its classifying space is
contractible. Let $D=C\setminus\{p\}$ be the full deletion.
There is no strict three-object chain in a comparability
tree. Hence the nerves of $D$ and that edge cover $NC$,
and their intersection is $NN$.

If $BN$ is contractible, this is a homotopy pushout along
CW cofibrations, and the actual inclusion $BD\to BC$ is
a homotopy equivalence. The deletion removes exactly the
identity of $p$ and its sole cross arrow, so $D$ has at
most sixteen arrows. Theorem~\ref{CP2:A-global-lower-seventeen}
excludes projective-plane type.

Otherwise $BN$ has the homotopy type of a nontrivial finite
group by Corollary~\ref{CP2:A-five-types}. The minimal ideal
of $N$ is a group: two rows or two columns with a nontrivial
local group would require at least four ideal elements,
and the outside identity would give order at least five.
Both trivial rational modules of $N$ are projective by
averaging over this group ideal.

All edges at $v$ have the same orientation, since any
incoming and outgoing pair would create a transitive
triangle. Apply the appropriately sided actual orbit
functor of Lemma~\ref{CP2:A-high-orbits}, replacing $N$
by the trivial corner and all its stalks by their orbit
sets. Its pullback
$H^*(B\overline C;\mathbb Q)\to H^*(BC;\mathbb Q)$
is a ring isomorphism. Orbit sets of nonempty stalks
are nonempty, so the comparability tree is unchanged.
All corners of $\overline C$, except possibly one of
order at most eleven, have order at most two. Under the
assumption $BC\simeq\CP^2$, rational homology of
$B\overline C$ is bounded. Proposition~\ref{CP2:A-tree-one-large-corner-eleven}
therefore kills all degree-two products there and, by
pullback, on $BC$. This contradicts the primitive rational
square of the projective plane.
\end{proof}


\subsection{The two remaining coned three-state star cases}

Write $P=P_5=(\{0,1\}\times\{0,1\})^1$, with
$e_{ij}e_{kl}=e_{il}$.  For a left-$N$, right-$M$ biset $S$,
put
\[
 E=B(N,S,M),\qquad Y=\operatorname{hocofib}(E\longrightarrow BN).
\]
This is the classifying-space model of the source star obtained
by adjoining to the $(M,N,S)$ edge a trivial leaf on a singleton
stalk: that additional edge cones off $BM$.
All homology in this subsection is rational unless a coefficient
ring is displayed.

\begin{lemma}[Proved: elementary coned-bridge comparisons]
\label{CP2:A-coned-bridge-comparisons}
The following assertions hold for the preceding model.
\begin{enumerate}
\item If a second class of $Y$ restricts to zero on $BN$, its
square is zero.  In particular all second squares vanish if
$H^2(BN)=0$, or if $E\to BN$ has a section on second cohomology.
\item If the left $N$ action on $S$ is trivial, then $Y$ is
homotopy equivalent to a suspension.
\item A proper biset retraction of $S$ can be used in a
projective-plane exclusion: a star of type $\CP^2$ would give
a star of the same type on the image biset.  In particular a
retraction onto at most two states contradicts
Lemma~\ref{CP2:A-star-coned-binary}.
\item Suppose $BM$ is rationally acyclic, the right action
space $Z=B(S,M)$ is connected, and some left $N$ actor on $S$
is constant.  Then $H_2(Y)=0$.
\end{enumerate}
\end{lemma}
\begin{proof}
The quotient $Y/BN$ is the suspension of $E$. A class with
zero restriction lifts from this quotient, whose positive
products vanish. The cofiber exact sequence gives the second
assertion in (1).

For (2), $E=BN\times Z$ and its map to $BN$ is projection.
Choose a vertex of $Z$. Restricting the projection to the
resulting section gives the identity $BN\to BN$; its mapping
cone is a contractible subcomplex of the full cofiber model.
Collapsing it identifies the quotient with
$\Sigma(BN_+\wedge Z)$, with the chosen vertex used as the
basepoint of $Z$. Collapsing that contractible subcomplex is
a homotopy equivalence.

For (3), use Theorem~\ref{CP2:A-cut-retract} for the cut
from the center to its two leaves, with the identity on the
singleton stalk. This statement concerns an actual retract,
so the integral projective-plane recognition in that theorem
applies.

For (4), regard $Z$ as a left $N$ space. The cone of the
augmented complex $C_*(Z)\to\mathbb Q$ has no homology in
degrees zero and one, and its degree-two homology is $H_1(Z)$.
The derived $N$ coinvariants of this cone compute the reduced
homology of $Y$. The first nonzero hyperhomology row therefore
gives $H_2(Y)=H_1(Z)_N$. A constant actor with value $u$ has
$u m=u$ for every $m\in M$, by commutation. Its map on $Z$
factors through $BM$, followed by the fixed-state inclusion
$BM\to Z$. Thus it acts by zero on $H_1(Z)$. Its coinvariant
relation kills every element, proving (4).
\end{proof}

\begin{proposition}[Proved: a P5 target on at most three states]
\label{CP2:A-coned-three-target-P5}
If $N=P$ and $1\leq |S|\leq3$, the coned bridge $Y$ cannot
have the homotopy type of $\CP^2$. There is no restriction on
the finite center monoid $M$.
\end{proposition}
\begin{proof}
A state fixed by the left $P$ action gives a map $BP\to E$
using empty $M$ strings. Its composite with the projection
to $BP$ is the identity; the state need not be fixed by $M$.
Lemma~\ref{CP2:A-coned-bridge-comparisons}(1) then applies.

Otherwise the left rectangular-band act normal form has
a unique two-point regular row core $U=\{c_0,c_1\}$ and at
most one transient point. The left band acts by
$e_{ij}s=c_i$. Every commuting right $M$ actor fixes $U$
pointwise. Swapping the two band columns defines an
automorphism $F$ of the star that fixes $M$ and $S$.

Use the opposite of the length-two projective resolution
in Lemma~\ref{CP2:A-order-five-corner}, and tensor it with
$\mathbb Q[S]$. Quotienting the acyclic subcomplex supported
on the two core points gives the right $\mathbb QM$ complex
\[
 \mathbb Q[0]\oplus T[1],\qquad T=\mathbb Q[S]/\mathbb Q[U],
\]
with zero differential. The column-interchange semilinear
lift, the opposite of the lift displayed in
Lemma~\ref{CP2:A-mixed-five-reductions}, acts as $+1$ on
$\mathbb Q[0]$ and as $-1$ on $T[1]$. Derived tensoring with
the trivial left $M$ module consequently gives the natural
decomposition
\[
 H_k(E)=H_k(BM)\ \oplus\
 \operatorname{Tor}^{\mathbb QM}_{k-1}(T,\mathbb Q),
\]
with respective $F$ signs $+1$ and $-1$. The second term is
absent in degree zero. The core bridge $E_U=B(P,U,M)$ maps
onto exactly the first summand.

Here the geometric core comparison is needed as well.
The left act $U$ is $Pe_{00}$, a retract of the regular left
act $P$. Its action category is therefore contractible: the
regular action category is contractible, and the act inclusion
and retraction give an actual categorical retract. As $M$
fixes every core point,
\[
 E_U=B(P,U)\times BM\simeq BM.
\]
Its projection to $BP$ factors through the contractible
first factor and is nullhomotopic. This nullhomotopy yields
a map $\Sigma E_U\to Y$ whose connecting homomorphism to
$H_*(E)$ is the core inclusion. Since $H_3(BP)=H_4(BP)=0$,
the cofiber sequence identifies $H_4(Y)$ with $H_3(E)$.
Thus this suspension map realizes the entire $+1$
eigenspace $H_3(BM)$ in $H_4(Y)$.

Suppose $Y\simeq\CP^2$ and let $\alpha$ generate $H^2(Y)$.
If its restriction to $BP$ vanishes, its square is zero by
Lemma~\ref{CP2:A-coned-bridge-comparisons}(1). Otherwise the
nonzero restriction into the one-dimensional $H^2(BP)$,
on which column interchange is $-1$, implies
$F^*\alpha=-\alpha$. Hence $F^*(\alpha^2)=\alpha^2$.
Its pairing with the $-1$ eigenspace of $H_4(Y)$ is zero.
Its pairing with the $+1$ eigenspace is also zero, because
that eigenspace is realized by $\Sigma E_U$, where every
second square vanishes. Over $\mathbb Q$ the two eigenspaces
span $H_4(Y)$; therefore $\alpha^2=0$, a contradiction.
\end{proof}

\begin{proposition}[Proved: a P5 center on at most three states]
\label{CP2:A-coned-three-center-P5}
If $M=P$ and $1\leq |S|\leq3$, the coned bridge $Y$ cannot
have the homotopy type of $\CP^2$. There is no restriction on
the finite leaf monoid $N$.
\end{proposition}
\begin{proof}
The right $P$ fixed states form an $N$ invariant subset.
If it has a single element, that element is also $N$ fixed
and gives a section $BN\to E$, so all second squares vanish.

Suppose there are at least two fixed states. Cardinality
three leaves no two-point regular core orbit in addition,
so the entire band core is the fixed set $F$. Put
$X=\mathbb Q[S]$ and $U=\mathbb Q[F]$. Tensoring the explicit
left $\mathbb QP$ resolution with $X$ gives the $\mathbb QN$
complex
\[
 U\oplus U \xrightarrow{(v,w)\mapsto v+w} X
 \xrightarrow{d_1} U,
 \qquad d_1(s)=s e_{10}-s e_{00},\qquad d_1|_U=0,
\]
in degrees two, one and zero. Quotient its acyclic subcomplex
given by the first $U$ in degree two mapping identically onto
$U\subset X$. The quotient is
\[
 U[2]\ \oplus\ \bigl[X/U\xrightarrow{\bar d_1}U\bigr],
\]
where the bracket occupies degrees one and zero. The bridge
projection to $BN$ uses only the degree-zero augmentation;
it is zero on the displayed $U[2]$ summand. After derived
$N$ coinvariants and passage to the cofiber, $H_3(Y)$ has
the direct summand $U_N$. This is nonzero because $F$ is a
nonempty finite $N$ set. Thus $Y$ does not have the homology
of $\CP^2$.

It remains that there are no fixed states. There is then
one regular two-point column core $U=\{c_0,c_1\}$ and at
most one transient point. Every $N$ actor fixes the core
pointwise. The acyclic-core quotient of the same projective
resolution is
\[
 \mathbb Q[0]\oplus T[1],\qquad
 T=\mathbb Q[S]/\mathbb Q[U].
\]
Under this quotient the projection $E\to BN$ is exactly the
projection onto the first summand. Therefore
$H_2(Y)=\mathbb Q\otimes_{\mathbb QN}T$. If $T=0$ this group
is zero. Otherwise $T$ is one-dimensional: an $N$ actor has
scalar one if it fixes the unique transient state and scalar
zero if it sends that state into the core. Nonzero coinvariants
force every scalar to be one. All states are then $N$ fixed,
so Lemma~\ref{CP2:A-coned-bridge-comparisons}(2) makes $Y$ a
suspension. This excludes the final case.
\end{proof}

\begin{proposition}[Proved: the other five-element centers]
\label{CP2:A-coned-three-other-five}
Let $|M|=5$, $M\not\cong P$, $|N|\leq11$, and
$1\leq |S|\leq3$. The coned bridge $Y$ cannot have type
$\CP^2$.
\end{proposition}
\begin{proof}
The center $BM$ is rationally acyclic by
Lemma~\ref{CP2:A-order-five-corner}. If its minimal ideal
has one row, the trivial left $\mathbb QM$ module is
projective, by normalized local-group averaging. The actual
right-orbit quotient $S\to S/M$ induces a rational homology
and cohomology-ring equivalence from $Y$ to the ordinary
left $N$ act cone on $S/M$. If $Y$ were a projective plane,
this latter cone would have bounded rational homology.
The left-act version of Theorem~\ref{CP2:A-bounded-eleven-acts}
kills every second product there. Pullback gives a
contradiction. This comparison requires only rational
boundedness, not integral projective-plane recognition.

Otherwise the minimal ideal has at least two rows. Since
$M\ne P$ and $|M|=5$, it has exactly one column. First suppose
its local group is trivial. Every core state is then globally
$M$ fixed: right multiplication fixes every element of a
one-column rectangular band. Conversely every fixed state
belongs to the core. Write $F$ for this nonempty fixed set.
If $|F|=1$, commutation makes its point $N$ fixed and gives
the section in Lemma~\ref{CP2:A-coned-bridge-comparisons}(1).
If $F=S$, the $M$ action is trivial. The bridge is
$B(N,S)\times BM$, and rational acyclicity of $BM$ again
gives a rational ring equivalence to the ordinary $N$ act
cone on $S$, excluded by bounded-eleven.

Only $S=\{u_0,u_1,w\}$ with $F=\{u_0,u_1\}$ remains.
The right transformation image fixes both core points and
sends $w$ to $w,u_0$, or $u_1$. If it has only one
nonidentity retraction $p$, then $p$ commutes with the whole
right image and with every $N$ actor. It is a biset retraction
onto two states, so the coned-binary lemma and the cut-retract
theorem exclude an actual projective-plane type.
Otherwise both retractions $p_0(w)=u_0$, $p_1(w)=u_1$ occur.
A map commuting with both is either the identity or a full
constant to $u_0$ or $u_1$: its value on $w$ determines its
values on both core points. If the $N$ image is the identity
alone, $Y$ is a suspension. If that image contains a constant,
the right action space is connected because $w$ reaches both
core points. Lemma~\ref{CP2:A-coned-bridge-comparisons}(4)
then gives $H_2(Y)=0$.

Finally suppose the one-column ideal has a nontrivial local
group. Its size and the outside identity force
$M=(\{0,1\}\times C_2)^1$, with
$(i,g)(j,h)=(i,gh)$. Its core is a disjoint union of $C_2$
orbits of sizes one or two. A single globally fixed state is
$N$ fixed, and the preceding section argument applies. If
all core orbits have one point, the transformation analysis
of the preceding paragraph applies verbatim, since the group
acts trivially on the core. The remaining possibility has a
single regular two-point orbit and at most one transient $w$.
With no transient the common core retraction already has two
states. With a transient, let $p_0,p_1$ be the two ideal
idempotent actors. They retract onto the same two-point core.
If $p_0(w)=p_1(w)$, their common retraction commutes with the
core transposition and all ideal actors, hence with $M$ and
$N$; apply coned-binary. If their values differ, a commuting
map sending $w$ into the core would identify the two core
points, contradicting commutation with the core transposition.
A commuting map fixing $w$ must fix both core points. Thus
the entire centralizer is the identity and the $N$ action is
trivial. Its cofiber is a suspension, completing the proof.
\end{proof}

\begin{corollary}[Proved: both outstanding large star profiles]
\label{CP2:A-star-coned-three-five-eight}
The source-star profiles
\[
 (\text{center},\text{leaf}_1,\text{leaf}_2,
   \text{stalk}_1,\text{stalk}_2)
 =(5,1,8,1,3),\qquad(8,1,5,1,3)
\]
and their leaf interchanges have no classifying space of type
$\CP^2$.
\end{corollary}
\begin{proof}
The singleton trivial leaf cones the center, giving precisely
the model $Y$. In the first profile use
Propositions~\ref{CP2:A-coned-three-center-P5} and
\ref{CP2:A-coned-three-other-five}. In the second profile a
non-$P$ five-element leaf has zero second rational cohomology,
so Lemma~\ref{CP2:A-coned-bridge-comparisons}(1) applies;
the remaining leaf $P$ is Proposition~\ref{CP2:A-coned-three-target-P5}.
\end{proof}

\begin{remark*}[Correction of a discarded route]
No graph model is asserted for a $P$ act with two globally
fixed core states and one transient state. In particular, when
the two row retractions send that transient to different core
states, the displayed projective complex has $H_0=\mathbb Q$,
$H_1=0$ and $H_2=\mathbb Q^2$; it is not an interval. An
initial proposed graph shortcut was withdrawn before use.
The degree-two module summand in the center-$P$ proof is the
replacement argument, and the regular-core proof for a target
$P$ uses the actual contractible left action category only.
\end{remark*}


\subsection{The fresh two-object scope at eighteen arrows}

\begin{lemma}[Proved: a pure categorical rectangular band]
\label{CP2:A-pure-categorical-band}
Let $C$ be a finite category consisting of a categorical
rectangular band with trivial local group and one additional
identity at each object. If $r_i,c_i$ are its positive local
row and column counts, the constant integral left module
has a projective resolution of length at most two. Thus
$H_j(BC;\mathbb Z)=0$ for $j\geq3$.
\end{lemma}
\begin{proof}
Put $R=\mathbb ZC$, choose a minimal idempotent $e$, and
let $P=Re$. At object $i$, $P$ is freely based by the $r_i$
minimal rows. Every band arrow sends every basis element
to its own target row. The augmentation $P\to\mathbf1$
therefore has kernel
\[
 V=\bigoplus_i S_i^{\,r_i-1},
\]
where $S_i$ is the rank-one module supported at $i$ and
annihilated by all band arrows. The representable $R1_i$
surjects onto $S_i$ by taking the coefficient of the extra
identity. Its kernel consists of the band arrows out of $i$.
For each of their $c_i$ source-column labels, the arrows
with that column form a copy of $P$. Hence
\[
 0\longrightarrow P^{c_i}\longrightarrow R1_i
   \longrightarrow S_i\longrightarrow0
\]
is exact over $\mathbb Z$. Splicing these sequences into
$0\to V\to P\to\mathbf1\to0$ gives
\[
 0\longrightarrow\bigoplus_i P^{c_i(r_i-1)}
 \longrightarrow\bigoplus_i(R1_i)^{r_i-1}
 \longrightarrow P\longrightarrow\mathbf1\longrightarrow0.
\]
All displayed terms preceding $\mathbf1$ are projective.
Tensoring with the constant right module proves the claimed
homology vanishing.
\end{proof}

\begin{proposition}[Proved: wholly strong two-object categories through eighteen]
\label{CP2:A-strong-two-eighteen}
A strongly connected two-object category with at most
eighteen arrows cannot have type $\CP^2$.
\end{proposition}
\begin{proof}
An object retract can be deleted, leaving a monoid with
fewer than twenty-three elements, excluded by
Theorem~\ref{A:monoid-minimum23-classification}. Otherwise
both identities are outside the categorical minimal ideal.
Write $r,c$ for its total row and column counts and $G$ for
its local group. If $r=2$ or $c=2$, the local counts on that
side are all one and Lemma~\ref{CP2:A-high-Rees} gives
finite-group homotopy type. In the other cases
$|G|rc+2\leq18$ and $r,c\geq3$, so $G=1$.
If one total is three, its two positive local counts are
one and two, and Lemma~\ref{CP2:A-baseline-binary-representable}
annihilates positive rational products. The sole remaining
possibility has $r=c=4$ and exactly eighteen arrows. Its
sixteen minimal-band arrows and two identities exhaust the
category. Lemma~\ref{CP2:A-pure-categorical-band} gives
$H_4=0$, completing the reduction.
\end{proof}

\begin{proposition}[Proved: rational three-state cones through order twelve]
\label{CP2:A-twelve-three-state-rational}
An ordinary act cone over a monoid of order at most twelve,
on at most three states, cannot have the rational homology
and cohomology ring of $\CP^2$.
\end{proposition}
\begin{proof}
Orders at most eleven follow from
Theorem~\ref{CP2:A-bounded-eleven-acts}. At order twelve,
the relative and two-row act obstructions force at least
three minimal rows and two minimal columns. The local group
is trivial, since a nontrivial one would already require
$3\cdot2\cdot2+1=13$ elements. The possible band shapes and
outside counts are
\[
 (3,2,6),\quad(4,2,4),\quad(5,2,2),\quad(3,3,3).
\]
The first case is excluded by
Proposition~\ref{CP2:A-three-row-six-outside-acts}; the next
two by Lemma~\ref{CP2:A-two-column-four-outside-bounded}.

In the last case there must be exactly three states, since
the binary act obstruction is unconditional. The vanishing
of first homology and a putative nonzero square imply
constant band action by Lemma~\ref{CP2:A-four-state-band}.
Actual constant-action compression gives a monoid with the
same outside algebra $A$, of dimension three, and two
augmentation modules of dimension two. The compressed
monoid has the same rational homology as the given cone,
so its Tor groups are bounded and its second Tor is
one-dimensional. If $A$ is local,
Lemma~\ref{CP2:A-local-three-small-module} excludes this.
If $A$ is nonlocal, Lemma~\ref{A:nonlocal-four} excludes it.
All comparisons just used are actual maps or homotopy
equivalences; no boundedness assertion is inferred from
an isolated degree-two comparison.
\end{proof}

\begin{proposition}[Proved: a corner of order at most two through eighteen]
\label{CP2:A-two-object-small-corner-eighteen}
A two-object collage with at most eighteen arrows and a
corner of order at most two cannot have type $\CP^2$.
\end{proposition}
\begin{proof}
Take the small corner as the left corner $N$. A trivial
corner is an ordinary act cone and is excluded by
Corollary~\ref{CP2:A-ordinary-act-eighteen}. Write $m=|M|$
and $s=|S|$ when $|N|=2$, so $m+s\leq16$.
If $m\leq11$, the trivial right rational $N$ module is
projective, and actual exact-orbit comparison gives a
rational ring equivalence to an ordinary $M$ act cone.
Its homology is bounded under the proposed target, so
Theorem~\ref{CP2:A-bounded-eleven-acts} applies. If
$m\geq13$, there are at most three connecting arrows,
excluded by the universal one- and two-cross obstructions
and Theorem~\ref{CP2:A-three-cross-reduction}.

Only $(m,|N|,s)=(12,2,4)$ remains. If $N=\{1,e\}$ with
$e^2=e$, its actor is a biset retraction onto $eS$. The
cut-retract theorem preserves an actual $\CP^2$ type on
the image. The $N$ action there is trivial and $BN$ is
contractible, so collapsing that corner gives an ordinary
$M$ act cone on at most four states by an actual homotopy
equivalence, a contradiction.

If $N=C_2$ and its action on $S$ is nontrivial, its orbit
set has at most three states. Exact rational orbit comparison
gives a rational ring and homology equivalence to the
ordinary $M$ cone on that set, contradicted by
Proposition~\ref{CP2:A-twelve-three-state-rational}. If the
$C_2$ action is trivial, collapsing $C_2$ gives an actual
categorical retract onto the four-state ordinary cone and
a rational homology equivalence. A retract of an integral
$\CP^2$ model has split integral homology and trivial
fundamental group. This rational equivalence forces its
second and fourth split summands to be $\mathbb Z$, and
pullback transfers the primitive square. Integral recognition
therefore makes the retract an actual $\CP^2$ model, again
contradicting the ordinary-act theorem. The rational and
actual comparisons have thus been kept separate in every
case.
\end{proof}

\begin{proposition}[Proved: the complete reduced eighteen-arrow scope]
\label{CP2:A-two-object-eighteen-scope}
To exclude every category with at most two objects and at
most eighteen arrows, it suffices to exclude the following
twenty-five positive-source collage profiles:
\[
\begin{gathered}
 (5,4,9),\ (5,5,8),\ (5,6,7),\ (5,7,6),\ (5,8,5),\ (5,9,4),\\
 (6,4,8),\ (6,5,7),\ (6,6,6),\ (6,7,5),\ (6,8,4),\\
 (7,3,8),\ (7,4,7),\ (7,5,6),\ (7,6,5),\ (7,7,4),\\
 (8,3,7),\ (8,4,6),\ (8,5,5),\ (8,6,4),\\
 (9,3,6),\ (9,4,5),\ (9,5,4),\\
 (10,3,5),\ (10,4,4).
\end{gathered}
\]
The first corner must have nonzero rational second homology.
This is a complete scope reduction, not a completion claim
for these finite profiles.
\end{proposition}
\begin{proof}
All budgets through seventeen are excluded by
Theorem~\ref{CP2:A-two-object-seventeen}; one-object models
are excluded by the monoid minimum. A remaining two-object
model is connected. Its wholly strong support is excluded
by Proposition~\ref{CP2:A-strong-two-eighteen}; otherwise it
is a collage. If both corner second cohomology groups vanish,
Lemma~\ref{CP2:A-collage-relative-cup} kills its second square.
By opposition choose a positive corner as the source and
write its order $m\geq5$, opposite order $n\geq1$ and
connecting-set size $s\geq1$. Then $m+n+s=18$.
There are exactly seventy-eight such ordered triples.

The preceding small-corner result removes $n\leq2$.
The universal one-, two- and three-cross results remove
$s\leq3$. The remaining inequalities $m\geq5,n\geq3,s\geq4$
leave twenty-eight triples. The cases $(5,3,10)$ and
$(6,3,9)$ are removed by
Theorem~\ref{CP2:A-baseline-right-zero-three-six} when the
opposite corner is the nonprojective-right three-element
monoid $R_3$, and by exact orbit comparison followed by
bounded-eleven for every other three-element corner.
The case $(11,3,4)$ is excluded by
Proposition~\ref{CP2:A-right-zero-four-states} for $R_3$,
and by the same projective endpoint comparison otherwise.
These are exactly the three removed triples, leaving the
displayed twenty-five. The independently executable integer
inventory in \path{monoid/code/two_object18_inventory.py}
records all seventy-eight budgets and their precise status.
\end{proof}

\begin{remark*}[Finite inputs still required at this frontier]
The frozen seventeen-arrow streams do not certify any of
the twenty-five new finite profiles. Positive source catalogs
through order nine are complete; source order ten is newly
required only for $(10,3,5)$ and $(10,4,4)$.

At order ten, the band possibilities are $2\times2$ with
six outside elements, $2\times3$ and $3\times2$ with four,
$2\times4$ and $4\times2$ with two, and $3\times3$ with
one. A nontrivial minimal group can only be $C_2$ with
shape $2\times2$ and two outside elements. Every table of
these shapes must be included before a justified homology
filter; the order-nine positive catalog does not cover them.

The opposite order-nine corner occurs only in $(5,9,4)$.
The acyclic-opposite P5 four-state lemma and exact
projective-right orbit comparison remove respectively
one-row and one-column minimal ideals. The remaining
two-sided order-nine catalog contains $21918$ band tables
and two $C_2$ Rees tables, all $21920$ of which are needed,
including tables of zero second rational homology.
Opposite order eight occurs in $(5,8,5)$ and $(6,8,4)$.
For $(5,8,5)$, Corollary~\ref{CP2:A-P5-eight-five-scope}
removes the one-row cases as well, so the complete $1205$
two-sided-band tables suffice. For $(6,8,4)$, those tables
alone do not yet cover the opposite-corner scope: the
omitted one-row cases require a further proof or their
complete catalogs.
No complete eighteen-arrow exclusion is asserted here.
\end{remark*}


\subsection{An acyclic opposite corner on five P5 states}

\begin{proposition}[Proved: the five-state acyclic-opposite P5 gate]
\label{CP2:A-P5-five-state-acyclic-opposite}
Let $N$ be any finite monoid with rationally acyclic
classifying space. A left-$N$, right-$P_5$ biset on at most
five states cannot give a collage of type $\CP^2$.
\end{proposition}
\begin{proof}
Write $P=P_5$, $X$ for the collage space, and $E=B(N,S,P)$
for its bridge. A right $P$ fixed state is excluded by
Lemma~\ref{CP2:A-mixed-five-reductions}(1).
In its absence the rectangular-band normal form has only
regular two-point core orbits. A unique such orbit is
excluded, at any state count, by part (3) of the same lemma.
With five states there can be at most two core orbits.
If there are exactly two and no transient, the full-regular-core
lemma applies. It remains to consider the four core states
$(a,j)$, with $a\in\{0,1\}$ indexing the orbits and
$j\in\{0,1\}$ the column, and one transient $w$.

For each band row $i$ there is an orbit $a_i$ such that
$we_{ij}=(a_i,j)$. If $a_0\ne a_1$, this right act is
the regular right $P$ act: send $w$ to $1$ and
$(a_i,j)$ to $e_{ij}$. Every commuting left $N$ actor
is consequently left multiplication by an element of $P$,
and the whole left action is specified by a homomorphism
$N\to P$. The projection $B(N,P,P)\to BN$ is a homotopy
equivalence. Indeed its right-action bar is the contractible
regular action category, and the equivariant map to a point
induces an equivalence after the other bar realization.
The homotopy pushout $BP\leftarrow E\to BN$ is therefore
homotopy equivalent to $BP$, which is not a projective plane.

Now suppose $a_0=a_1$. Row interchange
$\sigma(e_{ij})=e_{1-i,j}$ fixes the right action on every
state and defines a category automorphism $F$ fixing $N,S$.
Let $U$ be the four-point core and $A=\{0,1\}$ its orbit set,
with the left $N$ action induced on that set. Tensor the
length-two left projective $P$ resolution of
Lemma~\ref{CP2:A-order-five-corner} with $\mathbb Q[S]$.
Its first differential is zero, because the two row actors
agree. Its second differential maps isomorphically from
the two copies of the column-zero core onto $\mathbb Q[U]$.
Quotienting that acyclic subcomplex gives the $N$ complex
\[
 \mathbb Q[A][0]\oplus T[1],
 \qquad T=\mathbb Q[S]/\mathbb Q[U],
\]
with zero differential. The explicit semilinear lift in
Lemma~\ref{CP2:A-mixed-five-reductions} acts by $+1$ on
the first summand and $-1$ on the second. Thus the core
bridge $E_U=B(N,U,P)\to E$ induces exactly the $+1$
summand in rational homology, including in degree three.

Let $X_U$ be the full-core subcollage. The natural
nullhomotopy of its projection $E_U\to BP$ from
Lemma~\ref{CP2:A-P5-full-regular-core} gives
\[
 X_U\simeq BP\vee
 \operatorname{hocofib}\bigl(B(N,A)\longrightarrow BN\bigr).
\]
To justify the second factor, the equivariant map from the
right core-action category to its orbit set has contractible
fibers, each the action category of the regular column act
$e_{00}P$. It is a homotopy equivalence before and after
the $N$ bar. The second wedge summand is therefore the
ordinary left $N$ act cone on the two-element set $A$.
Its degree-two products vanish by the unconditional binary
act theorem. Those on $BP$ vanish as well, so every
degree-two product on $X_U$ is zero.

Since $BN$ is rationally acyclic and $BP$ has no third or
fourth rational homology, the two pushout sequences identify
$H_4(X_U)$ and $H_4(X)$ with $H_3(E_U)$ and $H_3(E)$,
respectively. Hence the map $H_4(X_U)\to H_4(X)$ realizes
the entire $+1$ eigenspace of $F$.

If $X$ had type $\CP^2$, let $\alpha$ generate its rational
second cohomology. A zero restriction to $BP$ would make
its square zero by the relative-cup lemma, since its
restriction to $BN$ already vanishes. Otherwise the
$-1$ action on $H^2(BP)$ forces $F^*\alpha=-\alpha$.
Its square is $F$ invariant and therefore pairs to zero
with the $-1$ eigenspace of $H_4(X)$. It also pairs to zero
with the $+1$ eigenspace: that space is in the image of
$H_4(X_U)$, and the restriction of $\alpha^2$ to $X_U$
vanishes. Thus $\alpha^2=0$, a contradiction.
\end{proof}

\begin{corollary}[Proved: the opposite-order-eight P5 scope at eighteen]
\label{CP2:A-P5-eight-five-scope}
For the profile $(5,8,5)$ with positive source corner, it
suffices to test all $1205$ ordinary order-eight monoids with
at least two minimal rows and two minimal columns as the
opposite corner. Every one of those tables is required,
including the ones with zero second rational homology.
\end{corollary}
\begin{proof}
The positive five-element source is $P_5$. A one-column
opposite ideal is removed by exact projective-right orbit
comparison and the bounded-eleven act theorem. A one-row
ideal makes the opposite classifying space rationally
acyclic, so the preceding proposition applies. The remaining
ideals have trivial group by the order-eight budget and
shapes $2\times2$, $2\times3$, $3\times2$. Their complete
catalogs contain $1181+12+12=1205$ tables, exactly as in
Proposition~\ref{CP2:A-P5-eight-corner-scope}. This is an
input reduction; no new five-state finite search is claimed
by this corollary.
\end{proof}


\subsection{Complete four- and five-object cardinality scopes at eighteen}

\begin{proposition}[Proved: the remaining small-corner finite scopes]
\label{CP2:A-four-five-eighteen-scopes}
For connected acyclic support on four or five objects at exactly eighteen
arrows, the following is a complete reduction to finite composition
problems. Profiles are identified under object isomorphism and
opposition; corner entries are retained as cardinalities.
\[
\begin{array}{c|rrrr}
\text{objects}&\text{all profiles}&\text{proved profiles}&
\text{finite profiles}&\text{largest remaining corner}\\\hline
4&61420&58696&2724&7\\
5&200737&199775&962&6
\end{array}
\]
The proposition does not exclude the remaining finite profiles.
\end{proposition}
\begin{proof}
Put every connected acyclic support in a topological order.
There are eighteen such labelled supports on four objects, in
seven isomorphism/opposition orbits, and 181 on five objects,
in twenty-five orbits. Assign positive integers to the diagonal
and the supported strict positions, with sum eighteen. The full
labelled decoration totals are respectively 320008 and 3525426.
Quotient each representative support's decorations by its full
automorphism and anti-automorphism group. The resulting numbers
of canonical profiles are 61420 and 200737.

For each profile, every invoked predicate is interpreted in
the preceding reduction proofs: support first homology,
singleton joins, relative two-block products, fixed stalks,
binary extreme stalks, the small excess and semisimple bounds,
actual singleton-comma deletion, and edge-tree products.
On four objects the relative semisimple bar has height at most
three; on five objects the separate eighteen-arrow budget in
Lemma~\ref{CP2:A-five18-semisimple-small} supplies the required
argument. The one-large-corner tree theorem is already part of
the original four-object partition. These predicates prove
58662 four-object profiles and 199743 five-object profiles.

For five objects, Proposition~\ref{CP2:A-tree-one-large-corner-eleven}
removes another twenty-eight original residual profiles.
Proposition~\ref{CP2:A-pendant-small-tree-eighteen} then removes
four more. For four objects the same pendant proposition removes
thirty-four profiles, including both remaining order-eight
corner cases. Every supplemental row retains its literal matrix,
previous index, and concrete pendant or tree witness. Thus
\[
61420=58662+34+2724,\qquad
200737=199743+28+4+962.
\]

The independent programs
\path{converse/scope18/check_four18_scope.py} and
\path{converse/scope18/check_five18_scope.py} import no producer
or profile-enumeration code. They reconstruct supports by moving
old vertices to new vertices, generate positive compositions by
separators, remove entire decoration orbits, and compare every
literal matrix, orbit weight, and complete proof-reason set.
The separate pendant checker verifies every supplemental witness
and exact predecessor-ledger binding. Their completed receipts
are \path{converse/scope18/four18_scope_independent_audit.json},
\path{converse/scope18/five18_scope_independent_audit.json}, and
\path{converse/scope18/pendant18_supplements_independent.json}.
Every residual corner therefore lies in the independently complete
ordinary monoid catalogs through order seven. Complete enumeration
of the connecting actions and compatible strict compositions,
with independently checked literal obstruction streams, is still
required for each residual profile.
\end{proof}


\subsection{Cyclic support on three and five objects at eighteen arrows}

\begin{lemma}[Proved: a binary collage into the five-element band monoid]
\label{CP2:A-P5-binary-source-block}
Let a finite category be a collage of a source category $S$ and
a sink corner $P_5$, the identity-adjoined $2\times2$ rectangular
band. Suppose every stalk $F(x)=C(x,*)$, for $x\in S$, is nonempty
and has at most two elements. If $H^2(BS;\mathbb Q)=0$, all rational
degree-two products on $BC$ vanish.
\end{lemma}
\begin{proof}
Write the band elements as $e_{ij}$, with
$e_{ij}e_{kl}=e_{il}$. If any stalk contains a point fixed by all
of $P_5$, that arrow gives a natural transformation from the
constant functor at its source object to the sink-corner inclusion.
The inclusion $BP_5\to BC$ is nullhomotopic. Every global second
class therefore restricts to zero on both blocks, and
Lemma~\ref{CP2:A-two-block} kills every product.

Assume no stalk has a $P_5$-fixed point. Every stalk then has two
elements. An idempotent transformation of a two-element set is
either the identity or a constant. If one band element acts as
the identity, the rectangular-band identities force all of them
to do so: in its row and column their products with that identity
give the identity, so those idempotents are bijective; then all
remaining entries are their products. This would give fixed
points. Thus every band element acts constantly. The relation
$e_{ij}e_{kl}=e_{il}$ shows that the constant value depends only
on $i$. The two row values must be distinct, since otherwise their
common value is fixed. Hence each stalk is canonically the
two-element row act $U=P_5e$ for a minimal idempotent $e$.

Every map between two stalks commutes with these two row constants,
and must preserve both row values. Under the specified
identifications all source-category arrows therefore act by the
identity on $U$. The cross-arrow bar is consequently
$B(P_5,U)\times BS$, and its projection to $BS$ is a homotopy
equivalence. Indeed $U=P_5e$ is an equivariant retract of the
regular left act $P_5$, by inclusion and $m\mapsto me$.
The action category of the regular left act has initial object
$1$, so its classifying space is contractible; its retract
$B(P_5,U)$ is contractible as well.

The actual collage homotopy pushout now identifies $BC\simeq BP_5$,
because the cross-arrow bar maps equivalently to the source block.
The latter space has type $S^2$, by
Corollary~\ref{CP2:A-five-types}, and its degree-two products are
zero. This proves the remaining case.
\end{proof}

\begin{theorem}[Proved: cyclic three-object support through eighteen]
\label{CP2:A-strong-three-eighteen}
A finite category on three objects with cyclic support and at most
eighteen arrows cannot have classifying space of type $\CP^2$.
\end{theorem}
\begin{proof}
Assume a connected counterexample. An object-retract deletion
saves at least three arrows and leaves at most fifteen; the
complete sixteen-arrow baseline excludes it. Thus there is no
object retract. If all three objects form a strong component,
Lemma~\ref{CP2:A-high-Rees} gives finite-group type unless both
total minimal row and column counts are at least four. In that
case the three identities are outside the ideal and there are
at least $4\cdot4+3=19$ arrows, a contradiction.

There remains a proper strong pair $S$ and one corner monoid $N$.
By opposition take $S$ to be the source component. Write
$q=|S|$, $m=|N|$, and $h$ for the total outgoing stalk cardinality.
Both stalks are nonempty, so
\[
 q+m+h\le18,\qquad q\ge6,\qquad h\ge2.
\]
If $h=2$, the full collage is a categorical join and has zero
positive rational products. Hence assume $h\ge3$. If both
$H^2(BS;\mathbb Q)$ and $H^2(BN;\mathbb Q)$ are zero, the
two-block lemma applies.

Suppose first that $H^2(BS;\mathbb Q)\ne0$. Its total minimal
row and column counts $r,c$ are both at least three; otherwise
one local count is one at each object and $S$ has finite-group
type. Since $q\le14$, the minimal group is trivial and
$rc+2\le q$ leaves only
\[
 (r,c)=(3,3),(3,4),(4,3).
\]
If $m\le2$ and each local row count is at most two, the trivial
right rational $N$-module is projective and
Corollary~\ref{CP2:A-baseline-two-row-instances} annihilates the
collage's second products. The only remaining local-row
distribution is $(1,3)$, in either order, with $(r,c)=(4,3)$.
It forces $q=14$, so the internal category consists exactly of
its twelve-element minimal band and its two identities.
At its one-row object $b$, a local minimal idempotent fixes every
incoming arrow from the other object, and the only unit is the
identity. There are no external entrances into $S$. Minimal cross
arrows supply the factorization of this idempotent required in
Lemma~\ref{CP2:A-high-one-sided}. Its one-sided object deletion
is an actual homotopy equivalence. It removes at least the two
corner arrows and two internal cross arrows, leaving at most
fourteen arrows, which the complete baseline excludes.

If $m\ge3$, then $q\le12$, so necessarily $r=c=3$. For $q=11$
the pair is again exactly its minimal band and its identities,
and the same one-row object deletion applies. For $q=12$ the
budget forces $m=3,h=3$. Its two outgoing stalks have sizes one
and two. Lemma~\ref{CP2:A-high-small-binary}, with source block
$S$ and the order-three sink, makes the square of its single
global second class zero. This exhausts the case of nonzero
second cohomology on $S$.

Finally suppose $H^2(BS;\mathbb Q)=0$ and
$H^2(BN;\mathbb Q)\ne0$. Then $m\ge5$. A singleton outgoing
stalk is fixed by all $N$ actors and the fixed-arrow corollary
applies. Therefore both stalks have size at least two, giving
$h\ge4$ and $q\le9$. The minimal group of $S$ is trivial:
a nontrivial group would already require
$2\cdot2\cdot2+2=10$ internal arrows. Thus
$rc+2\le9$ leaves either $r=2$, or $r=3,c=2$.

When $r=2$, the source-component collapse in
Lemma~\ref{CP2:A-high-one-sided} replaces $S$ by a trivial
object with the actual outgoing orbit act. The resulting space
is an ordinary act cone over $N$, with $m\le8$ from
$q\ge6,h\ge4$. Its homology is bounded, and the bounded-eleven
act theorem excludes its nonzero square. If $r=3,c=2,q=8$,
the pair is its six-element band plus identities; the one-row
object deletion above again applies.

The only new numerical case is $q=9,m=5,h=4$. The order-five
corner with nonzero second cohomology is necessarily $P_5$,
and the two stalks both have size two. Lemma~\ref{CP2:A-P5-binary-source-block}
applies since $H^2(BS;\mathbb Q)=0$. Its zero-product conclusion
contradicts the projective-plane square. All cyclic-support
possibilities have been excluded.
\end{proof}

\begin{corollary}[Proved: cyclic five-object support through eighteen]
\label{CP2:A-strong-five-eighteen}
A finite category on five objects with cyclic support and at
most eighteen arrows cannot have classifying space of type $\CP^2$.
\end{corollary}
\begin{proof}
Here $n\le18=2(5)+8$. Corollary~\ref{CP2:A-high-nine-support}
states that a projective-plane candidate in this range has acyclic
support. This contradicts the hypothesis directly.
\end{proof}


\subsection{Fresh order-ten positive source catalogs}

\begin{proposition}[Proved: the nontrivial-group part at order ten]
\label{CP2:A-C2-Rees-order10-sources}
Exactly fourteen ordinary isomorphism classes of order-ten
monoids have both nonzero rational second homology and a
nontrivial maximal subgroup in the minimal ideal. Their
minimal ideal has shape $2\times C_2\times2$. In the two
normalized sandwich families $b=0,1$, the full counts and
positive-second-homology counts are
\[
\begin{array}{c|r|r|r|r}
b&\text{compatible outside actors}&
\text{labelled tables}&\text{ordinary classes}&H_2(\mathbb Q)\ne0\\\hline
0&32&25&14&8\\
1&16&19&8&6.
\end{array}
\]
The twenty-two full classes in this table concern exactly
these two ideal shapes, not all order-ten monoids.
\end{proposition}
\begin{proof}
Positive rational second homology requires at least two
minimal rows and columns. At order ten a nontrivial local
group consequently forces two rows, two columns and group
$C_2$, with the identity and one further element $a$ outside
the eight-element ideal. Normalize that ideal as
\[
 I_b=\{(i,g,j):i,g,j\in\{0,1\}\},\qquad
 (i,g,j)(k,h,l)=(i,g+h+bjk,l),
\]
where group coordinates are taken modulo two. These are the
two normalized sandwich forms from
Proposition~\ref{CP2:A-positive-corners-nine}.
They remain distinct as ideal semigroups: the idempotents
are closed under products exactly when $b=0$.

Left and right multiplication by $a$ necessarily have the
forms
\[
 a(i,g,j)=(\sigma(i),g+\delta_i,j),\qquad
 (i,g,j)a=(i,g+\varepsilon_j,\tau(j)),
\]
for maps $\sigma,\tau:\{0,1\}\to\{0,1\}$ and two pairs of
group coordinates $\delta,\varepsilon$. The equation
$(x a)y=x(a y)$ is exactly
\[
 \varepsilon_j+b\,\tau(j)k
   =b\,j\sigma(k)+\delta_k\pmod2\qquad(j,k\in\{0,1\}).
\]
The other equations with one occurrence of $a$ follow from
these multiplier formulas. Try all $4^4=256$ parameter
tuples and then all ten possible values of $a^2$, retaining
precisely the associative literal tables. The new producer
\path{monoid/code/source10_c2_rees.py} performs exactly this
finite calculation. Canonicalization fixes the identity
and $a$ and ranges over every ideal automorphism, giving the
displayed counts.

The independent checker
\path{monoid/code/check_source10_c2_rees.py} does not import
the producer or assume its multiplier parameterization.
A left multiplier is determined by its images on the two
row generators $(i,0,0)$, because every ideal element is
$(i,0,0)(0,g,0)(0,0,j)$. The checker tries all $8^2$ such
image pairs and tests every literal left-multiplier equation.
It constructs right multipliers dually from the two column
generators. There are sixteen on each side. It then tries
every pair, testing the linking equation and their
commutation, and every square value in reverse order.
Every completed table is checked for full associativity.

For an independent isomorphism scope, it tries every
bijection on the four idempotents and every bijection on
their four-element complement, testing all ideal products.
Each ideal has exactly four automorphisms. The checker
therefore reaches every ordinary class without using the
producer's Rees-coordinate automorphism formula. Both
complete labelled-table sets, their exact digests, every
canonical class and every orbit multiplicity agree.

Finally the checker computes every principal two-sided
ideal, confirms that the displayed eight elements are
exactly the minimal ideal, and checks all one hundred
product equations for the quotient which forgets the group
coordinate. That six-element quotient is a rational
homology equivalence by Lemma~\ref{A:minimal-group-collapse}.
Its normalized bar boundaries are reduced by exact primitive
integer elimination. All forty-four labelled tables have
their rational second Betti numbers checked in this way.
These agree with the producer's rank-one balanced-form
condition $\sigma(1)-\sigma(0)=\tau(1)-\tau(0)$.
The independent receipt records $165868$ evaluated candidate
associativity equations, $4400$ quotient product equations,
and all forty-four exact bar checks. The literal catalog and
receipt are
\path{monoid/certificates/source10_c2_rees.json} and
\path{monoid/certificates/source10_c2_rees_verified.json}.
\end{proof}

\begin{remark*}[The separate rectangular-band source obligation]
The proposition completes only the nontrivial-group portion
of the order-ten positive source scope. The six rectangular
band shapes
$2\times2,2\times3,3\times2,2\times4,4\times2,3\times3$
remain separate input families. The fresh positive-only
producer \path{monoid/code/positive_corners_ten.cpp} computes
the exact balanced-form space of each outside actor multiset
before assigning any outside products. This space depends
only on the actors, so a zero space removes the whole family
without discarding a positive-second-homology monoid.
For a retained family every compatible product is still
enumerated and associativity enforced, and ordinary
canonicalization still ranges over all row, column and
outside permutations. A completed producer alone does not
certify the source catalog: an independently complete
enumeration and exact table/source join are still required.
No complete order-ten band count is asserted in this remark.
\end{remark*}


\subsection{Two binary stalks and five intermediate arrows}

\begin{proposition}[Proved: five labels cannot support a source cycle]
\label{A7:two-binary-five-labels}
Let $\mathcal C$ be a finite acyclic category with full support
$0<1<2<3<4$, and write $h_{ij}=|\mathcal C(i,j)|$.
Every rational degree-two cohomology product of $B\mathcal C$
is zero if
\[
 h_{01}=h_{02}=2,\qquad h_{13}\leq5.
\]
The same conclusion holds under the dual conditions
$h_{34}=h_{24}=2$ and $h_{13}\leq5$.
There is no restriction on the two long hom-sets or on $h_{04}$.
\end{proposition}
\begin{proof}
Use $A=\mathcal C(0,1)$, $U=\mathcal C(0,2)$,
$B=\mathcal C(1,2)$, $V=\mathcal C(1,3)$ and
$C=\mathcal C(2,3)$, and let
$P:A\times B\to U$, $Q:B\times C\to V$ be composition.
If every map $P(-,b):A\to U$ is constant, the existing
short-composite obstruction, Proposition~\ref{A7:short-composite-gates},
proves the assertion. Otherwise some $b_*\in B$ induces a
bijection, since both $A$ and $U$ have two elements.

Work first in the universal short-product category $F(P,Q,R)$.
Its $03$ hom-set $L$ is the quotient of
$(A\times V)\amalg(U\times C)$ by the relations
\[
 (a,Q(b,c))\sim(P(a,b),c).
\]
On the vertex set $L$, form the simple undirected graph $G$
whose edges are the distinct nonconstant image pairs of all
maps $A\to L$ indexed by $V$ and all maps $U\to L$ indexed
by $C$. Loops are discarded and parallel image pairs give one
edge. We show that a cycle in $G$ would require at least six
\emph{distinct} elements of $V$.

To keep track of the quotient, initially take one abstract
interval $I_v$ with endpoint set $A\times\{v\}$ for each
$v\in V$, and one interval $J_c$ with endpoint set
$U\times\{c\}$ for each $c\in C$.
First impose only the relations belonging to bijective columns
$P(-,b)$. Each such relation identifies both endpoints of
$J_c$ with the corresponding endpoints of $I_{Q(b,c)}$,
either preserving or reversing their order. Partition all these
intervals into classes under these whole-interval identifications.
Within each class the endpoint quotient has either two vertices
or one vertex. The latter case includes an orientation-reversing
identification around a chain of intervals. Regard a two-vertex
class as a single interval and a one-vertex class as a point;
no topological quotient with an extra midpoint is being used.
Different classes at this stage have disjoint vertex sets.
Every class contains an interval $I_v$, because the chosen
bijective column $b_*$ attaches every $J_c$ to such an interval.

The remaining columns are constant. If $P(-,b)$ has value
$u\in U$, their relations identify both endpoints of
$I_{Q(b,c)}$ with the $u$ endpoint of $J_c$. Mark every
interval class which contains an $I_v$ occurring on this
constant-column side. Each marked class is forced to collapse
to a single final vertex. After these collapses, every further
identification between different initial classes is incident to
the point of a marked class. All equivalences used here are
exactly the displayed defining endpoint relations for $L$.
In particular, the vertices of $G$ are their final vertex classes.

A nonloop edge of $G$ is the image of an uncollapsed interval
class. It therefore contains a label $v\in V$ belonging to
that class, and different edges can be represented by different
classes and hence different labels. A marked class cannot
represent a nonloop edge.
If a final vertex is incident to edges coming from two different
interval classes, its endpoint equivalence class must contain
a marked class: before the constant-column relations the two
classes were disjoint, and every relation joining different
classes is incident to a marked class. A marked class has just
one final vertex, even if it anchors arbitrarily many merges.
Consequently different such final vertices require different
marked classes and different labels from $V$.

Suppose now that $G$ has a simple cycle of length $m\geq3$.
Its $m$ distinct edges require $m$ distinct labels from
uncollapsed classes. At each of its $m$ distinct vertices,
the two incident cycle edges come from different interval
classes, so that vertex requires a label from a marked class.
Those $m$ marked labels are pairwise distinct and are disjoint
from the $m$ edge labels. Thus $|V|\geq2m\geq6$, a
contradiction. Reversed identifications were accounted for by
the possible point classes; parallel edges cannot create a
cycle in the simple graph relevant to the tree condition.
This proves that $G$ is a forest, including all its isolated
vertices.

Extend this forest to a tree on the full finite set $L$.
Use the unique two-vertex trees on $A$ and $U$, and the point
on the singleton extreme stalk $F(0,4)$. The maps $A\to U$
carry an edge to an edge or a vertex. The same holds for every
map from $A$ or $U$ to $L$ by the definition of $G$, and all
maps to the final singleton are constant. These are all the
source-act maps, so Lemma~\ref{A7:tree-act} proves that all
positive rational products on $BF$ vanish. The graph argument
did not use $R$. The universal transfer theorem,
Theorem~\ref{A7:universal-short-all-fields}, therefore proves
the desired degree-two product vanishing for every completion
and hence for $\mathcal C$. Reversing the object order gives
the stated dual condition.
\end{proof}


\subsection{Completion of the five-object eighteen-arrow case}

\begin{theorem}[Proved: five objects are excluded through eighteen arrows]
\label{CP2:A-five-eighteen-complete}
No finite category on five objects with at most eighteen
arrows has classifying space of type $\CP^2$.
\end{theorem}
\begin{proof}
Theorem~\ref{CP2:A-global-lower-eighteen} handles every smaller
budget. At eighteen arrows cyclic support is excluded by
Corollary~\ref{CP2:A-strong-five-eighteen}. For connected
acyclic support, Proposition~\ref{CP2:A-four-five-eighteen-scopes}
gives exactly 200737 canonical cardinality profiles, of which
199775 have the specified structural proofs. It remains to
check its unchanged ledger of 962 finite profiles, all with
corner order at most six.

Choose all corner monoids from the independently complete
ordinary-isomorphism catalogs. On each supported strict hom-set,
enumerate all commuting endpoint actions up to permutation of
that set. These permutations are independent between hom-sets,
so every category table admits this normalization. With all
actions fixed, enumerate every compatible product of two strict
arrows, imposing all typed associativity equations. No further
quotient of those composition maps is taken. The product of
the action-orbit weights recovers the labelled-table count with
the corner representatives fixed.

The sole enabled corner-table cut at five objects requires
every rational corner algebra to be semisimple and the strict
support height to be at most three. The relative bar complex
then has no fourth homology. Height four is explicitly guarded
against; no unproved five-object version of this cut is used.
All inherited projective and act-cone table cuts are disabled.
The exact producer scope and capacity bounds are recorded in
\path{code/multiobject18_engine_scope.md}.

All 962 profiles completed with successful producer exits and
closed literal streams. They cover 11322472 weighted tables and
413630 retained obstruction records. The independent scalar
checker verifies 410607 nonzero second cocycles with zero-square
primitives and 3023 positive third-homology obstructions, with
no unexcluded record. Each outcome contradicts the mod-two
homology ring of $\CP^2$. The checker evaluates 127338835 typed
associativity equations in these retained tables, in addition
to its normalized bar and cochain checks.

The separately written root finalizer
\path{code/check_five18_complete_join.py} imports no runner or
producer. It joins the full cardinality audit and both exact
supplements to the 962-row ledger, rehashes every bound source,
proof, catalog, binary and regression receipt, and requires all
962 individual completions. It checks each producer exit and
footer, checker exit and literal output, and both compressed
and uncompressed stream hashes. It independently reconstructs
the source-target cardinality matrix of every literal table.
Its full join has no missing profile and no survivor; the
compact receipt is \path{outputs/five18_complete_independent.json}.
Thus the proof partition and finite streams exhaust the entire
five-object scope, including every support type.
\end{proof}

\begin{corollary}[Proved: the object scope at eighteen]
\label{CP2:A-eighteen-at-most-four-objects}
Any eighteen-arrow category of type $\CP^2$ has two, three,
or four objects.
\end{corollary}
\begin{proof}
The monoid minimum excludes one object. The preceding theorem
excludes five. For $k\geq6$, we have
$18\leq\min\{21,2k+7\}$, so
Theorem~\ref{CP2:A-high-all-eight} excludes the category.
\end{proof}


\subsection{Completion of the order-ten positive source scope}

\begin{theorem}[Proved: the complete positive-second-homology catalog at order ten]
\label{CP2:A-positive-corners-ten}
Exactly $157145$ ordinary isomorphism classes of ten-element
monoids have nonzero rational second homology. The complete
partition by minimal-ideal shape is
\[
\begin{array}{c|r}
\text{minimal ideal}&\text{ordinary positive-second-homology classes}\\\hline
2\times2&151930\\
2\times3&2584\\
3\times2&2584\\
2\times4&16\\
4\times2&16\\
3\times3&1\\
2\times C_2\times2&14.
\end{array}
\]
These are ordinary isomorphism classes; opposition is not
identified. The completed catalog is a source-input theorem,
not a completion claim for an eighteen-arrow biset profile.
\end{theorem}
\begin{proof}
The minimal-ideal coordinates and rational local-group
collapse show that positive second homology requires at
least two minimal rows and two minimal columns. A group
monoid has no positive rational homology, so its identity
is outside the proper ideal. At order ten, the trivial-group
possibilities are exactly the six band shapes in the table.
A nontrivial group forces $2\times C_2\times2$ and is already
completely counted by
Proposition~\ref{CP2:A-C2-Rees-order10-sources}.

For each band normalize its elements as $e_{ij}$. Include
the identity actor and enumerate every multiset of row and
column transformation pairs for the other outside labels.
If $V,W$ are the row and column augmentation modules,
Lemma~\ref{A:band-Tor} identifies $H_2(M;\mathbb Q)$ with
$W\otimes_A V$, where $A$ is the contracted outside algebra.
The balancing relations depend only on the actor pairs,
not on the yet unassigned outside products. Their exact
rational rank is computed in dimension at most four.
A zero tensor removes the whole actor family; every
positive family remains. This exact test reduces the
$15504$ actor multisets in the largest $2\times2$ shape
to $252$ positive ones.

Every product involving the identity or band is forced.
For each remaining product, its possible values are exactly
the labels having the required composite actor pair.
Associativity is then enforced on all triples of nonidentity
outside labels. This suffices for full associativity:
when the band factor is on the right or left, the equation
is precisely the prescribed composite row or column actor;
when it is in the middle, the independent row and column
actions commute. Equations with two band factors follow
from the rectangular product, and identities were fixed
from the start.

The separately named producer
\path{monoid/code/positive_corners_ten_sparse.cpp} completes
every one of these finite scopes. It still canonicalizes
under all row, column and outside permutations. An exact
cache suppresses repeated canonicalization only, never a
product-search branch. Within one fixed actor family its
outside-product block has at most twenty-five entries,
each in $\{0,\ldots,9\}$. Encoding those entries by four
bits apiece is an injective one-hundred-bit word. A cache
hit therefore identifies an actual previously represented
table under a verified actor-family automorphism; no hash
collision can remove a distinct word. Ordinary class
canonicalization itself remains the full permutation test.

Two earlier newly authored producer versions remain as
controls: one has no cache, and one propagates every
associativity triple. On the completed order-nine
$2\times2$ scope, all three versions produce exactly the
same $7784$ table words, with the same $1190521$ DFS nodes
and $158631$ labelled completions. This is a regression
check, not the justification of order-ten completeness.

The independent checker
\path{monoid/code/check_positive_ten.cpp} uses fixed
reverse-order product assignments and never propagates
unknown cells. It independently reconstructs the exact
tensor, checks every completed input table for full
associativity and its displayed minimal ideal, and
constructs every band-preserving relabelling of each
reported ordinary class. For each actor family these
give an exact expected set of literal outside-product
words. Its reverse DFS must hit exactly that set.
Membership excludes an omitted class; complete coverage
excludes an extraneous class; intersecting relabelled sets
exclude distinct reports of the same ordinary class.
Thus its scope check also binds every literal table to
the claimed ordinary-isomorphism catalog.

The first attempt to apply that checker to the largest
order-ten family was killed with signal nine before
producing a receipt. It is not a premise here.
The separately named
\path{monoid/code/check_positive_ten_compact.cpp} preserves
the independent reverse DFS but stores input tables as
bytes and its exact completion words in sorted vectors,
with a bit vector recording visits. The full order-nine
independent results agree exactly, including every scope
count. The compact checker then completes all $252$
positive order-ten $2\times2$ families: $10981773$
literal tables, $537351915$ independent DFS nodes,
$151930000$ full input associativity equations and
$20010761607$ fully known outside associativity equations.
Its elapsed time was about $630$ seconds. The original
checker completes the other five families, and all six
producer/checker count and digest joins agree.

The independently checked band total is therefore
$157131$, and adding the fourteen local-group cases gives
$157145$. Five separate malformed, omitted or duplicate
catalog controls are rejected by both checker versions.
The final source join
\path{monoid/code/join_source10_catalog.py} rehashes every
finished catalog, producer receipt and independent receipt,
and binds them to their actual source versions. It also
recomputes every group-case representative from its actor
data and checks its ordinary canonical word, closing that
metadata-to-table link separately. The complete receipt is
\path{monoid/certificates/source10_complete_manifest.json}.
Its two derived TSVs contain the complete positive catalog
and the literal band coordinates used by later biset
enumerations. No interrupted checker prefix enters this
receipt.
\end{proof}

\begin{remark*}[Status relative to the preceding source note]
This theorem completes the six band obligations left open
in the preceding order-ten source note. The old source note
was preserved while a checkpoint PDF was building; the
separate completion theorem records the later finished
checks without changing the earlier frozen build inputs.
The twenty-five eighteen-arrow finite biset scopes remain
separate obligations.
\end{remark*}


\subsection{A six-element source and a one-row opposite corner}

\begin{lemma}[Proved: the positive six-element source shapes]
\label{CP2:A-six-positive-shapes}
Let $M$ have six elements and $H_2(BM;\mathbb Q)\ne0$.
Its minimal ideal is the $2\times2$ band
$K=\{e_{ij}:i,j\in\{0,1\}\}$, and
$M=K\amalg\{1,a\}$. Up to independent relabelling of the
rows and columns, the two actors of $a$ and its square are
exactly one of
\[
\begin{array}{c|c|c}
\sigma_a&\tau_a&a^2\\\hline
\mathrm{id}&\mathrm{id}&1\\
\mathrm{id}&\mathrm{id}&a\\
(01)&(01)&1\\
\mathrm{const}_0&\mathrm{const}_0&e_{00}\\
\mathrm{const}_0&\mathrm{const}_0&a.
\end{array}
\]
In each case $H_2(BM;\mathbb Q)$ is one-dimensional and
every rational degree-two square on $BM$ is zero.
\end{lemma}
\begin{proof}
The minimal-ideal and local-group reductions require at
least two rows and two columns. Six elements leave only
the displayed band and its two outside elements; a
nontrivial local group is too large. The two augmentation
modules are one-dimensional. Their balanced tensor is
nonzero exactly when the augmentation scalars of every
row and column actor agree. For a transformation of two
points the scalar is $1,-1,0$ for the identity, interchange
and a constant, respectively. Thus the actors of $a$ are
the three paired types in the table. The required actors
of $a^2$ then leave exactly the displayed products.
In the constant case the constant labels can independently
be normalized to zero. The balanced tensor has dimension
one, and Lemma~\ref{A:rank-one-square} gives the square
assertion.
\end{proof}

\begin{theorem}[Proved: a one-row opposite corner on at most four states]
\label{CP2:A-six-one-row-four-states}
Let $M$ be a six-element monoid, let $N$ be a finite monoid
whose minimal ideal has one row, and let $S$ be a nonempty
left-$N$, right-$M$ biset with at most four elements.
The associated two-object collage does not have
classifying space of type $\CP^2$.
\end{theorem}
\begin{proof}
Write $X$ for the collage space and $E=B(N,S,M)$ for its
two-sided action bar. The one-row hypothesis has two
consequences used separately below: $BN$ is rationally
acyclic, and the trivial \emph{left} $\mathbb QN$ module
is projective. For the latter, choose a minimal idempotent
$f$, with local group $G=fNf$. The average
$p=|G|^{-1}\sum_{g\in G}g$ obeys $np=p$ for every $n$,
because left multiplication permutes the group coordinate
in the unique row and the chosen column. Therefore
$\mathbb QN p\cong\mathbb Q$ is a projective summand.

If $H^2(BM;\mathbb Q)=0$, the relative-cup lemma already
excludes $X$. Otherwise use the preceding five source
types, and put $P=K^1\cong P_5$, $e=e_{00}$.
A state fixed by all band elements is fixed by $M$:
$sm=(se)m=s(em)=s$, since $em\in K$. Such a state makes
the source-corner inclusion nullhomotopic in $X$, so again
the relative-cup lemma applies.

Assume there are no fixed band states. The rectangular-band
act normal form then has regular two-point column cores.
If there are two cores, all four states are core states.
Write them $(b,j)$, where $j$ is the labelled column, and
put $h_{(b,j)}=e_{0j}$. For every $m\in M$ and $n\in N$,
\[
 h_s m=h_{sm},\qquad h_{ns}=h_s.
\]
These identities give the same equivariant natural
nullhomotopy of $E\to BM$ as in
Lemma~\ref{CP2:A-P5-full-regular-core}. Hence
$X\simeq BM\vee\operatorname{hocofib}(E\to BN)$.
The source summand already has one-dimensional second
cohomology and zero square. If $X$ were a projective plane,
its only second class would come from that summand and
would have zero square, a contradiction.

It remains to handle a unique core $\{c_0,c_1\}$.
Every band action on $S$ is then the constant map
$se_{ij}=c_j$, and every $N$ actor fixes $c_0,c_1$
individually. Write $A(s)=sa$.

\paragraph{Three split-algebra cases.}
For the first two source types, $a$ is central. In the
involution case put $p=(1+a)/2$; in the idempotent case
put $p=a$. The algebra $R=\mathbb QM$ splits as
\[
 R\cong\mathbb QP\times\mathbb Q,
 \qquad pR\cong\mathbb QP,
\]
and the trivial module lies in its $p$ summand. The monoid
map $\pi:M\to P$ fixes the band and sends $a$ to $1$.
Let $T$ be the actual quotient of $S$ by the $A$-orbits
in the involution case, or the image $AS$ in the idempotent
case. It is an $N$--$P$ biset, the quotient map is compatible
with $\pi$, and
$\mathbb Q[S]p\to\mathbb Q[T]$ is an isomorphism.

For the fifth source type, put $h=a-e$. Here $h$ is a
central idempotent, annihilates $\mathbb QK$ on both sides,
and $p=1-h$ gives the same algebra splitting, now with
$\pi(a)=e$. The set $AS$ is $N$ invariant, contains $c_0$
and does not contain $c_1$. Collapse $AS$ to $c_0$ to form
the actual $N$--$P$ quotient $T$. On permutation modules,
\[
 [s]p=[s]-[As]+[c_0].
\]
The vector $[c_0]$ and these vectors for $s\notin AS$
are a basis of $\mathbb Q[S]p$ and map to the evident
basis of $\mathbb Q[T]$.

In all three cases the resulting actual collage functor
induces a rational homology equivalence in every degree.
Indeed the trivial $M$ module is entirely in the $p$
block, and the displayed permutation-module isomorphism
identifies the derived right-action bar with that over
$P$. This comparison is $N$ linear, so it survives the
other derived tensor; the two corner maps and the
mapping-cylinder sequence give the assertion for the
collage. Its rational cohomology isomorphism respects
products because it is induced by an actual functor.
The target has at most four states. The proof of
Lemma~\ref{CP2:A-P5-four-state-acyclic-opposite} excludes
even a rational projective-plane cohomology ring: its
fixed-state, unique-core and full-core arguments use only
one-dimensional second cohomology and its nonzero square.
Thus these three cases are impossible.

\paragraph{The simultaneous-interchange case.}
Suppose $\sigma_a=\tau_a=(01)$ and $a^2=1$.
Set $v=e_{10}-e_{00}$ and $p_-=(1-a)/2$ in $R$.
An exact projective left resolution is
\[
 0\longrightarrow Re\xrightarrow{j}Rp_-
 \xrightarrow{d}Re\xrightarrow{\epsilon}\mathbb Q
 \longrightarrow0,
 \qquad d(p_-)=v,\qquad j(e_{i0})=e_{i0}p_-.
\]
The middle term has basis $p_-,e_{00}p_-,e_{10}p_-$;
$d$ has image $\mathbb Qv$ and kernel the last two
basis vectors, which are the displayed copy of $Re$.

The row interchange $\sigma(e_{ij})=e_{1-i,j}$,
$\sigma(a)=a$, fixes the right action on $S$ and gives
an automorphism $F$ of the whole category. Its semilinear
lift sends $e_{i0}$ in degree zero to $e_{1-i,0}$,
$rp_-$ in degree one to $-\sigma(r)p_-$, and $e_{i0}$
in degree two to $-e_{1-i,0}$. These maps commute with
both differentials. Tensoring with $\mathbb Q[S]$ gives
$\mathbb Qc_0$ in degrees zero and two, and
$\mathbb Q[S]p_-$ in degree one. The first differential
is zero, and the second sends $c_0$ to
$(c_0-c_1)/2$. This is an $N$-equivariant acyclic pair.
Its quotient is
\[
 \mathbb Q[0]\oplus T[1],
 \qquad T=\mathbb Q[S]p_-/\mathbb Q(c_0-c_1),
\]
with signs $+1,-1$ under $F$. Because $BN$ is rationally
acyclic, all positive rational homology of the bridge has
sign $-1$. The source resolution also gives homology only
in degrees zero and two, with sign $-1$ on $H_2(BM)$.
The mapping-cylinder sequence therefore gives sign $-1$
on $H_4(X)$.

If $X$ were a projective plane, a generator $\alpha$ in
degree two either restricts to zero on $BM$, when its
square vanishes by the relative-cup lemma, or has nonzero
restriction and hence $F^*\alpha=-\alpha$. Its square
then has sign $+1$, contrary to the sign on $H^4(X)$.

\paragraph{The nilpotent constant case: the exact relative complex.}
It remains to treat $\sigma_a=\tau_a=\mathrm{const}_0$
and $a^2=e$. Put
\[
 h=a-e,\qquad q=1-e,\qquad
 u_i=e_{i1}-e_{i0},\qquad v=e_{10}-e_{00}.
\]
Here $h$ is central, $h^2=0$, and $hK=Kh=0$.
The projectives $P_0=Re$ and $B=Rq$ have bases
$e_{00},e_{10}$ and $q,u_0,u_1,h$, respectively.
Let $j:Re\to B$ send $e_{i0}$ to $u_i$.
The following projective resolution continues in every
degree:
\[
 P_0=Re,\qquad P_1=B,\qquad P_k=Re\oplus B\quad(k\ge2),
\]
\[
 d_1(b)=bv,\qquad
 d_2(t,b)=j(t)+bh,\qquad
 d_k(t,b)=(0,j(t)+bh)\quad(k\ge3).
\]
To check exactness, $\ker\epsilon=\mathbb Qv$, while
$\ker d_1=\mathbb Q\{u_0,u_1,h\}$.
The map $j$ is an isomorphism onto the first two summands,
and right multiplication by $h$ sends $q$ to $h$ and
annihilates $u_0,u_1,h$. Thus the kernel of $d_2$, and
of each subsequent displayed map on its nonzero target
component, is exactly
$0\oplus\mathbb Q\{u_0,u_1,h\}$, the image of the
next differential. These assertions check exactness in
all degrees, not only through a finite prefix.

On the right permutation module put
\[
 D=\mathbb Q[S](1-e)
   =\operatorname{span}\{[s]-[c_0]:s\ne c_0\},
 \qquad d=[c_1]-[c_0],
 \qquad H(z)=zh.
\]
Then $H([s]-[c_0])=[As]-[c_0]$, $H^2=0$, and $Hd=0$.
Tensoring the resolution gives the $N$ complex
\[
 C_0=\mathbb Q,\quad C_1=D,\quad C_k=\mathbb Q\oplus D
 \ (k\ge2),
\]
with $d_1=0$, $d_2(t,z)=td+Hz$, and
$d_k(t,z)=(0,td+Hz)$ thereafter.

The augmentation $\mathbb Q[S]\to\mathbb Q$ models the
actual bridge projection to $BM$. On this resolution its
target complex is $\mathbb Q$ in degree zero and every
degree at least two, zero in degree one, and all target
differentials are zero. The chain map is the identity in
degree zero and the first-coordinate projection in all
degrees at least two. Its exact kernel is the $N$ complex
\[
 K_0=0,\qquad K_k=D\quad(k\ge1),\qquad
 d_K=H\quad(k\ge2).                                      \tag{*}
\]
Every target term is a trivial $N$ module. Since $BN$ is
rationally acyclic, its derived $N$ coinvariants compare
quasi-isomorphically to the target itself. Also collapsing
$BN$ in the collage is a rational homology equivalence.
The resulting cofiber is the cofiber of the actual
bridge projection to $BM$. The short exact kernel sequence
therefore gives
\[
 \widetilde C_*(X;\mathbb Q)\simeq
 \bigl(\mathbb Q\otimes_{\mathbb QN}^{\mathbf L}K\bigr)[1],
 \qquad
 H_2(X;\mathbb Q)\cong(D/HD)_N.                         \tag{**}
\]
The second assertion is the degree-one homology of the
derived tensor of a complex beginning in degree one.
This calculation uses the actual relative complex; it
does not presume that a vertical-homology summand survives
a spectral sequence.

The set map $A$ sends both core points to $c_0$ and has
$A^2$ equal to the constant $c_0$. With at most two other
states, its possibilities divide into the following three
exhaustive cases.

If $H=0$, the complex $(*)$ is
$\bigoplus_{k\ge1}D[k]$ with zero differential.
Every nonzero $\operatorname{Tor}^{\mathbb QN}_j(\mathbb Q,D)$
would consequently occur in arbitrarily high homological
degrees of $(**)$. Bounded homology forces all these Tor
groups to vanish, including coinvariants, and then
$H_2(X;\mathbb Q)=0$.

If $HD=\mathbb Qd=L$, this is a trivial left $N$ module
and hence projective. The surjection $H:D\to L$ has an
$N$-linear section with image $J\cong L$, so
$D=\ker H\oplus J$. Put $V=\ker H/L$.
There is an acyclic subcomplex of $(*)$ whose degree-one
term is $L$, whose terms in all higher degrees are
$L\oplus J$, and whose differential is the isomorphism
$J\to L$ on consecutive terms. The quotient is the
zero-differential complex
\[
 (D/L)[1]\ \oplus\ \bigoplus_{k\ge2}V[k],
 \qquad D/L\cong V\oplus J.
\]
Boundedness of $(**)$ forces every Tor group with $V$
to vanish because its shifts repeat indefinitely.
The remaining $J\cong\mathbb Q$ is projective on the
left, so it contributes only in degree two of $X$.
In particular $H_4(X;\mathbb Q)=0$.

In the last case $S=\{c_0,c_1,w,x\}$, after relabelling,
and
\[
 Ax=w,\qquad Aw=Ac_0=Ac_1=c_0,
 \qquad HD=\mathbb Q([w]-[c_0]).
\]
Every transformation commuting with $A$ and fixing the
two core points is determined by its value on $x$.
Its four possibilities are
\[
\begin{array}{c|cc}
 &x&w\\\hline
 I&x&w\\
 F&w&c_0\\
 P_0&c_0&c_0\\
 P_1&c_1&c_0.
\end{array}
\]
All fix $c_0,c_1$. With ordinary function composition,
$F^2=P_0$, $FP_i=P_i$, $P_iF=P_0$, and $P_iP_j=P_j$.
If both $P_0$ and $P_1$ occur in the image of $N$, the
module $D/HD$ has basis $d,[x]-[c_0]$ and its coinvariant
relations kill first $[x]-[c_0]$ using $P_0$, then $d$
using $P_1$. Formula $(**)$ gives $H_2(X;\mathbb Q)=0$.

If neither of these two idempotents occurs, only $I$
occurs, since $F^2=P_0$. The $N$ action is trivial.
Collapsing the rationally acyclic $BN$ gives an actual
functor to the ordinary $M$ act cone which is a rational
homology and cohomology-ring equivalence. Its homology
would be bounded, contrary to the bounded-small-act
theorem, for example
Theorem~\ref{CP2:A-bounded-eleven-acts}.

If exactly one of $P_0,P_1$ occurs, it is central in the
whole image of $N$. Indeed the only extra possible actor
is $F$, and it can occur only with $P_0$, with which it
commutes. The corresponding transformation is therefore
a proper $N$--$M$ biset idempotent with image
$\{c_0,c_1\}$. The actual cut-retraction theorem
\ref{CP2:A-cut-retract} would preserve projective-plane
type in that two-cross-arrow subcollage, contradicting
Lemma~\ref{CP2:A-high-two-cross}.
This exhausts the last case and proves the theorem.
\end{proof}

\begin{corollary}[Proved: the opposite-order-eight scope for $(6,8,4)$]
\label{CP2:A-six-eight-four-target-scope}
For the eighteen-arrow profile $(6,8,4)$ with a positive
source corner, it suffices to test the five sources of
Lemma~\ref{CP2:A-six-positive-shapes} against all $1205$
ordinary order-eight monoids having at least two minimal
rows and two minimal columns. Every table in that target
catalog is required, including those with zero rational
second homology.
\end{corollary}
\begin{proof}
A one-row opposite ideal is excluded by the theorem.
A one-column ideal gives exact projective-right
coinvariants, hence an actual orbit quotient inducing a
rational cohomology-ring isomorphism onto an ordinary
six-element-monoid act cone. Its bounded rational homology
forces zero products by the bounded-small-act theorem.
The remaining order-eight ideals have trivial group and
shapes $2\times2$, $2\times3$, or $3\times2$.
Their complete ordinary catalogs have
$1181+12+12=1205$ tables, as already proved and checked in
Proposition~\ref{CP2:A-P5-eight-corner-scope}.
This corollary closes the input reduction only; it makes
no claim that the new four-state finite biset search has
been completed.
\end{proof}


\subsection{The eight-state P5--P5 scope}

\begin{lemma}[Proved: a unique regular P5 core opposite P5]
\label{CP2:A-P5-double-unique-core}
Let $S$ be any finite nonempty left-$P_5$, right-$P_5$
biset. Suppose the right action has no fixed state and its
minimal-band core is a unique regular two-point column
orbit. The collage cannot have rational projective-plane
cohomology.
\end{lemma}
\begin{proof}
Write $P=P_5$, $E=B(P,S,P)$ and $X$ for the collage.
The semilinear resolution computation in
Lemma~\ref{CP2:A-mixed-five-reductions}(3) is equivariant
for any commuting opposite action. Before taking its
derived coinvariants it gives the complex
$\mathbb Q[0]\oplus T[1]$, with signs $+1,-1$ under
the source-row interchange $F$, respectively. Therefore
\[
 H_3(E;\mathbb Q)=H_3(BP;\mathbb Q)\oplus
 \operatorname{Tor}^{\mathbb QP}_2(\mathbb Q,T).
\]
The first summand is zero, so $F_*$ is $-1$ on this group.
Both corner spaces have zero rational homology in degrees
three and four by their length-two projective resolution.
The mapping-cylinder exact sequence consequently identifies
$H_4(X;\mathbb Q)$ with $H_3(E;\mathbb Q)$, with sign
$-1$. If $H^2(X;\mathbb Q)$ is one-dimensional, its
generator is an eigenvector of the involution $F^*$ with
eigenvalue $+1$ or $-1$. Its square has eigenvalue $+1$
in either case and must be zero. This contradicts the
projective-plane cohomology ring. No acyclicity assumption
on the opposite copy of $BP$ was used.
\end{proof}

\begin{proposition}[Proved: the complete finite right-act scope at $(5,5,8)$]
\label{CP2:A-P5-double-eight-scope}
To exclude the P5--P5 portion of the eighteen-arrow profile
$(5,5,8)$, it suffices to enumerate every left $P_5$ action
commuting with each of the twenty-nine specified eight-state
right acts. They are the right normal forms indexed
$1,2,\ldots,29$ in the checked eight-state catalog
\path{monoid/certificates/p5_small_targets/right_acts_8.tsv}.
State automorphisms may be divided out, with exact orbit
weights. No corner automorphism or opposition quotient is
needed in that finite enumeration.
\end{proposition}
\begin{proof}
If both actions have fixed states, the source and target
corner inclusions are separately nullhomotopic in the
collage, by the corresponding natural transformations.
Every second class therefore restricts to zero on both
corners and has zero square by the relative-cup lemma.
Otherwise at least one action has no fixed state.
Opposition and $P_5\cong P_5^{\mathrm{op}}$ put such an
action on the right, preserving the classifying-space
homotopy type. This step does not assert that both actions
are fixed-point free.

The independently checked eight-state normal-form catalog
has $1212$ ordinary right-act classes, of which exactly
the first $31$ have no fixed state. Index zero has a
unique two-point core and is excluded by the preceding
lemma. Index thirty consists of four regular two-point
cores and no transient states, so the full-core wedge
lemma applies. The other twenty-nine forms are the stated
scope. The cost inventory
\path{monoid/certificates/p5_double18_cost_inventory.json}
reconstructs their core types and counts every commuting
endomorphism from its core-orbit map and allowed transient
images; their largest centralizer has $2624$ elements.
The normal-form and Burnside receipt remains the earlier
complete checked eight-state receipt. A completed finite
search, its independent scope check and its literal cup
certificates are separate obligations.
\end{proof}


\begin{proposition}[Proved: the complete P5--P5 eight-state exclusion]
\label{CP2:A-P5-double-eight-complete}
No two-object category with both corner monoids $P_5$,
eight arrows in one direction and none in the other has
classifying space of type $\CP^2$.
\end{proposition}
\begin{proof}
Apply Proposition~\ref{CP2:A-P5-double-eight-scope}.
The separately named producer
\path{monoid/code/p5_double_eighteen.cpp} completes all
twenty-nine remaining right-act scopes. It constructs
every commuting transformation by propagation of point
images, then every homomorphism from the opposite $P_5$
into that full centralizer. Complete state orbits are
represented, with stabilizers and exact orbit-weight sums.
There are $74351$ labelled actions and $22352$ state orbits.
The necessary mod-two homology conditions exclude $1584$
orbits in degree one and $12362$ in degree two. All $8406$
remaining literal categories have an explicit nonzero
two-cocycle and a three-cochain whose coboundary is its
Alexander--Whitney square.

The independent scope checker
\path{monoid/code/p5_double_eighteen_check.cpp} constructs
centralizers from all core-orbit maps and compatible
transient images. It then assigns target images in fixed
order without propagation, independently recomputes the
first two homology groups, and matches every state orbit,
stabilizer and orbit weight. Its $157847$ search nodes
give exactly the stated complete counts. The separate
literal checker checks all $7145100$ associativity,
$4303872$ cocycle and $21519360$ square-primitive equations.
The exact stream join rebuilds each retained category
from its right and left actions and matches its complete
typed multiplication table and reason to the literal
certificate stream, through both ends of file.

A fresh seven-state control run passes the same three
independent checks and agrees, scope by scope, with every
included row of the earlier unrestricted seven-state
P5--P5 census. This control is additional to the complete
eight-state checks. Source, binary, input and closed-stream
hashes, commands and exact count joins are recorded in
\path{monoid/certificates/two18_p5_runs/p5_p5_s8.complete_verified.json}
under the checkpoint's external-file mapping. The three
analytic gates and the complete finite check exhaust the
scope. The mixed five-element corner cases are a separate
obligation.
\end{proof}


\subsection{The mixed five-element corner scope on eight states}

\begin{proposition}[Proved: the remaining mixed-five eighteen-arrow inputs]
\label{CP2:A-P5-mixed-eight-scope}
For the eighteen-arrow profile $(5,5,8)$, the portion not
having two copies of $P_5$ reduces, after opposition, to
the twenty-nine right $P_5$ acts indexed $1,\ldots,29$
in the complete eight-state normal-form catalog, against
the forty-two non-$P_5$ five-element opposite monoids
whose minimal ideal has more than one column.
Every commuting left action of each such monoid is required.
\end{proposition}
\begin{proof}
If neither corner is $P_5$, both have zero positive rational
homology by Lemma~\ref{CP2:A-order-five-corner}, so the
relative-cup lemma applies. If precisely one is $P_5$,
use opposition and its self-opposition isomorphism to put
it at the source. A one-column opposite ideal is excluded
by the exact-projective-right argument of
Lemma~\ref{CP2:A-mixed-five-reductions}(2).
The remaining opposite monoids are exactly the forty-two
previously classified and independently checked mixed-five
targets; all are rationally acyclic.

A fixed right state or a unique two-point right core is
excluded by parts (1) and (3) of that lemma. Among the
$1212$ eight-state right normal forms this leaves the
thirty forms indexed $1,\ldots,30$. The last has no
transients and is excluded by the full-regular-core wedge
lemma. Thus exactly the stated $29\cdot42=1218$ finite
scopes remain. This is an input-scope assertion; it does
not itself claim their finite enumeration has finished.
\end{proof}


\subsection{A row-independent P5 action across a cut}

\begin{proposition}[Proved: the relative sign obstruction]
\label{CP2:A-P5-row-independent}
Let $S$ be a finite nonempty left-$N$, right-$P_5$ biset
with no right fixed state. Suppose the two right band-row
actors agree on every state:
\[
 se_{0j}=se_{1j}\qquad(s\in S,\ j=0,1).
\]
Let $A$ be the finite left $N$ set of regular column-core
orbits. If the ordinary left $N$ act cone on $A$ has zero
rational degree-two products, then a collage on $S$ whose
second rational cohomology has dimension at most one has
zero degree-two squares.

The hypothesis on that act cone holds for every $N$ of
order at most seven, or for arbitrary finite $N$ when
$|A|\le2$. No bounded-homology hypothesis is required.
\end{proposition}
\begin{proof}
Let $U\subseteq S$ be the full band core, $X_U\subseteq X$
the subcollage, and $E_U\subseteq E$ their two-sided
action bars. Since there is no fixed state, every orbit
in $U$ is a regular two-point column orbit.
The full-core nullhomotopy gives
\[
 X_U\simeq BP_5\vee
 \operatorname{hocofib}\bigl(B(N,A)\longrightarrow BN\bigr).
\]
Thus every rational degree-two product on $X_U$ is zero
under the stated hypothesis.

The row interchange $\sigma(e_{ij})=e_{1-i,j}$, fixing
the other corner and every state, defines an involution
$F$ of the entire pair $(X,X_U)$. The row-independence
hypothesis is precisely the compatibility needed here.
We compute its relative action, retaining the actual
subcollage rather than replacing it by an abstract complex.

Tensor the length-two left projective $P_5$ resolution
from Lemma~\ref{CP2:A-order-five-corner} with the right
module $T=\mathbb Q[S]/\mathbb Q[U]$.
Every band element annihilates $T$. The degree-zero and
degree-two terms therefore vanish, and the degree-one
term is $T$. The semilinear lift used in
Lemma~\ref{CP2:A-mixed-five-reductions} acts by $-1$
on this surviving term. The relative right-action bar
is consequently $T[1]$, with $F=-1$, as an $N$ complex.
This computes the relative bar because each projective
resolution term is flat and the original permutation-module
sequence $0\to\mathbb Q[U]\to\mathbb Q[S]\to T\to0$
is exact. Taking derived $N$ coinvariants preserves this
relative comparison and the sign.

The two mapping-cylinder models share both corner spaces;
their quotient is $\Sigma(E/E_U)$. Hence $F=-1$ on
every rational relative homology group $H_*(X,X_U)$.
Over $\mathbb Q$ the eigenspaces of an involution split
exactly, so the homology sequence of the pair shows that
\[
 H_4(X_U;\mathbb Q)\longrightarrow H_4(X;\mathbb Q)
\]
surjects onto the $+1$ eigenspace. This conclusion does
not require $BN$ to be rationally acyclic.

If $H^2(X;\mathbb Q)$ has dimension one, a generator is
an eigenvector of $F^*$ with eigenvalue $+1$ or $-1$.
Its square has eigenvalue $+1$. It pairs to zero with
the negative fourth-homology eigenspace by its sign, and
with the positive eigenspace because that space comes
from $X_U$, where all second products vanish.
Its square is therefore zero. If the second cohomology
is zero the assertion is immediate.

Finally, all ordinary act cones over $|N|\le7$ have zero
second products by Corollary~\ref{CP2:A-high-seven-act}.
For $|A|\le2$ use the unconditional binary act theorem
(and the singleton cone). This proves both stated
unconditional instances.
\end{proof}

\begin{corollary}[Proved: scope refinement for mixed-five eight-state actions]
\label{CP2:A-P5-mixed-eight-row-scope}
In Proposition~\ref{CP2:A-P5-mixed-eight-scope}, right-act
indices $1,4,10,20,23$ may also be omitted. The remaining
twenty-four right acts against the forty-two mixed-five
targets give exactly $1008$ finite scopes.
\end{corollary}
\begin{proof}
Literal inspection of the checked normal-form maps shows
that the indicated five acts have equal row actors, and
the opposite corner has order five. The proposition
therefore applies. The other right indices retain their
previous scope. The new reduced enumeration must record
this gate explicitly; it does not alter or promote any
earlier uncompleted producer stream.
\end{proof}


\subsection{Orienting the fresh two-object search by a nonzero restriction}

\begin{lemma}[Proved: a fixed-point-free positive source orientation]
\label{CP2:A-two-object-positive-orientation}
Suppose a two-object collage $C$ has classifying space
of type $\CP^2$. After possibly taking the opposite
category, its source corner $M$ has nonzero second
rational homology and its connecting right $M$ act has
no global fixed state. In fact a generator of
$H^2(BC;\mathbb Q)$ restricts nontrivially to that source
corner.
\end{lemma}
\begin{proof}
Let $\alpha$ be the rational degree-two generator.
Its square is nonzero. By the relative-cup lemma,
$\alpha$ cannot restrict to zero on both corners.
Choose a corner on which its restriction is nonzero,
and if needed take opposition to make it the source.
Opposition preserves the homotopy type and the
nonvanishing of that restriction. A global fixed state
of the resulting right act would supply a natural
transformation from the source inclusion to the constant
functor at the other object, making its restriction zero.
This is impossible.
\end{proof}

\begin{corollary}[Proved: the joint fixed-point-free eighteen-arrow scope]
\label{CP2:A-two-object-eighteen-oriented-scope}
To close the fresh two-object eighteen-arrow lower-bound
scope, it suffices to check the twenty-five ordered
positive-source profiles of
Proposition~\ref{CP2:A-two-object-eighteen-scope}, using
only right acts without a global fixed state, provided
all otherwise required opposite-corner tables are included.
In particular, $P_5$ targets in profiles with source order
six through nine must be included under this orientation
convention. They cannot be omitted merely by citing the
oppositely ordered profile with a fixed-point-free $P_5$
right act.
\end{corollary}
\begin{proof}
Choose the orientation from the lemma before applying
the complete eighteen-arrow cardinality and analytic
reductions. The ordinary positive-source catalogs include
both ordinary opposition classes, so the oriented source
is present. The remaining ordered size triple is among
the same twenty-five profiles. Every retained analytic
input reduction continues to apply to this subfamily.
Conversely, the finite checks must now include every
allowed opposite table: exchanging corners again might
lose the specified nonzero restriction and fixed-point-free
right action. This explains the explicit $P_5$ target
requirement. Once all these ordered scopes are joined,
the lemma supplies a checked orientation for every
possible projective-plane collage.

For minimal-band right-act normal forms, a state fixed
by the whole band is automatically fixed by the entire
source monoid: $sm=(se)m=s(em)=s$ for a band idempotent
$e$, because $em$ remains in the band. Thus the literal
fixed-state metadata of the complete band catalogs
tests this condition before assigning outside actors.
For a nontrivial minimal group the same argument uses
a state fixed by every element of the minimal ideal.
\end{proof}

\begin{remark*}[Scope accounting]
The new convention is a joint orientation reduction.
An isolated fixed-point-free profile receipt does not
exclude all categories of that ordered profile by itself.
It becomes an unrestricted two-object exclusion only
after the complete ordered profile join and its opposite
dependencies are checked. Earlier complete seventeen-arrow
receipts and the unpruned mixed-five eighteen-arrow run
are unchanged.
\end{remark*}


\subsection{Small positive sources in the chosen orientation}

\begin{lemma}[Proved: a fixed-point-free P5 source on at most five states]
\label{CP2:A-P5-oriented-five}
For any finite opposite monoid $N$, a left-$N$,
right-$P_5$ biset on at most five states with no global
right fixed state cannot have a collage with the rational
cohomology ring of $\CP^2$.
\end{lemma}
\begin{proof}
There is either one regular two-point core orbit, or two.
With one orbit the two row actors agree on every state,
so Proposition~\ref{CP2:A-P5-row-independent} applies with
the one-element core-orbit set. With two orbits and no
transient, the full-core wedge lemma applies. The only
remaining case has four core states and one transient.
If the transient's two row images belong to the same
core orbit, the row actors agree and the row-independent
proposition applies with a two-element core-orbit set.
If those row images belong to different orbits, the right
act is the regular right $P_5$ act. The commuting left
action is left multiplication through a homomorphism
$N\to P_5$, and the bridge projection to $BN$ is an actual
homotopy equivalence. The collage is then homotopy
equivalent to $BP_5$. None of these arguments assumes
that $BN$ is rationally acyclic.
\end{proof}

\begin{lemma}[Proved: the nilpotent six-source low-degree comparison]
\label{CP2:A-six-nilpotent-low-degree-any-opposite}
Let $M$ be the six-element source with constant outside
actors and $a^2=e_{00}$. Suppose its right act has a
unique regular two-point band core $\{c_0,c_1\}$, with
any finite commuting opposite monoid $N$. Use
$D=\operatorname{span}\{[s]-[c_0]:s\ne c_0\}$,
$d=[c_1]-[c_0]$, and $H=A-E$ as in
Theorem~\ref{CP2:A-six-one-row-four-states}.
For its collage $X$ there is an isomorphism
\[
 H_2(X;\mathbb Q)\cong(D/HD)_N,
\]
under which the image of a generator of $H_2(BM;\mathbb Q)$
is the class of $d$, up to sign. No acyclicity or
projectivity assumption on $N$ is required here.
\end{lemma}
\begin{proof}
Use the exact all-degree source resolution displayed in
Theorem~\ref{CP2:A-six-one-row-four-states}. After tensoring
with $\mathbb Q[S]$, its degree-zero term is the trivial
$N$ module $\mathbb Qc_0$, its first differential is zero,
and its positive-degree subcomplex $C^+$ is
\[
 C^+_1=D,\qquad C^+_k=\mathbb Q\oplus D\ (k\ge2),
 \qquad d_2(t,z)=td+Hz.
\]
The right-action augmentation to a point is the projection
onto the degree-zero summand. It models the actual bridge
projection to $BN$ after taking the other bar. Thus, in
the chain homotopy pushout for the collage, that summand
cancels the $BN$ chains. What remains is the cone of
\[
 \mathbb Q\otimes_{\mathbb QN}^{\mathbf L}C^+
 \longrightarrow \widetilde C_*(BM;\mathbb Q).
\]
This is a comparison of the two actual projection maps:
the map to the source corner is induced by
$\mathbb Q[S]\to\mathbb Q$ and the augmentation of the
$N$ bar, while the map to the opposite corner is the
degree-zero projection just described.

The reduced source complex from that resolution has
$\mathbb Q$ in every degree at least two and zero
differentials. Its degree-two generator receives the
first coordinate of $C^+_2$. Using the free $N$ bar in
degrees zero and one, the degree-two group in the cone
is $\mathbb Q\oplus D$, its outgoing differential is zero,
and its degree-three boundaries span exactly
\[
 (t,-td-Hz),\qquad (0,nz-z)
 \quad(t\in\mathbb Q,\ z\in D,\ n\in N).
\]
There are no other degree-three terms contributing a
boundary in this degree. Eliminating the first coordinate
identifies the quotient with $(D/HD)_N$ and sends the
source vector $(1,0)$ to $[d]$. This proves both claims
without a high-degree comparison to $BN$.
\end{proof}

\begin{proposition}[Proved: a six-element source cannot support the chosen class on four states]
\label{CP2:A-six-oriented-four}
Let $M$ have six elements, let $N$ be any finite monoid,
and let $S$ be a left-$N$, right-$M$ biset on at most four
states. There is no projective-plane collage whose
degree-two generator restricts nontrivially to $BM$.
\end{proposition}
\begin{proof}
A global right fixed state would make that restriction
zero, so assume there is none. The positive-source
classification gives the five monoids of
Lemma~\ref{CP2:A-six-positive-shapes}.
If the band core has two regular orbits, all four states
are core states. The equivariant full-core nullhomotopy
from the proof of Theorem~\ref{CP2:A-six-one-row-four-states}
gives a $BM$ wedge summand, whose one-dimensional second
cohomology has zero square. This already excludes a
projective-plane collage.

Thus the core is the unique pair $U=\{c_0,c_1\}$.
The three split-algebra cases in that theorem give actual
rational cohomology-ring equivalences to $P_5$ collages
on quotient sets of size at most four. Those quotient
sets retain the two distinct core points and have no
right fixed state. The comparisons are $N$ linear before
taking its derived coinvariants, so they work for every
$N$. Lemma~\ref{CP2:A-P5-oriented-five} excludes them.

For the simultaneous-interchange source, use its exact
resolution $Re\to Rp_-\to Re$ and the same semilinear
row-swap lift. The core right act is $eM$, a retract of
the regular right act, so its action category is
contractible. Since $N$ fixes both core states, its core
bridge projects by a homotopy equivalence to $BN$.
The full-core nullhomotopy at the source consequently
identifies the core subcollage $X_U$ with $BM$.
The relative permutation module $\mathbb Q[S]/\mathbb Q[U]$
is annihilated by the band, so tensoring the source
resolution leaves only its $p_-$ part in degree one.
The row-swap lift acts there by $-1$.
As in Proposition~\ref{CP2:A-P5-row-independent}, all
relative homology of $(X,X_U)$ therefore has sign $-1$.
The positive fourth-homology eigenspace comes from
$H_4(BM;\mathbb Q)=0$. A square from one-dimensional
second cohomology has sign $+1$ and must vanish.

It remains to consider the nilpotent constant source.
The set-map classification from the six-source theorem
is exhaustive: $H=0$, $HD=\mathbb Qd$, or the four-state
chain $Ax=w$, $Aw=Ac_0=Ac_1=c_0$.
If $H=0$, the right transformation image is exactly the
identity and the two distinct constants, excluded by
Lemma~\ref{CP2:A-two-constant-image} for arbitrary $N$.
If $HD=\mathbb Qd$, the preceding low-degree lemma makes
$H_2(BM)\to H_2(X)$ zero. Duality contradicts the
specified nonzero restriction of the second generator.

In the chain case every opposite actor is one of
$I,F,P_0,P_1$ listed in the six-source proof. If both
$P_0,P_1$ occur, their relations kill both basis classes
in $(D/HD)_N$, so $H_2(X;\mathbb Q)=0$.
If exactly one occurs, it is a central proper biset
idempotent onto $U$. The cut-retraction theorem and the
unconditional two-cross-arrow lemma exclude this case.
If neither occurs, only the identity occurs and the
opposite action is trivial. The actual collapse to the
ordinary $M$ act cone induces an isomorphism onto second
rational cohomology by
Lemma~\ref{CP2:A-trivial-biset-side-H2}. All second
products on that cone vanish by the unconditional
order-seven act theorem. Thus its pullbacks exhaust
$H^2(X)$ and have zero square. This proves the proposition.
\end{proof}

\begin{corollary}[Proved: twenty-two remaining jointly oriented profiles]
\label{CP2:A-two-object-eighteen-oriented-twenty-two}
Under the nonzero-source-restriction orientation of
Lemma~\ref{CP2:A-two-object-positive-orientation}, the
ordered eighteen-arrow profiles $(5,8,5)$, $(5,9,4)$,
and $(6,8,4)$ are analytically excluded. The remaining
finite scope consists of the other twenty-two profiles
in Proposition~\ref{CP2:A-two-object-eighteen-scope}.
All of their opposite corners have order at most seven.
All permitted opposite tables, including the four reverse
$P_5$ target pieces, remain required.
\end{corollary}
\begin{proof}
Apply the five-state $P_5$ lemma to the first two profiles,
and the oriented six-source proposition to the third.
The complete twenty-five-profile list contains no other
opposite corner of order at least eight. These exclusions
use the chosen nonzero source restriction and are joined
through that orientation lemma; they do not assert that
an arbitrarily pre-oriented source has a nonzero
restriction.
\end{proof}


\subsection{A forest obstruction from a completed short map}

\begin{proposition}[Proved: a forest in the short composition map]
\label{A7:short-map-forest}
Let $\mathcal C$ be a finite acyclic category with full support
$0<1<2<3<4$, and write $h_{ij}=|\mathcal C(i,j)|$.
Suppose $h_{01}\leq2$ and $h_{03}\leq2$.
On the full vertex set $U=\mathcal C(0,2)$, form the simple
undirected graph $G_P$ whose edges are the two-element image
sets of the maps
\[
 \mathcal C(0,1)\longrightarrow U,\qquad a\longmapsto b\circ a,
 \qquad b\in\mathcal C(1,2).
\]
Repeated unordered pairs give one edge, and a one-element image
gives no edge. If $G_P$ is a forest, every rational degree-two
cohomology product of $B\mathcal C$ is zero.
There is no restriction on $h_{12}$, $h_{02}$ or $h_{04}$.

Dually, suppose $h_{34}\leq2$ and $h_{14}\leq2$.
On the full vertex set $W=\mathcal C(2,4)$, form the simple
graph $G_R$ from the two-element image sets of
\[
 \mathcal C(3,4)\longrightarrow W,\qquad d\longmapsto d\circ c,
 \qquad c\in\mathcal C(2,3).
\]
If $G_R$ is a forest, the same product-vanishing conclusion holds.
\end{proposition}
\begin{proof}
Collapse $\mathcal C(0,4)$ to a singleton. By
Proposition~\ref{A7:extreme-arrow-collapse}, the quotient map
induces a surjection on second rational cohomology and preserves
products. The short map defining $G_P$ is unchanged. It therefore
suffices to prove product vanishing in this quotient.

The quotient is the left-act cone of its full subcategory on
$1,2,3,4$, with source stalks
\[
 A=\mathcal C(0,1),\quad U=\mathcal C(0,2),\quad
 L=\mathcal C(0,3),\quad \{*\}.
\]
Give $A$ and $L$ their unique trees with at most two vertices.
Extend the forest $G_P$, including all isolated vertices of $U$,
to a tree on exactly $U$. Such an extension exists: connect its
distinct connected components by arbitrary additional edges.
Every act map $A\to U$ sends its only possible edge to an edge
of this chosen tree or to a vertex, by the definition of $G_P$.
Every other nonidentity act map has target $L$ or $\{*\}$, so
it also sends each tree edge to an edge or a vertex.
The compatible-tree lemma, Lemma~\ref{A7:tree-act}, annihilates
all positive rational cohomology products of this act cone.
Pullback proves the asserted degree-two vanishing for
$\mathcal C$. Applying this argument to the opposite category
gives the sink statement.
\end{proof}

\begin{corollary}[Proved: an entire cylinder of short-map completions]
\label{A7:short-map-forest-prefix}
Fix a positive hom-set profile and a partial assignment of the
three short maps
\[
 P:A\times B\to U,\qquad Q:B\times C\to V,
 \qquad R:C\times D\to W.
\]
If every value of $P$ is assigned and its graph $G_P$ satisfies
the source hypothesis of Proposition~\ref{A7:short-map-forest},
then no completion of this partial assignment to an associative
category of that profile can model $\CP^2$.
The dual assertion holds when every value of $R$ is assigned and
its graph satisfies the sink hypothesis. In particular, the
source obstruction can be applied before any value of $Q$ or $R$
has been assigned.
\end{corollary}
\begin{proof}
The graph in the proposition depends only on the completed map
$P$, respectively $R$, and on the fixed hom-set cardinalities.
Every completion has exactly the same graph and obeys the same
cardinality hypotheses. Its rational second products therefore
vanish, which is incompatible with the projective-plane ring.
\end{proof}

\begin{remark*}
Completeness of the relevant short map is essential to this
certificate. A forest made only from some assigned columns can
acquire a cycle when another column is assigned. The test used
here first verifies that every entry of $P$, respectively $R$,
is known. It then retains the full vertex set, discards loops,
and deduplicates unordered edges before testing for a cycle.
The resulting obstruction removes every literal completion of
the recorded global prefix. Its unrestricted completion count
is the product of the appropriate hom-set sizes over all unknown
short-map entries; this number is not a count of associative
categories or of first-use-normalized representatives.
\end{remark*}


\subsection{Published counts and the exact scopes of our catalogs}

\begin{proposition}[Proved: catalog completeness in the scopes actually used]
\label{CP2:A-catalog-scope-comparison-through-ten}
Our ordinary monoid catalogs are complete in every order from one
through seven. Our source catalogs in orders eight, nine and ten
are complete for the condition $H_2(BM;\mathbb Q)\ne0$; they are
not catalogs of all monoids in those orders. The exact counts are
\[
\begin{array}{c|r|r}
\text{order}&\text{complete ordinary catalog}&H_2(\mathbb Q)\ne0\\\hline
1&1&0\\
2&2&0\\
3&7&0\\
4&35&0\\
5&228&1\\
6&2237&5\\
7&31559&43
\end{array}
\]
\[
\begin{array}{c|r|r}
\text{order}&\text{published ordinary count}&\text{positive source catalog}\\\hline
8&1668997&510\\
9&3685886630&8046\\
10&\text{not listed}&157145
\end{array}
\]
Here ``not listed'' refers only to the opened ordinary-count source,
not to a claim that the number is unknown. The complete target
catalogs restricted to at least two minimal rows and two minimal
columns have $1205$ entries in order eight and $21920$ in order nine.
Those target catalogs include their rational-$H_2$-zero entries.
\end{proposition}
\begin{proof}
Theorem~\ref{CP2:A-monoid-census} gives the independent full table
enumeration, unrestricted-permutation checker and exact equality
with the supplied RESIDUE census in orders one through five.
Its counts agree in each order with the official OEIS sequence
\cite{OEIS_A058129}; the agreement supplements the exhaustion
proof and does not replace it. The positive subsets through order
seven are compared as complete canonical sets with those full
catalogs. Proposition~\ref{CP2:A-positive-corners-eight},
Proposition~\ref{CP2:A-positive-corners-nine} and
Theorem~\ref{CP2:A-positive-corners-ten} prove the complete
positive-source scopes in the larger orders, including the two
nontrivial-minimal-group classes in order nine and fourteen in
order ten. Their proof reductions are what justify omitting other
monoids from a source search.

The order-eight two-sided target count is
$1181+12+12=1205$. In order nine it is
$21534+191+191+1+1+2=21920$, by the complete minimal-ideal
families in the order-nine proposition. Positivity is not imposed
on either target list. The stronger jointly oriented
eighteen-arrow reduction in
Corollary~\ref{CP2:A-two-object-eighteen-oriented-twenty-two}
subsequently removes the profiles requiring those larger targets;
the earlier catalogs and their exact scopes remain valid.

For comparison, the opened Distler--Kelsey university repository
abstract reports the following counts up to \emph{isomorphism and
anti-isomorphism}, also listed in OEIS A058133:
\[
\begin{array}{c|r}
\text{order}&\text{isomorphism--opposition classes}\\\hline
8&858977\\
9&1844075697\\
10&52991253973742.
\end{array}
\]
We use that abstract only for its published counts; we do not claim
to have read the article's enumeration proof
\cite{CP2-DistlerKelsey2009,OEIS_A058133}.
If $a$ is the ordinary count, $b$ the opposition-orbit count and
$s$ the number of self-opposite ordinary classes, opposition is
an involution and $a=2b-s$. The opened A058132 values
$s_8=48957$ and $s_9=2264764$ give exactly the two ordinary
counts in the table above \cite{OEIS_A058132}. Neither the
ordinary nor the opposition total is equated with a positive-source
subset. The opened ordinary and self-opposite sequences do not
list an order-ten value, so we infer none from them.
\end{proof}

\begin{remark*}[Inherited covers and the completeness obligation]
The inherited order-seven list of $32823$ oriented tables is a
redundant cover, not an ordinary-isomorphism census of size
$31559$. The inherited order-eight positive-source cover has
$516$ entries, including six explicitly marked rational-$H_2$-zero
entries. Our exact positive count is $510$. We neither identify
those list lengths with the published all-monoid totals nor infer
completeness from their agreement. Each finite gate uses the full
catalog or the proved restricted catalog appropriate to its role;
missing historical output is never a premise of this comparison.
\end{remark*}


\subsection{The fresh seventeen-arrow linear search}
We independently enumerate the $4368$ ordered positive six-tuples of sum
seventeen and their $2240$ opposition orbits. The complete exclusion is
proved in Theorem~\ref{CP2:A-linear-seventeen-complete}: $2065$ original
proof cases, one new three-by-three triangle case, and $174$ complete
finite profiles. These comprise the original $150$ profiles and $24$
additional profiles identified by the scope audit. Their literal checks
cover $18{,}114{,}854{,}855$ weighted tables and $30{,}884{,}728$ retained
records, with no survivor. The complete independent join is recorded in
\path{outputs/linear17_complete_summary.json}.

Every corner input comes from our complete ordinary-isomorphism census
through order seven. The original row-major search tries every unforced
typed composition and propagates only required associativity equations.
Its finite table cache can repeat a record after eviction but cannot
lose an unseen case. The later canonical-biset enumeration uses the
proved exact orbit/stabilizer construction and all balanced equivariant
composition maps of Proposition~\ref{CP2:A-linear-biset-enumeration}.
Its integer weights restore labelled-table counts. Only the separately
proved corner, graph, singleton-long and projective-singleton cuts are
used, with exact source and proof digests in each receipt.

Every retained table is recorded literally with its nonboundary mod-two
second cocycle and either a three-cochain primitive for its square or
positive third homology. The separately written checker verifies all
identity, typing and associativity equations, the bar ranks and every
certificate equation. The complete-scope finalizer independently
regenerates all $2240$ profiles, distinguishes the two frozen profile
inputs, and checks each receipt's range, profile, source, proof,
certificate and completion-footer fields. An absent receipt remains a
pending profile. One historical proof statement was broadened to its
actual sided projectivity hypothesis; the exact old and current text
digests, checkpoint location and checked comparison are explicitly bound
in the join. No arbitrary proof-digest mismatch is accepted.

\section{Part B: two-dimensional resolutions}
\label{CP2:partB}
The three cardinal-gap questions and both parts of Lead 1 remain on the
active ledger. In accordance with the requested priority, the main
finite-category problem is being pursued first. We have not yet checked
the nonsingular-equations preprint against the precise S2COUNT and
INDLOCAL uses, so no conclusion from those uses is asserted here.

The outstanding first case is whether every perfect group admitting a
two-dimensional classifying space and having $H_2(G;\mathbb Z)\cong
\mathbb Z$ has a nontrivial countable quotient. The subsequent questions
ask for the exact cardinal intervals below the continuum and below
$\kappa^{\aleph_0}$ in infinite homology rank. Their statements are not
replaced by claims about only countable groups or only one cardinal.

\section{Part C: the independent converse and source record}
\label{CP2:partC}
% Independently authored for FINMON_CP2.  Include in a document with amsthm.
\subsection{Finite monoids with bounded homology}

We count identities among the arrows of a finite category.  In this
subsection, a finite monoid is regarded as a category with one object.
The universal group of a monoid $M$ is denoted by $U(M)$.
The checks below confirm the four assertions requested in C.1.  The
chronology of the independent derivations is recorded separately in the
check-event table; a later comparison with the supplied proofs is not
mislabelled as an independent derivation.

\begin{lemma}[Proved: nilpotence from a cohomological bound]
\label{CP2:C-Knilpotence}
Let $R=\mathbb Z$ or $R=\mathbb Z_{(p)}$, and let $X$ be a connected CW
complex such that $H^q(X;R)=0$ for all $q>n$.  The ideal of classes of
rank zero in $\mathrm{KU}_R^0(X)$ is nilpotent.
\end{lemma}
\begin{proof}
We give the convergence argument explicitly.  Let $Y$ be the zero
component of the zeroth infinite-loop space of $\mathrm{KU}_R$.
Its positive homotopy groups are $R$ in even degrees and zero in odd
degrees.  Choose an even integer $N>n+1$, and let $F$ be the homotopy
fiber of the Postnikov map $Y\to P_NY$.
The finite Postnikov sections of $F$ have their nonzero homotopy groups
in even degrees greater than $N$.  Obstruction theory therefore gives
\[
 [X,P_mF]=0,\qquad [\Sigma X,P_mF]=0
\]
for every $m$: the relevant obstruction and choice groups are ordinary
cohomology groups in degrees greater than $n$.  We use reduced
cohomology and based maps here.  Since a simply connected space is the
homotopy limit of its Postnikov sections, the mapping-space homotopy
limit sequence has zero inverse-limit and zero derived-limit terms in
these two degrees.  It follows that $[X,F]=0$ and hence that
\[
 [X,Y]\longrightarrow [X,P_NY]
\]
is injective.  This argument applies to arbitrary CW complexes $X$;
finite type is not assumed.

The product of $r$ rank-zero classes is induced by the reduced
multiplication map $Y^{\wedge r}\to Y$.  The factorization through the
smash product follows because a zero factor makes the product zero;
it can equally be expressed using the multiplication of the ring
spectrum $\mathrm{KU}_R$.  As $Y$ is simply connected,
$Y^{\wedge r}$ is $(2r-1)$-connected.  If $2r-1\ge N$, its map to the
$N$-type $P_NY$ is nullhomotopic.  The product is consequently zero by
the injectivity just proved.  The same $r$ works for every choice of
the rank-zero classes.
\end{proof}

\begin{theorem}[Proved: R9 Theorems 1.1 and 1.2]
\label{CP2:C-finite-group-maps}
Let $X$ be a connected CW complex with finite fundamental group
$\pi$, let $G$ be a finite group, and let $f:BG\to X$ be a continuous
map.  If $H^q(X;\mathbb Z)=0$ in all sufficiently high degrees, then
$f_*:G\to\pi$ is trivial.  More precisely, if
$H^q(X;\mathbb Z_{(p)})=0$ in all sufficiently high degrees, then
$p$ does not divide the order of $f_*(G)$.
\end{theorem}
\begin{proof}
Let $c:X\to B\pi$ classify the universal covering.  The regular
complex representation of $\pi$ gives a vector bundle $V\to X$ of
rank $d=|\pi|$.  By Lemma~\ref{CP2:C-Knilpotence},
$v=[V]-d$ is nilpotent in the corresponding K-theory ring.

Suppose first that integral cohomology is bounded and $f_*(G)$ is
nontrivial.  There is an element $g\in G$ of prime-power order
$p^a$ with $f_*(g)\ne1$: take a preimage of an element of order $p$
and then take its $p$-primary component.  Put $P=\langle g\rangle$.
The composite $BP\to BG\xrightarrow{f}X$ pulls $v$ back to the
virtual representation
\[
 w=\operatorname{Res}_{P}(\mathbb C[\pi])-d.
\]
Atiyah's Theorem~7.2 identifies the inverse limit of K-theory over
finite skeleta of $BP$ with the augmentation-ideal completion of
$R(P)$; Proposition~6.11 says that $R(P)$ injects into that
completion~\cite{CP2-Atiyah1961}.  A zero power of the pulled-back
class would therefore give a zero power of $w$ in $R(P)$.  Evaluation
of its character at $g$, however, gives $w(g)=-d$, so no power is
zero.  This is a contradiction.

For the prime-local assertion, suppose $p$ divides $|f_*(G)|$ and
choose $P$ as above for this prime.  We avoid commuting localization
with an infinite inverse limit.  Model $BP$ by the infinite lens
space and restrict to each finite odd-dimensional lens skeleton
$L^{2j+1}$.  Its reduced integral $K^0$ is a finite $p$-group: the
finite Atiyah--Hirzebruch spectral sequence has only finite
$p$-groups in the positive even positions contributing to $K^0$.
Thus
\[
 K^0(L^{2j+1})\longrightarrow
 K^0(L^{2j+1})\otimes\mathbb Z_{(p)}
\]
is injective.  On a finite CW complex the latter is
$\mathrm{KU}_{(p)}^0$.  Nilpotence of the localized pullback therefore
makes one fixed power of $w$ vanish integrally on every such skeleton.
Atiyah's Theorem~7.2 and Proposition~6.11 again force that power to
vanish in $R(P)$, contradicted by its character value $(-d)^r$.

Every map and bundle pullback in this proof is in the category of
based CW spaces up to based homotopy.  Removing basepoints only
conjugates the induced group homomorphism and does not change either
conclusion.
\end{proof}

\begin{remark}
\label{CP2:C-linear-extension}
The regular representation can be replaced by any finite-dimensional
complex representation of $\pi_1(X)$.  A finite-order complex matrix
is diagonalizable with eigenvalues on the unit circle, so its trace
equals its dimension only when it is the identity.  Consequently,
without assuming $\pi_1(X)$ finite, the image of $f_*$ lies in the
intersection of the kernels of all such representations.  In
particular the integral assertion holds when $\pi_1(X)$ is
residually finite or is a complex linear group.
\end{remark}

\begin{lemma}[Proved: the finite-group corner and its comparison map]
\label{CP2:C-minimal-corner}
For a finite monoid $M$ there is an idempotent $e$ such that $H=eMe$
is a finite group and a map $j_e:BH\to BM$ inducing a surjection
$H\to U(M)$ on fundamental groups.  In particular $U(M)$ is finite.
\end{lemma}
\begin{proof}
Choose a minimal nonempty two-sided ideal $I$ of $M$ and an idempotent
$e\in I$; a suitable positive power of any element of $I$ is
idempotent because its powers are eventually periodic.  For
$x\in eMe$, minimality gives $MxM=I$, hence
$e=axb$ for some $a,b\in M$.  Thus
\[
 e=(eae)x(ebe)
\]
inside the finite monoid $eMe$.  In a finite monoid, a product equal
to the identity forces both factors to be units: if $uv=1$, left
multiplication by $u$ is surjective and hence injective, so
$vu=1$.  Applying this twice to the displayed factorization shows
that every $x\in eMe$ is a unit.  Thus $H=eMe$ is a group.

The universal map $\eta:M\to U(M)$ sends idempotents to $1$, and
\[
 \eta(m)=\eta(eme) \qquad(m\in M).
\]
The subgroup $\eta(H)$ therefore contains every generator of
$U(M)$, proving surjectivity.

We spell out the topological comparison, since the inclusion
$eMe\subset M$ need not preserve the identity.  Let
$\operatorname{Kar}(M)$ have the idempotents as objects and
$\operatorname{Kar}(M)(a,b)=bMa$; composition is multiplication.
The full inclusion $i:M\to\operatorname{Kar}(M)$ at the object $1$
is a homotopy equivalence on classifying spaces.  Indeed, for an
idempotent $a$, the comma category $(i\downarrow a)$ has objects
$x\in aM$ and morphisms $m:x\to y$ satisfying $ym=x$.
The arrows $x:x\to a$ define a natural transformation from the
identity functor to the constant functor at $a$; naturality is
exactly $ym=x$.  This contracts its classifying space.  Quillen's
Theorem~A proves that $Bi$ is a homotopy equivalence.
The full inclusion of the object $e$ now defines the span
\[
 BH\longrightarrow B\operatorname{Kar}(M)
       \xleftarrow[\simeq]{Bi} BM.
\]
We define $j_e$ using a homotopy inverse of $Bi$.  Its effect on
fundamental groups is the homomorphism $\eta|_H$, since the two
arrows represented by $e$ identify the object $e$ with the object
$1$ after groupoid completion.  This proves the asserted typing and
surjectivity of the map.
\end{proof}

\begin{theorem}[Proved: R9 Theorem 2.1]
\label{CP2:C-unbounded-homology}
If a finite monoid $M$ has $H_k(BM;\mathbb Z)=0$ in all sufficiently
high degrees, then $U(M)=1$.  If a prime $p$ divides $|U(M)|$, then
$H_k(BM;\mathbb F_p)$ is nonzero in infinitely many degrees.
\end{theorem}
\begin{proof}
The integral universal-coefficient exact sequence converts eventual
vanishing of integral homology into eventual vanishing of integral
cohomology.  Apply Theorem~\ref{CP2:C-finite-group-maps} to the map
of Lemma~\ref{CP2:C-minimal-corner}; its induced homomorphism is both
surjective and trivial, so $U(M)=1$.

For the second assertion, suppose mod-$p$ homology vanished in all
large degrees.  The normalized bar chain group in every degree has
finite rank because $M$ is finite, so all integral homology groups
are finitely generated.  The homology universal-coefficient exact
sequence implies $H_k(BM;\mathbb Z)\otimes\mathbb F_p=0$ for all
large $k$.  A finitely generated abelian group with this property
has zero localization at $p$.  The cohomology universal-coefficient
sequence then gives eventual vanishing of
$H^q(BM;\mathbb Z_{(p)})$.  The prime-local conclusion of
Theorem~\ref{CP2:C-finite-group-maps}, applied to the same surjection,
implies $p\nmid|U(M)|$, a contradiction.
\end{proof}

\begin{theorem}[Proved: R9 Theorem 2.2]
\label{CP2:C-finite-equivalences}
For a finite monoid $M$, the following are equivalent:
\begin{enumerate}
\item $BM$ is homotopy equivalent to a finite CW complex;
\item $BM$ is homotopy equivalent to a finite-dimensional CW complex;
\item $H_k(BM;\mathbb Z)=0$ in all sufficiently high degrees.
\end{enumerate}
When these conditions hold, $BM$ is simply connected.  If homology
vanishes above degree $d$, a finite model can be chosen of dimension
at most $d+1$; if $H_d$ is torsion-free, dimension $d$ suffices.
\end{theorem}
\begin{proof}
The first implication and the implication from the second condition
to the third follow from cellular homology.  Under the third
condition Theorem~\ref{CP2:C-unbounded-homology} gives simple
connectivity.  The bar complex has finite rank in each degree, so
integral homology is finitely generated.  Hatcher's
Proposition~4C.1~\cite{CP2-Hatcher2002}, whose exact statement we
opened, supplies a homotopy-equivalent CW complex with one
generator cell for each infinite cyclic homology summand and a
generator--relator pair in successive dimensions for each finite
cyclic summand.  Only finitely many cells are required here.  The
dimension assertions follow from that same cell count.

An independent derivation recorded before reading the supplied
proof also constructs a finite model directly.  For $d\ge2$,
take $L=(BM)^{(d+1)}$.  Relative Hurewicz for $(BM,L)$ identifies
$H_{d+1}(L)$ with $\pi_{d+2}(BM,L)$.  The former is finite free.
Attach one $(d+2)$-cell for each basis element, using representative
relative disks in $BM$.  The resulting finite simply connected
complex maps to $BM$ by an integral homology isomorphism, and the
homological Whitehead theorem makes it a homotopy equivalence.
The Hatcher argument above gives the sharper stated dimension.
\end{proof}

\begin{corollary}[Proved using the stated outside theorem: R9 Corollary 2.5]
\label{CP2:C-characterization}
The homotopy types of classifying spaces of finite monoids having
bounded integral homology are exactly the homotopy types of finite
simply connected CW complexes.  The realization in the reverse
direction can be chosen to be aperiodic.
\end{corollary}
\begin{proof}
The forward implication is Theorem~\ref{CP2:C-finite-equivalences}.
For the reverse implication we use Ben-Bassat's Theorem~C
\cite{CP2-BenBassat2026}.  It realizes every finite simply connected
CW complex by a finite aperiodic monoid.  The resulting homology is
bounded because homotopy equivalence preserves homology.

We checked the proof at the steps used here.  Theorem~A identifies
the framed collage with the homotopy cofiber, respecting its map
from the old corner.  Lemma~5.3 replaces a finite connected
simplicial source by a reduced source through contractible
edge-replacement diagrams.  Theorem~B therefore applies to each
attaching sphere.  Theorem~C starts with a finite simply connected
CW model without one-cells and applies Theorem~B successively; each
attaching sphere is connected.  The aperiodicity check is the
stationary-power argument in Lemma~5.2.  We found no fault in these
steps.  We cite the unrefereed theorem as an outside result and do
not claim an independent re-proof of that whole paper.
\end{proof}

\begin{proposition}[Refuted: R8 Conjecture 7.1; proved replacement]
\label{CP2:C-conjecture-counterexample}
Finite fundamental group does not suffice for a connected finite CW
complex to be a finite-monoid classifying space.  The exact
replacement, using the outside theorem in
Corollary~\ref{CP2:C-characterization}, is simple connectivity.
In particular, $\mathbb{RP}^2$ is not homotopy equivalent to $BM$
for any finite monoid $M$.
\end{proposition}
\begin{proof}
The usual CW structure on $\mathbb{RP}^2$ has one cell in each of
dimensions $0,1,2$, with the two-cell attached along the square of
the one-cell loop.  It is finite and connected, its fundamental
group is $\langle a\mid a^2=1\rangle\cong C_2$, and its homology
vanishes above degree two.  Theorem~\ref{CP2:C-unbounded-homology}
therefore rules out every finite monoid.

For a completely enumerated finite counterexample we also give the
simplicial complex $K$ on vertices $1,\ldots,6$ with facets
\[
\begin{gathered}
123,\quad124,\quad135,\quad146,\quad156,\\
236,\quad245,\quad256,\quad345,\quad346.
\end{gathered}
\]
Its complete graph is connected.  Collapse the tree with edges
$12,13,14,15,16$.  Its fundamental-group presentation has the ten
remaining edges as generators and the ten triangular boundary
words as relators.  Nine elementary Tietze eliminations, recorded
literally in \texttt{rp2\_certificate.json}, leave one generator
with its square as the sole relator.  Thus $\pi_1(K)=C_2$.
The same certificate contains its cellular boundary matrix after
tree collapse; its determinant is $-2$, giving
$H_1(K;\mathbb Z)=\mathbb Z/2$ and $H_2(K;\mathbb Z)=0$.
All higher homology vanishes by dimension.  Our independent checker
verifies each Tietze substitution, the determinant by integer
elimination, the two-element group table, and the full composition
table of the face-poset category.  That category has 31 objects and
121 arrows including identities; its nerve is the barycentric
subdivision of $K$.  In particular this is a finite explicit
counterexample without any reliance on recognizing a surface.
Necessity and sufficiency of the corrected simple-connectivity
condition are precisely
Corollary~\ref{CP2:C-characterization}.
\end{proof}

\subsection{Scope of the source record}
The bibliography marks sources actually opened. The machine-readable
source records give exact URLs and hashes for downloaded files. The
report separately lists sources not yet obtained or not yet checked at
the needed theorem. Attached sources and other authors' papers are not
copied into the checkpoint. Theorems carried only as background are not
converted into proved results by their occurrence in a prior register.
\subsection{Comparison with the opened semigroup literature}
We opened Sweeney's Theorem 2.7 and Corollary 2.8, and the
integral projective-resolution method in Sections 3.1--3.5,
including Examples 3.1 and 3.3\cite{CP2-Sweeney2025}.
The idempotent-ideal projectives in our certificates implement
the same homological principle: exact embedded kernel lattices
are covered by direct summands of free modules, trivial
coinvariants turn the maps into integer matrices, and Smith
normal form records torsion. Our independent verifier checks
the full ideals, actions, differential equations and saturated
images, rather than assuming exactness from a truncated rank
calculation. The degree-two bar comparisons and cup primitives
are additional data needed to distinguish a projective plane
from a space having only the same additive homology. We have
not run Sweeney's external software package, and do not claim
equality with an unreceived output catalog.

Fiedorowicz's unnumbered lemma on pages 33--34 of his 2002
paper gives the five-element rectangular-band sphere monoid
and a projective right-module resolution\cite{CP2-Fiedorowicz2002}.
Our rectangular-band calculation contains this example, with
the left-module orientation used in Sweeney's Example 3.1.
We opened that entire three-page paper, including its separate
unnumbered simplicial-monoid counterexample theorem on pages
34--35. That theorem is not a minimality statement for finite
projective-plane monoids.

We also opened the introduction and Theorems A--C of
Martinian--Steinberg\cite{CP2-MartinianSteinberg2025}.
Their finite wedges of simply connected Moore spaces do not
supply the projective-plane cup square: positive-degree
products on such a wedge vanish. The proof-level comparison
with that paper remains explicitly unfinished. Exact source
versions, URLs, hashes, read locators and remaining obligations
are recorded in \path{notes/root_source_record_R11.json} and
the report. Availability of an attached source is not treated
as having checked all its proofs.

\bibliographystyle{plain}
\bibliography{FINMON_CP2_references}
\end{document}
